\documentclass[a4paper,12pt]{book}

%% ─────────────────────────────────────────────
%%  PACKAGES
%% ─────────────────────────────────────────────
\usepackage[utf8]{inputenc}
\usepackage[T1]{fontenc}
\usepackage[english]{babel}
\usepackage{amsmath,amssymb,amsthm,mathtools}
\usepackage{tcolorbox}
\tcbuselibrary{theorems,skins,breakable}
\usepackage{tikz,tikz-cd,pgfplots}
\pgfplotsset{compat=1.18}
\usetikzlibrary{patterns,arrows.meta,calc,decorations.markings}
\usepackage{graphicx,float,caption}
\usepackage{enumitem,booktabs,array,longtable}
\usepackage{hyperref}
\usepackage{cleveref}
\usepackage{makeidx}
\usepackage{minitoc}
\usepackage{esint}
\makeindex

%% ─────────────────────────────────────────────
%%  MATH OPERATORS
%% ─────────────────────────────────────────────
\DeclareMathOperator{\Hom}{Hom}
\DeclareMathOperator{\End}{End}
\DeclareMathOperator{\Aut}{Aut}
\DeclareMathOperator{\im}{im}
\DeclareMathOperator{\Id}{Id}
\DeclareMathOperator{\Ker}{Ker}
\DeclareMathOperator{\coker}{coker}
\DeclareMathOperator{\Spec}{Spec}
\DeclareMathOperator{\Tr}{Tr}
\DeclareMathOperator{\supp}{supp}
\DeclareMathOperator{\diam}{diam}
\DeclareMathOperator{\sgn}{sgn}
\DeclareMathOperator{\conv}{conv}
\DeclareMathOperator{\codim}{codim}
\DeclareMathOperator{\spn}{span}
\DeclareMathOperator{\dom}{dom}
\DeclareMathOperator{\ran}{ran}
\DeclareMathOperator{\dist}{dist}
\newcommand{\N}{\mathbb{N}}
\newcommand{\Z}{\mathbb{Z}}
\newcommand{\Q}{\mathbb{Q}}
\newcommand{\R}{\mathbb{R}}
\newcommand{\C}{\mathbb{C}}
\newcommand{\K}{\mathbb{K}}
\newcommand{\F}{\mathbb{F}}
\newcommand{\dd}{\mathrm{d}}
\newcommand{\norm}[1]{\left\| #1 \right\|}
\newcommand{\abs}[1]{\left| #1 \right|}
\newcommand{\inner}[2]{\left\langle #1, #2 \right\rangle}
\newcommand{\weakto}{\rightharpoonup}
\newcommand{\weakstarto}{\overset{*}{\rightharpoonup}}

%% ─────────────────────────────────────────────
%%  THEOREM ENVIRONMENTS
%% ─────────────────────────────────────────────
\theoremstyle{definition}
\newtheorem{definition}{Definition}[chapter]
\newtheorem{example}[definition]{Example}
\newtheorem{exercise}{Exercise}[chapter]
\theoremstyle{plain}
\newtheorem{theorem}[definition]{Theorem}
\newtheorem{proposition}[definition]{Proposition}
\newtheorem{lemma}[definition]{Lemma}
\newtheorem{corollary}[definition]{Corollary}
\theoremstyle{remark}
\newtheorem{remark}[definition]{Remark}
\newtheorem{notation}[definition]{Notation}

\tcolorboxenvironment{definition}{enhanced,breakable,colback=blue!5!white,colframe=blue!75!black,fonttitle=\bfseries}
\tcolorboxenvironment{theorem}{enhanced,breakable,colback=red!5!white,colframe=red!75!black,fonttitle=\bfseries}
\tcolorboxenvironment{proposition}{enhanced,breakable,colback=red!5!white,colframe=red!60!black,fonttitle=\bfseries}
\tcolorboxenvironment{lemma}{enhanced,breakable,colback=yellow!5!white,colframe=yellow!70!black,fonttitle=\bfseries}
\tcolorboxenvironment{corollary}{enhanced,breakable,colback=orange!5!white,colframe=orange!70!black,fonttitle=\bfseries}
\tcolorboxenvironment{example}{enhanced,breakable,colback=green!5!white,colframe=green!60!black,fonttitle=\bfseries}
\tcolorboxenvironment{exercise}{enhanced,breakable,colback=purple!5!white,colframe=purple!60!black,fonttitle=\bfseries}
\tcolorboxenvironment{remark}{enhanced,breakable,colback=gray!5!white,colframe=gray!50!black,fonttitle=\bfseries}

%% =============================================================
\begin{document}

\begin{titlepage}
\begin{tikzpicture}[remember picture, overlay]
  \fill[left color=blue!15!white, right color=blue!5!white]
    (current page.south west) rectangle (current page.north east);
  \fill[color=orange!70!yellow]
    ([yshift=-2cm]current page.north west) rectangle ([yshift=-2.3cm]current page.north east);
  \fill[color=blue!80!black, rounded corners=8pt]
    ([xshift=3cm, yshift=-4cm]current page.north west) rectangle
    ([xshift=-3cm, yshift=-10cm]current page.north east);
  \node[anchor=center, text=white, font=\Huge\bfseries]
    at ([yshift=-6cm]current page.north) {Functional Analysis};
  \node[anchor=center, text=orange!80!yellow, font=\Large]
    at ([yshift=-7.5cm]current page.north) {Notes de Cours / Lecture Notes};
  \node[anchor=center, text=white!80, font=\large]
    at ([yshift=-8.8cm]current page.north) {Graduate Level --- 2025--2026};
  \node[anchor=center, text=blue!80!black, font=\Large\itshape]
    at ([yshift=-12cm]current page.north) {Ya\'e Ulrich Gaba};
  \draw[orange!70!yellow, line width=2pt]
    ([xshift=5cm, yshift=-13.5cm]current page.north west) --
    ([xshift=-5cm, yshift=-13.5cm]current page.north east);
  \node[anchor=center, text width=12cm, align=center, text=blue!60!black, font=\itshape]
    at ([yshift=-16cm]current page.north) {From Banach spaces to operator algebras: a journey through modern functional analysis.};
  \node[anchor=south, text=gray, font=\small]
    at ([yshift=1.5cm]current page.south) {\today};
\end{tikzpicture}
\end{titlepage}
\dominitoc
\tableofcontents

%% =============================================================
%%  CHAPTER 1 : BANACH SPACES — FOUNDATIONS
%% =============================================================
\chapter{Banach Spaces --- Foundations}
\label{ch:banach}
\minitoc\bigskip

\noindent
Functional analysis extends the ideas of linear algebra and real analysis to
infinite-dimensional spaces.  The natural framework is that of a
\emph{normed vector space}\index{normed vector space}, and the single most
important property such a space can enjoy is \emph{completeness}.  A complete
normed space is called a \emph{Banach space}\index{Banach space}, after
Stefan Banach\index{Banach, Stefan}, who laid the foundations of the subject
in his 1932 monograph.

Throughout this chapter $\K$ denotes either $\R$ or $\C$, and all vector
spaces are over~$\K$ unless stated otherwise.

%% ─────────────────────────────────────────────
\section{Normed vector spaces}
\label{sec:normed-spaces}
%% ─────────────────────────────────────────────

\begin{definition}[Normed vector space]
\label{def:norm}
Let $E$ be a vector space over $\K$.\index{norm}  A \textbf{norm} on $E$ is a
function $\norm{\cdot} : E \to [0,+\infty)$ satisfying, for all
$x,y\in E$ and $\lambda\in\K$:
\begin{enumerate}[label=\textup{(N\arabic*)}]
  \item \label{it:norm-pos} $\norm{x} = 0 \iff x = 0$
        \hfill(positive definiteness)\index{positive definiteness};
  \item \label{it:norm-hom} $\norm{\lambda x} = \abs{\lambda}\,\norm{x}$
        \hfill(absolute homogeneity)\index{homogeneity};
  \item \label{it:norm-tri} $\norm{x+y} \le \norm{x}+\norm{y}$
        \hfill(triangle inequality)\index{triangle inequality}.
\end{enumerate}
The pair $(E,\norm{\cdot})$ is called a \textbf{normed vector space}
(or simply a \textbf{normed space}).
\end{definition}

Every norm induces a metric\index{metric} $d(x,y)=\norm{x-y}$ and hence a
topology on~$E$.  We write $B(x,r)=\{y\in E:\norm{y-x}<r\}$ for the open
ball\index{open ball} and $\overline{B}(x,r)=\{y\in E:\norm{y-x}\le r\}$
for the closed ball\index{closed ball}.  The \textbf{unit ball}\index{unit ball}
is $B_E = \overline{B}(0,1)$.

\begin{definition}[Seminorm]
\label{def:seminorm}
A function $p : E \to [0,+\infty)$ satisfying \ref{it:norm-hom} and
\ref{it:norm-tri} (but not necessarily \ref{it:norm-pos}) is called a
\textbf{seminorm}\index{seminorm}.
\end{definition}

\begin{remark}
\label{rem:seminorm-kernel}
If $p$ is a seminorm, then $N = \{x\in E : p(x)=0\}$ is a vector subspace,
and $p$ induces a norm on the quotient $E/N$.
\end{remark}

%% ─────────────────────────────────────────────
\subsection{Fundamental examples}
\label{ssec:norm-examples}

\begin{example}[Euclidean space $\R^n$]
\label{ex:Rn-norms}
\index{Euclidean space}\index{Rn@$\R^n$}
On $\R^n$ (or $\C^n$) the most common norms are, for
$x=(x_1,\dots,x_n)$:
\begin{align}
  \norm{x}_p &= \Bigl(\sum_{k=1}^n \abs{x_k}^p\Bigr)^{1/p},
               \qquad 1\le p<\infty, \label{eq:lp-finite}\\
  \norm{x}_\infty &= \max_{1\le k\le n} \abs{x_k}. \label{eq:linf-finite}
\end{align}
The case $p=2$ gives the standard Euclidean norm\index{Euclidean norm}.
That $\norm{\cdot}_p$ is indeed a norm for $p\ge 1$ follows from
\textbf{Minkowski's inequality}\index{Minkowski's inequality} (see

ef{prop:minkowski} below).
\end{example}

\begin{example}[Sequence spaces $\ell^p$]
\label{ex:ell-p}
\index{ell-p@$\ell^p$}\index{sequence space}
For $1 \le p < \infty$ define
\[
  \ell^p = \Bigl\{ x = (x_n)_{n\ge 1} \subset \K :
           \sum_{n=1}^{\infty}\abs{x_n}^p < \infty \Bigr\},
  \qquad
  \norm{x}_{\ell^p} = \Bigl(\sum_{n=1}^{\infty}\abs{x_n}^p\Bigr)^{1/p}.
\]
For $p=\infty$:
\[
  \ell^\infty = \bigl\{ x = (x_n)_{n\ge 1} \subset \K :
                \sup_n \abs{x_n}<\infty \bigr\},
  \qquad
  \norm{x}_{\ell^\infty} = \sup_{n\ge 1}\abs{x_n}.
\]
The subspace $c_0\subset\ell^\infty$\index{c0@$c_0$} consists of sequences
converging to zero; $c\subset\ell^\infty$\index{c@$c$} consists of all
convergent sequences.  Both are closed subspaces of~$\ell^\infty$.
\end{example}

\begin{example}[Space of continuous functions $C(K)$]
\label{ex:CK}
\index{C(K)@$C(K)$}\index{continuous functions, space of}
Let $K$ be a compact topological space\index{compact space}.  The space
\[
  C(K) = \{ f : K \to \K \mid f \text{ is continuous} \}
\]
equipped with the \textbf{supremum norm}\index{supremum norm}
$\norm{f}_\infty = \sup_{t\in K}\abs{f(t)}$ is a normed space.
For $K=[a,b]\subset\R$ we often write $C([a,b])$\index{C([a,b])@$C([a,b])$}.
\end{example}

\begin{example}[Preview: $L^p$ spaces]
\label{ex:Lp-preview}
\index{Lp@$L^p$}
Let $(\Omega,\mathcal{A},\mu)$ be a measure space\index{measure space}.
For $1\le p<\infty$ one defines
\[
  L^p(\Omega,\mu) = \Bigl\{ f : \Omega\to\K \text{ measurable} :
                     \int_\Omega \abs{f}^p\,\dd\mu < \infty \Bigr\}\Big/{\sim},
\]
where $f\sim g$ iff $f=g$ $\mu$-a.e.\index{almost everywhere}  The norm is
$\norm{f}_{L^p}=\bigl(\int\abs{f}^p\,\dd\mu\bigr)^{1/p}$.  These spaces
will be discussed in detail later; for now they serve as motivation.
\end{example}

%% ─────────────────────────────────────────────
\subsection{H\"older's and Minkowski's inequalities}
\label{ssec:holder-minkowski}

The fundamental tool for proving that $\norm{\cdot}_p$ is a norm is the
following pair of classical inequalities.

\begin{proposition}[{H\"older's inequality}]
\label{prop:holder}
\index{Holder's inequality@H\"older's inequality}
Let $1\le p\le\infty$ and $\frac{1}{p}+\frac{1}{q}=1$
(with $q=\infty$ when $p=1$).  For any sequences
$(a_n),(b_n)\in\K^{\N}$:
\[
  \sum_{n=1}^{\infty} \abs{a_n b_n}
  \le \Bigl(\sum_{n=1}^{\infty}\abs{a_n}^p\Bigr)^{1/p}
     \Bigl(\sum_{n=1}^{\infty}\abs{b_n}^q\Bigr)^{1/q}.
\]
In particular, if $a\in\ell^p$ and $b\in\ell^q$ then $ab\in\ell^1$.
\end{proposition}

\begin{proof}
The cases $p=1$ or $p=\infty$ are immediate.  Assume $1<p<\infty$.
Recall Young's inequality\index{Young's inequality}: for $\alpha,\beta\ge 0$,
\[
  \alpha\beta \le \frac{\alpha^p}{p} + \frac{\beta^q}{q}.
\]
If $A=\norm{a}_p=0$ or $B=\norm{b}_q=0$ the result is trivial.
Otherwise set $\alpha_n = \abs{a_n}/A$ and $\beta_n = \abs{b_n}/B$.  Then
\[
  \frac{\abs{a_n b_n}}{AB}
  = \alpha_n\beta_n
  \le \frac{\alpha_n^p}{p} + \frac{\beta_n^q}{q}.
\]
Summing over $n$:
\[
  \frac{1}{AB}\sum_n \abs{a_n b_n}
  \le \frac{1}{p}\sum_n\frac{\abs{a_n}^p}{A^p}
    + \frac{1}{q}\sum_n\frac{\abs{b_n}^q}{B^q}
  = \frac{1}{p}+\frac{1}{q} = 1. \qedhere
\]
\end{proof}

\begin{proposition}[{Minkowski's inequality}]
\label{prop:minkowski}
\index{Minkowski's inequality}
For $1\le p\le\infty$ and $a,b\in\ell^p$:
\[
  \norm{a+b}_p \le \norm{a}_p + \norm{b}_p.
\]
\end{proposition}

\begin{proof}
For $p=1$ and $p=\infty$ this is clear.  Let $1<p<\infty$.  We have
\begin{align*}
  \sum_n \abs{a_n+b_n}^p
  &\le \sum_n \abs{a_n+b_n}^{p-1}\bigl(\abs{a_n}+\abs{b_n}\bigr)\\
  &= \sum_n \abs{a_n+b_n}^{p-1}\abs{a_n}
    +\sum_n \abs{a_n+b_n}^{p-1}\abs{b_n}.
\end{align*}
Apply H\"older's inequality to each sum with exponents $q$ and $p$
(noting $(p-1)q=p$):
\[
  \sum_n\abs{a_n+b_n}^p
  \le \Bigl(\sum_n\abs{a_n+b_n}^p\Bigr)^{1/q}
      \bigl(\norm{a}_p+\norm{b}_p\bigr).
\]
Dividing both sides by $\bigl(\sum\abs{a_n+b_n}^p\bigr)^{1/q}$ gives the
result (the case when the sum is zero being trivial).
\end{proof}

%% ─────────────────────────────────────────────
\section{Convergence and series in normed spaces}
\label{sec:convergence-series}
%% ─────────────────────────────────────────────

\begin{definition}[Convergence]
\label{def:convergence}
\index{convergence!in normed spaces}
A sequence $(x_n)$ in a normed space $(E,\norm{\cdot})$ \textbf{converges}
to $x\in E$ if $\norm{x_n - x}\to 0$.  We write $x_n\to x$ or
$\lim_{n\to\infty}x_n = x$.
\end{definition}

\begin{definition}[Cauchy sequence]
\label{def:cauchy}
\index{Cauchy sequence}
A sequence $(x_n)$ in~$E$ is a \textbf{Cauchy sequence} if for every
$\varepsilon>0$ there exists $N\in\N$ such that $\norm{x_m-x_n}<\varepsilon$
for all $m,n\ge N$.
\end{definition}

Every convergent sequence is Cauchy.  The converse is the content of
completeness.

\begin{definition}[Series]
\label{def:series}
\index{series!in normed spaces}
Let $(x_n)_{n\ge 1}$ be a sequence in $E$.  The \textbf{series}
$\sum_{n=1}^\infty x_n$ is said to \textbf{converge} if the sequence of
partial sums $S_N = \sum_{n=1}^N x_n$ converges in $E$.  The series
\textbf{converges absolutely}\index{absolute convergence} if
$\sum_{n=1}^\infty\norm{x_n} < \infty$.
\end{definition}

\begin{remark}
In general, absolute convergence does \emph{not} imply convergence.
In fact, the converse characterises Banach spaces (see

ef{thm:abs-conv-characterization} below).
\end{remark}

%% ─────────────────────────────────────────────
\section{Banach spaces: definition and completeness}
\label{sec:banach-def}
%% ─────────────────────────────────────────────

\begin{definition}[Banach space]
\label{def:banach}
\index{Banach space!definition}
A normed vector space $(E,\norm{\cdot})$ is called a \textbf{Banach space}
if it is \emph{complete}\index{completeness}, i.e., every Cauchy sequence
in~$E$ converges in~$E$.
\end{definition}

\begin{theorem}[Completeness of $\ell^p$]
\label{thm:ell-p-complete}
\index{ell-p@$\ell^p$!completeness}
For every $1\le p\le\infty$, the space $\ell^p$ is a Banach space.
\end{theorem}

\begin{proof}
We give the proof for $1\le p<\infty$; the case $p=\infty$ is analogous
(and simpler).

Let $(x^{(k)})_{k\ge 1}$ be a Cauchy sequence in $\ell^p$, where
$x^{(k)}=\bigl(x_n^{(k)}\bigr)_{n\ge 1}$.  For each fixed~$n$,
\[
  \abs{x_n^{(k)}-x_n^{(j)}}
  \le \norm{x^{(k)}-x^{(j)}}_{\ell^p} \to 0
  \qquad\text{as }k,j\to\infty,
\]
so $(x_n^{(k)})_{k\ge 1}$ is Cauchy in $\K$, hence converges to some
$x_n\in\K$.  Set $x=(x_n)_{n\ge 1}$.

\textbf{Step 1:} $x\in\ell^p$.  Let $\varepsilon>0$ and choose $K$ such
that $\norm{x^{(k)}-x^{(j)}}_p<\varepsilon$ for all $k,j\ge K$.  For any
$N\in\N$:
\[
  \sum_{n=1}^N \abs{x_n^{(k)}-x_n^{(j)}}^p < \varepsilon^p
  \quad\text{for all }k,j\ge K.
\]
Letting $j\to\infty$ (with $N$ fixed):
$\sum_{n=1}^N\abs{x_n^{(k)}-x_n}^p \le \varepsilon^p$.
Since this holds for all~$N$,
\[
  \norm{x^{(k)}-x}_p \le \varepsilon
  \qquad\text{for all }k\ge K.
\]
In particular $x = x^{(K)}+(x-x^{(K)}) \in \ell^p$.

\textbf{Step 2:} $x^{(k)}\to x$ in $\ell^p$.  The inequality above
shows $\norm{x^{(k)}-x}_p\le\varepsilon$ for $k\ge K$, which is exactly
$x^{(k)}\to x$.
\end{proof}

\begin{theorem}[{Completeness of $C(K)$}]
\label{thm:CK-complete}
\index{C(K)@$C(K)$!completeness}
If $K$ is a compact topological space, then $(C(K),\norm{\cdot}_\infty)$
is a Banach space.
\end{theorem}

\begin{proof}
Let $(f_n)$ be a Cauchy sequence in $C(K)$.  For each $t\in K$,
$\abs{f_m(t)-f_n(t)}\le\norm{f_m-f_n}_\infty\to 0$, so $(f_n(t))$ converges
in $\K$.  Define $f(t)=\lim_{n\to\infty}f_n(t)$.

Given $\varepsilon>0$, choose $N$ so that $\norm{f_m-f_n}_\infty<\varepsilon$
for $m,n\ge N$.  Letting $m\to\infty$ we get
$\abs{f(t)-f_n(t)}\le\varepsilon$ for all $t\in K$ and $n\ge N$, i.e.,
the convergence is \emph{uniform}\index{uniform convergence}.  As a uniform
limit of continuous functions, $f$ is continuous, so $f\in C(K)$.  Moreover
$\norm{f-f_n}_\infty\le\varepsilon$ for $n\ge N$.
\end{proof}

\begin{remark}
\label{rem:incomplete-example}
\index{incomplete normed space}
The space $C([0,1])$ equipped with the $L^1$-norm
$\norm{f}_1=\int_0^1\abs{f(t)}\,\dd t$ is \emph{not} complete.  One can
construct a Cauchy sequence of continuous functions whose $L^1$-limit is a
discontinuous function (e.g., a step function).
\end{remark}

%% ─────────────────────────────────────────────
\section{Finite-dimensional normed spaces}
\label{sec:finite-dim}
%% ─────────────────────────────────────────────

Finite-dimensional spaces are the ``tame'' case: they enjoy several
remarkable properties that fail spectacularly in infinite dimensions.

\begin{theorem}[Equivalence of norms in finite dimension]
\label{thm:norms-equiv}
\index{equivalence of norms}\index{finite-dimensional space!equivalence of norms}
On a finite-dimensional vector space $E$, all norms are equivalent:
if $\norm{\cdot}_a$ and $\norm{\cdot}_b$ are two norms on~$E$ with
$\dim E = n < \infty$, there exist constants $0<\alpha\le\beta$ such that
\[
  \alpha\,\norm{x}_a \le \norm{x}_b \le \beta\,\norm{x}_a
  \qquad\forall\, x\in E.
\]
\end{theorem}

\begin{proof}
It suffices to show that every norm $\norm{\cdot}$ on $E$ is equivalent
to the $\ell^1$-norm relative to a basis.  Fix a basis $e_1,\dots,e_n$
and write $x=\sum_{i=1}^n \xi_i e_i$; set
$\norm{x}_1=\sum_{i=1}^n\abs{\xi_i}$.

\textbf{Upper bound.}  By the triangle inequality and homogeneity,
\[
  \norm{x} \le \sum_{i=1}^n\abs{\xi_i}\,\norm{e_i}
  \le M\norm{x}_1,
  \qquad M = \max_i \norm{e_i}.
\]

\textbf{Lower bound.}  The map $\varphi:(\K^n,\norm{\cdot}_1)\to\R$
defined by $\varphi(\xi)=\norm{\sum_i\xi_i e_i}$ is continuous (by the
upper bound just proved) and positive on the unit sphere
$S=\{x\in\K^n:\norm{x}_1=1\}$.  Since $S$ is compact in $\K^n$
(here we use finite dimension), $\varphi$ attains a minimum
$m=\min_{S}\varphi > 0$.  Hence for $x\ne 0$:
\[
  \norm{x}
  = \norm{x}_1\cdot\varphi\!\bigl(x/\norm{x}_1\bigr)
  \ge m\,\norm{x}_1.
\]
Thus $m\,\norm{x}_1\le\norm{x}\le M\,\norm{x}_1$.
\end{proof}

\begin{corollary}
\label{cor:finite-dim-banach}
\index{finite-dimensional space!completeness}
Every finite-dimensional normed space is a Banach space.
\end{corollary}

\begin{proof}
By 
ef{thm:norms-equiv}, any norm on $E\cong\K^n$ is equivalent to the
Euclidean norm.  Since $(\K^n,\norm{\cdot}_2)$ is complete, so is
$(E,\norm{\cdot})$.
\end{proof}

\begin{theorem}[{Compactness of the closed unit ball in finite dimension}]
\label{thm:unit-ball-compact}
\index{unit ball!compactness}\index{finite-dimensional space!compact unit ball}
A normed space $E$ has the property that $B_E=\overline{B}(0,1)$ is
compact if and only if $\dim E<\infty$.
\end{theorem}

\begin{proof}
$(\Rightarrow)$  If $\dim E<\infty$, then by equivalence of norms $B_E$
is homeomorphic to the closed unit ball in $(\K^n,\norm{\cdot}_2)$, which
is compact by the Heine--Borel theorem\index{Heine--Borel theorem}.

$(\Leftarrow)$  Suppose $\dim E=\infty$.  We construct a sequence in $B_E$
with no convergent subsequence, using Riesz's lemma (
ef{lem:riesz}
below).  Start with $x_1\in E$, $\norm{x_1}=1$.  Set $E_1=\spn\{x_1\}$.
By Riesz's lemma with $\theta=1/2$, there exists $x_2$ with
$\norm{x_2}=1$ and $\norm{x_2-y}\ge 1/2$ for all $y\in E_1$.  Set
$E_2=\spn\{x_1,x_2\}$ and repeat.  By induction we obtain
$(x_n)\subset B_E$ with $\norm{x_m-x_n}\ge 1/2$ for $m\ne n$, so no
subsequence can be Cauchy.
\end{proof}

\begin{lemma}[Riesz's lemma]
\label{lem:riesz}
\index{Riesz's lemma}
Let $E$ be a normed space, $F\subsetneq E$ a \emph{closed} proper subspace,
and $0<\theta<1$.  Then there exists $x_\theta\in E$ with
$\norm{x_\theta}=1$ and
\[
  \dist(x_\theta,F) = \inf_{y\in F}\norm{x_\theta-y} \ge \theta.
\]
\end{lemma}

\begin{proof}
Pick $z\in E\setminus F$.  Since $F$ is closed,
$d = \dist(z,F) > 0$.  Choose $y_0\in F$ with
$\norm{z-y_0}\le d/\theta$ (possible since $d/\theta > d$).
Set
\[
  x_\theta = \frac{z-y_0}{\norm{z-y_0}}.
\]
Then $\norm{x_\theta}=1$, and for any $y\in F$:
\[
  \norm{x_\theta - y}
  = \frac{\norm{z - y_0 - \norm{z-y_0}\,y}}{\norm{z-y_0}}
  = \frac{\norm{z - \underbrace{(y_0+\norm{z-y_0}\,y)}_{\in F}}}{\norm{z-y_0}}
  \ge \frac{d}{\norm{z-y_0}}
  \ge \frac{d}{d/\theta} = \theta. \qedhere
\]
\end{proof}

\begin{remark}
One cannot take $\theta=1$ in general.  However, when $F$ has finite
codimension\index{codimension}, the infimum $d=\dist(z,F)$ is attained and
one can find $x$ with $\dist(x,F)=1$.
\end{remark}

%% ─────────────────────────────────────────────
\section{Continuous linear maps and the operator norm}
\label{sec:continuous-linear}
%% ─────────────────────────────────────────────

\begin{definition}[Bounded linear map]
\label{def:bounded-linear}
\index{bounded linear map}\index{linear map!bounded}
Let $(E,\norm{\cdot}_E)$ and $(F,\norm{\cdot}_F)$ be normed spaces.
A linear map $T:E\to F$ is \textbf{bounded} if there exists $C\ge 0$ such
that
\[
  \norm{Tx}_F \le C\,\norm{x}_E \qquad\forall\,x\in E.
\]
\end{definition}

\begin{proposition}[Continuity $=$ boundedness]
\label{prop:cont-iff-bounded}
\index{bounded linear map!equivalence with continuity}
For a linear map $T:E\to F$ between normed spaces, the following are
equivalent:
\begin{enumerate}[label=\textup{(\roman*)}]
  \item $T$ is continuous;
  \item $T$ is continuous at $0$;
  \item $T$ is bounded;
  \item $T$ is Lipschitz\index{Lipschitz map}.
\end{enumerate}
\end{proposition}

\begin{proof}
(i)$\Rightarrow$(ii) is trivial.

(ii)$\Rightarrow$(iii):  Continuity at $0$ gives $\delta>0$ with
$\norm{Tx}_F<1$ whenever $\norm{x}_E<\delta$.  For $x\ne 0$, set
$y = \frac{\delta}{2\norm{x}_E}\,x$; then $\norm{y}_E<\delta$, so
$\norm{Ty}_F<1$, i.e., $\norm{Tx}_F < \frac{2}{\delta}\norm{x}_E$.

(iii)$\Rightarrow$(iv):
$\norm{Tx-Ty}_F = \norm{T(x-y)}_F \le C\norm{x-y}_E$.

(iv)$\Rightarrow$(i) is clear.
\end{proof}

\begin{definition}[Operator norm]
\label{def:operator-norm}
\index{operator norm}
For $T\in\mathcal{L}(E,F)$ (the space of bounded linear maps $E\to F$), the
\textbf{operator norm} is
\[
  \norm{T}_{\mathcal{L}(E,F)}
  = \sup_{\substack{x\in E\\ x\ne 0}} \frac{\norm{Tx}_F}{\norm{x}_E}
  = \sup_{\norm{x}_E\le 1}\norm{Tx}_F
  = \sup_{\norm{x}_E= 1}\norm{Tx}_F.
\]
\end{definition}

\begin{notation}
We write $\mathcal{L}(E,F)$\index{L(E,F)@$\mathcal{L}(E,F)$} for the space
of bounded linear maps $E\to F$, and $E^* = \mathcal{L}(E,\K)$ for the
\textbf{(continuous) dual space}\index{dual space}\index{E*@$E^*$}.
When $E=F$ we write $\mathcal{L}(E)=\mathcal{L}(E,E)$.
\end{notation}

\begin{theorem}[$\mathcal{L}(E,F)$ is Banach when $F$ is]
\label{thm:LEF-banach}
\index{L(E,F)@$\mathcal{L}(E,F)$!completeness}
If $F$ is a Banach space then $\mathcal{L}(E,F)$ is a Banach space.
In particular the dual $E^*$ is always a Banach space.
\end{theorem}

\begin{proof}
Let $(T_n)$ be a Cauchy sequence in $\mathcal{L}(E,F)$.  For each $x\in E$,
\[
  \norm{T_m x - T_n x}_F \le \norm{T_m-T_n}\,\norm{x}_E \to 0,
\]
so $(T_n x)$ is Cauchy in $F$, hence converges to some element we call $Tx$.

\textbf{Linearity.}  For $\alpha,\beta\in\K$ and $x,y\in E$:
$T(\alpha x+\beta y) = \lim_n T_n(\alpha x+\beta y)
= \alpha\lim_n T_n x + \beta\lim_n T_n y = \alpha Tx+\beta Ty$.

\textbf{Boundedness.}  Since $(T_n)$ is Cauchy, it is bounded: say
$\norm{T_n}\le M$ for all $n$.  Then
$\norm{Tx}_F = \lim_n\norm{T_n x}_F \le M\norm{x}_E$, so
$T\in\mathcal{L}(E,F)$.

\textbf{Convergence.}  Given $\varepsilon>0$, pick $N$ such that
$\norm{T_m-T_n}<\varepsilon$ for $m,n\ge N$.  For any $x$ with
$\norm{x}\le 1$ and $n\ge N$:
\[
  \norm{T_m x - T_n x}_F \le \varepsilon.
\]
Letting $m\to\infty$: $\norm{Tx-T_n x}_F\le\varepsilon$.  Taking the sup
over $\norm{x}\le 1$: $\norm{T-T_n}\le\varepsilon$.
\end{proof}

\begin{corollary}
\label{cor:dual-banach}
\index{dual space!completeness}
For any normed space $E$, the dual space $E^*$ is a Banach space.
\end{corollary}

%% ─────────────────────────────────────────────
\section{Absolutely convergent series and completeness}
\label{sec:abs-conv}
%% ─────────────────────────────────────────────

The following elegant characterisation of Banach spaces reduces
completeness to a statement about series.

\begin{theorem}[Characterisation of completeness]
\label{thm:abs-conv-characterization}
\index{absolute convergence!characterization of completeness}
\index{Banach space!characterization via series}
A normed space $E$ is a Banach space if and only if every absolutely
convergent series in $E$ converges.
\end{theorem}

\begin{proof}
$(\Rightarrow)$  Suppose $E$ is Banach and $\sum\norm{x_n}<\infty$.
The partial sums $S_N=\sum_{n=1}^N x_n$ form a Cauchy sequence since
for $M>N$:
\[
  \norm{S_M-S_N} = \Bigl\|\sum_{n=N+1}^M x_n\Bigr\|
  \le \sum_{n=N+1}^M\norm{x_n} \to 0 \quad\text{as }N,M\to\infty.
\]
Hence $(S_N)$ converges.

$(\Leftarrow)$  Let $(y_n)$ be a Cauchy sequence.  Choose a subsequence
$(y_{n_k})$ with $\norm{y_{n_{k+1}}-y_{n_k}}<2^{-k}$.  Set
$x_k = y_{n_{k+1}}-y_{n_k}$.  Then $\sum\norm{x_k}<\sum 2^{-k}=1<\infty$,
so by hypothesis $\sum x_k$ converges, i.e., the telescoping partial sums
$y_{n_{k+1}}-y_{n_1}$ converge.  Thus $(y_{n_k})$ converges, and since
$(y_n)$ is Cauchy with a convergent subsequence, $(y_n)$ itself converges.
\end{proof}

%% ─────────────────────────────────────────────
\section{Quotient spaces}
\label{sec:quotient}
%% ─────────────────────────────────────────────

\begin{definition}[Quotient norm]
\label{def:quotient-norm}
\index{quotient space}\index{quotient norm}
Let $E$ be a normed space and $F\subset E$ a \emph{closed} subspace.
On the quotient vector space $E/F$ we define
\[
  \norm{\dot{x}}_{E/F} = \inf_{y\in F}\norm{x-y}_E
  = \dist(x,F),
  \qquad \dot{x} = x+F.
\]
\end{definition}

\begin{proposition}
\label{prop:quotient-norm-is-norm}
The quotient norm is indeed a norm on $E/F$ (closedness of $F$ is essential
for positive definiteness).  Moreover, if $E$ is a Banach space, so is $E/F$.
\end{proposition}

\begin{proof}
\textbf{Norm axioms.}
\begin{itemize}
  \item \emph{Positive definiteness:}
    $\norm{\dot x}=0$ means $\dist(x,F)=0$, so $x\in\overline{F}=F$
    (since $F$ is closed), i.e., $\dot x = \dot 0$.
  \item \emph{Homogeneity:}
    $\norm{\lambda\dot x}=\inf_{y\in F}\norm{\lambda x-y}
    =\abs{\lambda}\inf_{y\in F}\norm{x-y/\lambda}
    =\abs{\lambda}\,\norm{\dot x}$ for $\lambda\ne 0$.
  \item \emph{Triangle inequality:}
    For any $y_1,y_2\in F$:
    $\norm{(x_1+x_2)-(y_1+y_2)}\le\norm{x_1-y_1}+\norm{x_2-y_2}$.
    Taking infima over $y_1,y_2\in F$ gives
    $\norm{\dot x_1+\dot x_2}\le\norm{\dot x_1}+\norm{\dot x_2}$.
\end{itemize}

\textbf{Completeness.}  We use 
ef{thm:abs-conv-characterization}.
Let $\sum_n\norm{\dot x_n}_{E/F}<\infty$.  For each $n$, pick a
representative $z_n\in\dot x_n$ with $\norm{z_n}_E<\norm{\dot x_n}_{E/F}+2^{-n}$.
Then $\sum\norm{z_n}_E<\infty$, so since $E$ is Banach,
$S=\sum_n z_n$ converges in $E$.  The partial sums of $\sum\dot x_n$
equal $\dot{S}_N$ (the class of $\sum_{k=1}^N z_k$), and
$\norm{\dot{S}_N-\dot{S}}\le\norm{S_N-S}_E\to 0$.
\end{proof}

The canonical projection\index{canonical projection}
$\pi:E\to E/F$, $\pi(x)=\dot x$ is a bounded linear map with
$\norm{\pi}=1$ (when $F\ne E$).

%% ─────────────────────────────────────────────
\section{Completion of a normed space}
\label{sec:completion}
%% ─────────────────────────────────────────────

\begin{theorem}[Completion]
\label{thm:completion}
\index{completion}\index{normed space!completion}
Let $(E,\norm{\cdot})$ be a normed space.  There exists a Banach space
$\widehat{E}$ and an isometric linear injection $\iota:E\hookrightarrow\widehat{E}$
such that $\iota(E)$ is dense in $\widehat{E}$.  The pair
$(\widehat{E},\iota)$ is unique up to isometric isomorphism.
\end{theorem}

\begin{proof}[Proof (sketch)]
\index{completion!construction}
Consider the set $\mathcal{C}$ of all Cauchy sequences in $E$.  Define an
equivalence relation $(x_n)\sim(y_n)$ iff $\norm{x_n-y_n}\to 0$.
Set $\widehat{E}=\mathcal{C}/{\sim}$.  Vector space operations are
defined componentwise, and the norm
$\norm{[(x_n)]}_{\widehat{E}}=\lim_n\norm{x_n}_E$ (which exists since
$(\norm{x_n})$ is Cauchy in $\R$) is well-defined.  The map
$\iota:E\to\widehat{E}$ sends $x$ to the class of the constant sequence
$(x,x,x,\dots)$.  One verifies that $\widehat{E}$ is complete and
$\iota(E)$ is dense; we leave the details as an exercise.

For uniqueness, if $(\widetilde{E},j)$ is another completion, the map
$j\circ\iota^{-1}:\iota(E)\to\widetilde{E}$ is an isometry between dense
subsets of Banach spaces; by the BLT theorem\index{BLT theorem}
(Bounded Linear Transformation extension), it extends to an isometric
isomorphism $\widehat{E}\to\widetilde{E}$.
\end{proof}

The following commutative diagram summarises the universal property of
the completion:
\[
\begin{tikzcd}
  E \arrow[r,"\iota",hook] \arrow[dr,"T"'] & \widehat{E} \arrow[d,dashed,"\widehat{T}"] \\
  & F
\end{tikzcd}
\]
Given any bounded linear map $T:E\to F$ into a Banach space $F$, there
exists a unique $\widehat{T}:\widehat{E}\to F$ with
$\widehat{T}\circ\iota=T$ and $\norm{\widehat{T}}=\norm{T}$.

%% ─────────────────────────────────────────────
\section{Exercises for Chapter 1}
\label{sec:ex-ch1}
%% ─────────────────────────────────────────────

\begin{exercise}[$\star$]
\label{exo:norm-continuous}
\index{norm!continuity}
Show that the norm $\norm{\cdot}:E\to\R$ is a (Lipschitz) continuous
function: $\bigl|\norm{x}-\norm{y}\bigr|\le\norm{x-y}$.
\end{exercise}

\begin{exercise}[$\star$]
\label{exo:closed-subspace-banach}
\index{closed subspace!Banach}
Let $E$ be a Banach space and $F\subset E$ a subspace.  Prove that $F$ is
a Banach space (with the induced norm) if and only if $F$ is closed in $E$.
\end{exercise}

\begin{exercise}[$\star$]
\label{exo:c0-closed}
Prove that $c_0$ is a closed subspace of $\ell^\infty$, hence a Banach
space.\index{c0@$c_0$!closedness}
\end{exercise}

\begin{exercise}[$\star\star$]
\label{exo:ell-p-separable}
\index{ell-p@$\ell^p$!separability}\index{separable space}
Show that $\ell^p$ is separable for $1\le p<\infty$, but
$\ell^\infty$ is not separable.
\end{exercise}

\begin{exercise}[$\star\star$]
\label{exo:operator-norm-properties}
\index{operator norm!properties}
Let $E,F,G$ be normed spaces, $T\in\mathcal{L}(E,F)$ and
$S\in\mathcal{L}(F,G)$.  Prove that $S\circ T\in\mathcal{L}(E,G)$ and
$\norm{S\circ T}\le\norm{S}\,\norm{T}$.
\end{exercise}

\begin{exercise}[$\star\star$]
\label{exo:neumann-series}
\index{Neumann series}\index{invertibility}
Let $E$ be a Banach space and $T\in\mathcal{L}(E)$ with $\norm{T}<1$.
Prove that $\Id-T$ is invertible in $\mathcal{L}(E)$ and
\[
  (\Id-T)^{-1} = \sum_{n=0}^{\infty}T^n,
\]
with $\norm{(\Id-T)^{-1}}\le\frac{1}{1-\norm{T}}$.
\textit{Hint:} use the characterisation of completeness via absolutely
convergent series.
\end{exercise}

\begin{exercise}[$\star\star$]
\label{exo:GL-open}
\index{invertible operators!openness}
Deduce from the previous exercise that the set $GL(E)$ of invertible
elements of $\mathcal{L}(E)$ is open, and that the map $T\mapsto T^{-1}$
is continuous on $GL(E)$.
\end{exercise}

\begin{exercise}[$\star\star$]
\label{exo:equivalent-norms-open-map}
Let $\norm{\cdot}_a$ and $\norm{\cdot}_b$ be two norms on a vector
space~$E$.  Show they are equivalent if and only if the identity map
$\Id:(E,\norm{\cdot}_a)\to(E,\norm{\cdot}_b)$ is a homeomorphism.
\end{exercise}

\begin{exercise}[$\star\star\star$]
\label{exo:ell-infty-dual}
\index{dual space!of $\ell^1$}
Prove that the dual of $\ell^1$ is isometrically isomorphic to
$\ell^\infty$.  \textit{Hint:} given $\varphi\in(\ell^1)^*$, define
$a_n=\varphi(e_n)$ and show $a\in\ell^\infty$ with
$\norm{a}_\infty=\norm{\varphi}$.
\end{exercise}

\begin{exercise}[$\star\star\star$]
\label{exo:quotient-completion}
\index{quotient space!and completion}
Let $E$ be a normed space and $F$ a closed subspace.  Prove that the
completion of $E/F$ is isometrically isomorphic to $\widehat{E}/\overline{F}$,
where $\overline{F}$ denotes the closure of $\iota(F)$ in $\widehat{E}$.
\end{exercise}

\begin{exercise}[$\star$]
\label{exo:parallelogram}
\index{parallelogram law}
Show that a norm $\norm{\cdot}$ on a vector space $E$ satisfies the
\textbf{parallelogram law}
\[
  \norm{x+y}^2 + \norm{x-y}^2 = 2\bigl(\norm{x}^2+\norm{y}^2\bigr)
\]
for all $x,y\in E$ if and only if it arises from an inner product via
$\norm{x}=\sqrt{\inner{x}{x}}$.  This is the
\textbf{Jordan--von Neumann theorem}\index{Jordan--von Neumann theorem}.
\end{exercise}

\begin{exercise}[$\star\star$]
\label{exo:banach-not-hilbert}
Show that $\ell^p$ for $p\ne 2$ does not satisfy the parallelogram law,
hence is not a Hilbert space.\index{ell-p@$\ell^p$!not Hilbert for $p\ne 2$}
\end{exercise}

\begin{exercise}[$\star\star$]
\label{exo:riesz-lemma-not-one}
\index{Riesz's lemma!$\theta=1$ impossible}
Give an example of a Banach space $E$ and a closed subspace $F\subsetneq E$
such that there is no $x$ with $\norm{x}=1$ and $\dist(x,F)=1$.
\textit{Hint:} consider $c_0$ inside $\ell^\infty$.
\end{exercise}

\begin{exercise}[$\star\star\star$]
\label{exo:completion-Lp}
\index{Lp@$L^p$!as completion}
Let $C_c(\R)$\index{Cc@$C_c(\R)$} be the space of continuous functions
with compact support, equipped with the $L^p$-norm.  Show that $C_c(\R)$
is not complete for any $1\le p<\infty$, and that its completion is
$L^p(\R)$.
\end{exercise}

%% =============================================================
%%  CHAPTER 2 : HILBERT SPACES
%% =============================================================
\chapter{Hilbert Spaces}
\label{ch:hilbert}
\minitoc\bigskip

\noindent
Hilbert spaces\index{Hilbert space} are Banach spaces whose norm arises from
an inner product.  This extra structure is enormously powerful: it gives rise
to orthogonal projections\index{orthogonal projection}, the Riesz
representation theorem\index{Riesz representation theorem}, and the theory of
orthonormal bases\index{orthonormal basis}, bringing the subject much closer
to the intuition of finite-dimensional linear algebra.  The archetypal
Hilbert space is $L^2$\index{L2@$L^2$}, and the theory developed here
underpins quantum mechanics\index{quantum mechanics}, signal
processing\index{signal processing}, and the modern theory of PDEs.

%% ─────────────────────────────────────────────
\section{Inner products and the Cauchy--Schwarz inequality}
\label{sec:inner-products}
%% ─────────────────────────────────────────────

\begin{definition}[Inner product]
\label{def:inner-product}
\index{inner product}
Let $H$ be a vector space over $\K$ ($=\R$ or $\C$).  An \textbf{inner
product} (or \textbf{scalar product}) is a map
$\inner{\cdot}{\cdot}:H\times H\to\K$ satisfying:
\begin{enumerate}[label=\textup{(IP\arabic*)}]
  \item \label{ip:linear}
    $\inner{\alpha x_1+\beta x_2}{y}
    =\alpha\inner{x_1}{y}+\beta\inner{x_2}{y}$
    for all $\alpha,\beta\in\K$, $x_1,x_2,y\in H$
    \hfill(linearity in the first variable\index{linearity});
  \item \label{ip:conj}
    $\inner{x}{y}=\overline{\inner{y}{x}}$
    \hfill(conjugate symmetry\index{conjugate symmetry});
  \item \label{ip:pos}
    $\inner{x}{x}\ge 0$, with equality iff $x=0$
    \hfill(positive definiteness).
\end{enumerate}
The pair $(H,\inner{\cdot}{\cdot})$ is called a
\textbf{pre-Hilbert space}\index{pre-Hilbert space}
(or \textbf{inner product space}\index{inner product space}).
\end{definition}

\begin{remark}
\label{rem:convention}
We adopt the convention that the inner product is linear in the
\emph{first} argument (the ``physicist's convention''\index{physicist's convention}
is the opposite).  Note that \ref{ip:linear} and \ref{ip:conj} together
give conjugate-linearity in the second argument:
$\inner{x}{\alpha y}=\bar\alpha\inner{x}{y}$.
\end{remark}

\begin{notation}
Every inner product induces a norm\index{norm!from inner product}:
\[
  \norm{x} = \sqrt{\inner{x}{x}}.
\]
We write $x\perp y$\index{orthogonality} (``$x$ is orthogonal to $y$'')
when $\inner{x}{y}=0$.
\end{notation}

\begin{theorem}[Cauchy--Schwarz inequality]
\label{thm:cauchy-schwarz}
\index{Cauchy--Schwarz inequality}
For all $x,y$ in a pre-Hilbert space $H$:
\[
  \abs{\inner{x}{y}} \le \norm{x}\,\norm{y},
\]
with equality if and only if $x$ and $y$ are linearly dependent.
\end{theorem}

\begin{proof}
If $y=0$ both sides vanish.  Assume $y\ne 0$ and let $\lambda\in\K$.
Then
\[
  0 \le \norm{x-\lambda y}^2
  = \norm{x}^2 - \lambda\inner{y}{x} - \bar\lambda\inner{x}{y}
    + \abs{\lambda}^2\norm{y}^2.
\]
Choose $\lambda = \inner{x}{y}/\norm{y}^2$.  Then
\[
  0 \le \norm{x}^2
  - \frac{\abs{\inner{x}{y}}^2}{\norm{y}^2}
  - \frac{\abs{\inner{x}{y}}^2}{\norm{y}^2}
  + \frac{\abs{\inner{x}{y}}^2}{\norm{y}^4}\norm{y}^2
  = \norm{x}^2 - \frac{\abs{\inner{x}{y}}^2}{\norm{y}^2}.
\]
Rearranging gives $\abs{\inner{x}{y}}^2\le\norm{x}^2\norm{y}^2$.
Equality holds iff $\norm{x-\lambda y}=0$, i.e., $x=\lambda y$.
\end{proof}

\begin{corollary}
The function $\norm{x}=\sqrt{\inner{x}{x}}$ is indeed a norm.
\end{corollary}

\begin{proof}
Only the triangle inequality needs verification:
\begin{align*}
  \norm{x+y}^2
  &= \norm{x}^2 + 2\operatorname{Re}\inner{x}{y} + \norm{y}^2 \\
  &\le \norm{x}^2 + 2\norm{x}\,\norm{y} + \norm{y}^2
  = \bigl(\norm{x}+\norm{y}\bigr)^2. \qedhere
\end{align*}
\end{proof}

\begin{proposition}[Polarization identity]
\label{prop:polarization}
\index{polarization identity}
The inner product can be recovered from the norm.  Over $\R$:
\[
  \inner{x}{y} = \frac{1}{4}\bigl(\norm{x+y}^2 - \norm{x-y}^2\bigr).
\]
Over $\C$:
\[
  \inner{x}{y} = \frac{1}{4}\sum_{k=0}^{3}i^k\norm{x+i^k y}^2.
\]
\end{proposition}

\begin{proof}
Direct expansion using $\norm{x+y}^2=\norm{x}^2+2\operatorname{Re}\inner{x}{y}+\norm{y}^2$.
For the complex case, the four terms $\norm{x+y}^2$, $i\norm{x+iy}^2$,
$-\norm{x-y}^2$, $-i\norm{x-iy}^2$ yield $4\inner{x}{y}$ when summed.
\end{proof}

\begin{proposition}[Parallelogram law]
\label{prop:parallelogram}
\index{parallelogram law}
In any inner product space:
\[
  \norm{x+y}^2 + \norm{x-y}^2 = 2\bigl(\norm{x}^2 + \norm{y}^2\bigr).
\]
Conversely, a norm satisfying this identity comes from an inner product
(Jordan--von Neumann theorem\index{Jordan--von Neumann theorem}).
\end{proposition}

\begin{proof}
Expanding both sides:
$\norm{x+y}^2+\norm{x-y}^2 = (\norm{x}^2+2\operatorname{Re}\inner{x}{y}+\norm{y}^2)
+(\norm{x}^2-2\operatorname{Re}\inner{x}{y}+\norm{y}^2) = 2\norm{x}^2+2\norm{y}^2$.
\end{proof}

%% ─────────────────────────────────────────────
\section{Hilbert spaces: definition and examples}
\label{sec:hilbert-def}
%% ─────────────────────────────────────────────

\begin{definition}[Hilbert space]
\label{def:hilbert}
\index{Hilbert space!definition}
A \textbf{Hilbert space} is a pre-Hilbert space that is complete with
respect to the norm induced by its inner product.
\end{definition}

\begin{example}[$\K^n$]
\label{ex:Kn-hilbert}
\index{Kn@$\K^n$!Hilbert space}
The space $\K^n$ with $\inner{x}{y}=\sum_{k=1}^n x_k\bar{y}_k$ is a
Hilbert space of dimension~$n$.
\end{example}

\begin{example}[$\ell^2$]
\label{ex:ell2}
\index{ell-2@$\ell^2$}
The space $\ell^2$ with inner product
$\inner{x}{y}=\sum_{n=1}^\infty x_n\bar{y}_n$ is a Hilbert space.
The series converges absolutely by the Cauchy--Schwarz inequality for
sequences.  Completeness was proved in 
ef{thm:ell-p-complete}.
\end{example}

\begin{example}[$L^2(\Omega,\mu)$]
\label{ex:L2}
\index{L2@$L^2$!Hilbert space}
For a measure space $(\Omega,\mathcal{A},\mu)$, the space
$L^2(\Omega,\mu)$ with inner product
\[
  \inner{f}{g} = \int_\Omega f\,\bar{g}\,\dd\mu
\]
is a Hilbert space (the Fischer--Riesz theorem,

ef{thm:fischer-riesz}).
\end{example}

\begin{example}[Sobolev space $H^1(\Omega)$ --- preview]
\label{ex:sobolev-preview}
\index{Sobolev space!$H^1$}
Let $\Omega\subset\R^n$ be open.  The Sobolev space
\[
  H^1(\Omega) = \bigl\{f\in L^2(\Omega) :
  \partial_i f\in L^2(\Omega) \text{ for }i=1,\dots,n\bigr\}
\]
(derivatives in the distributional sense) is a Hilbert space with inner
product
\[
  \inner{f}{g}_{H^1}
  = \int_\Omega f\bar{g}\,\dd x
    + \sum_{i=1}^n\int_\Omega (\partial_i f)\overline{(\partial_i g)}\,\dd x.
\]
We will return to Sobolev spaces in later chapters.
\end{example}

%% ─────────────────────────────────────────────
\section{Projection onto closed convex sets}
\label{sec:projection-convex}
%% ─────────────────────────────────────────────

The following theorem is the cornerstone of Hilbert space geometry.

\begin{theorem}[Projection onto a closed convex set]
\label{thm:projection-convex}
\index{projection!onto closed convex set}\index{closed convex set}
Let $H$ be a Hilbert space and $C\subset H$ a non-empty, closed, convex
set.  For every $x\in H$ there exists a unique $p\in C$ such that
\[
  \norm{x-p} = \dist(x,C) := \inf_{y\in C}\norm{x-y}.
\]
Moreover, $p$ is characterised by:
\begin{equation}\label{eq:projection-char}
  p\in C \quad\text{and}\quad
  \operatorname{Re}\inner{x-p}{y-p}\le 0 \quad\forall\,y\in C.
\end{equation}
\end{theorem}

\begin{proof}
Set $d=\dist(x,C)$ and choose a minimising sequence $(y_n)\subset C$
with $\norm{x-y_n}\to d$.

\textbf{Existence (via the parallelogram law).}
For $m,n\in\N$, since $C$ is convex,
$\frac{y_m+y_n}{2}\in C$, so
$\norm{x-\tfrac{y_m+y_n}{2}}\ge d$.  The parallelogram law gives:
\begin{align*}
  \norm{y_m-y_n}^2
  &= \norm{(y_m-x)-(y_n-x)}^2 \\
  &= 2\norm{y_m-x}^2+2\norm{y_n-x}^2-\norm{(y_m-x)+(y_n-x)}^2 \\
  &= 2\norm{y_m-x}^2+2\norm{y_n-x}^2-4\Bigl\|\frac{y_m+y_n}{2}-x\Bigr\|^2 \\
  &\le 2\norm{y_m-x}^2+2\norm{y_n-x}^2 - 4d^2.
\end{align*}
As $m,n\to\infty$, the right side tends to $2d^2+2d^2-4d^2=0$.
Hence $(y_n)$ is Cauchy, and since $H$ is complete and $C$ is closed,
$y_n\to p\in C$ with $\norm{x-p}=d$.

\textbf{Uniqueness.}  If $p,p'$ both achieve the minimum, apply the
parallelogram law to $y_m=p$, $y_n=p'$ to get $\norm{p-p'}^2\le 0$.

\textbf{Characterisation.}  ($\Rightarrow$) Let $y\in C$ and $t\in(0,1]$.
Since $C$ is convex, $p+t(y-p)\in C$, so
\[
  d^2 \le \norm{x-p-t(y-p)}^2
  = d^2 - 2t\operatorname{Re}\inner{x-p}{y-p}+t^2\norm{y-p}^2.
\]
Hence $2\operatorname{Re}\inner{x-p}{y-p}\le t\norm{y-p}^2$.  Letting
$t\to 0^+$ gives $\operatorname{Re}\inner{x-p}{y-p}\le 0$.

($\Leftarrow$)  If $p$ satisfies \eqref{eq:projection-char}, then for
any $y\in C$:
\[
  \norm{x-y}^2 = \norm{(x-p)-(y-p)}^2
  = \norm{x-p}^2 - 2\operatorname{Re}\inner{x-p}{y-p}+\norm{y-p}^2
  \ge \norm{x-p}^2. \qedhere
\]
\end{proof}

We denote the projection of $x$ onto $C$ by $P_C(x)$\index{PC@$P_C$}.

\begin{proposition}[Projection is a contraction]
\label{prop:projection-contraction}
\index{projection!contraction}
The map $P_C:H\to C$ is a contraction:
$\norm{P_C(x)-P_C(y)}\le\norm{x-y}$ for all $x,y\in H$.
\end{proposition}

\begin{proof}
Set $p=P_C(x)$, $q=P_C(y)$.  By the characterisation:
$\operatorname{Re}\inner{x-p}{q-p}\le 0$ and
$\operatorname{Re}\inner{y-q}{p-q}\le 0$.  Adding:
\[
  \operatorname{Re}\inner{(x-p)-(y-q)}{q-p}\le 0,
\]
i.e., $\operatorname{Re}\inner{(x-y)-(p-q)}{q-p}\le 0$, hence
$\operatorname{Re}\inner{x-y}{q-p}\le\operatorname{Re}\inner{p-q}{q-p}=-\norm{p-q}^2$.
Therefore
$\norm{p-q}^2\le\operatorname{Re}\inner{x-y}{p-q}\le\norm{x-y}\,\norm{p-q}$.
\end{proof}

%% ─────────────────────────────────────────────
\subsection{Projection onto a closed subspace}
\label{ssec:projection-subspace}

When $C=F$ is a \emph{closed subspace}\index{closed subspace}, the
characterisation simplifies considerably.

\begin{theorem}[Projection onto a closed subspace]
\label{thm:projection-subspace}
\index{projection!onto closed subspace}\index{orthogonal complement}
Let $F$ be a closed subspace of a Hilbert space $H$.  For every $x\in H$:
\begin{enumerate}[label=\textup{(\roman*)}]
  \item There exists a unique $p\in F$ nearest to $x$, characterised by
    \begin{equation}\label{eq:proj-subspace}
      p\in F \quad\text{and}\quad x-p\in F^\perp,
    \end{equation}
    where $F^\perp = \{z\in H:\inner{z}{y}=0\ \forall y\in F\}$ is the
    \textbf{orthogonal complement}\index{orthogonal complement!definition} of $F$.
  \item $H = F \oplus F^\perp$ (algebraic and topological direct sum).
  \item The projection $P_F:H\to F$ is a bounded linear operator with
    $\norm{P_F}=1$ (when $F\ne\{0\}$), $P_F^2=P_F$, and
    $\inner{P_Fx}{y}=\inner{x}{P_Fy}$ for all $x,y\in H$ (i.e., $P_F$
    is self-adjoint\index{self-adjoint operator}).
\end{enumerate}
\end{theorem}

\begin{proof}
(i) Since $F$ is a closed convex set, 
ef{thm:projection-convex}
gives a unique $p\in F$ minimising $\norm{x-p}$.  The characterisation
\eqref{eq:projection-char} gives
$\operatorname{Re}\inner{x-p}{y-p}\le 0$ for all $y\in F$.
Since $F$ is a subspace, for any $z\in F$ and $t\in\R$ we may take
$y=p+tz$ to get
$t\operatorname{Re}\inner{x-p}{z}\le 0$ for all $t\in\R$, forcing
$\operatorname{Re}\inner{x-p}{z}=0$.
Replacing $z$ by $iz$ (in the complex case) gives
$\operatorname{Im}\inner{x-p}{z}=0$ too, so $\inner{x-p}{z}=0$
for all $z\in F$, i.e., $x-p\in F^\perp$.

(ii)  Every $x\in H$ decomposes as $x = P_F x + (x-P_F x)$ with
$P_F x\in F$ and $x-P_F x\in F^\perp$.  If $z\in F\cap F^\perp$ then
$\inner{z}{z}=0$, so $z=0$.  Hence the sum is direct.

(iii)  Linearity: let $x_1,x_2\in H$ and $\alpha,\beta\in\K$.
Set $p_i=P_Fx_i$.  Then $\alpha p_1+\beta p_2\in F$ and
$(\alpha x_1+\beta x_2)-(\alpha p_1+\beta p_2)=\alpha(x_1-p_1)+\beta(x_2-p_2)\in F^\perp$,
so $P_F(\alpha x_1+\beta x_2)=\alpha p_1+\beta p_2$.

Boundedness: $\norm{P_Fx}^2=\inner{P_Fx}{P_Fx}=\inner{P_Fx}{x}\le\norm{P_Fx}\,\norm{x}$,
giving $\norm{P_Fx}\le\norm{x}$.  For $x\in F$, $P_Fx=x$, so $\norm{P_F}=1$.

Idempotence: $P_F^2=P_F$ since $P_Fx\in F$ implies $P_F(P_Fx)=P_Fx$.

Self-adjointness:
$\inner{P_Fx}{y}=\inner{P_Fx}{P_Fy+(y-P_Fy)}=\inner{P_Fx}{P_Fy}$
(since $P_Fx\in F$ and $y-P_Fy\in F^\perp$).  By symmetry,
$\inner{x}{P_Fy}=\inner{P_Fx}{P_Fy}$.
\end{proof}

\begin{corollary}
\label{cor:perp-perp}
\index{orthogonal complement!double}
For any closed subspace $F$ of $H$: $(F^\perp)^\perp = F$.
\end{corollary}

\begin{proof}
$F\subset(F^\perp)^\perp$ is clear.  If $x\in(F^\perp)^\perp$, decompose
$x=p+q$ with $p\in F$, $q\in F^\perp$.  Then
$q=x-p\in(F^\perp)^\perp$ (since both $x$ and $p$ lie there), so
$q\in F^\perp\cap(F^\perp)^\perp$, giving $q=0$ and $x=p\in F$.
\end{proof}

\begin{corollary}
\label{cor:dense-perp}
A subspace $F$ is dense in $H$ if and only if $F^\perp=\{0\}$.
\end{corollary}

\begin{proof}
$\overline{F}$ is a closed subspace and $\overline{F}^\perp=F^\perp$.
If $F^\perp=\{0\}$, then $\overline{F}=(\overline{F}^\perp)^\perp=\{0\}^\perp=H$.
Conversely, if $\overline{F}=H$ and $z\in F^\perp$, then $\inner{z}{y}=0$
for all $y\in F$; by continuity, $\inner{z}{y}=0$ for all $y\in H$,
forcing $z=0$.
\end{proof}

%% ─────────────────────────────────────────────
\section{The Riesz representation theorem}
\label{sec:riesz-representation}
%% ─────────────────────────────────────────────

\begin{theorem}[{Riesz--Fr\'echet representation theorem}]
\label{thm:riesz-frechet}
\index{Riesz representation theorem}\index{Riesz--Fr\'echet theorem}
Let $H$ be a Hilbert space.  For every continuous linear functional
$\varphi\in H^*$ there exists a unique $u\in H$ such that
\[
  \varphi(x) = \inner{x}{u} \qquad\forall\, x\in H.
\]
Moreover, $\norm{\varphi}_{H^*}=\norm{u}_H$.  The map
$J:H\to H^*$ defined by $Ju = \inner{\cdot}{u}$ is a
conjugate-linear\index{conjugate-linear map} isometric isomorphism
(an isometric isomorphism in the real case).
\end{theorem}

\begin{proof}
\textbf{Existence.}
If $\varphi=0$, take $u=0$.  Assume $\varphi\ne 0$.
Set $F=\Ker\varphi$\index{kernel}; this is a closed hyperplane
(i.e., $\codim F=1$).  Since $F\ne H$, we have $F^\perp\ne\{0\}$;
pick $v\in F^\perp\setminus\{0\}$.

For any $x\in H$, the element $x-\frac{\varphi(x)}{\varphi(v)}v$ lies in
$F$ (check: $\varphi$ applied to it gives $0$).  Hence
\[
  0 = \inner{x-\frac{\varphi(x)}{\varphi(v)}v}{v}
  = \inner{x}{v} - \frac{\varphi(x)}{\varphi(v)}\norm{v}^2.
\]
Solving: $\varphi(x) = \inner{x}{\frac{\overline{\varphi(v)}}{\norm{v}^2}\,v}$.
Set $u = \frac{\overline{\varphi(v)}}{\norm{v}^2}\,v$.

\textbf{Uniqueness.}  If $\inner{x}{u}=\inner{x}{u'}$ for all $x$,
then $\inner{x}{u-u'}=0$ for all $x$; taking $x=u-u'$ gives $u=u'$.

\textbf{Isometry.}  On one hand,
$\abs{\varphi(x)}=\abs{\inner{x}{u}}\le\norm{x}\,\norm{u}$, so
$\norm{\varphi}\le\norm{u}$.  On the other hand,
$\varphi(u)=\norm{u}^2$, so
$\norm{\varphi}\ge\abs{\varphi(u/\norm{u})}=\norm{u}$.

\textbf{Conjugate-linearity.}
$J(\alpha u_1+\beta u_2)(x)
=\inner{x}{\alpha u_1+\beta u_2}
=\bar\alpha\inner{x}{u_1}+\bar\beta\inner{x}{u_2}
=\bar\alpha\,Ju_1(x)+\bar\beta\,Ju_2(x)$.

\textbf{Surjectivity} was established in the existence part.
\end{proof}

\begin{remark}
\label{rem:riesz-self-dual}
\index{self-dual}
The Riesz theorem says that a Hilbert space is (conjugate-linearly)
isometrically isomorphic to its own dual: $H\cong H^*$.  This
\textbf{self-duality}\index{Hilbert space!self-duality} is a key
distinction from general Banach spaces.
\end{remark}

\begin{corollary}[{Lax--Milgram theorem --- a preview}]
\label{cor:lax-milgram-preview}
\index{Lax--Milgram theorem}
Let $a:H\times H\to\K$ be a continuous, coercive sesquilinear form on a
Hilbert space $H$.  Then for every $\varphi\in H^*$ there exists a unique
$u\in H$ with $a(u,v)=\varphi(v)$ for all $v\in H$.
\end{corollary}

This will be proved in full generality later; it is the foundation of
the variational approach to elliptic PDEs\index{elliptic PDE}.

%% ─────────────────────────────────────────────
\section{Orthonormal systems and Hilbert bases}
\label{sec:hilbert-bases}
%% ─────────────────────────────────────────────

\begin{definition}[Orthonormal system]
\label{def:ONS}
\index{orthonormal system}
A family $(e_i)_{i\in I}$ in a Hilbert space $H$ is called an
\textbf{orthonormal system} (ONS) if
\[
  \inner{e_i}{e_j} = \delta_{ij}
  \qquad\forall\, i,j\in I,
\]
where $\delta_{ij}$\index{Kronecker delta} is the Kronecker delta.
\end{definition}

\begin{definition}[Hilbert basis]
\label{def:hilbert-basis}
\index{Hilbert basis}\index{orthonormal basis}
An orthonormal system $(e_i)_{i\in I}$ is a \textbf{Hilbert basis}
(or \textbf{complete orthonormal system}) if it is \emph{maximal},
i.e., no proper orthonormal system contains it.  Equivalently,
$\spn\{e_i:i\in I\}$ is dense in $H$.
\end{definition}

\begin{definition}[Fourier coefficients]
\label{def:fourier-coeff}
\index{Fourier coefficients}
Given an orthonormal system $(e_i)_{i\in I}$ and $x\in H$, the scalars
$\hat{x}_i = \inner{x}{e_i}$ are called the \textbf{Fourier coefficients}
of $x$.
\end{definition}

\begin{lemma}[Best approximation]
\label{lem:best-approx}
\index{best approximation}
Let $e_1,\dots,e_n$ be orthonormal and $x\in H$.  Then
\[
  \min_{\alpha_1,\dots,\alpha_n\in\K}
  \Bigl\|x - \sum_{k=1}^n\alpha_k e_k\Bigr\|^2
  = \norm{x}^2 - \sum_{k=1}^n\abs{\inner{x}{e_k}}^2,
\]
and the minimum is attained at $\alpha_k=\inner{x}{e_k}$.
\end{lemma}

\begin{proof}
Set $s=\sum_{k=1}^n\alpha_k e_k$ and $\hat x_k=\inner{x}{e_k}$.  Then
\begin{align*}
  \norm{x-s}^2
  &= \norm{x}^2 - \sum_k\bar\alpha_k\inner{x}{e_k}
     - \sum_k\alpha_k\overline{\inner{x}{e_k}}
     + \sum_k\abs{\alpha_k}^2 \\
  &= \norm{x}^2 - \sum_k\abs{\hat x_k}^2
     + \sum_k\abs{\alpha_k-\hat x_k}^2.
\end{align*}
This is minimised when $\alpha_k=\hat x_k$, and the minimum value is
$\norm{x}^2-\sum_k\abs{\hat x_k}^2$.
\end{proof}

\begin{theorem}[Bessel's inequality]
\label{thm:bessel}
\index{Bessel's inequality}
Let $(e_i)_{i\in I}$ be an orthonormal system in $H$.  For every $x\in H$:
\[
  \sum_{i\in I}\abs{\inner{x}{e_i}}^2 \le \norm{x}^2.
\]
In particular, at most countably many Fourier coefficients are nonzero.
\end{theorem}

\begin{proof}
For any finite subset $J\subset I$, 
ef{lem:best-approx} gives
$\sum_{i\in J}\abs{\inner{x}{e_i}}^2\le\norm{x}^2$.  Taking the
supremum over all finite $J$ gives the result.  The set
$\{i\in I:\inner{x}{e_i}\ne 0\}=\bigcup_{n\ge 1}\{i:\abs{\inner{x}{e_i}}>1/n\}$;
each set in the union is finite (by Bessel), hence the union is countable.
\end{proof}

\begin{theorem}[{Parseval's theorem}]
\label{thm:parseval}
\index{Parseval's theorem}\index{Parseval's identity}
Let $(e_i)_{i\in I}$ be an orthonormal system in a Hilbert space $H$.
The following are equivalent:
\begin{enumerate}[label=\textup{(\roman*)}]
  \item $(e_i)_{i\in I}$ is a Hilbert basis;
  \item $\overline{\spn}\{e_i:i\in I\}=H$;
  \item For every $x\in H$:
    $x = \sum_{i\in I}\inner{x}{e_i}\,e_i$\quad
    (convergence in norm);
  \item \textup{(Parseval's identity)}
    $\norm{x}^2 = \sum_{i\in I}\abs{\inner{x}{e_i}}^2$
    for every $x\in H$;
  \item \textup{(Generalised Parseval)}
    $\inner{x}{y}=\sum_{i\in I}\inner{x}{e_i}\overline{\inner{y}{e_i}}$
    for all $x,y\in H$;
  \item $\inner{x}{e_i}=0$ for all $i\in I$ implies $x=0$.
\end{enumerate}
\end{theorem}

\begin{proof}
We prove the cycle
(i)$\Rightarrow$(vi)$\Rightarrow$(ii)$\Rightarrow$(iii)$\Rightarrow$(iv)$\Rightarrow$(v)$\Rightarrow$(ii),
plus (ii)$\Leftrightarrow$(i).

\textbf{(i)$\Rightarrow$(vi):}  If $\inner{x}{e_i}=0$ for all $i$ and
$x\ne 0$, then $e_i\cup\{x/\norm{x}\}$ would be a strictly larger ONS,
contradicting maximality.

\textbf{(vi)$\Rightarrow$(ii):}  If $F=\overline{\spn}\{e_i\}\ne H$,
pick $x\in F^\perp\setminus\{0\}$.  Then $\inner{x}{e_i}=0$ for all $i$,
contradicting (vi).

\textbf{(ii)$\Rightarrow$(iii):}
Let $x\in H$ and $\varepsilon>0$.  By (ii), there exist finitely many
indices $i_1,\dots,i_n$ and scalars $\alpha_k$ with
$\norm{x-\sum_k\alpha_k e_{i_k}}<\varepsilon$.  By the best approximation
lemma, replacing $\alpha_k$ by $\inner{x}{e_{i_k}}$ can only decrease the
error.  Since Bessel's inequality guarantees
$\sum_{i\in I}\abs{\inner{x}{e_i}}^2<\infty$, we can arrange the nonzero
Fourier coefficients in a sequence and the partial sums converge to $x$.

\textbf{(iii)$\Rightarrow$(iv):}
$\norm{x}^2 = \inner{x}{x}
= \inner{\sum_i\inner{x}{e_i}e_i}{\sum_j\inner{x}{e_j}e_j}
= \sum_i\abs{\inner{x}{e_i}}^2$
(the last step uses continuity of the inner product).

\textbf{(iv)$\Rightarrow$(v):}
Apply the polarization identity\index{polarization identity} to
$\norm{x+y}^2$, $\norm{x-y}^2$, etc.

\textbf{(v)$\Rightarrow$(ii):}
If $x\perp\overline{\spn}\{e_i\}$, then $\inner{x}{e_i}=0$ for all $i$,
so (v) gives $\inner{x}{x}=0$, i.e., $x=0$.  By 
ef{cor:dense-perp},
$\overline{\spn}\{e_i\}=H$.

\textbf{(ii)$\Leftrightarrow$(i):}
$\overline{\spn}\{e_i\}=H$ iff $\{e_i\}^\perp=\{0\}$ iff $(e_i)$ is
maximal (any vector orthogonal to all $e_i$ must be zero, so no ONS
properly contains $(e_i)$).
\end{proof}

%% ─────────────────────────────────────────────
\section{Separable spaces and countable bases}
\label{sec:separable}
%% ─────────────────────────────────────────────

\begin{theorem}[Existence of Hilbert bases]
\label{thm:ONB-exists}
\index{Hilbert basis!existence}
Every Hilbert space admits a Hilbert basis.
\end{theorem}

\begin{proof}
Apply Zorn's lemma\index{Zorn's lemma} to the partially ordered set
(by inclusion) of orthonormal systems.  Every chain has an upper bound
(its union), so there exists a maximal element.
\end{proof}

\begin{theorem}
\label{thm:separable-countable-basis}
\index{separable Hilbert space}\index{Hilbert basis!countable}
A Hilbert space $H$ is separable if and only if it admits a countable
(or finite) Hilbert basis.
\end{theorem}

\begin{proof}
$(\Leftarrow)$  If $(e_n)_{n\ge 1}$ is a countable Hilbert basis, the
set $D=\bigl\{\sum_{k=1}^N q_k e_k : N\in\N,\, q_k\in\Q+i\Q\bigr\}$
is countable and dense.

$(\Rightarrow)$  If $H$ is separable, every ONS is countable (distinct
elements are distance $\sqrt{2}$ apart, so each lives in a distinct
member of a countable open cover).  By 
ef{thm:ONB-exists}, a Hilbert
basis exists, and it must be countable.
\end{proof}

\begin{theorem}[Isomorphism theorem]
\label{thm:hilbert-isomorphism}
\index{Hilbert space!isomorphism theorem}
Every separable infinite-dimensional Hilbert space is isometrically
isomorphic to $\ell^2$.
\end{theorem}

\begin{proof}
Let $(e_n)_{n\ge 1}$ be a Hilbert basis.  Define
$U:H\to\ell^2$ by $Ux=\bigl(\inner{x}{e_n}\bigr)_{n\ge 1}$.
By Parseval, $\norm{Ux}_{\ell^2}=\norm{x}_H$, so $U$ is an isometry
(hence injective).  For surjectivity, given
$(\alpha_n)\in\ell^2$, the series $\sum_n\alpha_n e_n$ converges in $H$
(partial sums are Cauchy by Bessel/Parseval) to some $x$ with $Ux=(\alpha_n)$.
\end{proof}

%% ─────────────────────────────────────────────
\section{The Fischer--Riesz theorem}
\label{sec:fischer-riesz}
%% ─────────────────────────────────────────────

\begin{theorem}[{Fischer--Riesz}]
\label{thm:fischer-riesz}
\index{Fischer--Riesz theorem}\index{L2@$L^2$!completeness}
For any measure space $(\Omega,\mathcal{A},\mu)$, the space
$L^2(\Omega,\mu)$ is a Hilbert space.
\end{theorem}

\begin{proof}
We need to show completeness.  Let $(f_n)$ be a Cauchy sequence in $L^2$.
Choose a subsequence $(f_{n_k})$ with
$\norm{f_{n_{k+1}}-f_{n_k}}_2 < 2^{-k}$.  Define
\[
  g_N = \sum_{k=1}^N \abs{f_{n_{k+1}}-f_{n_k}},\qquad
  g = \sum_{k=1}^{\infty}\abs{f_{n_{k+1}}-f_{n_k}}.
\]
By Minkowski's inequality in $L^2$:
\[
  \norm{g_N}_2 \le \sum_{k=1}^N 2^{-k} < 1.
\]
By monotone convergence\index{monotone convergence theorem},
$\norm{g}_2\le 1$, so $g<\infty$ a.e.  This means the telescoping series
\[
  f_{n_1} + \sum_{k=1}^{\infty}(f_{n_{k+1}}-f_{n_k})
\]
converges absolutely a.e.\ to a function $f$.  We have $f_{n_k}\to f$
a.e.  Moreover
$\abs{f_{n_k}-f_{n_1}}^2\le g^2\in L^1$, so by dominated
convergence\index{dominated convergence theorem},
$\norm{f-f_{n_k}}_2\to 0$.  Since $(f_n)$ is Cauchy and a subsequence
converges to $f$, the entire sequence converges to $f$ in $L^2$.
\end{proof}

\begin{remark}
The same argument (with $p$ replacing $2$) proves that $L^p$ is complete
for all $1\le p<\infty$.  The case $p=\infty$ requires a different (simpler)
argument.
\end{remark}

%% ─────────────────────────────────────────────
\section{Hilbert direct sum}
\label{sec:hilbert-sum}
%% ─────────────────────────────────────────────

\begin{definition}[Hilbert direct sum]
\label{def:hilbert-sum}
\index{Hilbert direct sum}\index{direct sum!Hilbert}
Let $(H_i)_{i\in I}$ be a family of Hilbert spaces.  The
\textbf{Hilbert direct sum} is
\[
  \bigoplus_{i\in I}^{2} H_i
  = \Bigl\{ (x_i)_{i\in I} : x_i\in H_i,\;
    \sum_{i\in I}\norm{x_i}_{H_i}^2 < \infty \Bigr\}
\]
with inner product $\inner{(x_i)}{(y_i)} = \sum_{i\in I}\inner{x_i}{y_i}_{H_i}$.
\end{definition}

\begin{proposition}
\label{prop:hilbert-sum-complete}
\index{Hilbert direct sum!completeness}
The Hilbert direct sum is a Hilbert space.
\end{proposition}

\begin{proof}
The inner product is well-defined by the Cauchy--Schwarz inequality
applied to $\ell^2(I)$.  Completeness follows from the completeness of
each $H_i$ and of $\ell^2(I)$: if $(x^{(n)})$ is Cauchy in
$\bigoplus^2 H_i$, then for each $i$, $(x_i^{(n)})$ is Cauchy in $H_i$,
hence converges to some $x_i$.  A standard $\varepsilon/3$ argument (as
in the proof of completeness of $\ell^p$) shows
$x=(x_i)\in\bigoplus^2 H_i$ and $x^{(n)}\to x$.
\end{proof}

\begin{example}
$\ell^2(I) = \bigoplus_{i\in I}^2 \K$ is the Hilbert direct sum of
copies of the scalar field.
\end{example}

%% ─────────────────────────────────────────────
\section{Operators on Hilbert spaces: the adjoint}
\label{sec:adjoint}
%% ─────────────────────────────────────────────

\begin{definition}[Adjoint operator]
\label{def:adjoint}
\index{adjoint operator}\index{operator!adjoint}
Let $H,K$ be Hilbert spaces and $T\in\mathcal{L}(H,K)$.  The
\textbf{adjoint} of $T$ is the unique operator $T^*\in\mathcal{L}(K,H)$
satisfying
\[
  \inner{Tx}{y}_K = \inner{x}{T^* y}_H
  \qquad\forall\, x\in H,\; y\in K.
\]
\end{definition}

\begin{proof}[Proof of existence and uniqueness]
For fixed $y\in K$, the map $x\mapsto\inner{Tx}{y}_K$ is a continuous
linear functional on $H$ (with norm $\le\norm{T}\,\norm{y}$).  By the
Riesz representation theorem (
ef{thm:riesz-frechet}), there exists a
unique $z\in H$ with $\inner{Tx}{y}=\inner{x}{z}$ for all $x$.  Define
$T^*y=z$.  Linearity and boundedness of $T^*$ follow from uniqueness
and the estimate $\norm{T^*y}\le\norm{T}\,\norm{y}$.
\end{proof}

\begin{proposition}[Properties of the adjoint]
\label{prop:adjoint-props}
\index{adjoint operator!properties}
Let $S,T\in\mathcal{L}(H,K)$ and $\alpha\in\K$.
\begin{enumerate}[label=\textup{(\roman*)}]
  \item $(S+T)^*=S^*+T^*$ and $(\alpha T)^*=\bar\alpha\, T^*$.
  \item $(T^*)^*=T$.
  \item $\norm{T^*}=\norm{T}$.
  \item $\norm{T^*T}=\norm{T}^2$.
  \item If $R\in\mathcal{L}(K,L)$, then $(RT)^*=T^*R^*$.
  \item $\Ker T^* = (\ran T)^\perp$ and $\overline{\ran T}=(\Ker T^*)^\perp$.
\end{enumerate}
\end{proposition}

\begin{proof}
We prove the most important parts.

(ii) For all $x\in H$, $y\in K$:
$\inner{x}{T^{**}y}=\inner{T^*x}{y}=\overline{\inner{y}{T^*x}}
=\overline{\inner{Ty}{x}}=\inner{x}{Ty}$, so $T^{**}=T$ (here we
identify $H$ with $H^{**}$ via the inner product).

Wait --- let us be more careful.  We have $T^*\in\mathcal{L}(K,H)$,
so $T^{**}=(T^*)^*\in\mathcal{L}(H,K)$.  For $x\in H$, $y\in K$:
$\inner{T^{**}x}{y}_K=\inner{x}{T^*y}_H=\inner{Tx}{y}_K$.
Since this holds for all $y$, $T^{**}x=Tx$.

(iii) $\norm{T^*}\le\norm{T}$ since
$\abs{\inner{x}{T^*y}}=\abs{\inner{Tx}{y}}\le\norm{T}\,\norm{x}\,\norm{y}$.
By (ii), $\norm{T}=\norm{T^{**}}\le\norm{T^*}$.

(iv) $\norm{T^*T}\le\norm{T^*}\,\norm{T}=\norm{T}^2$.  Conversely,
$\norm{Tx}^2=\inner{Tx}{Tx}=\inner{T^*Tx}{x}\le\norm{T^*Tx}\,\norm{x}
\le\norm{T^*T}\,\norm{x}^2$, so $\norm{T}^2\le\norm{T^*T}$.

(vi) $y\in\Ker T^*$ iff $T^*y=0$ iff $\inner{x}{T^*y}=0$ for all $x$
iff $\inner{Tx}{y}=0$ for all $x$ iff $y\in(\ran T)^\perp$.
By 
ef{cor:perp-perp}, $\overline{\ran T}=((\ran T)^\perp)^\perp=(\Ker T^*)^\perp$.
\end{proof}

\begin{definition}[Self-adjoint, normal, unitary]
\label{def:special-operators}
\index{self-adjoint operator}\index{normal operator}\index{unitary operator}
Let $T\in\mathcal{L}(H)$.
\begin{itemize}
  \item $T$ is \textbf{self-adjoint} (or \textbf{Hermitian}\index{Hermitian operator}) if $T^*=T$.
  \item $T$ is \textbf{normal} if $T^*T=TT^*$.
  \item $T$ is \textbf{unitary} if $T^*T=TT^*=\Id$ (equivalently, $T$
        is a surjective isometry\index{isometry}).
\end{itemize}
\end{definition}

\begin{proposition}
\label{prop:self-adjoint-real-spectrum}
\index{self-adjoint operator!real inner products}
If $T=T^*$, then $\inner{Tx}{x}\in\R$ for all $x\in H$.  Conversely,
if $\K=\C$ and $\inner{Tx}{x}\in\R$ for all $x$, then $T=T^*$.
\end{proposition}

\begin{proof}
If $T=T^*$: $\overline{\inner{Tx}{x}}=\inner{x}{Tx}=\inner{T^*x}{x}=\inner{Tx}{x}$,
so $\inner{Tx}{x}\in\R$.

Conversely, if $\inner{Tx}{x}\in\R$ for all $x$, then
$\inner{Tx}{x}=\overline{\inner{Tx}{x}}=\inner{x}{Tx}=\inner{T^*x}{x}$.
Hence $\inner{(T-T^*)x}{x}=0$ for all $x$.  Over $\C$, the polarization
identity implies $T-T^*=0$.
\end{proof}

\begin{remark}
\label{rem:real-case-warning}
\index{self-adjoint operator!real case}
Over $\R$, the condition $\inner{Tx}{x}=0$ for all $x$ does \emph{not}
imply $T=0$ (consider a $90^\circ$ rotation in $\R^2$).  The implication
relies on the full complex polarization identity.
\end{remark}

%% ─────────────────────────────────────────────
\section{Exercises for Chapter 2}
\label{sec:ex-ch2}
%% ─────────────────────────────────────────────

\begin{exercise}[$\star$]
\label{exo:inner-product-continuous}
\index{inner product!continuity}
Show that the inner product $\inner{\cdot}{\cdot}:H\times H\to\K$ is
(jointly) continuous.  Specifically, if $x_n\to x$ and $y_n\to y$,
then $\inner{x_n}{y_n}\to\inner{x}{y}$.
\end{exercise}

\begin{exercise}[$\star$]
\label{exo:orthogonal-pythagorean}
\index{Pythagorean theorem}
Prove the \textbf{Pythagorean theorem}: if $x\perp y$ then
$\norm{x+y}^2=\norm{x}^2+\norm{y}^2$.  Generalise to finite sums.
\end{exercise}

\begin{exercise}[$\star$]
\label{exo:closed-subspace-perp}
Show that if $S\subset H$ is any subset, then $S^\perp$ is a closed
subspace of $H$.
\end{exercise}

\begin{exercise}[$\star\star$]
\label{exo:projection-formula}
\index{projection!explicit formula}
Let $F=\spn\{e_1,\dots,e_n\}$ where $e_1,\dots,e_n$ are orthonormal.
Show that the orthogonal projection onto $F$ is given by
$P_F x = \sum_{k=1}^n\inner{x}{e_k}\,e_k$.
\end{exercise}

\begin{exercise}[$\star\star$]
\label{exo:gram-schmidt}
\index{Gram--Schmidt process}
Let $(v_n)_{n\ge 1}$ be a linearly independent sequence in a Hilbert
space $H$.  Describe the \textbf{Gram--Schmidt orthonormalization} process
and prove it produces an orthonormal system $(e_n)$ with
$\spn\{e_1,\dots,e_n\}=\spn\{v_1,\dots,v_n\}$ for each $n$.
\end{exercise}

\begin{exercise}[$\star\star$]
\label{exo:L2-basis-trig}
\index{trigonometric system}\index{Fourier series}
Show that the system $\bigl(\frac{1}{\sqrt{2\pi}}e^{int}\bigr)_{n\in\Z}$
is an orthonormal system in $L^2([-\pi,\pi])$.  Proving completeness
requires the Stone--Weierstrass theorem or Fej\'er's theorem; state this
as a fact and deduce Parseval's identity for Fourier series:
\[
  \frac{1}{2\pi}\int_{-\pi}^{\pi}\abs{f(t)}^2\,\dd t
  = \sum_{n\in\Z}\abs{\hat{f}(n)}^2.
\]
\end{exercise}

\begin{exercise}[$\star\star$]
\label{exo:riesz-application}
\index{Riesz representation theorem!application}
Let $H=L^2([0,1])$ and define $\varphi(f)=\int_0^1 t\,f(t)\,\dd t$.
Show that $\varphi\in H^*$ and find the element $u\in H$ given by the
Riesz representation theorem such that $\varphi(f)=\inner{f}{u}$.
Compute $\norm{\varphi}$.
\end{exercise}

\begin{exercise}[$\star\star$]
\label{exo:adjoint-matrix}
\index{adjoint operator!matrix representation}
Let $T:\K^n\to\K^m$ be the linear map represented by the matrix $A$.
Show that $T^*$ is represented by $A^*=\bar{A}^T$ (the conjugate
transpose).
\end{exercise}

\begin{exercise}[$\star\star$]
\label{exo:normal-characterization}
\index{normal operator!characterization}
Show that $T\in\mathcal{L}(H)$ is normal if and only if
$\norm{Tx}=\norm{T^*x}$ for all $x\in H$.
\end{exercise}

\begin{exercise}[$\star\star$]
\label{exo:unitary-preserves-inner}
\index{unitary operator!characterization}
Let $U\in\mathcal{L}(H)$.  Prove the following are equivalent:
(a)~$U$ is unitary;
(b)~$U$ is surjective and $\inner{Ux}{Uy}=\inner{x}{y}$ for all $x,y$;
(c)~$U$ is surjective and $\norm{Ux}=\norm{x}$ for all $x$.
\end{exercise}

\begin{exercise}[$\star\star\star$]
\label{exo:weak-convergence-hilbert}
\index{weak convergence!in Hilbert space}
A sequence $(x_n)$ in a Hilbert space $H$ \textbf{converges weakly} to
$x$, written $x_n\weakto x$, if $\inner{x_n}{y}\to\inner{x}{y}$ for all
$y\in H$.
\begin{enumerate}[label=(\alph*)]
  \item Show that strong convergence implies weak convergence but not
        conversely ($e_n\weakto 0$ in $\ell^2$ but $\norm{e_n}=1$).
  \item Prove that a weakly convergent sequence is bounded.
        \textit{Hint:} use the Riesz theorem to identify $H$ with $H^*$,
        then apply the uniform boundedness principle (accept it for now).
  \item Show that $x_n\weakto x$ and $\norm{x_n}\to\norm{x}$ together
        imply $x_n\to x$ (strong convergence).
\end{enumerate}
\end{exercise}

\begin{exercise}[$\star\star\star$]
\label{exo:hilbert-sum-universal}
\index{Hilbert direct sum!universal property}
Let $(H_i)_{i\in I}$ be a family of closed, mutually orthogonal subspaces
of a Hilbert space $H$ with $H=\overline{\bigoplus_{i\in I}H_i}$.  Show
that $H$ is isometrically isomorphic to $\bigoplus_{i\in I}^2 H_i$.
\end{exercise}

\begin{exercise}[$\star\star\star$]
\label{exo:compact-operator-preview}
\index{compact operator!preview}
Let $H$ be a separable Hilbert space with Hilbert basis $(e_n)$.
Define $T:H\to H$ by $Te_n = \frac{1}{n}e_n$.
\begin{enumerate}[label=(\alph*)]
  \item Show that $T$ is bounded, self-adjoint, and $\norm{T}=1$.
  \item Show that $T$ maps bounded sets to sets with compact closure.
        (Such an operator is called \textbf{compact}\index{compact operator}.)
  \item Show that $T$ is \emph{not} surjective.
  \item Find $\Ker T$ and $\overline{\ran T}$.
\end{enumerate}
\end{exercise}

\begin{exercise}[$\star$]
\label{exo:perp-complement-sum}
Let $F,G$ be closed subspaces of a Hilbert space.  Show that
$(F+G)^\perp = F^\perp\cap G^\perp$.  Is it always true that
$(F\cap G)^\perp = F^\perp + G^\perp$?
\end{exercise}

\begin{exercise}[$\star\star$]
\label{exo:nearest-point-strictly-convex}
\index{strictly convex norm}
A normed space is \textbf{strictly convex} if
$\norm{x}=\norm{y}=1$ and $x\ne y$ imply $\norm{x+y}<2$.  Show that
Hilbert spaces are strictly convex, and deduce uniqueness of nearest
points in closed convex sets without using the parallelogram law directly.
\end{exercise}

\begin{exercise}[$\star\star\star$]
\label{exo:lax-milgram}
\index{Lax--Milgram theorem!proof}
Prove the Lax--Milgram theorem (
ef{cor:lax-milgram-preview}):
if $a:H\times H\to\K$ is a sesquilinear form satisfying
$\abs{a(u,v)}\le M\norm{u}\,\norm{v}$ (continuity) and
$\operatorname{Re}\, a(u,u)\ge\alpha\norm{u}^2$ (coercivity) for some
$\alpha>0$, then for every $\varphi\in H^*$ there exists a unique $u\in H$
with $a(u,v)=\varphi(v)$ for all $v\in H$, and
$\norm{u}\le\frac{1}{\alpha}\norm{\varphi}$.

\textit{Hint:} By Riesz, write $a(u,v)=\inner{Au}{v}$ for some
$A\in\mathcal{L}(H)$.  Show $A$ is injective with closed range, then
surjective.
\end{exercise}
%%%%%%%%%%%%%%%%%%%%%%%%%%%%%%%%%%%%%%%%%%%%%%%%%%%%%%%%%%%%
%%  CHAPTER 3 — Bounded Linear Operators
%%%%%%%%%%%%%%%%%%%%%%%%%%%%%%%%%%%%%%%%%%%%%%%%%%%%%%%%%%%%
\chapter{Bounded Linear Operators}\label{ch:bounded-operators}
\minitoc

\vspace{1em}

The study of bounded linear operators between Banach spaces lies at the heart of
functional analysis. In this chapter we develop the algebraic and topological
theory of the space $\mathcal{L}(E,F)$, introduce compact operators, and build
the spectral theory of bounded operators. The chapter culminates with the rich
additional structure available when the underlying spaces are Hilbert spaces:
the adjoint, and the classes of normal, self-adjoint, and unitary operators.

%-------------------------------------------------------------------
\section{The space \texorpdfstring{$\mathcal{L}(E,F)$}{L(E,F)} revisited}
\index{bounded linear operator}\index{operator!bounded linear}
%-------------------------------------------------------------------

Throughout this section, $(E,\norm{\cdot}_E)$ and $(F,\norm{\cdot}_F)$ denote
normed spaces over the field $\K\in\{\R,\C\}$.

\begin{definition}[Bounded linear operator]\label{def:bounded-op}
\index{operator!bounded}\index{bounded linear operator!definition}
A linear map $T\colon E\to F$ is \emph{bounded} if there exists $C\geq 0$ such
that
\[
  \norm{Tx}_F \leq C\,\norm{x}_E \qquad \text{for all } x\in E.
\]
The infimum of all such constants is the \emph{operator norm}\index{operator norm}
\[
  \norm{T} \;=\; \sup_{\norm{x}_E\leq 1}\norm{Tx}_F
           \;=\; \sup_{\norm{x}_E = 1}\norm{Tx}_F
           \;=\; \sup_{x\neq 0}\frac{\norm{Tx}_F}{\norm{x}_E}.
\]
We denote by $\mathcal{L}(E,F)$\index{L(E,F)@$\mathcal{L}(E,F)$} the vector
space of all bounded linear operators from $E$ to $F$, equipped with the
operator norm. When $E=F$ we write $\mathcal{L}(E)=\mathcal{L}(E,E)$\index{L(E)@$\mathcal{L}(E)$}.
\end{definition}

\begin{proposition}[Completeness of $\mathcal{L}(E,F)$]
\label{prop:LEF-complete}
\index{L(E,F)@$\mathcal{L}(E,F)$!completeness}
If $F$ is a Banach space, then $\mathcal{L}(E,F)$ is a Banach space.
\end{proposition}

\begin{proof}
Let $(T_n)$ be a Cauchy sequence in $\mathcal{L}(E,F)$. For each $x\in E$,
\[
  \norm{T_n x - T_m x}_F \leq \norm{T_n - T_m}\,\norm{x}_E \to 0,
\]
so $(T_n x)$ is Cauchy in $F$, hence convergent. Define $Tx = \lim_{n\to\infty}T_n x$.
Linearity of $T$ is clear. Since $(T_n)$ is Cauchy, it is bounded:
$\sup_n\norm{T_n}=M<\infty$. Then $\norm{Tx}_F = \lim_n\norm{T_n x}_F \leq M\norm{x}_E$,
so $T\in\mathcal{L}(E,F)$. Finally, given $\varepsilon>0$, choose $N$ so that
$\norm{T_n-T_m}<\varepsilon$ for $n,m\geq N$. Letting $m\to\infty$,
$\norm{T_n - T}\leq\varepsilon$ for $n\geq N$.
\end{proof}

\begin{example}[Multiplication operators]\label{ex:mult-op}
\index{multiplication operator}\index{operator!multiplication}
Let $\varphi\in L^\infty(\mu)$ for a measure space $(\Omega,\Sigma,\mu)$.
The operator $M_\varphi\colon L^p(\mu)\to L^p(\mu)$ defined by
$M_\varphi f = \varphi f$ satisfies $\norm{M_\varphi} = \norm{\varphi}_\infty$.
\end{example}

\begin{example}[Integral operators]\label{ex:integral-op}
\index{integral operator}\index{operator!integral}
Let $K\in L^2([0,1]^2)$. The operator $T_K\colon L^2([0,1])\to L^2([0,1])$
defined by
\[
  (T_K f)(x) = \int_0^1 K(x,y)\,f(y)\,\dd y
\]
satisfies $\norm{T_K}\leq \norm{K}_{L^2([0,1]^2)}$ by the Cauchy--Schwarz
inequality\index{Cauchy--Schwarz inequality}.
\end{example}

\begin{proposition}[Composition]\label{prop:composition-ops}
\index{operator!composition}\index{composition of operators}
Let $E,F,G$ be normed spaces and $T\in\mathcal{L}(E,F)$, $S\in\mathcal{L}(F,G)$.
Then $ST\in\mathcal{L}(E,G)$ and $\norm{ST}\leq\norm{S}\,\norm{T}$.
\end{proposition}

\begin{proof}
For any $x\in E$,
$\norm{STx}_G \leq \norm{S}\,\norm{Tx}_F \leq \norm{S}\,\norm{T}\,\norm{x}_E$.
\end{proof}

\begin{corollary}\label{cor:LE-algebra}
\index{Banach algebra}\index{operator algebra}
If $E$ is a Banach space, then $\mathcal{L}(E)$ is a unital Banach algebra with
product given by composition and unit~$\Id_E$.
\end{corollary}

%-------------------------------------------------------------------
\section{Finite-rank and compact operators}
\index{finite-rank operator}\index{compact operator}
%-------------------------------------------------------------------

\begin{definition}[Finite-rank operator]\label{def:finite-rank}
\index{finite-rank operator!definition}\index{operator!finite-rank}
An operator $T\in\mathcal{L}(E,F)$ is of \emph{finite rank} if $\dim(\ran T)<\infty$.
We write $\mathcal{F}(E,F)$\index{F(E,F)@$\mathcal{F}(E,F)$} for the set of
finite-rank operators, which is a linear subspace of $\mathcal{L}(E,F)$.
The \emph{rank}\index{rank!of an operator} of $T$ is $\dim(\ran T)$.
\end{definition}

\begin{example}\label{ex:rank-one}
\index{rank-one operator}
Let $f\in E'$ and $y\in F$. The operator $T = y\otimes f$\index{tensor product!rank-one operator}
defined by $Tx = f(x)\,y$ has rank at most one and $\norm{T}=\norm{f}\,\norm{y}$.
Every finite-rank operator is a finite sum of such rank-one operators.
\end{example}

\begin{definition}[Compact operator]\label{def:compact-op}
\index{compact operator!definition}\index{operator!compact}
An operator $T\in\mathcal{L}(E,F)$ is \emph{compact} if $T(B_E)$ is relatively
compact in $F$, where $B_E=\{x\in E:\norm{x}\leq 1\}$. Equivalently, for every
bounded sequence $(x_n)$ in $E$, the sequence $(Tx_n)$ has a convergent
subsequence in $F$. We write $\mathcal{K}(E,F)$\index{K(E,F)@$\mathcal{K}(E,F)$}
for the set of compact operators.
\end{definition}

\begin{proposition}[Basic properties of compact operators]
\label{prop:compact-basic}
\index{compact operator!properties}
Let $E,F,G$ be Banach spaces.
\begin{enumerate}[label=(\roman*)]
  \item Every finite-rank operator is compact.
  \item $\mathcal{K}(E,F)$ is a closed linear subspace of $\mathcal{L}(E,F)$.
  \item If $T\in\mathcal{K}(E,F)$ and $S\in\mathcal{L}(F,G)$, then
        $ST\in\mathcal{K}(E,G)$.
  \item If $T\in\mathcal{L}(E,F)$ and $S\in\mathcal{K}(F,G)$, then
        $ST\in\mathcal{K}(E,G)$.
\end{enumerate}
Thus $\mathcal{K}(E,F)$ is a closed two-sided ideal in $\mathcal{L}(E)$ when
$E=F=G$.
\end{proposition}

\begin{proof}
\emph{(i)} If $T$ has finite rank, then $T(B_E)$ is a bounded subset of a
finite-dimensional subspace, hence relatively compact.

\emph{(ii)} \textsc{Linearity.} If $T_1,T_2\in\mathcal{K}(E,F)$ and
$\alpha\in\K$, take any bounded sequence $(x_n)$. Extract a subsequence
$(x_{n_k})$ so that $(T_1 x_{n_k})$ converges, then a further subsequence so
that $(T_2 x_{n_{k_l}})$ converges. Then $((T_1+\alpha T_2)x_{n_{k_l}})$
converges.

\textsc{Closedness.}\index{compact operator!closedness}
Let $T_n\to T$ in $\mathcal{L}(E,F)$ with each $T_n$ compact. Let $(x_k)$ be
a bounded sequence in $E$ with $\norm{x_k}\leq M$. By a diagonal argument, we
extract a subsequence $(x_{k_j})$ such that $(T_n x_{k_j})_j$ converges for
every $n$. Given $\varepsilon>0$, choose $n_0$ with
$\norm{T-T_{n_0}}<\varepsilon/(3M)$. Choose $J$ so that
$\norm{T_{n_0}x_{k_i}-T_{n_0}x_{k_j}}<\varepsilon/3$ for $i,j\geq J$. Then
\[
  \norm{Tx_{k_i}-Tx_{k_j}}
  \leq 2\norm{T-T_{n_0}}\,M + \norm{T_{n_0}x_{k_i}-T_{n_0}x_{k_j}}
  < \varepsilon.
\]
So $(Tx_{k_j})$ is Cauchy, hence convergent.

\emph{(iii)} If $(x_n)$ is bounded, extract a subsequence with
$(Tx_{n_k})$ convergent. Then $(STx_{n_k})$ converges by continuity of $S$.

\emph{(iv)} If $(x_n)$ is bounded, then $(Tx_n)$ is bounded in $F$, so
compactness of $S$ yields a convergent subsequence of $(STx_n)$.
\end{proof}

\begin{example}[Hilbert--Schmidt operators are compact]
\label{ex:HS-compact}
\index{Hilbert--Schmidt operator}\index{compact operator!Hilbert--Schmidt}
The integral operator $T_K$ from Example~\ref{ex:integral-op} is compact: it is
the limit (in operator norm) of finite-rank operators obtained by truncating the
singular value decomposition (or an orthonormal expansion of $K$).
\end{example}

\begin{remark}\label{rem:compact-not-finite-rank}
\index{compact operator!vs.\ finite-rank}
Not every compact operator has finite rank. The diagonal operator
$T\colon\ell^2\to\ell^2$ defined by $T e_n = \frac{1}{n}e_n$ is compact (being
the norm-limit of the finite-rank operators $T_N$ that agree with $T$ on
$e_1,\dots,e_N$ and vanish on $e_n$ for $n>N$) but has infinite-dimensional
range.
\end{remark}

%-------------------------------------------------------------------
\section{Spectrum and resolvent}
\index{spectrum}\index{resolvent}
%-------------------------------------------------------------------

In this section, $E$ is a Banach space over $\C$ and $T\in\mathcal{L}(E)$.

\begin{definition}[Resolvent set and spectrum]\label{def:spectrum}
\index{resolvent set}\index{spectrum!definition}
The \emph{resolvent set} of $T$ is
\[
  \rho(T) = \bigl\{\lambda\in\C : \lambda\Id - T \text{ is bijective in }
  \mathcal{L}(E)\bigr\}.
\]
The \emph{spectrum}\index{spectrum} of $T$ is $\sigma(T) = \C\setminus\rho(T)$.
For $\lambda\in\rho(T)$, the \emph{resolvent operator}\index{resolvent operator}
is $R(\lambda,T) = (\lambda\Id - T)^{-1}\in\mathcal{L}(E)$.
\end{definition}

\begin{definition}[Classification of the spectrum]\label{def:spectrum-parts}
\index{spectrum!point}\index{eigenvalue}\index{spectrum!continuous}
\index{spectrum!residual}
The spectrum is partitioned as $\sigma(T) = \sigma_p(T)\cup\sigma_c(T)\cup\sigma_r(T)$
where:
\begin{enumerate}[label=(\roman*)]
  \item $\sigma_p(T)$: the \emph{point spectrum} (eigenvalues): $\lambda\Id-T$ is
        not injective.
  \item $\sigma_c(T)$: the \emph{continuous spectrum}: $\lambda\Id-T$ is injective,
        $\ran(\lambda\Id-T)$ is dense but not all of $E$.
  \item $\sigma_r(T)$: the \emph{residual spectrum}: $\lambda\Id-T$ is injective,
        $\ran(\lambda\Id-T)$ is not dense.
\end{enumerate}
\end{definition}

\begin{remark}\label{rem:bdd-inverse-auto}
By the open mapping theorem (Theorem~\ref{thm:omt} or the relevant
theorem from Chapter~2), if $\lambda\Id-T$ is bijective, then
$(\lambda\Id-T)^{-1}$ is automatically bounded. Hence the definition of
$\rho(T)$ simply requires bijectivity.
\end{remark}

%-------------------------------------------------------------------
\section{The Neumann series}\label{sec:neumann-series}
\index{Neumann series}
%-------------------------------------------------------------------

\begin{theorem}[Neumann series]\label{thm:neumann}
\index{Neumann series!theorem}\index{invertibility!Neumann series}
Let $E$ be a Banach space and $T\in\mathcal{L}(E)$ with $\norm{T}<1$. Then
$\Id-T$ is invertible in $\mathcal{L}(E)$ and
\[
  (\Id-T)^{-1} = \sum_{n=0}^{\infty} T^n,
\]
with $\norm{(\Id-T)^{-1}} \leq \dfrac{1}{1-\norm{T}}$.
\end{theorem}

\begin{proof}
Since $\norm{T^n}\leq\norm{T}^n$ and $\norm{T}<1$, the series
$S_N = \sum_{n=0}^{N}T^n$ is absolutely convergent in the Banach space
$\mathcal{L}(E)$:
\[
  \sum_{n=0}^{\infty}\norm{T^n} \leq \sum_{n=0}^{\infty}\norm{T}^n
  = \frac{1}{1-\norm{T}} < \infty.
\]
Let $S = \sum_{n=0}^{\infty}T^n$. Then
\[
  (\Id - T)S_N = \sum_{n=0}^{N}T^n - \sum_{n=1}^{N+1}T^n = \Id - T^{N+1},
\]
and $\norm{T^{N+1}}\leq\norm{T}^{N+1}\to 0$. So $(\Id-T)S = \Id$. Similarly
$S(\Id-T)=\Id$, whence $(\Id-T)^{-1}=S$.
\end{proof}

\begin{example}[Application: Fredholm integral equations of the second kind]
\label{ex:neumann-fredholm}
\index{Fredholm integral equation}\index{Neumann series!application}
Consider the equation $f - \lambda K f = g$ in $L^2([0,1])$, where $K$ is the
integral operator with kernel $k\in L^2([0,1]^2)$ and $\lambda\in\C$. If
$\abs{\lambda}<\norm{K}^{-1}$, the Neumann series gives the unique solution
\[
  f = \sum_{n=0}^{\infty}\lambda^n K^n g = g + \lambda Kg + \lambda^2 K^2 g + \cdots.
\]
The partial sums $f_N = \sum_{n=0}^{N}\lambda^n K^n g$ converge to $f$ in
$L^2([0,1])$ with rate $\norm{f-f_N}\leq\frac{(\abs{\lambda}\norm{K})^{N+1}}{1-\abs{\lambda}\norm{K}}\norm{g}$.
\end{example}

\begin{corollary}[Invertible operators form an open set]
\label{cor:inv-open}
\index{invertibility!open set}\index{GL(E)@$\mathrm{GL}(E)$}
The set $\mathrm{GL}(E)$ of invertible elements of $\mathcal{L}(E)$ is open in
the operator-norm topology, and the map $T\mapsto T^{-1}$ is continuous (in fact,
analytic).
\end{corollary}

\begin{proof}
Let $S\in\mathrm{GL}(E)$ and let $T\in\mathcal{L}(E)$ with
$\norm{T-S}<\norm{S^{-1}}^{-1}$. Write
\[
  T = S\bigl(\Id - S^{-1}(S-T)\bigr).
\]
Since $\norm{S^{-1}(S-T)}\leq\norm{S^{-1}}\,\norm{S-T}<1$, the Neumann series
shows that $\Id - S^{-1}(S-T)$ is invertible, hence $T$ is invertible. Moreover,
\[
  \norm{T^{-1}-S^{-1}}
  = \norm{T^{-1}(S-T)S^{-1}}
  \leq \norm{T^{-1}}\,\norm{S-T}\,\norm{S^{-1}},
\]
which tends to zero as $T\to S$ (since $\norm{T^{-1}}$ remains bounded).
\end{proof}

\begin{corollary}[Spectrum is compact]\label{cor:spectrum-compact}
\index{spectrum!compactness}
For every $T\in\mathcal{L}(E)$, the spectrum $\sigma(T)$ is a non-empty compact
subset of $\{\lambda\in\C : \abs{\lambda}\leq\norm{T}\}$.
\end{corollary}

\begin{proof}
\textsc{Boundedness.}\index{spectrum!bounded} If $\abs{\lambda}>\norm{T}$, write
$\lambda\Id - T = \lambda(\Id - \lambda^{-1}T)$. Since
$\norm{\lambda^{-1}T}<1$, the Neumann series shows $\lambda\in\rho(T)$.

\textsc{Closedness.} The resolvent set is the preimage of $\mathrm{GL}(E)$ under
the continuous map $\lambda\mapsto\lambda\Id - T$. Since $\mathrm{GL}(E)$ is
open (Corollary~\ref{cor:inv-open}), $\rho(T)$ is open, so $\sigma(T)$ is
closed. Being closed and bounded, it is compact.

\textsc{Non-emptiness.}\index{spectrum!non-empty} Suppose $\sigma(T)=\emptyset$.
Then $R(\lambda,T)$ is an entire $\mathcal{L}(E)$-valued function. For any
$f\in E'$ and $x\in E$, the function $\lambda\mapsto f(R(\lambda,T)x)$ is
entire. For $\abs{\lambda}>2\norm{T}$,
\[
  \norm{R(\lambda,T)} \leq \frac{1}{\abs{\lambda}-\norm{T}}
  \leq \frac{1}{\abs{\lambda}/2} \to 0 \quad\text{as }\abs{\lambda}\to\infty.
\]
By Liouville's theorem\index{Liouville's theorem}, $f(R(\lambda,T)x)=0$ for all
$\lambda$. Since $f$ and $x$ were arbitrary, $R(\lambda,T)=0$, a contradiction
since $R(\lambda,T)(\lambda\Id-T)=\Id$.
\end{proof}

%-------------------------------------------------------------------
\section{Spectral radius}\label{sec:spectral-radius}
\index{spectral radius}
%-------------------------------------------------------------------

\begin{definition}[Spectral radius]\label{def:spectral-radius}
\index{spectral radius!definition}
The \emph{spectral radius} of $T\in\mathcal{L}(E)$ is
\[
  r(T) = \sup\bigl\{\abs{\lambda} : \lambda\in\sigma(T)\bigr\}.
\]
\end{definition}

\begin{theorem}[Spectral radius formula / Gelfand's formula]
\label{thm:spectral-radius}
\index{spectral radius!formula}\index{Gelfand's formula}
For every $T\in\mathcal{L}(E)$,
\[
  r(T) = \lim_{n\to\infty}\norm{T^n}^{1/n} = \inf_{n\geq 1}\norm{T^n}^{1/n}.
\]
\end{theorem}

\begin{proof}
Write $\alpha = \inf_{n\geq 1}\norm{T^n}^{1/n}$ and
$\beta = \limsup_{n\to\infty}\norm{T^n}^{1/n}$.

\textsc{Step 1: $r(T)\leq\alpha$.}\index{spectral radius!lower bound}
We first show $\alpha=\inf$ equals $\lim\inf$, using the submultiplicativity
$\norm{T^{m+n}}\leq\norm{T^m}\,\norm{T^n}$. The sequence $a_n = \log\norm{T^n}$
is subadditive, so by Fekete's lemma\index{Fekete's lemma},
$\lim_{n}a_n/n = \inf_n a_n/n$. Exponentiating, $\lim_n\norm{T^n}^{1/n}=\alpha$.
Hence $\beta = \alpha$ and the limit exists.

\textsc{Step 2: $r(T)\leq\alpha$.}
If $\lambda\in\sigma(T)$, then $\lambda^n\in\sigma(T^n)$ (since
$\lambda^n\Id - T^n = (\lambda\Id-T)\bigl(\sum_{k=0}^{n-1}\lambda^{n-1-k}T^k\bigr)$;
if $\lambda^n\Id-T^n$ were invertible, so would $\lambda\Id-T$ be). Thus
$\abs{\lambda}^n\leq\norm{T^n}$, so $\abs{\lambda}\leq\norm{T^n}^{1/n}$ for
all $n$. Taking the infimum over $n$ gives $\abs{\lambda}\leq\alpha$, hence
$r(T)\leq\alpha$.

\textsc{Step 3: $\alpha\leq r(T)$.}\index{spectral radius!upper bound}
For $\abs{\lambda}>r(T)$, $\lambda\in\rho(T)$ and the resolvent has the Laurent
expansion
\[
  R(\lambda,T) = \sum_{n=0}^{\infty}\frac{T^n}{\lambda^{n+1}},
\]
convergent in $\mathcal{L}(E)$ for $\abs{\lambda}>\norm{T}$. By analytic
continuation\index{analytic continuation}, this series converges for all
$\abs{\lambda}>r(T)$.

For any $f\in E'$ and $x\in E$, the scalar power series
$\sum_{n=0}^{\infty}f(T^n x)/\lambda^{n+1}$ converges for $\abs{\lambda}>r(T)$.
By the root test, $\limsup_n\abs{f(T^n x)}^{1/n}\leq r(T)$ for all $f,x$.

By the uniform boundedness principle\index{uniform boundedness principle}, for
each $x$ the sequence $(\norm{T^n x}^{1/n})$ satisfies
$\limsup_n\norm{T^n x}^{1/n}\leq r(T)$. Applying the uniform boundedness
principle again (to the operators $T^n/r^n$ for any $r>r(T)$) shows
$\sup_n\norm{T^n}/r^n<\infty$, hence $\limsup_n\norm{T^n}^{1/n}\leq r$.
Since $r>r(T)$ was arbitrary, $\alpha = \beta\leq r(T)$.
\end{proof}

\begin{example}\label{ex:spectral-radius-nilpotent}
\index{nilpotent operator}\index{spectral radius!nilpotent}
Let $T\colon\C^n\to\C^n$ be nilpotent ($T^n=0$). Then $\norm{T^n}^{1/n}=0$ for
all large $n$, so $r(T)=0$ and $\sigma(T)=\{0\}$.
\end{example}

\begin{example}[Shift operator]\label{ex:spectral-radius-shift}
\index{shift operator!spectral radius}
The right shift $R\colon\ell^2\to\ell^2$, $(x_1,x_2,\dots)\mapsto(0,x_1,x_2,\dots)$
satisfies $\norm{R^n}=1$ for all $n$, so $r(R)=1$. One verifies
$\sigma(R)=\overline{\mathbb{D}}$ (the closed unit disc).
\end{example}

%-------------------------------------------------------------------
\section{Perturbation of invertible operators}
\index{perturbation!of invertible operators}
%-------------------------------------------------------------------

\begin{theorem}[Perturbation bound]\label{thm:perturbation}
\index{perturbation!bound}\index{invertibility!perturbation}
Let $S\in\mathrm{GL}(E)$ and $T\in\mathcal{L}(E)$ with
$\norm{T-S}<\norm{S^{-1}}^{-1}$. Then $T\in\mathrm{GL}(E)$ and
\[
  \norm{T^{-1}} \leq \frac{\norm{S^{-1}}}{1-\norm{S^{-1}}\,\norm{T-S}},
  \qquad
  \norm{T^{-1}-S^{-1}} \leq
  \frac{\norm{S^{-1}}^2\,\norm{T-S}}{1-\norm{S^{-1}}\,\norm{T-S}}.
\]
\end{theorem}

\begin{proof}
Write $T = S(\Id - S^{-1}(S-T))$ and set $U = S^{-1}(S-T)$. Then
$\norm{U}\leq\norm{S^{-1}}\,\norm{S-T}<1$, so $\Id-U$ is invertible by the
Neumann series with $\norm{(\Id-U)^{-1}}\leq(1-\norm{U})^{-1}$. Hence
$T^{-1}=(\Id-U)^{-1}S^{-1}$ and
\[
  \norm{T^{-1}} \leq \frac{\norm{S^{-1}}}{1-\norm{S^{-1}}\,\norm{S-T}}.
\]
For the second inequality,
\[
  T^{-1}-S^{-1} = T^{-1}(S-T)S^{-1},
\]
so $\norm{T^{-1}-S^{-1}}\leq\norm{T^{-1}}\,\norm{S-T}\,\norm{S^{-1}}$, and
we substitute the bound on $\norm{T^{-1}}$.
\end{proof}

\begin{example}[Perturbation of the identity]\label{ex:perturbation-identity}
\index{perturbation!of the identity}
If $T\in\mathcal{L}(E)$ with $\norm{T}<1$, then $\Id-T$ is invertible and
\[
  \norm{(\Id-T)^{-1}-\Id} = \norm{(\Id-T)^{-1}-(\Id-T)^{-1}(\Id-T)}
  = \norm{(\Id-T)^{-1}T} \leq \frac{\norm{T}}{1-\norm{T}}.
\]
This shows that invertibility is ``stable'': operators close to the identity
remain invertible, with a quantitative bound on how far the inverse is from the
identity.
\end{example}

\begin{remark}[Analytic dependence of the resolvent]\label{rem:resolvent-analytic}
\index{resolvent!analytic dependence}
The resolvent map $\lambda\mapsto R(\lambda,T)$ is analytic on $\rho(T)$ as a
function valued in $\mathcal{L}(E)$. Indeed, for $\lambda_0\in\rho(T)$, writing
$\lambda\Id-T = (\lambda_0\Id-T)\bigl(\Id+(\lambda-\lambda_0)R(\lambda_0,T)\bigr)$,
the Neumann series gives a power series expansion in $(\lambda-\lambda_0)$
convergent for $\abs{\lambda-\lambda_0}<\norm{R(\lambda_0,T)}^{-1}$.
\end{remark}

\begin{corollary}[Resolvent identity]\label{cor:resolvent-identity}
\index{resolvent identity}
For $\lambda,\mu\in\rho(T)$,
\[
  R(\lambda,T) - R(\mu,T) = (\mu-\lambda)\,R(\lambda,T)\,R(\mu,T).
\]
In particular, $R(\lambda,T)$ and $R(\mu,T)$ commute.
\end{corollary}

\begin{proof}
Multiply both sides on the left by $(\lambda\Id-T)$ and on the right by
$(\mu\Id-T)$:
\begin{align*}
  (\lambda\Id-T)\bigl[R(\lambda,T)-R(\mu,T)\bigr](\mu\Id-T)
  &= (\mu\Id-T) - (\lambda\Id-T) \\
  &= (\mu-\lambda)\Id. \qedhere
\end{align*}
\end{proof}

%-------------------------------------------------------------------
\section{Operators on Hilbert spaces: the adjoint}
\index{adjoint!operator}\index{Hilbert space!operators on}
%-------------------------------------------------------------------

Throughout this section, $H$ and $K$ denote Hilbert spaces over $\K\in\{\R,\C\}$.

\begin{theorem}[Existence of the adjoint]\label{thm:adjoint-existence}
\index{adjoint!existence}\index{Riesz representation theorem!and adjoint}
For every $T\in\mathcal{L}(H,K)$ there exists a unique operator
$T^*\in\mathcal{L}(K,H)$ such that
\[
  \inner{Tx}{y}_K = \inner{x}{T^*y}_H
  \qquad\text{for all } x\in H,\; y\in K.
\]
Moreover, $\norm{T^*}=\norm{T}$.
\end{theorem}

\begin{proof}
Fix $y\in K$. The map $\varphi_y\colon H\to\K$ defined by $\varphi_y(x)=\inner{Tx}{y}_K$
is a bounded linear functional with $\norm{\varphi_y}\leq\norm{T}\,\norm{y}$.
By the Riesz representation theorem\index{Riesz representation theorem}, there
is a unique $z_y\in H$ with $\varphi_y(x)=\inner{x}{z_y}_H$ for all $x$, and
$\norm{z_y}=\norm{\varphi_y}\leq\norm{T}\,\norm{y}$. Define $T^*y = z_y$.

\textsc{Linearity.} For $y_1,y_2\in K$ and $\alpha\in\K$,
\[
  \inner{x}{T^*(y_1+\alpha y_2)}_H
  = \inner{Tx}{y_1+\alpha y_2}_K
  = \inner{Tx}{y_1}_K + \overline{\alpha}\inner{Tx}{y_2}_K
  = \inner{x}{T^*y_1 + \overline{\alpha}T^*y_2}_H.
\]
Wait---in a complex Hilbert space, $\inner{Tx}{\alpha y} = \overline{\alpha}\inner{Tx}{y}$,
so
\[
  \inner{x}{T^*(\alpha y)}_H = \overline{\alpha}\inner{x}{T^*y}_H
  = \inner{x}{\alpha T^*y}_H.
\]
Hence $T^*(\alpha y)=\alpha T^*y$ and $T^*$ is linear.

\textsc{Boundedness and norm.} We have $\norm{T^*}\leq\norm{T}$ from the
bound $\norm{T^*y}\leq\norm{T}\,\norm{y}$. For the reverse inequality,
\[
  \norm{Tx}_K = \sup_{\norm{y}\leq 1}\abs{\inner{Tx}{y}_K}
  = \sup_{\norm{y}\leq 1}\abs{\inner{x}{T^*y}_H}
  \leq \norm{x}\,\norm{T^*},
\]
whence $\norm{T}\leq\norm{T^*}$.
\end{proof}

\begin{proposition}[Properties of the adjoint]\label{prop:adjoint-props}
\index{adjoint!properties}
Let $S,T\in\mathcal{L}(H,K)$, $U\in\mathcal{L}(K,L)$, and $\alpha\in\K$.
\begin{enumerate}[label=(\roman*)]
  \item $(T+\alpha S)^* = T^* + \overline{\alpha}S^*$.
        \index{adjoint!linearity}
  \item $(UT)^* = T^*U^*$.
        \index{adjoint!composition}
  \item $T^{**} = T$.
        \index{adjoint!double}
  \item $\norm{T^*T} = \norm{TT^*} = \norm{T}^2$.
        \index{adjoint!norm identity}\index{C*-identity@$C^*$-identity}
\end{enumerate}
\end{proposition}

\begin{proof}
\emph{(i)} Direct from the inner product definition.

\emph{(ii)} For all $x\in H$ and $z\in L$,
\[
  \inner{UTx}{z}_L = \inner{Tx}{U^*z}_K = \inner{x}{T^*U^*z}_H.
\]
By uniqueness, $(UT)^*=T^*U^*$.

\emph{(iii)} For all $x\in H$ and $y\in K$,
$\inner{y}{T^{**}x}_K = \inner{T^*y}{x}_H = \overline{\inner{x}{T^*y}_H}
= \overline{\inner{Tx}{y}_K} = \inner{y}{Tx}_K$, so $T^{**}=T$.

\emph{(iv)} We have $\norm{T^*T}\leq\norm{T^*}\,\norm{T}=\norm{T}^2$. Conversely,
\[
  \norm{Tx}^2 = \inner{Tx}{Tx}_K = \inner{T^*Tx}{x}_H
  \leq \norm{T^*Tx}\,\norm{x}
  \leq \norm{T^*T}\,\norm{x}^2.
\]
Taking the supremum over $\norm{x}\leq 1$ gives $\norm{T}^2\leq\norm{T^*T}$.
The equality $\norm{TT^*}=\norm{T}^2$ follows by replacing $T$ with $T^*$.
\end{proof}

\begin{remark}[The $C^*$-identity]\label{rem:Cstar}
\index{C*-identity@$C^*$-identity}
Property (iv) is the \emph{$C^*$-identity}. It shows that $\mathcal{L}(H)$
with the adjoint operation is a $C^*$-algebra\index{C*-algebra@$C^*$-algebra}.
\end{remark}

%-------------------------------------------------------------------
\section{Normal, self-adjoint, and unitary operators}
\index{normal operator}\index{self-adjoint operator}\index{unitary operator}
%-------------------------------------------------------------------

\begin{definition}\label{def:special-ops}
Let $T\in\mathcal{L}(H)$.
\begin{enumerate}[label=(\roman*)]
  \item $T$ is \emph{normal}\index{normal operator!definition} if $T^*T = TT^*$.
  \item $T$ is \emph{self-adjoint}\index{self-adjoint operator!definition}
        (or \emph{Hermitian}\index{Hermitian operator}) if $T^*=T$.
  \item $T$ is \emph{unitary}\index{unitary operator!definition} if
        $T^*T = TT^* = \Id$.
  \item $T$ is a \emph{(orthogonal) projection}\index{orthogonal projection!operator}
        if $T^2=T$ and $T^*=T$.
\end{enumerate}
\end{definition}

\begin{proposition}[Characterizations of normality]\label{prop:normal-char}
\index{normal operator!characterization}
The following are equivalent for $T\in\mathcal{L}(H)$:
\begin{enumerate}[label=(\roman*)]
  \item $T$ is normal.
  \item $\norm{Tx}=\norm{T^*x}$ for all $x\in H$.
  \item $T = A + iB$ where $A,B$ are self-adjoint and $AB=BA$.
\end{enumerate}
\end{proposition}

\begin{proof}
(i)$\Rightarrow$(ii):
$\norm{Tx}^2 = \inner{T^*Tx}{x} = \inner{TT^*x}{x} = \norm{T^*x}^2$.

(ii)$\Rightarrow$(i): Polarization\index{polarization identity} gives
$\inner{T^*Tx}{y}=\inner{TT^*x}{y}$ for all $x,y$, so $T^*T=TT^*$.

(i)$\Leftrightarrow$(iii): Set $A=(T+T^*)/2$, $B=(T-T^*)/(2i)$, both
self-adjoint. Then $T=A+iB$ and
\[
  TT^*-T^*T = 2i(AB-BA),
\]
so normality is equivalent to $AB=BA$.
\end{proof}

\begin{proposition}[Properties of self-adjoint operators]
\label{prop:self-adjoint-props}
\index{self-adjoint operator!properties}
Let $T\in\mathcal{L}(H)$ be self-adjoint. Then:
\begin{enumerate}[label=(\roman*)]
  \item $\inner{Tx}{x}\in\R$ for all $x\in H$.
  \item $\norm{T} = \sup_{\norm{x}\leq 1}\abs{\inner{Tx}{x}}$.
  \item If $\K=\C$ and $\inner{Tx}{x}=0$ for all $x$, then $T=0$.
\end{enumerate}
\end{proposition}

\begin{proof}
\emph{(i)} $\overline{\inner{Tx}{x}} = \inner{x}{Tx} = \inner{T^*x}{x}
= \inner{Tx}{x}$.

\emph{(ii)} Let $M = \sup_{\norm{x}\leq 1}\abs{\inner{Tx}{x}}$. Clearly
$M\leq\norm{T}$. For the converse, the parallelogram-type identity
\[
  4\,\mathrm{Re}\,\inner{Tx}{y}
  = \inner{T(x+y)}{x+y} - \inner{T(x-y)}{x-y}
\]
(using $T^*=T$) gives $\abs{\mathrm{Re}\,\inner{Tx}{y}}\leq M\cdot\frac{1}{4}
(\norm{x+y}^2+\norm{x-y}^2)=\frac{M}{2}(\norm{x}^2+\norm{y}^2)$.
Taking $x,y$ on the unit sphere and choosing the phase of $y$ to make
$\inner{Tx}{y}$ real and positive, $\abs{\inner{Tx}{y}}\leq M$. Hence
$\norm{Tx}=\sup_{\norm{y}=1}\abs{\inner{Tx}{y}}\leq M$ and $\norm{T}\leq M$.

\emph{(iii)} By polarization, $\inner{Tx}{y}=0$ for all $x,y$, so $T=0$.
\end{proof}

\begin{theorem}[Spectrum of a self-adjoint operator is real]
\label{thm:spectrum-selfadj-real}
\index{self-adjoint operator!spectrum}\index{spectrum!self-adjoint operator}
Let $H$ be a complex Hilbert space and $T\in\mathcal{L}(H)$ self-adjoint. Then
$\sigma(T)\subset\R$. More precisely,
$\sigma(T)\subset[m,M]$ where $m=\inf_{\norm{x}=1}\inner{Tx}{x}$ and
$M=\sup_{\norm{x}=1}\inner{Tx}{x}$, and both $m,M\in\sigma(T)$.
\end{theorem}

\begin{proof}
\textsc{Step 1: $\sigma(T)\subset\R$.}\index{spectrum!real}
Let $\lambda = \alpha + i\beta$ with $\beta\neq 0$. For $x\in H$ with $\norm{x}=1$,
\[
  \norm{(\lambda\Id-T)x}^2
  = \norm{(\alpha\Id-T)x}^2 + \beta^2\norm{x}^2
  \geq \beta^2 > 0.
\]
(Here we used $\mathrm{Re}\inner{(\alpha\Id-T)x}{i\beta x}
= \beta\,\mathrm{Im}\inner{(\alpha\Id-T)x}{x}=0$ since
$\inner{(\alpha\Id-T)x}{x}\in\R$.)
Hence $\lambda\Id-T$ is bounded below. Similarly, $\overline{\lambda}\Id-T^*
= \overline{\lambda}\Id-T$ is bounded below, so
$\ran(\lambda\Id-T)$ is dense (if $y\perp\ran(\lambda\Id-T)$, then
$(\overline{\lambda}\Id-T)y=0$, contradicting $\overline{\lambda}\Id-T$
bounded below). By the open mapping theorem, $\lambda\Id-T$ is invertible.

\textsc{Step 2: $\sigma(T)\subset[m,M]$.}
If $\lambda>M$, then $\inner{(\lambda\Id-T)x}{x}\geq(\lambda-M)\norm{x}^2$,
so $\lambda\Id-T$ is bounded below and, since self-adjoint, invertible. Similarly
for $\lambda<m$.

\textsc{Step 3: $m,M\in\sigma(T)$.}\index{spectrum!endpoints}
We show $M\in\sigma(T)$; the argument for $m$ is analogous. By definition of
$M$, there exists a sequence $(x_n)$ with $\norm{x_n}=1$ and
$\inner{Tx_n}{x_n}\to M$. Then
\[
  \norm{(M\Id-T)x_n}^2
  = M^2 - 2M\inner{Tx_n}{x_n} + \norm{Tx_n}^2.
\]
Since $\inner{Tx_n}{x_n}\to M$ and
$\norm{Tx_n}^2 = \inner{T^2 x_n}{x_n}\leq\norm{T^2}\leq\norm{T}^2$, while
also $\norm{Tx_n}^2\geq\inner{Tx_n}{x_n}^2\to M^2$, we deduce
$\norm{Tx_n}^2\to M^2$. Hence $\norm{(M\Id-T)x_n}^2\to M^2-2M\cdot M+M^2=0$.
So $M\Id-T$ is not bounded below, hence $M\in\sigma(T)$.
\end{proof}

\begin{proposition}[Spectral properties of unitary operators]
\label{prop:unitary-spectrum}
\index{unitary operator!spectrum}\index{spectrum!unitary operator}
If $U\in\mathcal{L}(H)$ is unitary, then $\sigma(U)\subset\{z\in\C:\abs{z}=1\}$.
\end{proposition}

\begin{proof}
Since $\norm{U}=1$, $\sigma(U)\subset\overline{\mathbb{D}}$. Since $U^{-1}=U^*$
with $\norm{U^*}=1$, $\sigma(U^{-1})\subset\overline{\mathbb{D}}$. But
$\lambda\in\sigma(U)$ iff $\lambda^{-1}\in\sigma(U^{-1})$, so
$\abs{\lambda}\leq 1$ and $\abs{\lambda}^{-1}\leq 1$, giving $\abs{\lambda}=1$.
\end{proof}

\begin{proposition}[Orthogonal projections]\label{prop:orth-proj-char}
\index{orthogonal projection!characterization}
Let $P\in\mathcal{L}(H)$. Then $P$ is an orthogonal projection if and only if
there exists a closed subspace $M\subset H$ such that $P$ is the orthogonal
projection onto $M$. In that case, $\sigma(P)\subset\{0,1\}$.
\end{proposition}

\begin{proof}
If $P^2=P=P^*$, set $M=\ran P = \Ker(\Id-P)$. Since $P$ is continuous, $M$ is
closed. For $x\in H$, write $x = Px + (x-Px)$. We have $Px\in M$ and
$P(x-Px)=Px-P^2 x=0$, so $x-Px\in\Ker P$. Moreover
$\inner{Px}{x-Px}=\inner{Px}{x-Px}=\inner{P^2 x}{x-Px}=\inner{Px}{P(x-Px)}=0$.
So $P$ is the orthogonal projection onto $M$.

Conversely, if $P$ is the orthogonal projection onto a closed subspace $M$, then
$P^2=P$ and $\inner{Px}{y}=\inner{Px}{Py+(y-Py)}=\inner{Px}{Py}=\inner{x}{Py}$
(since $Px\perp(y-Py)$ and $x-Px\perp Py$), so $P^*=P$.

For the spectrum: if $\lambda\notin\{0,1\}$, then
$(\lambda\Id-P)^{-1}=\frac{1}{\lambda}\Id + \frac{1}{\lambda(\lambda-1)}P$.
\end{proof}

\begin{proposition}[Kernel and range of the adjoint]\label{prop:ker-ran-adjoint}
\index{adjoint!kernel and range}\index{kernel!of adjoint}\index{range!of adjoint}
Let $T\in\mathcal{L}(H,K)$. Then:
\begin{enumerate}[label=(\roman*)]
  \item $\Ker T^* = (\ran T)^\perp$.
  \item $\Ker T = (\ran T^*)^\perp$.
  \item $\overline{\ran T} = (\Ker T^*)^\perp$.
  \item $\overline{\ran T^*} = (\Ker T)^\perp$.
\end{enumerate}
\end{proposition}

\begin{proof}
\emph{(i)} $y\in\Ker T^*$ iff $T^*y=0$ iff $\inner{x}{T^*y}=0$ for all
$x\in H$ iff $\inner{Tx}{y}=0$ for all $x\in H$ iff $y\in(\ran T)^\perp$.

\emph{(ii)} Apply (i) to $T^*$, noting $T^{**}=T$.

\emph{(iii)} and \emph{(iv)} follow from (i) and (ii) using $M^{\perp\perp}=\overline{M}$.
\end{proof}

\begin{example}[Matrix operators]\label{ex:adjoint-matrix}
\index{adjoint!matrix}\index{matrix!adjoint}
If $H=\C^n$ and $T$ is represented by a matrix $A=(a_{ij})$, then $T^*$
corresponds to $A^*=(\overline{a_{ji}})$, the conjugate transpose. Self-adjoint
means $A=A^*$ (Hermitian matrix), unitary means $A^*A=I$ (unitary matrix), and
normal means $A^*A=AA^*$.
\end{example}

\begin{example}[Adjoint of a multiplication operator]\label{ex:adjoint-mult}
\index{multiplication operator!adjoint}\index{adjoint!multiplication operator}
On $L^2(\mu)$, the multiplication operator $M_\varphi$ with $\varphi\in L^\infty(\mu)$
satisfies $M_\varphi^* = M_{\overline{\varphi}}$. Indeed,
\[
  \inner{M_\varphi f}{g} = \int \varphi f\,\overline{g}\,\dd\mu
  = \int f\,\overline{\overline{\varphi}g}\,\dd\mu = \inner{f}{M_{\overline{\varphi}}g}.
\]
Thus $M_\varphi$ is self-adjoint iff $\varphi$ is real-valued a.e., normal always
(since multiplication operators commute), and unitary iff $\abs{\varphi}=1$ a.e.
\end{example}

\begin{example}[Adjoint of the right shift]\label{ex:adjoint-shift}
\index{shift operator!adjoint}\index{adjoint!shift operator}
The right shift $R\colon\ell^2\to\ell^2$, $R(x_1,x_2,\dots)=(0,x_1,x_2,\dots)$,
has adjoint $R^*=L$, the left shift: $L(x_1,x_2,x_3,\dots)=(x_2,x_3,\dots)$.
Indeed, $\inner{Rx}{y}=\sum_{n=1}^\infty x_n\overline{y_{n+1}}=\inner{x}{Ly}$.
Note that $R^*R=\Id$ but $RR^*\neq\Id$ (it is the projection onto $\{0\}\oplus\ell^2$),
so $R$ is an isometry but not unitary.
\end{example}

%-------------------------------------------------------------------
\section{Exercises}
%-------------------------------------------------------------------

\begin{exercise}[$\star$]\label{exo:op-norm-matrix}
\index{operator norm!matrix}
Let $A\colon\R^n\to\R^n$ be a linear map represented by a matrix $(a_{ij})$.
Show that $\norm{A}\leq\bigl(\sum_{i,j}a_{ij}^2\bigr)^{1/2}$ (the Frobenius
norm\index{Frobenius norm}).
\end{exercise}

\begin{exercise}[$\star$]\label{exo:compact-diagonal}
\index{compact operator!diagonal}
Let $(d_n)\subset\C$ be a bounded sequence and define $T\colon\ell^2\to\ell^2$
by $T e_n = d_n e_n$. Show that $T$ is compact if and only if $d_n\to 0$.
\end{exercise}

\begin{exercise}[$\star\star$]\label{exo:spectrum-shift}
\index{shift operator!spectrum}
Let $R$ be the unilateral right shift on $\ell^2(\N)$. Compute $\sigma(R)$,
$\sigma_p(R)$, $\sigma_c(R)$, and $\sigma_r(R)$. Compute the same for the
left shift $L=R^*$.
\end{exercise}

\begin{exercise}[$\star$]\label{exo:neumann-application}
\index{Neumann series!application}
Let $T\in\mathcal{L}(E)$ be nilpotent of order $N$ ($T^N=0$ but $T^{N-1}\neq 0$).
Compute $(\Id-T)^{-1}$ explicitly.
\end{exercise}

\begin{exercise}[$\star\star$]\label{exo:spectral-radius-compact}
\index{spectral radius!compact operator}
Let $T\in\mathcal{K}(E)$ be a compact operator on an infinite-dimensional Banach
space. Show that $0\in\sigma(T)$.
\end{exercise}

\begin{exercise}[$\star\star$]\label{exo:adjoint-kernel-range}
\index{adjoint!kernel and range}
Let $T\in\mathcal{L}(H)$. Prove that:
\begin{enumerate}[label=(\alph*)]
  \item $\Ker T^* = (\ran T)^\perp$.
  \item $\overline{\ran T} = (\Ker T^*)^\perp$.
  \item $T$ is injective if and only if $\ran T^*$ is dense.
\end{enumerate}
\end{exercise}

\begin{exercise}[$\star\star$]\label{exo:normal-eigenspaces}
\index{normal operator!eigenspaces}
Let $T\in\mathcal{L}(H)$ be normal. Show that eigenspaces corresponding to
distinct eigenvalues are orthogonal. Show that $\Ker(T-\lambda\Id)=\Ker(T^*-\overline{\lambda}\Id)$.
\end{exercise}

\begin{exercise}[$\star\star\star$]\label{exo:volterra-spectrum}
\index{Volterra operator!spectrum}
Consider the Volterra operator\index{Volterra operator}
$V\colon L^2([0,1])\to L^2([0,1])$, $(Vf)(x)=\int_0^x f(t)\,\dd t$.
\begin{enumerate}[label=(\alph*)]
  \item Show that $V$ is compact.
  \item Show that $V$ has no eigenvalues.
  \item Conclude $\sigma(V)=\{0\}$ and compute $r(V)$.
  \item Compute $\norm{V}$. (\emph{Hint:} use $V^*$ and the $C^*$-identity.)
\end{enumerate}
\end{exercise}

\begin{exercise}[$\star$]\label{exo:self-adj-positive}
\index{positive operator}\index{self-adjoint operator!positive}
An operator $T\in\mathcal{L}(H)$ is \emph{positive} (written $T\geq 0$) if
$T=T^*$ and $\inner{Tx}{x}\geq 0$ for all $x$. Show that $T^*T\geq 0$ for
any $T\in\mathcal{L}(H)$.
\end{exercise}

\begin{exercise}[$\star\star\star$]\label{exo:approximate-eigenvalue}
\index{approximate eigenvalue}\index{spectrum!approximate point spectrum}
Let $T\in\mathcal{L}(H)$ be self-adjoint. Show that every
$\lambda\in\sigma(T)$ is an \emph{approximate eigenvalue}: there exists
$(x_n)$ with $\norm{x_n}=1$ and $\norm{Tx_n-\lambda x_n}\to 0$.
\end{exercise}

\begin{exercise}[$\star\star$]\label{exo:unitary-char}
\index{unitary operator!characterization}
Let $U\in\mathcal{L}(H)$. Show the following are equivalent:
(i) $U$ is unitary;
(ii) $U$ is surjective and $\norm{Ux}=\norm{x}$ for all $x$;
(iii) $U$ is surjective and $\inner{Ux}{Uy}=\inner{x}{y}$ for all $x,y$.
\end{exercise}

\begin{exercise}[$\star$]\label{exo:resolvent-bound}
\index{resolvent!bound}
Show that for any $T\in\mathcal{L}(E)$ and $\lambda\in\rho(T)$,
$\norm{R(\lambda,T)}\geq\frac{1}{\mathrm{dist}(\lambda,\sigma(T))}$.
\end{exercise}

\begin{exercise}[$\star\star$]\label{exo:spectrum-compact-op}
\index{compact operator!spectrum}
Let $T\in\mathcal{K}(E)$ where $E$ is an infinite-dimensional Banach space.
\begin{enumerate}[label=(\alph*)]
  \item Show that every nonzero $\lambda\in\sigma(T)$ is an eigenvalue.
  \item Show that $\sigma(T)\setminus\{0\}$ is at most countable, with $0$ as the
        only possible accumulation point.
\end{enumerate}
(\emph{Hint:} use the Riesz lemma for part (a) and the fact that eigenspaces for
distinct eigenvalues of a compact operator are finite-dimensional.)
\end{exercise}

\begin{exercise}[$\star\star\star$]\label{exo:numerical-range}
\index{numerical range}
The \emph{numerical range} of $T\in\mathcal{L}(H)$ is
$W(T)=\{\inner{Tx}{x}:\norm{x}=1\}$.
\begin{enumerate}[label=(\alph*)]
  \item Show that $W(T)$ is convex (Toeplitz--Hausdorff theorem; this is hard).
  \item Show that $\sigma(T)\subset\overline{W(T)}$.
  \item Show that if $T$ is normal, $\overline{W(T)}=\conv(\sigma(T))$.
\end{enumerate}
\end{exercise}

\begin{exercise}[$\star\star$]\label{exo:compact-selfadj-eigenvalues}
\index{compact operator!self-adjoint}
Let $T\in\mathcal{K}(H)$ be self-adjoint on a separable Hilbert space $H$.
Show that either $\norm{T}$ or $-\norm{T}$ is an eigenvalue of $T$.
(\emph{Hint:} use the fact that $\norm{T}=\sup_{\norm{x}=1}\abs{\inner{Tx}{x}}$
and extract a convergent subsequence.)
\end{exercise}


%%%%%%%%%%%%%%%%%%%%%%%%%%%%%%%%%%%%%%%%%%%%%%%%%%%%%%%%%%%%
%%  CHAPTER 4 — The Hahn--Banach Theorem and Consequences
%%%%%%%%%%%%%%%%%%%%%%%%%%%%%%%%%%%%%%%%%%%%%%%%%%%%%%%%%%%%
\chapter{The Hahn--Banach Theorem and Consequences}\label{ch:hahn-banach}
\minitoc

\vspace{1em}

The Hahn--Banach theorem is one of the three pillars of functional analysis (together
with the uniform boundedness principle and the open mapping/closed graph theorems).
It guarantees that normed spaces have a rich supply of continuous linear
functionals, and its geometric formulations provide powerful convex separation
results. This chapter presents the analytic and geometric forms, and develops
fundamental consequences: the characterisation of the dual space, reflexivity,
and the Banach--Alaoglu theorem.

%-------------------------------------------------------------------
\section{Analytic form: real case}
\index{Hahn--Banach theorem!analytic form}\index{Hahn--Banach theorem!real case}
%-------------------------------------------------------------------

\begin{definition}[Sublinear functional]\label{def:sublinear}
\index{sublinear functional}
A function $p\colon E\to\R$ on a real vector space $E$ is \emph{sublinear} if
\begin{enumerate}[label=(\roman*)]
  \item $p(\lambda x) = \lambda\, p(x)$ for all $\lambda\geq 0$ and $x\in E$
        (positive homogeneity), and
  \item $p(x+y)\leq p(x)+p(y)$ for all $x,y\in E$ (subadditivity).
\end{enumerate}
\end{definition}

\begin{theorem}[Hahn--Banach, real analytic form]\label{thm:HB-real}
\index{Hahn--Banach theorem!real analytic form}
Let $E$ be a real vector space, $p\colon E\to\R$ a sublinear functional,
$G\subset E$ a linear subspace, and $g\colon G\to\R$ a linear functional
satisfying $g(x)\leq p(x)$ for all $x\in G$. Then there exists a linear
functional $f\colon E\to\R$ extending $g$ (i.e., $f|_G=g$) such that
$f(x)\leq p(x)$ for all $x\in E$.
\end{theorem}

\begin{proof}
\textsc{Step 1: Extension by one dimension.}\index{Hahn--Banach theorem!one-step extension}
Let $x_0\in E\setminus G$ and set $G_1 = G\oplus\R x_0$. We seek $c\in\R$ such
that $f_1(x+tx_0)=g(x)+tc$ satisfies $f_1\leq p$ on $G_1$, i.e.,
\[
  g(x)+tc \leq p(x+tx_0) \qquad\text{for all } x\in G,\; t\in\R.
\]
For $t>0$: $c\leq \frac{1}{t}\bigl(p(x+tx_0)-g(x)\bigr) = p(x/t+x_0)-g(x/t)$,
so $c\leq\inf_{u\in G}\bigl(p(u+x_0)-g(u)\bigr)$.

For $t<0$ (set $s=-t>0$): $-c\leq p(x-sx_0)/s-g(x/s) = p(v-x_0)-g(v)$ for
$v=x/s$, so $c\geq\sup_{v\in G}\bigl(g(v)-p(v-x_0)\bigr)$.

We must verify the sup $\leq$ the inf. For any $u,v\in G$:
\[
  g(v)-p(v-x_0) \leq g(u)+p(u+x_0)-p(v-x_0)-g(u)
\]
holds if $g(v)+g(u) \leq p(v-x_0)+p(u+x_0)$. But $g(v)+g(u) = g(v+u)
\leq p(v+u) \leq p(v-x_0)+p(u+x_0)$ by subadditivity. So the choice of $c$ is
possible and the one-step extension works.

\textsc{Step 2: Application of Zorn's lemma.}\index{Zorn's lemma!in Hahn--Banach proof}
Consider the set
\[
  \mathcal{P} = \bigl\{(H,h) : G\subset H\subset E \text{ subspace},\;
  h\colon H\to\R \text{ linear},\; h|_G=g,\; h\leq p\text{ on }H\bigr\},
\]
partially ordered by $(H_1,h_1)\leq(H_2,h_2)$ iff $H_1\subset H_2$ and
$h_2|_{H_1}=h_1$.

Every chain $\{(H_i,h_i)\}_{i\in I}$ has an upper bound: set
$H=\bigcup_i H_i$ and $h(x)=h_i(x)$ if $x\in H_i$. This is well-defined,
linear, extends $g$, and satisfies $h\leq p$.

By Zorn's lemma\index{Zorn's lemma}, $\mathcal{P}$ has a maximal element $(H_0,f)$.
If $H_0\neq E$, then Step~1 would yield a strict extension, contradicting
maximality. Hence $H_0=E$.
\end{proof}

%-------------------------------------------------------------------
\section{Analytic form: complex case}
\index{Hahn--Banach theorem!complex case}
%-------------------------------------------------------------------

\begin{theorem}[Hahn--Banach, complex analytic form]\label{thm:HB-complex}
\index{Hahn--Banach theorem!complex analytic form}
Let $E$ be a complex vector space, $p\colon E\to[0,\infty)$ a seminorm, $G\subset E$
a linear subspace, and $g\colon G\to\C$ a linear functional satisfying
$\abs{g(x)}\leq p(x)$ for all $x\in G$. Then there exists a linear functional
$f\colon E\to\C$ extending $g$ such that $\abs{f(x)}\leq p(x)$ for all $x\in E$.
\end{theorem}

\begin{proof}
\index{Hahn--Banach theorem!complex from real}
Write $g = g_1 + i\,g_2$ where $g_1=\mathrm{Re}\,g$ and $g_2=\mathrm{Im}\,g$
are real-linear functionals on $G$ (viewed as a real vector space). Note that
$g(ix) = ig(x)$ gives $g_1(ix)+ig_2(ix) = -g_2(x)+ig_1(x)$, so
\begin{equation}\label{eq:HB-complex-relation}
  g_2(x) = -g_1(ix).
\end{equation}
Moreover, $g_1(x) = \mathrm{Re}\,g(x)\leq\abs{g(x)}\leq p(x)$, and $p$ is a
sublinear functional on $E_{\R}$ (the real vector space underlying $E$).

By the real Hahn--Banach theorem (Theorem~\ref{thm:HB-real}), extend $g_1$ to a
real-linear functional $f_1\colon E\to\R$ with $f_1\leq p$ on $E$. Define
\[
  f(x) = f_1(x) - i\,f_1(ix).
\]
Then $f$ is complex-linear: $f(ix) = f_1(ix)-if_1(-x) = f_1(ix)+if_1(x) = if(x)$.
Also $f|_G = g$ by \eqref{eq:HB-complex-relation}.

For the bound, write $f(x) = \abs{f(x)}\,e^{i\theta}$. Then
$\abs{f(x)} = f(e^{-i\theta}x) = f_1(e^{-i\theta}x) \leq p(e^{-i\theta}x) = p(x)$.
\end{proof}

%-------------------------------------------------------------------
\section{Consequences: extension of functionals, separation of points}
\index{Hahn--Banach theorem!consequences}
%-------------------------------------------------------------------

\begin{corollary}[Extension of bounded linear functionals]
\label{cor:HB-extension}
\index{Hahn--Banach theorem!extension of functionals}
\index{bounded linear functional!extension}
Let $E$ be a normed space, $G\subset E$ a subspace, and $g\in G'$. Then there
exists $f\in E'$ with $f|_G = g$ and $\norm{f}_{E'} = \norm{g}_{G'}$.
\end{corollary}

\begin{proof}
Apply Theorem~\ref{thm:HB-complex} (or Theorem~\ref{thm:HB-real} in the real
case) with $p(x) = \norm{g}\,\norm{x}$. The extension satisfies
$\abs{f(x)}\leq\norm{g}\,\norm{x}$, so $\norm{f}\leq\norm{g}$. Since
$f|_G=g$, $\norm{f}\geq\norm{g}$.
\end{proof}

\begin{corollary}[Separation of points by functionals]\label{cor:HB-separation-points}
\index{Hahn--Banach theorem!separation of points}\index{separation!of points}
Let $E$ be a normed space and $x_0\in E$ with $x_0\neq 0$. Then there exists
$f\in E'$ with $\norm{f}=1$ and $f(x_0)=\norm{x_0}$.
\end{corollary}

\begin{proof}
On $G=\K x_0$, define $g(\alpha x_0)=\alpha\norm{x_0}$. Then $\norm{g}=1$.
Extend by Corollary~\ref{cor:HB-extension}.
\end{proof}

\begin{corollary}[The dual separates points]\label{cor:dual-separates}
\index{dual space!separates points}
If $x,y\in E$ and $f(x)=f(y)$ for all $f\in E'$, then $x=y$.
Equivalently, $\norm{x}=\sup\{\abs{f(x)}:f\in E',\,\norm{f}\leq 1\}$.
\end{corollary}

\begin{proof}
Apply Corollary~\ref{cor:HB-separation-points} to $x-y$.
\end{proof}

\begin{corollary}[Functionals that vanish on a subspace]\label{cor:HB-annihilator}
\index{annihilator}
Let $G$ be a closed subspace of a normed space $E$ and $x_0\notin G$. Then there
exists $f\in E'$ with $f|_G=0$, $\norm{f}=1$, and $f(x_0)=\mathrm{dist}(x_0,G)>0$.
\end{corollary}

\begin{proof}
Set $d=\mathrm{dist}(x_0,G)>0$ (since $G$ is closed). On $H=G\oplus\K x_0$,
define $g(y+\alpha x_0)=\alpha d$. Then $\abs{g(y+\alpha x_0)}=\abs{\alpha}\,d
\leq\abs{\alpha}\,\norm{x_0-(-y/\alpha)}=\norm{\alpha x_0+y}$ when $\alpha\neq 0$,
and $g(y)=0$ when $\alpha=0$. So $\norm{g}\leq 1$, and in fact $\norm{g}=1$
since for $y_n\in G$ with $\norm{x_0-y_n}\to d$,
$\abs{g(x_0-y_n)}/\norm{x_0-y_n}=d/\norm{x_0-y_n}\to 1$. Extend by
Corollary~\ref{cor:HB-extension}.
\end{proof}

%-------------------------------------------------------------------
\section{Geometric forms: separation of convex sets}
\index{Hahn--Banach theorem!geometric form}\index{separation!of convex sets}
%-------------------------------------------------------------------

\begin{definition}[Hyperplane]\label{def:hyperplane}
\index{hyperplane}\index{closed hyperplane}
A \emph{(closed) hyperplane} in $E$ is a set of the form
$H = \{x\in E : f(x) = \alpha\}$ where $f\in E'\setminus\{0\}$ and $\alpha\in\R$.
A \emph{supporting hyperplane}\index{supporting hyperplane} at a point $x_0$ of
a convex set $C$ is a hyperplane $H$ containing $x_0$ such that $C$ lies in one
of the closed half-spaces determined by $H$.
\end{definition}

\begin{definition}[Minkowski functional]\label{def:minkowski-functional}
\index{Minkowski functional}\index{gauge function}
Let $C$ be a convex subset of a real vector space $E$ with $0\in\mathrm{int}(C)$
(the algebraic interior, or core). The \emph{Minkowski functional} (or
\emph{gauge}\index{gauge}) of $C$ is
\[
  p_C(x) = \inf\{t>0 : x\in tC\} = \inf\bigl\{t>0 : x/t\in C\bigr\}.
\]
\end{definition}

\begin{proposition}[Properties of the Minkowski functional]
\label{prop:minkowski-props}
\index{Minkowski functional!properties}
Let $C$ be a convex set in a real normed space $E$ with $0\in\mathrm{int}(C)$.
Then:
\begin{enumerate}[label=(\roman*)]
  \item $p_C$ is sublinear (positively homogeneous and subadditive).
  \item If $C$ is open, then $C = \{x : p_C(x)<1\}$.
  \item If $C$ is symmetric ($-C=C$), then $p_C$ is a seminorm.
  \item If $C$ is open, bounded, and symmetric, then $p_C$ is a norm equivalent
        to $\norm{\cdot}$.
\end{enumerate}
\end{proposition}

\begin{proof}
\emph{(i)} Positive homogeneity is clear. For subadditivity: if $x\in sC$ and
$y\in tC$, then $x+y = (s+t)\bigl(\frac{s}{s+t}\cdot\frac{x}{s}+\frac{t}{s+t}
\cdot\frac{y}{t}\bigr)\in(s+t)C$ by convexity. So $p_C(x+y)\leq s+t$; taking
infima gives $p_C(x+y)\leq p_C(x)+p_C(y)$.

\emph{(ii)} If $p_C(x)<1$, there exists $t<1$ with $x\in tC\subset C$ (since
$0\in C$ and $C$ is convex: $x = t(x/t)+(1-t)\cdot 0\in C$). Conversely, if
$x\in C$ and $C$ is open, there exists $\varepsilon>0$ with
$(1+\varepsilon)x\in C$, so $p_C(x)\leq(1+\varepsilon)^{-1}<1$.

\emph{(iii)} $p_C(-x) = \inf\{t>0:-x\in tC\}=\inf\{t>0:x\in t(-C)\}=\inf\{t>0:x\in tC\}=p_C(x)$.

\emph{(iv)} Since $C$ is bounded and open with $0\in C$, there exist $r,R>0$ with
$B(0,r)\subset C\subset B(0,R)$. Then $\norm{x}/R\leq p_C(x)\leq\norm{x}/r$.
\end{proof}

\begin{theorem}[Separation of a point from a convex set]
\label{thm:separation-point-convex}
\index{separation!point and convex set}\index{Hahn--Banach theorem!separation}
Let $E$ be a real normed space, $C\subset E$ a non-empty open convex set, and
$x_0\in E\setminus C$. Then there exists $f\in E'\setminus\{0\}$ and
$\alpha\in\R$ such that
\[
  f(x) < \alpha \leq f(x_0) \qquad\text{for all } x\in C.
\]
\end{theorem}

\begin{proof}
Without loss of generality assume $0\in C$ (translate if necessary). Let
$p=p_C$ be the Minkowski functional. Since $x_0\notin C$ and $C$ is open,
$p(x_0)\geq 1$ (by Proposition~\ref{prop:minkowski-props}(ii)).

On $G=\R x_0$, define $g(tx_0)=t\,p(x_0)$. For $t\geq 0$,
$g(tx_0)=tp(x_0)=p(tx_0)$. For $t<0$, $g(tx_0)=tp(x_0)<0\leq p(tx_0)$.
So $g\leq p$ on $G$.

By the real Hahn--Banach theorem, extend $g$ to $f\colon E\to\R$ with $f\leq p$.
For $x\in C$, $f(x)\leq p(x)<1$ (since $C$ is open). Also
$f(x_0)=g(x_0)=p(x_0)\geq 1$. Set $\alpha=1$.

It remains to verify $f$ is continuous. Since $f\leq p$ and $p(x)\leq\norm{x}/r$
for some $r>0$ (as $B(0,r)\subset C$), $f(x)\leq\norm{x}/r$. Also
$f(-x)\leq p(-x)$; if $C$ is not symmetric we still have $-C$ containing a ball,
giving $f(-x)\leq\norm{x}/r'$. In general, since $0\in\mathrm{int}(C)$, there
exists $r>0$ with $B(0,r)\subset C$, so $p(x)\leq\norm{x}/r$ and
$f(x)\leq\norm{x}/r$. Similarly $-f(x)=f(-x)\leq p(-x)\leq\norm{-x}/r$, giving
$\abs{f(x)}\leq\norm{x}/r$. (For the bound on $p(-x)$: since $0\in C$ and $C$ is
open, $B(0,r)\subset C$ for some $r>0$, and also $-r'\cdot x/\norm{x}\in C$ for
some $r'>0$, but more directly, for $x$ with $\norm{x}\leq r$, $-x\in B(0,r)
\subset C$ so $p(-x)<1\leq\norm{-x}/r\cdot r/\norm{x}$ gives what we need.)
Hence $f\in E'$.
\end{proof}

\begin{theorem}[Strict separation of convex sets]\label{thm:strict-separation}
\index{separation!strict}\index{Hahn--Banach theorem!strict separation}
Let $E$ be a real normed space, $A\subset E$ compact and convex, $B\subset E$
closed and convex, with $A\cap B=\emptyset$. Then there exist $f\in E'\setminus\{0\}$
and $\alpha,\beta\in\R$ with $\alpha<\beta$ such that
\[
  f(a)\leq\alpha < \beta\leq f(b) \qquad\text{for all } a\in A,\; b\in B.
\]
\end{theorem}

\begin{proof}
\textsc{Step 1.} The set $A-B = \{a-b:a\in A,\,b\in B\}$ is closed and convex
(closed because $A$ is compact and $B$ is closed: if $a_n-b_n\to z$, extract
$a_{n_k}\to a\in A$, then $b_{n_k}\to a-z\in B$, so $z=a-(a-z)\in A-B$).
Also $0\notin A-B$ since $A\cap B=\emptyset$.

\textsc{Step 2.} Since $A-B$ is closed and $0\notin A-B$, there exists $r>0$
with $B(0,r)\cap(A-B)=\emptyset$. Then $C = (A-B)+B(0,r)$ is an open convex set
not containing $0$ (if $0\in C$, then there exist $a\in A$, $b\in B$, $u$ with
$\norm{u}<r$ and $a-b+u=0$, so $a-b=-u$ and $\norm{a-b}<r$, contradicting
$\mathrm{dist}(A,B)\geq r$). Actually, let us re-examine: the distance
$d=\inf\{\norm{a-b}:a\in A,b\in B\}>0$ by compactness of $A$ and closedness
of $B$. Set $C=A-B+B(0,d/2)$; this is open, convex, and $0\notin C$.

\textsc{Step 3.} By Theorem~\ref{thm:separation-point-convex} applied to $C$
and $x_0=0$: there exists $f\in E'\setminus\{0\}$ and $\gamma$ with
$f(x)<\gamma\leq f(0)=0$ for all $x\in C$. In particular, for all $a\in A$
and $b\in B$, $f(a-b)<0$, i.e., $f(a)<f(b)$. Then
$\sup_{a\in A}f(a)\leq\inf_{b\in B}f(b)$. Since $A$ is compact, the sup is a
max; set $\alpha=\max_A f$ and $\beta=\inf_B f$. We need $\alpha<\beta$. Indeed,
$f(a-b+u)<0$ for all $\norm{u}<d/2$, so taking $u\to 0$,
$f(a)\leq f(b)$ with equality only if $f(a-b)=0$. But then
$f(a-b+u)=f(u)<0$ for suitable $u$, contradicting that $f$ can take positive
values on $B(0,d/2)$ (since $f\neq 0$). More precisely: since $f\neq 0$,
$\sup_{\norm{u}<d/2}f(u)>0$, so for any $a\in A$, $b\in B$:
$f(a)+\sup f(u)\leq f(b)$, giving $f(a)<f(b)$. Hence $\alpha<\beta$.
\end{proof}

\begin{remark}\label{rem:separation-complex}
\index{separation!complex case}
In a complex normed space, the geometric Hahn--Banach theorem applies to the
real parts: if $f\in E'$ (complex), then $\mathrm{Re}\,f$ separates, and $f$
is determined by $\mathrm{Re}\,f$.
\end{remark}

\begin{corollary}[Supporting hyperplane theorem]\label{cor:supporting-hyperplane}
\index{supporting hyperplane!theorem}
Let $C$ be a closed convex subset of a normed space $E$ and $x_0\in\partial C$.
Then there exists a supporting hyperplane to $C$ at $x_0$: a functional
$f\in E'\setminus\{0\}$ with $\mathrm{Re}\,f(x)\leq\mathrm{Re}\,f(x_0)$
for all $x\in C$.
\end{corollary}

\begin{proof}
Since $x_0\in\partial C$, there exist $y_n\notin C$ with $y_n\to x_0$.
For each $n$, by Theorem~\ref{thm:separation-point-convex} (applied to
$\mathrm{int}(C)$ if it is nonempty, or a slight adaptation otherwise),
there exists $f_n\in E'$ with $\norm{f_n}=1$ and
$\mathrm{Re}\,f_n(x)\leq\mathrm{Re}\,f_n(y_n)$ for all $x\in C$.
By weak*-compactness of the unit ball (Banach--Alaoglu, Theorem~\ref{thm:banach-alaoglu}
below), extract a weak*-convergent subnet $f_{n_\alpha}\weakstarto f$ with
$\norm{f}\leq 1$. Then $\mathrm{Re}\,f(x)\leq\mathrm{Re}\,f(x_0)$ for all
$x\in C$, and $f\neq 0$ since $\mathrm{Re}\,f(x_0)=\lim\mathrm{Re}\,f_{n_\alpha}(y_{n_\alpha})
\geq\limsup\mathrm{Re}\,f_{n_\alpha}(x)$ for any interior point $x$.
\end{proof}

%-------------------------------------------------------------------
\section{Applications to dual spaces}
\index{dual space!applications of Hahn--Banach}
%-------------------------------------------------------------------

\begin{theorem}[Dual of $\ell^p$]\label{thm:dual-ellp}
\index{dual space!of $\ell^p$}\index{lp spaces@$\ell^p$ spaces!dual}
For $1\leq p<\infty$, the map $\Phi\colon\ell^q\to(\ell^p)'$ defined by
\[
  \Phi(y)(x) = \sum_{n=1}^{\infty}x_n y_n, \qquad y=(y_n)\in\ell^q,\;
  x=(x_n)\in\ell^p,
\]
where $\frac{1}{p}+\frac{1}{q}=1$ (with $q=\infty$ when $p=1$), is an
isometric isomorphism.
\end{theorem}

\begin{proof}
\textsc{$\Phi$ is isometric.}\index{H\"older's inequality}
By H\"older's inequality, $\abs{\Phi(y)(x)}\leq\norm{x}_p\norm{y}_q$, so
$\norm{\Phi(y)}\leq\norm{y}_q$. For the reverse, when $1<p<\infty$,
take $x_n = \abs{y_n}^{q/p}\sgn(\overline{y_n})$. Then
$\norm{x}_p^p = \sum\abs{y_n}^q = \norm{y}_q^q$ and
$\Phi(y)(x) = \sum\abs{y_n}^{q/p+1}=\norm{y}_q^q$, so
$\norm{\Phi(y)}\geq\norm{y}_q^q/\norm{y}_q^{q/p}=\norm{y}_q$.
The case $p=1$ is similar.

\textsc{$\Phi$ is surjective.}\index{dual space!surjectivity of identification}
Let $f\in(\ell^p)'$. Set $y_n = f(e_n)$ where $e_n$ is the $n$-th standard
basis vector. For any finite set $F\subset\N$, the choice
$x = \sum_{n\in F}\abs{y_n}^{q/p}\sgn(\overline{y_n})\,e_n$ gives
\[
  \sum_{n\in F}\abs{y_n}^q = f(x) \leq \norm{f}\,\norm{x}_p
  = \norm{f}\Bigl(\sum_{n\in F}\abs{y_n}^q\Bigr)^{1/p},
\]
so $\bigl(\sum_{n\in F}\abs{y_n}^q\bigr)^{1/q}\leq\norm{f}$. Thus
$y\in\ell^q$ with $\norm{y}_q\leq\norm{f}$. By density of finitely supported
sequences, $f=\Phi(y)$.
\end{proof}

\begin{theorem}[Riesz--Markov representation, statement]\label{thm:riesz-markov}
\index{Riesz--Markov theorem}\index{dual space!of $C([0,1])$}
The dual of $C([0,1])$ is isometrically isomorphic to $\mathcal{M}([0,1])$, the
space of finite signed (or complex) Radon measures on $[0,1]$, via the pairing
\[
  \mu \mapsto \Bigl(f\mapsto\int_0^1 f\,\dd\mu\Bigr).
\]
\end{theorem}

\begin{remark}
The proof of the Riesz--Markov theorem requires measure-theoretic tools and will
be presented in Chapter~9 when we study the theory of distributions and measures
on locally compact spaces.
\end{remark}

%-------------------------------------------------------------------
\section{Bidual and canonical embedding}\label{sec:bidual}
\index{bidual}\index{canonical embedding}
%-------------------------------------------------------------------

\begin{definition}[Bidual]\label{def:bidual}
\index{bidual!definition}
The \emph{bidual} of a normed space $E$ is $E'' = (E')'$, the dual of the dual
space.
\end{definition}

\begin{theorem}[Canonical embedding is isometric]\label{thm:canonical-embedding}
\index{canonical embedding!isometry}\index{bidual!canonical embedding}
The map $J\colon E\to E''$ defined by
\[
  J(x)(f) = f(x), \qquad x\in E,\; f\in E',
\]
is a linear isometric embedding (i.e., $\norm{J(x)}_{E''}=\norm{x}_E$ for all
$x\in E$).
\end{theorem}

\begin{proof}
Linearity is clear. For the norm: $\abs{J(x)(f)} = \abs{f(x)}\leq\norm{f}\,\norm{x}$,
so $\norm{J(x)}\leq\norm{x}$. Conversely, by Corollary~\ref{cor:HB-separation-points},
there exists $f_0\in E'$ with $\norm{f_0}=1$ and $f_0(x)=\norm{x}$. Then
$\abs{J(x)(f_0)}=\norm{x}$, so $\norm{J(x)}\geq\norm{x}$.
\end{proof}

\begin{definition}[Reflexive space]\label{def:reflexive}
\index{reflexive space}\index{reflexive space!definition}
A Banach space $E$ is \emph{reflexive} if $J\colon E\to E''$ is surjective (and
hence an isometric isomorphism).
\end{definition}

\begin{example}\label{ex:reflexive-examples}
\index{reflexive space!examples}
\begin{enumerate}[label=(\roman*)]
  \item Every Hilbert space is reflexive (by the Riesz representation theorem).
  \item $\ell^p$ and $L^p(\mu)$ are reflexive for $1<p<\infty$.
  \item $\ell^1$, $\ell^\infty$, $c_0$, $L^1(\mu)$, $L^\infty(\mu)$, and
        $C([0,1])$ are not reflexive.
\end{enumerate}
\end{example}

\begin{proposition}\label{prop:reflexive-closed-subspace}
\index{reflexive space!closed subspace}
A closed subspace of a reflexive Banach space is reflexive.
\end{proposition}

\begin{proof}
Let $F$ be a closed subspace of a reflexive space $E$. Let $\xi\in F''$. Define
$\Phi\in E''$ by $\Phi(f) = \xi(f|_F)$ for $f\in E'$. Since $E$ is reflexive,
$\Phi=J_E(x)$ for some $x\in E$, i.e., $f(x)=\xi(f|_F)$ for all $f\in E'$.

We claim $x\in F$. If $x\notin F$, by Corollary~\ref{cor:HB-annihilator} there
exists $f\in E'$ with $f|_F=0$ and $f(x)\neq 0$. But then
$f(x)=\xi(f|_F)=\xi(0)=0$, a contradiction.

So $x\in F$ and for any $g\in F'$, extend $g$ to $f\in E'$ (by
Corollary~\ref{cor:HB-extension}). Then $\xi(g) = \xi(f|_F) = f(x) = g(x) = J_F(x)(g)$.
Hence $\xi = J_F(x)$ and $J_F$ is surjective.
\end{proof}

\begin{proposition}\label{prop:reflexive-quotient}
\index{reflexive space!quotient}
A quotient of a reflexive Banach space by a closed subspace is reflexive.
\end{proposition}

\begin{proof}
Let $E$ be reflexive and $G\subset E$ closed. Let $\pi\colon E\to E/G$ be the
quotient map. Given $\xi\in(E/G)''$, define $\Psi\in E''$ by $\Psi(f)=\xi(f\circ\pi^{-1})$
--- more precisely, note that $\pi'\colon(E/G)'\to E'$ defined by
$\pi'(g)=g\circ\pi$ is an isometric embedding with range $G^\perp$
(Exercise~\ref{exo:annihilator-props}). Define $\Psi\in E''$ by
$\Psi(f)=\xi((\pi')^{-1}(f))$ for $f\in G^\perp$ and extend to $E'$.
Since $E$ is reflexive, $\Psi=J_E(x)$ for some $x\in E$. Then for any
$g\in(E/G)'$, $\xi(g)=\Psi(\pi'g)=(\pi'g)(x)=g(\pi(x))=J_{E/G}(\pi(x))(g)$.
Hence $\xi=J_{E/G}(\pi(x))$ and the quotient is reflexive.
\end{proof}

\begin{remark}[Characterizations of reflexivity]\label{rem:reflexive-char}
\index{reflexive space!characterizations}
To summarize, for a Banach space $E$ the following are equivalent:
\begin{enumerate}[label=(\roman*)]
  \item $E$ is reflexive (i.e., $J$ is surjective).
  \item $B_E$ is weakly compact (Corollary~\ref{cor:reflexive-alaoglu}).
  \item Every bounded sequence in $E$ has a weakly convergent subsequence
        (Eberlein--\v{S}mulian theorem, Chapter~8).
  \item Every $f\in E'$ attains its norm on $B_E$ (James' theorem).
\end{enumerate}
\end{remark}

%-------------------------------------------------------------------
\section{The subdifferential}
\index{subdifferential}\index{convex function!subdifferential}
%-------------------------------------------------------------------

\begin{definition}[Subdifferential]\label{def:subdifferential}
\index{subdifferential!definition}
Let $E$ be a real normed space and $\varphi\colon E\to\R\cup\{+\infty\}$ a convex
function. The \emph{subdifferential} of $\varphi$ at a point $x_0\in E$ where
$\varphi(x_0)<+\infty$ is
\[
  \partial\varphi(x_0) = \bigl\{f\in E' : \varphi(x)-\varphi(x_0)
  \geq f(x-x_0) \text{ for all } x\in E\bigr\}.
\]
Elements of $\partial\varphi(x_0)$ are called \emph{subgradients}\index{subgradient}.
\end{definition}

\begin{proposition}[Non-emptiness of the subdifferential]
\label{prop:subdiff-nonempty}
\index{subdifferential!non-emptiness}
If $\varphi\colon E\to\R\cup\{+\infty\}$ is convex, lower semicontinuous, and
$\varphi(x_0)\in\R$, then $\partial\varphi(x_0)\neq\emptyset$.
\end{proposition}

\begin{proof}
The epigraph $\mathrm{epi}(\varphi)=\{(x,t):t\geq\varphi(x)\}$\index{epigraph}
is a closed convex subset of $E\times\R$, and $(x_0,\varphi(x_0))\in\partial\,\mathrm{epi}(\varphi)$.
By the supporting hyperplane theorem (Corollary~\ref{cor:supporting-hyperplane}),
there exist $f\in E'$ and $\beta\in\R$, not both zero, such that
\[
  f(x)+\beta t \leq f(x_0)+\beta\varphi(x_0)
  \qquad\text{for all } (x,t)\in\mathrm{epi}(\varphi).
\]
Since $(x_0,t)\in\mathrm{epi}(\varphi)$ for all $t\geq\varphi(x_0)$, we need
$\beta t\leq\beta\varphi(x_0)$ for large $t$, so $\beta\leq 0$. If $\beta=0$,
then $f(x)\leq f(x_0)$ for all $x$ in $\dom\varphi$, which (if $\dom\varphi$
contains a neighborhood of $x_0$) forces $f=0$. Since $\varphi$ is continuous
at $x_0$ (being convex and finite, it is continuous on the interior of its
effective domain), $\beta<0$. Dividing by $-\beta>0$:
\[
  \varphi(x) \geq \varphi(x_0) + (-f/\beta)(x-x_0),
\]
so $-f/\beta\in\partial\varphi(x_0)$.
\end{proof}

\begin{example}[Subdifferential of the norm]\label{ex:subdiff-norm}
\index{subdifferential!of the norm}\index{norm!subdifferential}
For $\varphi(x)=\norm{x}$ on a normed space $E$:
\begin{enumerate}[label=(\roman*)]
  \item If $x\neq 0$: $\partial\varphi(x) = \{f\in E':\norm{f}=1,\, f(x)=\norm{x}\}$.
  \item $\partial\varphi(0) = \{f\in E':\norm{f}\leq 1\} = B_{E'}$.
\end{enumerate}
\end{example}

%-------------------------------------------------------------------
\section{The Banach--Alaoglu theorem}\label{sec:banach-alaoglu}
\index{Banach--Alaoglu theorem}
%-------------------------------------------------------------------

Recall that the \emph{weak* topology}\index{weak* topology} on $E'$ is the
coarsest topology making all evaluation maps $f\mapsto f(x)$ ($x\in E$)
continuous.

\begin{notation}\label{not:weak-star}
\index{weak* convergence}
We write $f_\alpha\weakstarto f$ to denote convergence in the weak* topology:
$f_\alpha(x)\to f(x)$ for all $x\in E$.
\end{notation}

\begin{theorem}[Banach--Alaoglu]\label{thm:banach-alaoglu}
\index{Banach--Alaoglu theorem!statement}
Let $E$ be a normed space. Then the closed unit ball
$B_{E'} = \{f\in E':\norm{f}\leq 1\}$ is compact in the weak* topology.
\end{theorem}

\begin{remark}\label{rem:alaoglu-proof-deferred}
The proof uses Tychonoff's theorem and will be given in Chapter~8 (Weak
Topologies). Here we state it and develop some immediate consequences.
\end{remark}

\begin{corollary}\label{cor:alaoglu-bounded}
\index{Banach--Alaoglu theorem!consequences}
Every bounded net $(f_\alpha)$ in $E'$ has a weak*-convergent subnet.
\end{corollary}

\begin{corollary}[Goldstine's theorem]\label{cor:goldstine}
\index{Goldstine's theorem}
Let $E$ be a normed space. Then $J(B_E)$ is weak*-dense in $B_{E''}$.
\end{corollary}

\begin{proof}
$J(B_E)$ is a convex subset of $E''$. If it were not weak*-dense in $B_{E''}$,
there would exist $\xi_0\in B_{E''}$ not in the weak*-closure of $J(B_E)$.
By the geometric Hahn--Banach theorem in the weak* topology, there would exist
$f\in E'$ and $\alpha$ such that $\mathrm{Re}\,J(x)(f)\leq\alpha<\mathrm{Re}\,\xi_0(f)$
for all $x\in B_E$. But $\mathrm{Re}\,J(x)(f)=\mathrm{Re}\,f(x)$, so
$\sup_{\norm{x}\leq 1}\mathrm{Re}\,f(x)=\norm{f}\leq\alpha<\mathrm{Re}\,\xi_0(f)
\leq\norm{\xi_0}\,\norm{f}\leq\norm{f}$, a contradiction.
\end{proof}

\begin{corollary}\label{cor:reflexive-alaoglu}
\index{reflexive space!weak compactness}
A Banach space $E$ is reflexive if and only if $B_E$ is weakly compact.
\end{corollary}

\begin{proof}
If $E$ is reflexive, then $J(B_E)=B_{E''}$ is weak*-compact (Banach--Alaoglu
applied to $E'$), and since $J$ is a homeomorphism from the weak topology on
$E$ to the weak* topology on $E''$, $B_E$ is weakly compact.

Conversely, if $B_E$ is weakly compact, then $J(B_E)$ is weak*-compact in
$E''$, hence weak*-closed. By Goldstine's theorem, $J(B_E)$ is weak*-dense in
$B_{E''}$, so $J(B_E)=B_{E''}$ and $J$ is surjective.
\end{proof}

%-------------------------------------------------------------------
\section{Exercises}
%-------------------------------------------------------------------

\begin{exercise}[$\star$]\label{exo:HB-dominated}
\index{Hahn--Banach theorem!exercise}
Let $E$ be a real vector space and $p\colon E\to\R$ a sublinear functional.
Show that there exists a linear functional $f\colon E\to\R$ with $f\leq p$.
(\emph{Hint:} apply Theorem~\ref{thm:HB-real} with $G=\{0\}$.)
\end{exercise}

\begin{exercise}[$\star$]\label{exo:HB-finite-dim}
\index{Hahn--Banach theorem!finite-dimensional}
Give a proof of the Hahn--Banach theorem in finite dimensions without Zorn's
lemma (by induction on codimension).
\end{exercise}

\begin{exercise}[$\star\star$]\label{exo:dual-c0}
\index{dual space!of $c_0$}
Show that $(c_0)'\cong\ell^1$ isometrically, where $c_0$ is the space of
sequences converging to zero with the sup norm.
\end{exercise}

\begin{exercise}[$\star$]\label{exo:annihilator-props}
\index{annihilator!properties}
Let $E$ be a normed space and $G\subset E$ a subspace. The \emph{annihilator}
is $G^\perp = \{f\in E' : f|_G=0\}$.
\begin{enumerate}[label=(\alph*)]
  \item Show $G^\perp$ is a closed subspace of $E'$.
  \item Show $(E/\overline{G})' \cong G^\perp$ isometrically.
  \item Show $(\overline{G})' \cong E'/G^\perp$ isometrically.
\end{enumerate}
\end{exercise}

\begin{exercise}[$\star\star$]\label{exo:separation-exercise}
\index{separation!exercise}
Let $C$ be a closed convex subset of a normed space $E$ and $x_0\notin C$.
Show that there exists $f\in E'$ such that
$\mathrm{Re}\,f(x_0) > \sup_{x\in C}\mathrm{Re}\,f(x)$.
\end{exercise}

\begin{exercise}[$\star\star$]\label{exo:not-reflexive}
\index{reflexive space!non-examples}
\begin{enumerate}[label=(\alph*)]
  \item Show that $c_0$ is not reflexive by showing that $J(c_0)\neq(\ell^1)'
        \cong\ell^\infty$.
  \item Show that $\ell^1$ is not reflexive.
\end{enumerate}
\end{exercise}

\begin{exercise}[$\star\star\star$]\label{exo:james-reflexive}
\index{James' theorem (exercise)}
Let $E$ be a Banach space. Show that if every $f\in E'$ attains its norm on
$B_E$ (i.e., there exists $x$ with $\norm{x}=1$ and $f(x)=\norm{f}$), then
$E$ is reflexive. (This is James' theorem; proving it is quite hard---outline
the argument using the characterisation of reflexivity via weak compactness.)
\end{exercise}

\begin{exercise}[$\star\star$]\label{exo:minkowski-functional-exercise}
\index{Minkowski functional!exercise}
Let $C=\{x\in\R^n:\inner{a_i}{x}\leq b_i,\;i=1,\dots,m\}$ be a convex polytope
containing the origin in its interior. Compute the Minkowski functional $p_C$
explicitly.
\end{exercise}

\begin{exercise}[$\star\star$]\label{exo:subdiff-exercise}
\index{subdifferential!exercise}
Compute $\partial\varphi(x)$ for $\varphi(x)=\norm{x}_1$ on $\ell^1$ (the
$\ell^1$-norm). Characterize the subgradients at the origin and at a point
$x=(x_n)$ with all $x_n\neq 0$.
\end{exercise}

\begin{exercise}[$\star$]\label{exo:HB-quotient}
\index{Hahn--Banach theorem!quotient space}
Use the Hahn--Banach theorem to show that for any closed subspace $G$ of a
Banach space $E$ and any $x\notin G$,
\[
  \mathrm{dist}(x,G) = \max\bigl\{f(x) : f\in G^\perp,\;\norm{f}\leq 1\bigr\}.
\]
\end{exercise}

\begin{exercise}[$\star\star$]\label{exo:weak-star-metriz}
\index{weak* topology!metrizability}
Let $E$ be a separable normed space. Show that the weak* topology on $B_{E'}$
is metrizable. (\emph{Hint:} if $(x_n)$ is dense in $B_E$, use
$d(f,g)=\sum_{n=1}^{\infty}2^{-n}\abs{f(x_n)-g(x_n)}$.)
\end{exercise}

\begin{exercise}[$\star\star\star$]\label{exo:hahn-banach-uniqueness}
\index{Hahn--Banach theorem!uniqueness of extension}
Let $E$ be a normed space and $G$ a subspace. Show that every $g\in G'$
has a \emph{unique} Hahn--Banach extension to $E$ if and only if the dual
norm on $E'/G^\perp$ is strictly convex. (A norm is \emph{strictly convex}
\index{strictly convex norm} if $\norm{x}=\norm{y}=\norm{(x+y)/2}=1$
implies $x=y$.)
\end{exercise}

\begin{exercise}[$\star\star$]\label{exo:bidual-l1}
\index{bidual!of $\ell^1$}
Identify $(\ell^1)''$ explicitly. Show that the canonical embedding
$J\colon\ell^1\to(\ell^1)''$ corresponds to the natural inclusion
$\ell^1\hookrightarrow\ell^\infty\cong(\ell^1)'$ followed by the canonical
embedding. Conclude again that $\ell^1$ is not reflexive.
\end{exercise}

\begin{exercise}[$\star$]\label{exo:reflexive-dual}
\index{reflexive space!dual of reflexive}
Show that if $E$ is reflexive, then $E'$ is also reflexive.
\end{exercise}

\begin{exercise}[$\star\star$]\label{exo:mazur-separation}
\index{Mazur's theorem (exercise)}
Let $E$ be a normed space and $C\subset E$ a convex set. Using the Hahn--Banach
theorem, show that the closure of $C$ in the norm topology equals the closure
of $C$ in the weak topology. (This is Mazur's theorem.)
\end{exercise}

\begin{exercise}[$\star\star$]\label{exo:separation-two-convex}
\index{separation!two convex sets}
Let $A$ and $B$ be non-empty disjoint convex subsets of a normed space $E$,
with $A$ open. Show that there exist $f\in E'\setminus\{0\}$ and $\alpha\in\R$
such that $\mathrm{Re}\,f(a)<\alpha\leq\mathrm{Re}\,f(b)$ for all $a\in A$
and $b\in B$. Give an example showing that strict separation may fail without
compactness of one of the sets.
\end{exercise}

\begin{exercise}[$\star\star$]\label{exo:banach-limits}
\index{Banach limit}\index{Hahn--Banach theorem!Banach limits}
A \emph{Banach limit} is a linear functional $L\colon\ell^\infty\to\R$
satisfying: (i) $L\geq 0$ (i.e., $x_n\geq 0$ for all $n$ implies $L(x)\geq 0$);
(ii) $L((x_{n+1}))=L((x_n))$ (shift-invariance);
(iii) $L((1,1,1,\dots))=1$.
\begin{enumerate}[label=(\alph*)]
  \item Use the Hahn--Banach theorem to show that Banach limits exist.
        (\emph{Hint:} on the subspace of convergent sequences, take $L = \lim$,
        and use the sublinear functional $p(x) = \limsup_{n\to\infty}\frac{1}{n}
        \sum_{k=1}^{n}x_k$.)
  \item Show that if $(x_n)$ converges, then $L((x_n))=\lim_{n\to\infty}x_n$
        for every Banach limit $L$.
\end{enumerate}
\end{exercise}

\begin{exercise}[$\star\star\star$]\label{exo:total-subsets}
\index{total subset}\index{weak* topology!total subset}
A subset $S\subset E'$ is \emph{total} (or \emph{separating}) if $f(x)=0$ for
all $f\in S$ implies $x=0$. Show that $S$ is total if and only if $\spn(S)$ is
weak*-dense in $E'$. (\emph{Hint:} use the Hahn--Banach theorem in the
weak*-topology, identifying $(E',w^*)' \cong E$ via the canonical embedding.)
\end{exercise}

\begin{exercise}[$\star$]\label{exo:norm-attaining}
\index{norm-attaining functional}
Let $E=c_0$ and $f=(f_n)\in\ell^1=(c_0)'$. Show that $f$ attains its norm
on $B_{c_0}$ if and only if the supremum $\norm{f}_1=\sum\abs{f_n}$ is
attained in a suitable sense. Give an explicit $f\in(c_0)'$ that attains its
norm and one that does not (impossible in this case---show all functionals in
$\ell^1$ attain their norm on $B_{c_0}$).
\end{exercise}

\begin{remark}[Historical note]\label{rem:HB-history}
\index{Hahn--Banach theorem!history}
The Hahn--Banach theorem was proved independently by Hans Hahn (1927) for
separable spaces and by Stefan Banach (1929) in full generality. The geometric
forms were developed by various authors, including Mazur, Eidelheit, and
Dieudonn\'e. The use of Zorn's lemma (or equivalently, the axiom of choice)
is essential: there exist models of set theory (without full choice) in which
the Hahn--Banach theorem fails. However, weaker forms of choice such as the
Boolean prime ideal theorem suffice.
\end{remark}
%%%%%%%%%%%%%%%%%%%%%%%%%%%%%%%%%%%%%%%%%%%%%%%%%%%%%%%%%%%%%%%%%%%%%%%
%% CHAPTER 5 — The Uniform Boundedness Principle
%%%%%%%%%%%%%%%%%%%%%%%%%%%%%%%%%%%%%%%%%%%%%%%%%%%%%%%%%%%%%%%%%%%%%%%
\chapter{The Uniform Boundedness Principle}\label{ch:ubp}
\minitoc

\vspace{1em}
\noindent
The \emph{uniform boundedness principle}\index{uniform boundedness principle}, also known as the
\emph{Banach--Steinhaus theorem}\index{Banach--Steinhaus theorem}, is one of the cornerstones
of functional analysis. It asserts that a family of bounded linear operators that is pointwise
bounded must in fact be uniformly bounded in the operator norm. The proof rests on the
\emph{Baire category theorem}\index{Baire category theorem}, a fundamental topological result
concerning complete metric spaces. In this chapter we develop the Baire category theorem from
scratch, prove the uniform boundedness principle in its full generality, and present a wealth
of applications ranging from Fourier analysis to the theory of divergent series.

%-----------------------------------------------------------------------
\section{The Baire category theorem}\label{sec:baire}
%-----------------------------------------------------------------------

\subsection{Meagre and residual sets}\label{subsec:meagre}

\begin{definition}[Nowhere dense set]\label{def:nowhere-dense}
\index{nowhere dense set}
Let $(X,d)$ be a metric space. A subset $A \subset X$ is called
\textbf{nowhere dense} if its closure $\overline{A}$ has empty interior:
\[
  \mathrm{int}(\overline{A}) = \varnothing.
\]
Equivalently, $A$ is nowhere dense if and only if for every nonempty open set
$U \subset X$, there exists a nonempty open set $V \subset U$ such that
$V \cap A = \varnothing$.
\end{definition}

\begin{example}\label{ex:nowhere-dense-examples}
\index{nowhere dense set!examples}
\begin{enumerate}[label=(\roman*)]
  \item The set $\Z$ is nowhere dense in $\R$ with the usual metric.
  \item Any finite subset of $\R$ is nowhere dense.
  \item The Cantor set\index{Cantor set} $\mathcal{C} \subset [0,1]$ is closed with
    empty interior, hence nowhere dense.
  \item In any normed space $X$, a proper closed subspace $Y \subsetneq X$ is
    nowhere dense. Indeed, if $\mathrm{int}(\overline{Y}) = \mathrm{int}(Y) \neq \varnothing$,
    then $Y$ contains an open ball, and by the linearity of $Y$ this forces $Y = X$.
\end{enumerate}
\end{example}

\begin{definition}[Meagre and residual sets]\label{def:meagre}
\index{meagre set}\index{set of first category|see{meagre set}}
\index{residual set}\index{set of second category}
Let $(X,d)$ be a metric space.
\begin{enumerate}[label=(\roman*)]
  \item A set $M \subset X$ is \textbf{meagre} (or \textbf{of the first category})
    if it can be written as a countable union of nowhere dense sets:
    \[
      M = \bigcup_{n=1}^{\infty} A_n, \qquad \text{each } A_n \text{ nowhere dense.}
    \]
  \item A set $R \subset X$ is \textbf{residual} (or \textbf{comeagre}) if its
    complement $X \setminus R$ is meagre.
  \item A set $S \subset X$ is \textbf{of the second category} if it is not meagre.
\end{enumerate}
\end{definition}

\begin{remark}\label{rem:category-terminology}
\index{Baire category!terminology}
The terminology ``first'' and ``second'' category is due to Baire. Meagre sets are
``topologically small,'' analogous to measure-zero sets in measure theory. However,
these notions do not coincide: there exist meagre sets of full Lebesgue measure and
residual sets of measure zero.
\end{remark}

\subsection{Statement and proof of the Baire category theorem}

\begin{theorem}[Baire category theorem]\label{thm:baire}
\index{Baire category theorem!complete metric spaces}
Let $(X,d)$ be a complete metric space. Then:
\begin{enumerate}[label=(\alph*)]
  \item $X$ is of the second category in itself, i.e., $X$ is not meagre.
  \item Every countable intersection of dense open sets is dense in $X$.
  \item Equivalently, a countable union of closed sets with empty interior has
    empty interior.
\end{enumerate}
\end{theorem}

\begin{proof}
We prove statement (b); the equivalence with (a) and (c) follows by taking complements.

Let $\{U_n\}_{n=1}^{\infty}$ be a sequence of dense open subsets of $X$. We must show
that $G = \bigcap_{n=1}^{\infty} U_n$ is dense in $X$. Let $W \subset X$ be an
arbitrary nonempty open set; we need to show $G \cap W \neq \varnothing$.

\textbf{Step 1.}\index{Baire category theorem!proof} Since $U_1$ is dense and open,
$U_1 \cap W$ is a nonempty open set. Choose $x_1 \in U_1 \cap W$ and $r_1 > 0$
such that
\[
  \overline{B(x_1, r_1)} \subset U_1 \cap W, \qquad 0 < r_1 < 1.
\]
This is possible because $U_1 \cap W$ is open.

\textbf{Step 2.} Since $U_2$ is dense, $U_2 \cap B(x_1, r_1)$ is nonempty and open.
Choose $x_2 \in U_2 \cap B(x_1, r_1)$ and $r_2 > 0$ such that
\[
  \overline{B(x_2, r_2)} \subset U_2 \cap B(x_1, r_1), \qquad 0 < r_2 < \tfrac{1}{2}.
\]

\textbf{Induction.} Proceeding inductively, suppose we have chosen $x_n$ and $r_n$ with
$0 < r_n < 1/n$ and
\[
  \overline{B(x_n, r_n)} \subset U_n \cap B(x_{n-1}, r_{n-1}).
\]
Since $U_{n+1}$ is dense, $U_{n+1} \cap B(x_n, r_n)$ is nonempty and open. Choose
$x_{n+1} \in U_{n+1} \cap B(x_n, r_n)$ and $0 < r_{n+1} < 1/(n+1)$ with
\[
  \overline{B(x_{n+1}, r_{n+1})} \subset U_{n+1} \cap B(x_n, r_n).
\]

\textbf{Step 3 (Cauchy sequence).} The sequence $(x_n)_{n \geq 1}$ is Cauchy: for
$m > n$, we have $x_m \in B(x_n, r_n)$, so $d(x_m, x_n) < r_n < 1/n \to 0$. Since
$X$ is complete, $x_n \to x^*$ for some $x^* \in X$.

\textbf{Step 4 (Limit lies in the intersection).} For each fixed $n$, we have
$x_m \in \overline{B(x_n, r_n)}$ for all $m \geq n$, and since
$\overline{B(x_n, r_n)}$ is closed, $x^* \in \overline{B(x_n, r_n)} \subset U_n$.
Hence $x^* \in U_n$ for every $n$, so $x^* \in G$. Also,
$x^* \in \overline{B(x_1, r_1)} \subset W$, so $x^* \in G \cap W$.

Since $W$ was an arbitrary nonempty open set, $G$ is dense in $X$.

For (a) $\Leftrightarrow$ (b): suppose $X = \bigcup_n F_n$ with each $F_n$ nowhere
dense. Then $U_n = X \setminus \overline{F_n}$ is dense and open, but
$\bigcap_n U_n = X \setminus \bigcup_n \overline{F_n}
\subset X \setminus \bigcup_n F_n = \varnothing$,
contradicting (b). For (b) $\Leftrightarrow$ (c): if $X = \bigcup_n F_n$ with each
$F_n$ closed and $\mathrm{int}(F_n) = \varnothing$, then $U_n = X \setminus F_n$ is
dense and open and $\bigcap_n U_n = \varnothing$, contradicting (b).
\end{proof}

\begin{remark}\label{rem:baire-lcs}
\index{Baire category theorem!locally compact spaces}
The Baire category theorem also holds for locally compact Hausdorff spaces. The proof
is essentially the same, replacing completeness by local compactness: one chooses
the closed balls $\overline{B(x_n,r_n)}$ to be compact, and the nested intersection
of nonempty compact sets is nonempty by the finite intersection property.
\end{remark}

\begin{proposition}\label{prop:baire-consequences}
Let $X$ be a complete metric space.
\begin{enumerate}[label=(\roman*)]
  \item\index{Baire category theorem!consequences} If $X = \bigcup_{n=1}^{\infty} F_n$
    with each $F_n$ closed, then at least one $F_n$ has nonempty interior.
  \item Every residual subset of $X$ is dense.
  \item A nonempty complete metric space with no isolated points is uncountable.
\end{enumerate}
\end{proposition}

\begin{proof}
\begin{enumerate}[label=(\roman*)]
  \item This is the contrapositive of Theorem~\ref{thm:baire}(c).
  \item If $R$ is residual, then $X \setminus R$ is meagre, so $R$ contains a dense
    $G_\delta$ set by (b) of the theorem.
  \item If $X = \{x_1, x_2, \ldots\}$ were countable, then each singleton $\{x_n\}$
    is closed with empty interior (since $x_n$ is not isolated), so
    $X = \bigcup_n \{x_n\}$ is meagre, contradicting Theorem~\ref{thm:baire}(a).
    \qedhere
\end{enumerate}
\end{proof}

\begin{exercise}[$\Q$ is not a $G_\delta$; $\star$]\label{ex:baire-Q-not-Gdelta}
\index{$G_\delta$ set}\index{rational numbers!not a $G_\delta$}
Show that $\Q$ is \emph{not} a $G_\delta$ set in $\R$; that is, $\Q$ cannot be written
as a countable intersection of open sets.

\noindent\textit{Hint:} If $\Q = \bigcap_n U_n$ with each $U_n$ open and dense, then
$V_n = U_n \setminus \{q_n\}$ for an enumeration $\{q_n\}$ of $\Q$ would also be open
and dense, but $\bigcap_n V_n \cap \bigcap_n (X\setminus\{q_n\}) = \varnothing$,
which contradicts Baire.
\end{exercise}

\begin{exercise}[Osgood's theorem; $\star\star$]\label{ex:osgood}
\index{Osgood's theorem}
Let $f_n \colon [a,b] \to \R$ be a sequence of continuous functions converging
pointwise to $f \colon [a,b] \to \R$. Use the Baire category theorem to show that
$f$ is continuous on a dense $G_\delta$ subset of $[a,b]$.

\noindent\textit{Hint:} For each $k \geq 1$, let
$F_{k,m} = \{x \in [a,b] : \abs{f_n(x) - f_p(x)} \leq 1/k \text{ for all }
n,p \geq m\}$. Show that $\bigcup_m F_{k,m} = [a,b]$ and apply Baire.
\end{exercise}

%-----------------------------------------------------------------------
\section{The uniform boundedness principle}\label{sec:ubp-statement}
%-----------------------------------------------------------------------

We now turn to the main theorem of this chapter. Throughout, $X$ denotes a Banach
space\index{Banach space} and $Y$ a normed space.

\begin{theorem}[Banach--Steinhaus / Uniform Boundedness Principle]%
\label{thm:banach-steinhaus}
\index{Banach--Steinhaus theorem!statement}
\index{uniform boundedness principle!statement}
Let $X$ be a Banach space, $Y$ a normed space, and
$\{T_\alpha\}_{\alpha \in \mathcal{A}} \subset \mathcal{L}(X,Y)$\index{bounded
linear operator} a family of bounded linear operators. Suppose the family is
\textbf{pointwise bounded}\index{pointwise bounded family}: for every $x \in X$,
\[
  \sup_{\alpha \in \mathcal{A}} \norm{T_\alpha x}_Y < \infty.
\]
Then the family is \textbf{uniformly bounded}\index{uniformly bounded family}:
\[
  \sup_{\alpha \in \mathcal{A}} \norm{T_\alpha}_{\mathcal{L}(X,Y)} < \infty.
\]
\end{theorem}

\begin{proof}
\index{Banach--Steinhaus theorem!proof}
For each $n \in \N$, define the set
\[
  F_n = \bigl\{ x \in X : \sup_{\alpha \in \mathcal{A}} \norm{T_\alpha x}_Y
  \leq n \bigr\}
  = \bigcap_{\alpha \in \mathcal{A}} \{ x \in X : \norm{T_\alpha x}_Y \leq n \}.
\]
Each set $\{x : \norm{T_\alpha x}_Y \leq n\}$ is closed (it is the preimage of
$[0,n]$ under the continuous map $x \mapsto \norm{T_\alpha x}$), so $F_n$ is an
intersection of closed sets, hence closed.

The pointwise boundedness hypothesis gives $X = \bigcup_{n=1}^{\infty} F_n$.
Since $X$ is a Banach space (hence a complete metric space), the Baire category
theorem (Proposition~\ref{prop:baire-consequences}(i)) guarantees that some $F_{n_0}$
has nonempty interior. Hence there exist $x_0 \in X$ and $r > 0$ such that
$B(x_0, r) \subset F_{n_0}$, meaning
\[
  \norm{x - x_0} < r \implies \sup_{\alpha} \norm{T_\alpha x} \leq n_0.
\]

Now let $x \in X$ with $\norm{x} \leq 1$. Then $x_0 + rx/2 \in B(x_0, r)
\subset F_{n_0}$, and $x_0 \in F_{n_0}$. Therefore,
\[
  \norm{T_\alpha\!\bigl(\tfrac{r}{2} x\bigr)}
  = \norm{T_\alpha(x_0 + \tfrac{r}{2}x) - T_\alpha x_0}
  \leq \norm{T_\alpha(x_0 + \tfrac{r}{2}x)} + \norm{T_\alpha x_0}
  \leq n_0 + n_0 = 2n_0.
\]
Hence $\norm{T_\alpha x} \leq 4n_0/r$ for all $\alpha$ and all $\norm{x} \leq 1$.
This gives
\[
  \sup_{\alpha \in \mathcal{A}} \norm{T_\alpha} \leq \frac{4 n_0}{r} < \infty.
  \qedhere
\]
\end{proof}

\begin{remark}\label{rem:ubp-completeness-necessary}
\index{uniform boundedness principle!completeness is essential}
The completeness of $X$ is essential. Consider $X = c_{00}$\index{$c_{00}$ (finitely
supported sequences)}, the space of finitely
supported sequences with the $\ell^\infty$ norm. Define $T_n \colon c_{00} \to \K$ by
$T_n(x) = n\, x_n$. Then $\sup_n \abs{T_n(x)} < \infty$ for each $x \in c_{00}$
(since $x$ has finite support), but $\norm{T_n} = n \to \infty$.
\end{remark}

\begin{corollary}\label{cor:ubp-pointwise-limit}
\index{pointwise limit of operators}
Let $X$ be a Banach space, $Y$ a normed space, and $(T_n)_{n \geq 1} \subset
\mathcal{L}(X,Y)$ a sequence such that the limit
\[
  Tx = \lim_{n \to \infty} T_n x
\]
exists for every $x \in X$. Then:
\begin{enumerate}[label=(\roman*)]
  \item $\sup_n \norm{T_n} < \infty$.
  \item The limit operator $T \colon X \to Y$ is bounded, and
    $\norm{T} \leq \liminf_{n\to\infty} \norm{T_n}$.
\end{enumerate}
\end{corollary}

\begin{proof}
(i) A convergent sequence in $Y$ is bounded, so $\sup_n \norm{T_n x} < \infty$
for each $x$. The Banach--Steinhaus theorem gives $\sup_n \norm{T_n} = M < \infty$.

(ii) $T$ is linear (being the pointwise limit of linear maps). For any $x \in X$,
\[
  \norm{Tx} = \lim_{n\to\infty} \norm{T_n x}
  \leq \liminf_{n\to\infty} \norm{T_n}\, \norm{x},
\]
so $T$ is bounded with $\norm{T} \leq \liminf_n \norm{T_n}$.
\end{proof}

%-----------------------------------------------------------------------
\section{Applications to sequences of functionals}\label{sec:ubp-functionals}
%-----------------------------------------------------------------------

\begin{proposition}[Pointwise bounded functionals]\label{prop:ptwise-bounded-functionals}
\index{pointwise bounded!functionals}
Let $X$ be a Banach space and $(f_n)_{n \geq 1} \subset X^*$ a sequence of bounded
linear functionals. If $\sup_n \abs{f_n(x)} < \infty$ for every $x \in X$, then
$\sup_n \norm{f_n}_{X^*} < \infty$.
\end{proposition}

\begin{proof}
This is the special case $Y = \K$ of Theorem~\ref{thm:banach-steinhaus}.
\end{proof}

\begin{example}[Weak convergence implies boundedness]\label{ex:weak-conv-bounded}
\index{weak convergence!implies norm boundedness}
Let $X$ be a Banach space and suppose $x_n \weakto x$ in $X$ (i.e.,
$f(x_n) \to f(x)$ for all $f \in X^*$). Then $\sup_n \norm{x_n} < \infty$.

To see this, consider the canonical images $\hat{x}_n \in X^{**}$ defined by
$\hat{x}_n(f) = f(x_n)$. The hypothesis says that for each $f \in X^*$, the sequence
$(\hat{x}_n(f))_n$ is convergent, hence bounded. Since $X^*$ is a Banach space,
the uniform boundedness principle applied to $\{\hat{x}_n\} \subset
\mathcal{L}(X^*, \K)$ gives
\[
  \sup_n \norm{\hat{x}_n}_{X^{**}} = \sup_n \norm{x_n}_X < \infty,
\]
where the equality uses the isometric embedding $X \hookrightarrow X^{**}$.
\end{example}

\begin{proposition}[Weak$^*$ convergence implies boundedness]%
\label{prop:weakstar-conv-bounded}
\index{weak$^*$ convergence!implies norm boundedness}
Let $X$ be a Banach space and suppose $f_n \weakstarto f$ in $X^*$ (i.e.,
$f_n(x) \to f(x)$ for all $x \in X$). Then $\sup_n \norm{f_n}_{X^*} < \infty$
and $\norm{f} \leq \liminf_n \norm{f_n}$.
\end{proposition}

\begin{proof}
Apply Corollary~\ref{cor:ubp-pointwise-limit} with $T_n = f_n$ and $Y = \K$.
\end{proof}

%-----------------------------------------------------------------------
\section{Strong convergence versus uniform convergence}%
\label{sec:strong-vs-uniform}
%-----------------------------------------------------------------------

\begin{definition}\label{def:strong-uniform-convergence}
\index{strong operator convergence}\index{uniform operator convergence}
Let $X,Y$ be normed spaces and $(T_n)_{n \geq 1} \subset \mathcal{L}(X,Y)$.
\begin{enumerate}[label=(\roman*)]
  \item $(T_n)$ converges \textbf{strongly} (or \textbf{pointwise}) to
    $T \in \mathcal{L}(X,Y)$ if $T_n x \to Tx$ for every $x \in X$.
    We write $T_n \xrightarrow{s} T$.
  \item $(T_n)$ converges \textbf{uniformly} (or \textbf{in norm}) to $T$ if
    $\norm{T_n - T} \to 0$. We write $T_n \to T$.
\end{enumerate}
\end{definition}

\begin{remark}\label{rem:uniform-implies-strong}
Uniform convergence implies strong convergence, since
$\norm{T_n x - Tx} \leq \norm{T_n - T}\norm{x}$. The converse is false in general.
\end{remark}

\begin{example}\label{ex:strong-not-uniform}
\index{strong operator convergence!vs.\ uniform convergence}
On $X = \ell^2(\N)$, let $P_n$ be the orthogonal projection onto
$\spn\{e_1, \ldots, e_n\}$. Then $P_n \xrightarrow{s} \Id$ (since every $x \in
\ell^2$ is the sum of its Fourier series), but $\norm{P_n - \Id} = 1$ for all $n$,
so the convergence is not uniform.
\end{example}

\begin{theorem}[Banach--Steinhaus for strong convergence]%
\label{thm:bs-strong-convergence}
\index{Banach--Steinhaus theorem!strong convergence}
Let $X$ be a Banach space, $Y$ a normed space, and $(T_n) \subset \mathcal{L}(X,Y)$
a sequence converging strongly to $T$. Then:
\begin{enumerate}[label=(\roman*)]
  \item $\sup_n \norm{T_n} < \infty$ and $T \in \mathcal{L}(X,Y)$.
  \item If $x_n \to x$ in $X$, then $T_n x_n \to Tx$ in $Y$.
\end{enumerate}
\end{theorem}

\begin{proof}
(i) follows from Corollary~\ref{cor:ubp-pointwise-limit}. Let $M = \sup_n \norm{T_n}$.

(ii) We write
\[
  T_n x_n - Tx = T_n(x_n - x) + (T_n x - Tx).
\]
For the first term: $\norm{T_n(x_n - x)} \leq M \norm{x_n - x} \to 0$.
For the second term: $T_n x \to Tx$ by hypothesis. Hence $T_n x_n \to Tx$.
\end{proof}

\begin{exercise}[Strong convergence and equicontinuity; $\star$]%
\label{ex:strong-equicontinuous}
\index{equicontinuity}
Let $X$ be a Banach space, $Y$ a normed space, and $(T_n) \subset \mathcal{L}(X,Y)$
with $\sup_n \norm{T_n} \leq M$. Show that $T_n \xrightarrow{s} T$ if and only if
$T_n x \to Tx$ for all $x$ in a dense subset of $X$.
\end{exercise}

%-----------------------------------------------------------------------
\section{Application: divergent Fourier series}\label{sec:dubois-reymond}
%-----------------------------------------------------------------------

One of the most striking applications of the uniform boundedness principle is the
proof that there exist continuous functions whose Fourier series diverges at a
given point. This result, due to du Bois-Reymond\index{du Bois-Reymond} (1876),
was one of the motivations for the development of the Baire category approach.

\subsection{The Dirichlet kernel and partial sums}\label{subsec:dirichlet-kernel}

Let $f \colon [-\pi, \pi] \to \C$ be a $2\pi$-periodic integrable function.
The $N$-th partial sum\index{Fourier series!partial sums} of its Fourier series at
the point $t = 0$ is
\[
  S_N f(0) = \sum_{k=-N}^{N} \hat{f}(k),
  \qquad \hat{f}(k) = \frac{1}{2\pi} \int_{-\pi}^{\pi} f(t)\, e^{-ikt}\, dt.
\]
One computes
\[
  S_N f(0) = \frac{1}{2\pi} \int_{-\pi}^{\pi} f(t)\, D_N(t)\, dt,
\]
where $D_N$\index{Dirichlet kernel} is the \textbf{Dirichlet kernel}:
\[
  D_N(t) = \sum_{k=-N}^{N} e^{ikt} = \frac{\sin\bigl((N+\tfrac{1}{2})t\bigr)}%
  {\sin(t/2)}.
\]

\begin{lemma}\label{lem:dirichlet-L1-norm}
\index{Dirichlet kernel!$L^1$ norm}
The $L^1$ norm of the Dirichlet kernel satisfies
\[
  \frac{1}{2\pi}\int_{-\pi}^{\pi} \abs{D_N(t)}\, dt
  = \frac{4}{\pi^2} \ln N + O(1) \quad \text{as } N \to \infty.
\]
In particular, $\norm{D_N}_{L^1} \to \infty$.
\end{lemma}

\begin{proof}
We have
\begin{align*}
  \frac{1}{2\pi}\int_{-\pi}^{\pi} \abs{D_N(t)}\, dt
  &= \frac{1}{\pi}\int_{0}^{\pi}
    \frac{\abs{\sin\bigl((N+\tfrac{1}{2})t\bigr)}}{\sin(t/2)}\, dt \\
  &\geq \frac{1}{\pi}\int_{0}^{\pi}
    \frac{\abs{\sin\bigl((N+\tfrac{1}{2})t\bigr)}}{t/2}\, dt
    \quad (\text{since } \sin(t/2) \leq t/2) \\
  &= \frac{2}{\pi}\int_{0}^{(N+1/2)\pi}
    \frac{\abs{\sin u}}{u}\, du
    \quad (u = (N+\tfrac{1}{2})t) \\
  &\geq \frac{2}{\pi}\sum_{k=1}^{N} \int_{(k-1)\pi}^{k\pi}
    \frac{\abs{\sin u}}{u}\, du
  \geq \frac{2}{\pi}\sum_{k=1}^{N} \frac{1}{k\pi}
    \int_{0}^{\pi} \sin u\, du
  = \frac{4}{\pi^2}\sum_{k=1}^{N} \frac{1}{k}.
\end{align*}
Since $\sum_{k=1}^N 1/k = \ln N + O(1)$, we get the lower bound. A similar
calculation gives a matching upper bound.
\end{proof}

\subsection{The du Bois-Reymond theorem}

Consider the Banach space $X = C([-\pi,\pi])$\index{$C([-\pi,\pi])$} of continuous
$2\pi$-periodic functions with the supremum norm. For each $N$, the $N$-th partial
sum operator at $t = 0$ defines a bounded linear functional:
\[
  \Lambda_N \colon C([-\pi,\pi]) \to \C,
  \qquad \Lambda_N(f) = S_N f(0) = \frac{1}{2\pi}
  \int_{-\pi}^{\pi} f(t)\, D_N(t)\, dt.
\]
Its operator norm is
\[
  \norm{\Lambda_N} = \frac{1}{2\pi}\int_{-\pi}^{\pi} \abs{D_N(t)}\, dt
  = \frac{4}{\pi^2}\ln N + O(1) \to \infty.
\]

\begin{theorem}[du Bois-Reymond]\label{thm:dubois-reymond}
\index{du Bois-Reymond theorem}\index{Fourier series!divergence}
There exists a continuous $2\pi$-periodic function $f$ whose Fourier series
diverges at $t = 0$, i.e., $\sup_N \abs{S_N f(0)} = \infty$.

More precisely, the set of functions $f \in C([-\pi,\pi])$ for which
$\sup_N \abs{S_N f(0)} = \infty$ is a dense $G_\delta$ subset of
$C([-\pi,\pi])$; in particular, it is residual.
\end{theorem}

\begin{proof}
\index{du Bois-Reymond theorem!proof}
We apply the \emph{contrapositive} of the Banach--Steinhaus theorem
(Theorem~\ref{thm:banach-steinhaus}). Consider the functionals
$\Lambda_N \in C([-\pi,\pi])^*$ defined above. We have shown that
$\norm{\Lambda_N} \to \infty$.

If $\sup_N \abs{\Lambda_N(f)} < \infty$ were to hold for \emph{every}
$f \in C([-\pi,\pi])$, then the uniform boundedness principle would give
$\sup_N \norm{\Lambda_N} < \infty$, a contradiction. Hence there exists
$f \in C([-\pi,\pi])$ with $\sup_N \abs{S_N f(0)} = \infty$.

For the second statement, define $G_k = \{f \in C([-\pi,\pi]) :
\sup_N \abs{S_N f(0)} > k\}$. Each $G_k$ is open (as a union of the open sets
$\{\abs{\Lambda_N(f)} > k\}$ over $N$), and we need to show each $G_k$ is dense.
If $G_k$ were not dense for some $k$, there would exist a ball
$\overline{B(g,\varepsilon)}$ with $\sup_N \abs{\Lambda_N(f)} \leq k$ for all
$f \in \overline{B(g,\varepsilon)}$. Then for $\norm{h} \leq 1$:
\[
  \abs{\Lambda_N(\varepsilon h)} \leq \abs{\Lambda_N(g + \varepsilon h)}
  + \abs{\Lambda_N(g)} \leq 2k,
\]
giving $\norm{\Lambda_N} \leq 2k/\varepsilon$ for all $N$, contradicting
$\norm{\Lambda_N} \to \infty$. Hence each $G_k$ is open and dense, and
$\bigcap_k G_k = \{f : \sup_N \abs{S_N f(0)} = \infty\}$ is a dense $G_\delta$
by the Baire category theorem.
\end{proof}

\begin{remark}\label{rem:dubois-reymond-generic}
\index{Fourier series!generic divergence}
The theorem shows that ``most'' continuous functions (in the Baire category sense)
have divergent Fourier series at any given point. This stands in sharp contrast to
Carleson's celebrated theorem\index{Carleson's theorem} (1966) that the Fourier
series of any $L^2$ function converges almost everywhere.
\end{remark}

%-----------------------------------------------------------------------
\section{The Gibbs phenomenon}\label{sec:gibbs}
%-----------------------------------------------------------------------

\index{Gibbs phenomenon|(}
The Gibbs phenomenon describes the overshoot of partial Fourier sums near a jump
discontinuity. While not directly a consequence of the uniform boundedness principle,
it illustrates the behavior of the Dirichlet kernel and complements the discussion
of Fourier series.

\begin{example}[Gibbs phenomenon for the square wave]\label{ex:gibbs-square}
Consider the $2\pi$-periodic function
\[
  f(t) = \begin{cases} 1 & \text{if } 0 < t < \pi, \\
  -1 & \text{if } -\pi < t < 0, \\ 0 & \text{if } t \in \{0, \pm\pi\}.\end{cases}
\]
Its Fourier series is
$f(t) \sim \frac{4}{\pi}\sum_{k=0}^{\infty} \frac{\sin(2k+1)t}{2k+1}$.
The $N$-th partial sum at $t = \pi/(2N+1)$ satisfies
\[
  S_N f\!\left(\frac{\pi}{2N+1}\right) \to
  \frac{2}{\pi}\int_0^{\pi} \frac{\sin u}{u}\, du
  = \frac{2}{\pi}\,\mathrm{Si}(\pi) \approx 1.17898,
\]
where $\mathrm{Si}$\index{sine integral} is the sine integral. This represents an
overshoot of approximately $8.95\%$ above the function value $1$, and this overshoot
does not diminish as $N \to \infty$ --- it merely moves closer to the discontinuity.
\end{example}

\begin{proposition}\label{prop:gibbs-overshoot}
\index{Gibbs phenomenon!overshoot}
Let $f$ be a piecewise smooth $2\pi$-periodic function with a jump discontinuity at
$t_0$. Then the partial sums $S_N f$ exhibit an overshoot near $t_0$ whose magnitude
approaches
\[
  \frac{f(t_0^+) - f(t_0^-)}{2} \cdot
  \left(\frac{2}{\pi}\int_0^{\pi}\frac{\sin u}{u}\,du - 1\right)
  \approx 0.0895 \cdot \bigl(f(t_0^+) - f(t_0^-)\bigr).
\]
\end{proposition}
\index{Gibbs phenomenon|)}

%-----------------------------------------------------------------------
\section{The Banach--Steinhaus condensation theorem}%
\label{sec:condensation}
%-----------------------------------------------------------------------

The following refinement of the uniform boundedness principle describes the structure
of the ``exceptional set'' where pointwise boundedness fails.

\begin{theorem}[Condensation of singularities / Banach--Steinhaus]%
\label{thm:condensation}
\index{condensation of singularities}\index{Banach--Steinhaus theorem!condensation}
Let $X$ be a Banach space, $Y$ a normed space, and let
$(T_{n,m})_{n,m \geq 1} \subset \mathcal{L}(X,Y)$ be a double sequence of bounded
linear operators. Suppose that for each $n \geq 1$,
\[
  \sup_{m \geq 1} \norm{T_{n,m}} = \infty.
\]
Then the set
\[
  \mathcal{R} = \bigl\{ x \in X : \sup_{m \geq 1} \norm{T_{n,m} x} = \infty
  \text{ for all } n \geq 1 \bigr\}
\]
is a dense $G_\delta$ subset of $X$ (hence residual).
\end{theorem}

\begin{proof}
\index{condensation of singularities!proof}
For each $n, k \geq 1$, define
\[
  G_{n,k} = \bigl\{ x \in X : \exists\, m \text{ with }
  \norm{T_{n,m} x} > k \bigr\}
  = \bigcup_{m=1}^{\infty} \bigl\{ x \in X : \norm{T_{n,m} x} > k \bigr\}.
\]
Each set $\{x : \norm{T_{n,m}x} > k\}$ is open (preimage of $(k,\infty)$ under
a continuous map), so $G_{n,k}$ is open.

\textbf{Claim:} $G_{n,k}$ is dense for each $n, k$.

Suppose not: there exist $n_0$, $k_0$, $x_0 \in X$, and $r > 0$ with
$B(x_0, r) \cap G_{n_0,k_0} = \varnothing$. This means $\norm{T_{n_0,m}x} \leq k_0$
for all $m$ and all $x \in B(x_0, r)$. For $\norm{h} < r$:
\[
  \norm{T_{n_0,m} h} \leq \norm{T_{n_0,m}(x_0 + h)} + \norm{T_{n_0,m} x_0}
  \leq 2 k_0,
\]
so $\norm{T_{n_0,m}} \leq 2k_0/r$ for all $m$, contradicting
$\sup_m \norm{T_{n_0,m}} = \infty$.

Now observe that
\[
  \mathcal{R} = \bigcap_{n=1}^{\infty} \bigcap_{k=1}^{\infty} G_{n,k}.
\]
This is a countable intersection of dense open sets. By the Baire category theorem
(Theorem~\ref{thm:baire}), $\mathcal{R}$ is a dense $G_\delta$.
\end{proof}

\begin{corollary}\label{cor:condensation-fourier}
\index{Fourier series!condensation of divergence}
There exists a continuous $2\pi$-periodic function whose Fourier series diverges at
every rational multiple of $\pi$. More generally, given any countable set
$S = \{t_n\}_{n \geq 1} \subset [-\pi, \pi]$, the set of functions
$f \in C([-\pi,\pi])$ whose Fourier series diverges simultaneously at every
$t_n \in S$ is residual in $C([-\pi,\pi])$.
\end{corollary}

\begin{proof}
For each $n$, let $T_{n,m}(f) = S_m f(t_n)$ be the $m$-th partial Fourier sum at
$t_n$. Then $\norm{T_{n,m}} = \frac{1}{2\pi}\int_{-\pi}^{\pi}\abs{D_m(t)}\,dt
\to \infty$ as $m \to \infty$. The condensation theorem
(Theorem~\ref{thm:condensation}) gives the result.
\end{proof}

%-----------------------------------------------------------------------
\section{Uniform boundedness for families of operators}%
\label{sec:ubp-families}
%-----------------------------------------------------------------------

The uniform boundedness principle extends naturally to arbitrary (not necessarily
countable) families of operators.

\begin{theorem}[Uniform boundedness for operator families]%
\label{thm:ubp-families}
\index{uniform boundedness principle!for operator families}
Let $X$ be a Banach space, $Y$ a Banach space, and let
$\{T_\alpha\}_{\alpha \in \mathcal{A}} \subset \mathcal{L}(X,Y)$ and
$\{S_\beta\}_{\beta \in \mathcal{B}} \subset \mathcal{L}(Y,\K)$ be families of
operators such that
\[
  \sup_{\alpha \in \mathcal{A},\; \beta \in \mathcal{B}}
  \abs{S_\beta(T_\alpha x)} < \infty
  \qquad \text{for every } x \in X.
\]
Then $\sup_{\alpha,\beta} \norm{S_\beta \circ T_\alpha} < \infty$.
\end{theorem}

\begin{proof}
For each fixed $x \in X$, the family $\{S_\beta \circ T_\alpha\}$ is pointwise
bounded on $X$. First, for each $\alpha$, the family $\{S_\beta(T_\alpha x)\}_\beta$
is bounded. Since $Y$ is a Banach space, the Banach--Steinhaus theorem
applied to $\{S_\beta\}_\beta \subset Y^*$ gives
$C_\alpha := \sup_\beta \norm{S_\beta(T_\alpha x)} < \infty$ for each $x$, and
moreover $\sup_\beta \norm{S_\beta} < \infty$ (call this bound $K$).

Now $\sup_\alpha \norm{T_\alpha x} \leq \sup_\alpha \sup_\beta
\abs{S_\beta(T_\alpha x)} / \inf_{\norm{S_\beta}\neq 0} \ldots$ --- but more
directly, since for each $x$, $\sup_\alpha \sup_\beta \abs{S_\beta(T_\alpha x)}
< \infty$, we have in particular $\sup_\alpha \norm{T_\alpha x}_Y < \infty$ (by
Hahn--Banach applied in $Y$). The Banach--Steinhaus theorem then gives
$\sup_\alpha \norm{T_\alpha} < \infty$, say $\leq M$. Hence
$\norm{S_\beta \circ T_\alpha} \leq \norm{S_\beta}\norm{T_\alpha} \leq KM$.
\end{proof}

\begin{proposition}[Principle of condensation for bilinear forms]%
\label{prop:bilinear-ubp}
\index{bilinear form!uniform boundedness}
Let $X, Y$ be Banach spaces and $B \colon X \times Y \to \K$ a separately
continuous bilinear form. Then $B$ is jointly bounded:
\[
  \abs{B(x,y)} \leq M \norm{x}\norm{y}
  \qquad \text{for some } M > 0 \text{ and all } x \in X,\; y \in Y.
\]
\end{proposition}

\begin{proof}
\index{bilinear form!joint boundedness proof}
For each $x \in X$, the map $y \mapsto B(x,y)$ is a bounded linear functional on
$Y$, so there exists $T_x \in Y^*$ with $T_x(y) = B(x,y)$ and
$\norm{T_x} = \sup_{\norm{y}\leq 1} \abs{B(x,y)}$. The map $x \mapsto T_x$ is
a linear operator $T \colon X \to Y^*$.

For each $y$, the map $x \mapsto B(x,y)$ is bounded, so $\abs{B(x,y)} \leq
C_y \norm{x}$. In particular, $\sup_{\norm{x}\leq 1}\abs{T_x(y)} \leq C_y < \infty$.
Since $Y$ is a Banach space, the uniform boundedness principle (applied to
$\{T_x\}_{\norm{x}\leq 1} \subset Y^*$) gives that $x \mapsto T_x$ maps the unit
ball of $X$ to a bounded subset of $Y^*$. But this means $T$ is bounded, say
$\norm{T} \leq M$, and
\[
  \abs{B(x,y)} = \abs{T_x(y)} \leq \norm{T_x}\norm{y} \leq M\norm{x}\norm{y}.
  \qedhere
\]
\end{proof}

%-----------------------------------------------------------------------
\section{Exercises for Chapter~\ref{ch:ubp}}\label{sec:ubp-exercises}
%-----------------------------------------------------------------------

\begin{exercise}[Nowhere dense subsets; $\star$]\label{ex:nowhere-dense-union}
\index{nowhere dense set!exercises}
Show that a finite union of nowhere dense sets is nowhere dense. Give an example
showing that a countable union of nowhere dense sets need not be nowhere dense.
\end{exercise}

\begin{exercise}[Baire category in $\ell^p$; $\star$]\label{ex:baire-ell-p}
\index{Baire category theorem!in $\ell^p$}
Let $1 \leq p < \infty$. Show that $\ell^q \subset \ell^p$ is meagre in $\ell^p$
for $q < p$. \textit{Hint:} Write $\ell^q = \bigcup_n \{x \in \ell^p :
\norm{x}_q \leq n\}$.
\end{exercise}

\begin{exercise}[Pointwise but not uniform convergence; $\star$]\label{ex:ptwise-not-uniform}
\index{strong operator convergence!exercises}
Let $T_n \colon \ell^2 \to \ell^2$ be the right shift applied $n$ times:
$T_n(x_1, x_2, \ldots) = (0, \ldots, 0, x_1, x_2, \ldots)$ with $n$ leading zeros.
Show that $T_n \xrightarrow{s} 0$ but $\norm{T_n} = 1$ for all $n$. Reconcile this
with the Banach--Steinhaus theorem.
\end{exercise}

\begin{exercise}[Divergence on a dense set; $\star\star$]\label{ex:divergence-dense}
\index{Fourier series!exercises}
Let $X$ be a Banach space and $(T_n) \subset X^*$ with $\sup_n \norm{T_n} = \infty$.
Show that the set $\{x \in X : \sup_n \abs{T_n(x)} = \infty\}$ is dense in $X$.
Can you show it is residual?
\end{exercise}

\begin{exercise}[Unbounded Toeplitz matrix; $\star\star$]\label{ex:toeplitz}
\index{Toeplitz matrix}
A matrix $A = (a_{nk})_{n,k \geq 0}$ defines a \emph{summability method}: given
a sequence $(s_k)$, we set $\sigma_n = \sum_k a_{nk} s_k$ (when convergent). Show
that if the Toeplitz conditions hold ($\lim_n a_{nk} = 0$ for each $k$,
$\sup_n \sum_k \abs{a_{nk}} < \infty$, $\lim_n \sum_k a_{nk} = 1$), then
$\sigma_n \to s$ whenever $s_k \to s$. \textit{Hint:} Apply the Banach--Steinhaus
theorem to functionals on $c$ (the space of convergent sequences).
\end{exercise}

\begin{exercise}[Principle of condensation application; $\star\star$]\label{ex:condensation-power-series}
\index{condensation of singularities!exercises}
Let $(a_n)_{n \geq 0}$ be a sequence of complex numbers such that $\sum_n a_n z^n$
has radius of convergence $1$. Use the condensation theorem to show that the set
of points $e^{i\theta}$ on the unit circle where $\sum_n a_n e^{in\theta}$ diverges
is either empty or residual in the unit circle.
\end{exercise}

\begin{exercise}[Continuous nowhere differentiable functions; $\star\star\star$]%
\label{ex:nowhere-diff}
\index{nowhere differentiable function}\index{Weierstrass function}
Use the Baire category theorem to show that the set of continuous functions
$f \in C([0,1])$ that are nowhere differentiable is residual in $C([0,1])$.

\noindent\textit{Hint:} For each $n \geq 1$, let $A_n$ be the set of $f \in C([0,1])$
such that there exists $x \in [0,1]$ with
$\abs{f(x+h) - f(x)} \leq n\abs{h}$ for all sufficiently small $h$. Show that each
$A_n$ is closed and has empty interior.
\end{exercise}

\begin{exercise}[Banach--Steinhaus for Hilbert spaces; $\star$]\label{ex:bs-hilbert}
\index{Banach--Steinhaus theorem!Hilbert space version}
Let $H$ be a Hilbert space and $(T_n) \subset \mathcal{L}(H)$ a sequence of
self-adjoint operators with $\sup_n \abs{\inner{T_n x}{x}} < \infty$ for every
$x \in H$. Show that $\sup_n \norm{T_n} < \infty$. \textit{Hint:} Use the
polarization identity and the fact that $\norm{T} = \sup_{\norm{x}=1}
\abs{\inner{Tx}{x}}$ for self-adjoint $T$.
\end{exercise}

\begin{exercise}[Weak$^*$ sequential completeness; $\star\star$]\label{ex:weakstar-complete}
\index{weak$^*$ topology!sequential completeness}
Let $X$ be a Banach space and $(f_n) \subset X^*$ a weak$^*$ Cauchy sequence (i.e.,
$(f_n(x))$ is Cauchy in $\K$ for each $x \in X$). Show that there exists $f \in X^*$
such that $f_n \weakstarto f$. Is the analogous statement true for weak Cauchy
sequences in $X$?
\end{exercise}

\begin{exercise}[Superdense orbits; $\star\star\star$]\label{ex:superdense}
\index{hypercyclic operator}
Let $X$ be a separable Banach space and $T \in \mathcal{L}(X)$ a
\emph{hypercyclic}\index{hypercyclic operator} operator (there exists $x \in X$ such
that $\{T^n x : n \geq 0\}$ is dense in $X$). Use the Baire category theorem to
show that the set of hypercyclic vectors for $T$ is either empty or a dense $G_\delta$.
\end{exercise}

\begin{exercise}[Uniform boundedness and integration; $\star\star$]%
\label{ex:ubp-integration}
\index{uniform boundedness principle!integration}
Let $(f_n)$ be a sequence in $L^1([0,1])$ such that for every $g \in L^\infty([0,1])$,
the sequence $\left(\int_0^1 f_n g\, dx\right)$ is bounded. Show that
$\sup_n \norm{f_n}_{L^1} < \infty$.
\end{exercise}


%%%%%%%%%%%%%%%%%%%%%%%%%%%%%%%%%%%%%%%%%%%%%%%%%%%%%%%%%%%%%%%%%%%%%%%
%% CHAPTER 6 — The Open Mapping Theorem and Closed Graph Theorem
%%%%%%%%%%%%%%%%%%%%%%%%%%%%%%%%%%%%%%%%%%%%%%%%%%%%%%%%%%%%%%%%%%%%%%%
\chapter{The Open Mapping Theorem and Closed Graph Theorem}\label{ch:omt-cgt}
\minitoc

\vspace{1em}
\noindent
This chapter is devoted to two further pillars of functional analysis: the
\emph{open mapping theorem}\index{open mapping theorem} (also called the
\emph{Banach--Schauder theorem}\index{Banach--Schauder theorem}) and the
\emph{closed graph theorem}\index{closed graph theorem}. Like the uniform boundedness
principle, both rest upon the Baire category theorem and require completeness in an
essential way. Together with the Hahn--Banach theorem and the uniform boundedness
principle, these form the ``big four'' theorems\index{big four theorems} that
underpin much of linear functional analysis.

%-----------------------------------------------------------------------
\section{The open mapping theorem}\label{sec:omt}
%-----------------------------------------------------------------------

\subsection{Preliminary: the open mapping lemma}

\begin{definition}\label{def:open-map}
\index{open mapping}
A map $T \colon X \to Y$ between topological spaces is called \textbf{open} if
it maps open sets to open sets: $U \subset X$ open implies $T(U) \subset Y$ open.
\end{definition}

The key technical step is the following lemma, which uses the Baire category theorem.

\begin{lemma}[Open mapping lemma]\label{lem:open-mapping-lemma}
\index{open mapping lemma}
Let $X$ and $Y$ be Banach spaces and $T \in \mathcal{L}(X,Y)$ a surjective
bounded linear operator. Then there exists $\delta > 0$ such that
\[
  B_Y(0, \delta) \subset \overline{T(B_X(0,1))}.
\]
That is, the closure of the image of the open unit ball of $X$ contains an open
ball in $Y$.
\end{lemma}

\begin{proof}
\index{open mapping lemma!proof}
Since $T$ is surjective, $Y = T(X) = \bigcup_{n=1}^{\infty} T(B_X(0,n))
= \bigcup_{n=1}^{\infty} n\, T(B_X(0,1))$. Taking closures,
$Y = \bigcup_{n=1}^{\infty} n\, \overline{T(B_X(0,1))}$.

Since $Y$ is a Banach space (hence a complete metric space), the Baire category
theorem (Proposition~\ref{prop:baire-consequences}(i)) implies that some
$n_0\, \overline{T(B_X(0,1))}$ has nonempty interior. Since multiplication by $n_0$
is a homeomorphism, $\overline{T(B_X(0,1))}$ itself has nonempty interior. Hence
there exist $y_0 \in Y$ and $r > 0$ with $B_Y(y_0, r) \subset
\overline{T(B_X(0,1))}$.

Since $\overline{T(B_X(0,1))}$ is symmetric (if $y \in \overline{T(B_X(0,1))}$ then
$-y \in \overline{T(B_X(0,1))}$, because $B_X(0,1)$ is symmetric) and convex
(since $T(B_X(0,1))$ is convex), we can average:
\[
  B_Y(0, r) = \tfrac{1}{2}\bigl(B_Y(y_0,r) + B_Y(-y_0,r)\bigr)
  \subset \tfrac{1}{2}\bigl(\overline{T(B_X(0,1))} + \overline{T(B_X(0,1))}\bigr)
  \subset \overline{T(B_X(0,1))}.
\]
The last inclusion uses the convexity: if $u, v \in \overline{T(B_X(0,1))}$, then
$\frac{u+v}{2} \in \overline{T(B_X(0,1))}$ since $B_X(0,1)$ is convex and $T$ is
linear. Set $\delta = r$.
\end{proof}

\subsection{The main theorem}

\begin{theorem}[Open Mapping Theorem / Banach--Schauder]\label{thm:omt}
\index{open mapping theorem!statement and proof}
\index{Banach--Schauder theorem|see{open mapping theorem}}
Let $X$ and $Y$ be Banach spaces and $T \in \mathcal{L}(X,Y)$ surjective.
Then $T$ is an open mapping. More precisely, there exists $\delta > 0$ such that
\[
  B_Y(0, \delta) \subset T(B_X(0,1)).
\]
\end{theorem}

\begin{proof}
\index{open mapping theorem!proof}
\textbf{Step 1 (Closure contains a ball).}
By Lemma~\ref{lem:open-mapping-lemma}, there exists $\delta > 0$ such that
$B_Y(0, \delta) \subset \overline{T(B_X(0,1))}$. By scaling, for every $\varepsilon > 0$,
\begin{equation}\label{eq:omt-closure}
  B_Y(0, \varepsilon \delta) \subset \overline{T(B_X(0,\varepsilon))}.
\end{equation}

\textbf{Step 2 (Image itself contains a ball --- the approximation argument).}
We show that $B_Y(0, \delta) \subset T(B_X(0, 2))$. Let $y \in B_Y(0, \delta)$.
We construct a sequence $(x_n)$ in $X$ such that
\[
  \left\| y - \sum_{k=0}^{n} T x_k \right\| < \frac{\delta}{2^{n+1}},
  \qquad \norm{x_n} < \frac{1}{2^n}.
\]

\textit{Base case:} By~\eqref{eq:omt-closure} with $\varepsilon = 1$, the element
$y \in B_Y(0,\delta) \subset \overline{T(B_X(0,1))}$, so there exists $x_0 \in X$
with $\norm{x_0} < 1$ and $\norm{y - Tx_0} < \delta/2$.

\textit{Inductive step:} Having chosen $x_0, \ldots, x_n$, the residual
$y - \sum_{k=0}^n Tx_k \in B_Y(0, \delta/2^{n+1})$. By~\eqref{eq:omt-closure}
with $\varepsilon = 1/2^{n+1}$, there exists $x_{n+1}$ with
$\norm{x_{n+1}} < 1/2^{n+1}$ and
$\norm{y - \sum_{k=0}^{n+1} Tx_k} < \delta/2^{n+2}$.

\textbf{Step 3 (Convergence).}
Set $s_n = \sum_{k=0}^n x_k$. Since $\sum_{k=0}^\infty \norm{x_k}
< \sum_{k=0}^\infty 2^{-k} = 2$, the series converges absolutely, so $(s_n)$ is
Cauchy in the Banach space $X$. Let $x = \sum_{k=0}^\infty x_k$; then
$\norm{x} < 2$. By continuity of $T$,
\[
  Tx = T\!\left(\sum_{k=0}^\infty x_k\right) = \lim_{n\to\infty}
  \sum_{k=0}^n Tx_k = y.
\]
Hence $y = Tx$ with $x \in B_X(0,2)$, proving $B_Y(0,\delta) \subset T(B_X(0,2))$.

Replacing by $B_X(0,1)$ via scaling: $B_Y(0, \delta/2) \subset T(B_X(0,1))$.

\textbf{Step 4 (Openness).}
Let $U \subset X$ be open and $y_0 = Tx_0 \in T(U)$. Choose $r > 0$ with
$B_X(x_0, r) \subset U$. Then
\[
  T(U) \supset T(B_X(x_0,r)) = Tx_0 + T(B_X(0,r))
  \supset y_0 + B_Y(0, r\delta/2),
\]
so $T(U)$ contains an open ball around each of its points, hence $T(U)$ is open.
\end{proof}

\begin{remark}\label{rem:omt-completeness}
\index{open mapping theorem!completeness hypotheses}
The completeness of \emph{both} $X$ and $Y$ is used: completeness of $Y$ in the
Baire category argument (Step 1), and completeness of $X$ for the convergence of
the series in Step 3. Without completeness, the theorem can fail: consider the
identity map $\Id \colon (C([0,1]), \norm{\cdot}_{C^1}) \to (C([0,1]),
\norm{\cdot}_\infty)$, which is continuous and bijective but not open.
\end{remark}

%-----------------------------------------------------------------------
\section{The Banach isomorphism theorem}\label{sec:banach-iso}
%-----------------------------------------------------------------------

\begin{corollary}[Banach isomorphism theorem]\label{cor:banach-iso}
\index{Banach isomorphism theorem}\index{inverse mapping theorem}
Let $X$ and $Y$ be Banach spaces and $T \in \mathcal{L}(X,Y)$ a continuous
linear bijection. Then $T^{-1} \in \mathcal{L}(Y,X)$; that is, $T^{-1}$ is
automatically continuous. Hence $T$ is a topological isomorphism.
\end{corollary}

\begin{proof}
Since $T$ is a bijection, $T^{-1} \colon Y \to X$ is a well-defined linear map.
By the open mapping theorem (Theorem~\ref{thm:omt}), $T$ is open, so for any
open $U \subset X$, $T(U)$ is open in $Y$. But $(T^{-1})^{-1}(U) = T(U)$, so
$T^{-1}$ is continuous.
\end{proof}

\begin{example}\label{ex:equivalent-norms-banach}
\index{equivalent norms!Banach space criterion}
Let $X$ be a vector space equipped with two norms $\norm{\cdot}_1$ and
$\norm{\cdot}_2$ such that $(X, \norm{\cdot}_1)$ and $(X, \norm{\cdot}_2)$ are
both Banach spaces. If there exists $C > 0$ with $\norm{x}_2 \leq C\norm{x}_1$
for all $x$, then the norms are equivalent, i.e., there exists $c > 0$ with
$c\norm{x}_1 \leq \norm{x}_2$.

Indeed, the identity $\Id \colon (X, \norm{\cdot}_1) \to (X, \norm{\cdot}_2)$ is
continuous and bijective. By the Banach isomorphism theorem, $\Id^{-1}$ is
continuous, giving the reverse inequality.
\end{example}

\begin{example}\label{ex:banach-iso-ell1}
\index{Banach isomorphism theorem!applications}
Consider $\ell^1(\N)$ with its standard norm and the map $T \colon \ell^1 \to c_0$
defined by $T(x)_n = \sum_{k=n}^{\infty} x_k$ (partial tail sums). One can verify
that $T$ is a bounded linear bijection. The Banach isomorphism theorem guarantees
that $T^{-1}$ is bounded, even though an explicit formula for $T^{-1}$ may be
complicated: $(T^{-1}y)_n = y_n - y_{n+1}$.
\end{example}

\begin{proposition}[Bounded below implies open range]%
\label{prop:bounded-below-open}
\index{bounded below operator}\index{open range}
Let $X$ be a Banach space, $Y$ a normed space, and $T \in \mathcal{L}(X,Y)$.
The following are equivalent:
\begin{enumerate}[label=(\roman*)]
  \item $T$ is bounded below: there exists $c > 0$ with $\norm{Tx} \geq c\norm{x}$
    for all $x \in X$.
  \item $T$ is injective and $T(X)$ is closed in $Y$, and $T^{-1} \colon T(X) \to X$
    is bounded.
\end{enumerate}
If in addition $Y$ is a Banach space and $T$ is surjective, then (i) and (ii)
are equivalent to $T$ being a topological isomorphism.
\end{proposition}

\begin{proof}
(i) $\Rightarrow$ (ii): Boundedness below gives injectivity ($Tx = 0 \Rightarrow
x = 0$) and $\norm{T^{-1}} \leq 1/c$. To see $T(X)$ is closed, let
$Tx_n \to y$. Then $(Tx_n)$ is Cauchy, so $\norm{x_n - x_m} \leq
c^{-1}\norm{Tx_n - Tx_m} \to 0$, hence $(x_n)$ is Cauchy. Since $X$ is complete,
$x_n \to x$ and $y = Tx \in T(X)$.

(ii) $\Rightarrow$ (i): $T(X)$ with the norm inherited from $Y$ is a Banach space
(closed subspace of a --- possibly incomplete --- normed space, but $T(X) \cong X/\Ker T = X$
since $T$ is injective). The boundedness of $T^{-1}$ gives
$\norm{x} = \norm{T^{-1}(Tx)} \leq \norm{T^{-1}}\norm{Tx}$, so
$\norm{Tx} \geq c\norm{x}$ with $c = 1/\norm{T^{-1}}$.
\end{proof}

%-----------------------------------------------------------------------
\section{The closed graph theorem}\label{sec:cgt}
%-----------------------------------------------------------------------

\subsection{The graph of a linear operator}

\begin{definition}\label{def:graph}
\index{graph of an operator}
Let $X$ and $Y$ be normed spaces and $T \colon D(T) \to Y$ a (possibly unbounded)
linear operator defined on a subspace $D(T) \subset X$ (the \textbf{domain}\index{domain
of an operator} of $T$). The \textbf{graph} of $T$ is
\[
  \Gamma(T) = \{(x, Tx) : x \in D(T)\} \subset X \times Y.
\]
We equip $X \times Y$ with the product norm
$\norm{(x,y)} = \norm{x}_X + \norm{y}_Y$ (or equivalently
$\max\{\norm{x}, \norm{y}\}$, as these are equivalent).
\end{definition}

\begin{definition}[Closed operator]\label{def:closed-operator}
\index{closed operator}
An operator $T \colon D(T) \to Y$ is \textbf{closed} if its graph $\Gamma(T)$
is a closed subset of $X \times Y$. Equivalently, $T$ is closed if and only if:
whenever $(x_n) \subset D(T)$ with $x_n \to x$ in $X$ and $Tx_n \to y$ in $Y$,
then $x \in D(T)$ and $Tx = y$.
\end{definition}

\begin{remark}\label{rem:closed-vs-continuous}
\index{closed operator!vs.\ continuous operator}
A continuous operator with closed domain is always closed (if $x_n \to x$ and $D(T)$
is closed, then $x \in D(T)$ and $Tx_n \to Tx$ by continuity). However, a closed
operator need not be continuous --- this is the key distinction when dealing with
unbounded operators.
\end{remark}

\begin{example}\label{ex:differentiation-closed}
\index{differentiation operator!closed}
On $X = C([0,1])$ with the supremum norm, consider the operator
$T = d/dx$ with domain $D(T) = C^1([0,1])$. Then $T$ is closed: if $f_n \to f$
uniformly and $f_n' \to g$ uniformly, then $f$ is differentiable and $f' = g$.
However, $T$ is not bounded on $(D(T), \norm{\cdot}_\infty)$.
\end{example}

\subsection{Statement and proof of the closed graph theorem}

\begin{theorem}[Closed Graph Theorem]\label{thm:cgt}
\index{closed graph theorem!statement and proof}
Let $X$ and $Y$ be Banach spaces and $T \colon X \to Y$ a linear operator
(defined on all of $X$). If $T$ is closed (i.e., $\Gamma(T)$ is closed in
$X \times Y$), then $T$ is bounded.
\end{theorem}

\begin{proof}
\index{closed graph theorem!proof}
Since $X$ and $Y$ are Banach spaces, the product $X \times Y$ (with the norm
$\norm{(x,y)} = \norm{x} + \norm{y}$) is a Banach space. The graph $\Gamma(T)$
is a closed linear subspace of $X \times Y$, hence itself a Banach space.

Consider the projections:
\begin{align*}
  \pi_1 &\colon \Gamma(T) \to X, \quad \pi_1(x, Tx) = x, \\
  \pi_2 &\colon \Gamma(T) \to Y, \quad \pi_2(x, Tx) = Tx.
\end{align*}
Both are bounded: $\norm{\pi_1(x,Tx)} = \norm{x} \leq \norm{x} + \norm{Tx}$
and similarly for $\pi_2$.

The projection $\pi_1$ is a bounded linear bijection from the Banach space
$\Gamma(T)$ to the Banach space $X$. By the Banach isomorphism theorem
(Corollary~\ref{cor:banach-iso}), $\pi_1^{-1} \colon X \to \Gamma(T)$ is bounded.
Therefore $T = \pi_2 \circ \pi_1^{-1}$ is bounded, being the composition of two
bounded operators.
\end{proof}

\begin{remark}\label{rem:cgt-elegant}
The proof of the closed graph theorem via the Banach isomorphism theorem is
remarkably elegant: the entire argument reduces to the observation that two natural
projections and one inverse are bounded.
\end{remark}

\begin{theorem}[Closed graph theorem --- equivalent formulation]\label{thm:cgt-equiv}
\index{closed graph theorem!equivalent formulation}
Let $X$ and $Y$ be Banach spaces and $T \colon X \to Y$ a linear operator. The
following are equivalent:
\begin{enumerate}[label=(\roman*)]
  \item $T$ is bounded (i.e., continuous).
  \item $T$ is closed.
  \item For every sequence $(x_n)$ in $X$: if $x_n \to 0$ and $Tx_n \to y$,
    then $y = 0$.
\end{enumerate}
\end{theorem}

\begin{proof}
(i) $\Rightarrow$ (ii): If $T$ is continuous and $x_n \to x$, $Tx_n \to y$, then
$Tx_n \to Tx$ by continuity, so $y = Tx$ and $\Gamma(T)$ is closed.

(ii) $\Rightarrow$ (iii): Immediate from the definition, with $x = 0$.

(iii) $\Rightarrow$ (ii): Suppose $x_n \to x$ and $Tx_n \to y$. Set
$z_n = x_n - x$; then $z_n \to 0$ and $Tz_n = Tx_n - Tx \to y - Tx$. By (iii),
$y - Tx = 0$, so $y = Tx$.

(ii) $\Rightarrow$ (i): This is Theorem~\ref{thm:cgt}.
\end{proof}

\begin{example}[Automatic continuity]\label{ex:automatic-continuity}
\index{automatic continuity}
Let $H$ be a Hilbert space and $T \colon H \to H$ a linear operator satisfying
$\inner{Tx}{y} = \inner{x}{Ty}$ for all $x, y \in H$ (i.e., $T$ is formally
symmetric and defined on all of $H$). Then $T$ is bounded.

\textit{Proof:} We verify condition (iii). Suppose $x_n \to 0$ and $Tx_n \to y$.
For any $z \in H$:
\[
  \inner{y}{z} = \lim_n \inner{Tx_n}{z} = \lim_n \inner{x_n}{Tz} = 0.
\]
Since this holds for all $z$, $y = 0$. By the closed graph theorem, $T$ is bounded.
This is the \emph{Hellinger--Toeplitz theorem}\index{Hellinger--Toeplitz theorem}.
\end{example}

%-----------------------------------------------------------------------
\section{Applications of the closed graph theorem}\label{sec:cgt-applications}
%-----------------------------------------------------------------------

\subsection{Unbounded operators and their domains}\label{subsec:unbounded-ops}

\begin{definition}\label{def:unbounded-operator}
\index{unbounded operator}
An \textbf{unbounded operator} from $X$ to $Y$ is a linear map
$T \colon D(T) \to Y$ where $D(T)$ is a (not necessarily closed) linear subspace
of $X$. The operator is \textbf{densely defined}\index{densely defined operator}
if $D(T)$ is dense in $X$.
\end{definition}

\begin{proposition}\label{prop:closed-bounded-domain}
\index{closed operator!with closed domain}
Let $X$ and $Y$ be Banach spaces. If $T \colon D(T) \to Y$ is a closed operator
and $D(T)$ is a closed subspace of $X$, then $T$ is bounded.
\end{proposition}

\begin{proof}
Since $D(T)$ is a closed subspace of a Banach space, it is itself a Banach space.
The operator $T \colon D(T) \to Y$ is closed and defined on the entire Banach space
$D(T)$. By the closed graph theorem, $T$ is bounded.
\end{proof}

\begin{corollary}\label{cor:unbounded-closed-not-everywhere}
\index{unbounded operator!domain}
If $T \colon D(T) \to Y$ is closed and unbounded, then $D(T)$ cannot be a closed
subspace of $X$. In particular, if $T$ is densely defined, closed, and unbounded,
then $D(T) \neq X$.
\end{corollary}

\begin{example}\label{ex:unbounded-laplacian}
\index{Laplacian!as unbounded operator}
The Laplacian $\Delta = \sum_{k=1}^n \partial^2/\partial x_k^2$ on $L^2(\R^n)$
is a densely defined closed operator with domain
$D(\Delta) = H^2(\R^n)$\index{Sobolev space} (the Sobolev space of functions with
two weak derivatives in $L^2$). By the corollary, $H^2(\R^n)$ is not closed in
$L^2(\R^n)$ (which is obvious since $H^2 \neq L^2$ but $H^2$ is dense).
\end{example}

\begin{proposition}[Graph norm]\label{prop:graph-norm}
\index{graph norm}
Let $T \colon D(T) \to Y$ be a closed operator. The \textbf{graph norm} on $D(T)$
is defined by
\[
  \norm{x}_T = \norm{x}_X + \norm{Tx}_Y.
\]
Then $(D(T), \norm{\cdot}_T)$ is a Banach space, and $T \colon (D(T), \norm{\cdot}_T)
\to Y$ is bounded with $\norm{T}_{\norm{\cdot}_T \to Y} \leq 1$.
\end{proposition}

\begin{proof}
The map $\Phi \colon (D(T), \norm{\cdot}_T) \to \Gamma(T)$ defined by
$\Phi(x) = (x, Tx)$ is an isometric isomorphism. Since $T$ is closed, $\Gamma(T)$
is a closed subspace of the Banach space $X \times Y$, hence complete. Thus
$(D(T), \norm{\cdot}_T)$ is complete. The boundedness of $T$ follows from
$\norm{Tx} \leq \norm{x} + \norm{Tx} = \norm{x}_T$.
\end{proof}

\subsection{Application: automatic continuity of homomorphisms}

\begin{theorem}\label{thm:auto-cont-algebra}
\index{automatic continuity!Banach algebra homomorphisms}
Let $A$ and $B$ be Banach algebras\index{Banach algebra} and let
$\varphi \colon A \to B$ be an algebra homomorphism that is also a surjection.
If $B$ is semisimple\index{semisimple Banach algebra} (i.e., its Jacobson
radical\index{Jacobson radical} is zero), then $\varphi$ is continuous.
\end{theorem}

\begin{remark}
This deep result, due to B.~E.~Johnson\index{Johnson, B.~E.} (1967), illustrates
how the closed graph theorem can be used to derive automatic continuity in algebraic
settings. The proof uses the closed graph theorem together with properties of the
spectrum. We omit the full proof and refer to~\cite{dales2000} for details.
\end{remark}

%-----------------------------------------------------------------------
\section{The closed range theorem}\label{sec:closed-range}
%-----------------------------------------------------------------------

\begin{definition}\label{def:range-annihilator}
\index{annihilator}\index{preannihilator}
Let $X$ be a Banach space. For $S \subset X$, the \textbf{annihilator} of $S$ is
\[
  S^\perp = \{ f \in X^* : f(x) = 0 \text{ for all } x \in S \}.
\]
For $W \subset X^*$, the \textbf{preannihilator} of $W$ is
\[
  {}^\perp W = \{ x \in X : f(x) = 0 \text{ for all } f \in W \}.
\]
\end{definition}

\begin{theorem}[Closed range theorem]\label{thm:closed-range}
\index{closed range theorem}
Let $X$ and $Y$ be Banach spaces and $T \in \mathcal{L}(X,Y)$. The following
conditions are equivalent:
\begin{enumerate}[label=(\roman*)]
  \item $\ran(T) = T(X)$ is closed in $Y$.\index{range!closed}
  \item $\ran(T^*) = T^*(Y^*)$ is closed in $X^*$.
  \item $\ran(T) = {}^\perp(\Ker T^*)$.
  \item $\ran(T^*) = (\Ker T)^\perp$.
\end{enumerate}
\end{theorem}

\begin{proof}
\index{closed range theorem!proof}
We prove the cycle (i) $\Rightarrow$ (iv) $\Rightarrow$ (ii) $\Rightarrow$
(iii) $\Rightarrow$ (i).

\textbf{(i) $\Rightarrow$ (iv):}
The inclusion $\ran(T^*) \subset (\Ker T)^\perp$ is clear: if $f = T^*g$ and
$x \in \Ker T$, then $f(x) = g(Tx) = 0$.

For the reverse, let $f \in (\Ker T)^\perp$. Since $\ran(T)$ is closed, the
quotient map induces an isomorphism $\tilde{T} \colon X/\Ker T \to \ran(T)$
defined by $\tilde{T}([x]) = Tx$, and $\tilde{T}$ is bounded below (by the open
mapping theorem applied to $\tilde{T}$, since both $X/\Ker T$ and $\ran(T)$ are
Banach spaces). Define $\tilde{f} \colon X/\Ker T \to \K$ by $\tilde{f}([x]) = f(x)$
(well-defined since $f \in (\Ker T)^\perp$). Then $\tilde{f}$ is bounded. Define
$h \colon \ran(T) \to \K$ by $h(Tx) = f(x)$. Then $h$ is well-defined and bounded
(since $\abs{h(Tx)} = \abs{f(x)} \leq \norm{f}\norm{x}$ and $\norm{x} \leq
C\norm{Tx}$ by the bounded below property on $X/\Ker T$, we get
$\abs{h(y)} \leq \norm{f}C\norm{y}$ for $y \in \ran(T)$). Extend $h$ to
$g \in Y^*$ by Hahn--Banach. Then $T^*g = f$, so $f \in \ran(T^*)$.

\textbf{(iv) $\Rightarrow$ (ii):} Immediate, since $(\Ker T)^\perp$ is a closed
subspace of $X^*$.

\textbf{(ii) $\Rightarrow$ (iii):}
The inclusion $\ran(T) \subset {}^\perp(\Ker T^*)$ always holds. For the reverse,
suppose $y \in {}^\perp(\Ker T^*)$, i.e., $g(y) = 0$ for all $g \in \Ker T^*$.
We must show $y \in \overline{\ran(T)} = \ran(T)$ (we need to show the range is
closed). Define $\Phi \colon \ran(T^*) \to \K$ by $\Phi(T^*g) = g(y)$. This is
well-defined: if $T^*g_1 = T^*g_2$, then $g_1 - g_2 \in \Ker T^*$, so
$(g_1 - g_2)(y) = 0$, i.e., $g_1(y) = g_2(y)$.

Since $\ran(T^*)$ is closed (hypothesis (ii)), it is a Banach space. The functional
$\Phi$ is bounded: by the open mapping theorem, the induced map
$\tilde{T}^* \colon Y^*/\Ker T^* \to \ran(T^*)$ is an isomorphism, so there exists
$C > 0$ with $\inf_{h \in \Ker T^*} \norm{g - h} \leq C \norm{T^*g}$. Then
$\abs{\Phi(T^*g)} = \abs{g(y)} = \inf_{h \in \Ker T^*}\abs{(g-h)(y)} \leq
\norm{y}\inf_{h}\norm{g-h} \leq C\norm{y}\norm{T^*g}$.

By Hahn--Banach, extend $\Phi$ to $F \in (X^*)^* = X^{**}$. Under the canonical
embedding $X \hookrightarrow X^{**}$, we have $F(T^*g) = g(y)$ for all $g$.
However, we need $y \in \ran(T)$, not just $y \in {}^\perp(\Ker T^*)$.

Consider the restriction $\hat{T} \colon X \to \overline{\ran(T)}$. Then (ii) for
$T$ gives $\ran(T^*) = \ran(\hat{T}^*)$ is closed. Repeating the argument of
(i) $\Rightarrow$ (iv) with $\overline{\ran(T)}$ in place of $Y$: since
$\ran(\hat{T}^*)$ is closed, the induced map is bounded below, and the construction
shows $y \in \ran(T)$. Hence $\ran(T) = {}^\perp(\Ker T^*)$ and in particular
$\ran(T)$ is closed.

\textbf{(iii) $\Rightarrow$ (i):} ${}^\perp(\Ker T^*)$ is an intersection of
kernels of continuous functionals, hence closed. So $\ran(T) = {}^\perp(\Ker T^*)$
is closed.
\end{proof}

\begin{corollary}\label{cor:surjective-iff-bounded-below}
\index{surjective operator!characterization}
Let $T \in \mathcal{L}(X,Y)$ with $X, Y$ Banach spaces. Then $T$ is surjective
if and only if $T^*$ is bounded below, i.e., there exists $c > 0$ with
$\norm{T^* g} \geq c \norm{g}$ for all $g \in Y^*$.
\end{corollary}

\begin{proof}
If $T$ is surjective, then $\ran(T) = Y$ is closed, and the closed range theorem
gives $\ran(T^*) = (\Ker T)^\perp$. Moreover, $\Ker T^* = (\ran T)^\perp
= Y^\perp = \{0\}$, so $T^*$ is injective with closed range, hence bounded below
by Proposition~\ref{prop:bounded-below-open}.

Conversely, if $T^*$ is bounded below, then $T^*$ is injective and $\ran(T^*)$
is closed. By the closed range theorem, $\ran(T) = {}^\perp(\Ker T^*)
= {}^\perp \{0\} = Y$.
\end{proof}

%-----------------------------------------------------------------------
\section{Characterization of operators with closed range}%
\label{sec:closed-range-char}
%-----------------------------------------------------------------------

\begin{theorem}[Characterization via minimum modulus]\label{thm:closed-range-modulus}
\index{minimum modulus}\index{closed range!characterization}
Let $X$ and $Y$ be Banach spaces and $T \in \mathcal{L}(X,Y)$. Define the
\textbf{minimum modulus} of $T$ as
\[
  \gamma(T) = \inf\!\left\{ \frac{\norm{Tx}}{\mathrm{dist}(x, \Ker T)} :
  x \notin \Ker T \right\}.
\]
Then $\ran(T)$ is closed if and only if $\gamma(T) > 0$.
\end{theorem}

\begin{proof}
\index{minimum modulus!and closed range}
The quotient space $X/\Ker T$ is a Banach space (since $\Ker T$ is closed). The
induced map $\hat{T} \colon X/\Ker T \to Y$ defined by $\hat{T}([x]) = Tx$ is a
well-defined injective bounded operator with $\ran(\hat{T}) = \ran(T)$. Moreover,
\[
  \gamma(T) = \inf_{\norm{[x]} = 1} \norm{\hat{T}([x])}.
\]

If $\gamma(T) > 0$, then $\hat{T}$ is bounded below by $\gamma(T)$, so $\ran(\hat{T})$
is closed (Proposition~\ref{prop:bounded-below-open}).

Conversely, if $\ran(T)$ is closed, then $\hat{T} \colon X/\Ker T \to \ran(T)$ is
a bounded bijection between Banach spaces. By the Banach isomorphism theorem,
$\hat{T}^{-1}$ is bounded, say $\norm{\hat{T}^{-1}} = C$. Then
$\norm{[x]} \leq C\norm{Tx}$ for all $x$, giving $\gamma(T) \geq 1/C > 0$.
\end{proof}

\begin{proposition}\label{prop:closed-range-perturbation}
\index{closed range!stability under perturbation}
Let $T \in \mathcal{L}(X,Y)$ have closed range with $\gamma(T) > 0$. If
$S \in \mathcal{L}(X,Y)$ with $\norm{S} < \gamma(T)$, then $T + S$ also has
closed range.
\end{proposition}

\begin{proof}
For $x \notin \Ker(T+S)$:
\begin{align*}
  \norm{(T+S)x} &\geq \norm{Tx} - \norm{Sx} \geq \gamma(T)\,\mathrm{dist}(x, \Ker T)
  - \norm{S}\norm{x} \\
  &\geq \gamma(T)\,\mathrm{dist}(x, \Ker(T+S)) - \norm{S}\norm{x}.
\end{align*}
A more careful argument using the quotient space shows that $\gamma(T+S) \geq
\gamma(T) - \norm{S} > 0$, hence $\ran(T+S)$ is closed.
\end{proof}

%-----------------------------------------------------------------------
\section{Factorization principles}\label{sec:factorization}
%-----------------------------------------------------------------------

\begin{theorem}[Factorization through a closed subspace]%
\label{thm:factorization}
\index{factorization principle}
Let $X, Y, Z$ be Banach spaces and $T \in \mathcal{L}(X,Y)$,
$S \in \mathcal{L}(X,Z)$. Suppose $\Ker T \subset \Ker S$. Then there exists a
unique linear operator $R \colon \ran(T) \to Z$ such that $S = R \circ T$ on $X$.
If $\ran(T)$ is closed, then $R$ is bounded.
\end{theorem}

\begin{proof}
Define $R(Tx) = Sx$. This is well-defined: if $Tx_1 = Tx_2$, then
$x_1 - x_2 \in \Ker T \subset \Ker S$, so $Sx_1 = Sx_2$. Linearity is clear.

If $\ran(T)$ is closed, then $\ran(T)$ is a Banach space. We verify $R$ is closed:
suppose $Tx_n \to y \in \ran(T)$ and $R(Tx_n) = Sx_n \to z$. Since $\ran(T)$ is
closed, $y = Tx_0$ for some $x_0$. By the open mapping theorem applied to
$\hat{T} \colon X/\Ker T \to \ran(T)$, there exist $\tilde{x}_n$ with
$T\tilde{x}_n = Tx_n$ and $\norm{\tilde{x}_n - \tilde{x}_0} \to 0$ (choosing
coset representatives). Since $\Ker T \subset \Ker S$, $S\tilde{x}_n = Sx_n \to z$
and $S\tilde{x}_n \to S\tilde{x}_0$, so $z = Sx_0 = R(Tx_0) = R(y)$. By the
closed graph theorem, $R$ is bounded.
\end{proof}

\begin{corollary}[Open mapping theorem --- alternative formulation]%
\label{cor:omt-factor}
\index{open mapping theorem!factorization}
Let $T \in \mathcal{L}(X,Y)$ be surjective. Then $T$ factors as
$X \xrightarrow{\pi} X/\Ker T \xrightarrow{\hat{T}} Y$, where $\pi$ is the
quotient map and $\hat{T}$ is a topological isomorphism.
\end{corollary}

\begin{proof}
$\hat{T}$ is a bijective bounded operator from $X/\Ker T$ (a Banach space) to $Y$
(a Banach space). By the Banach isomorphism theorem, $\hat{T}$ is an isomorphism.
\end{proof}

%-----------------------------------------------------------------------
\section{Fr\'echet spaces and the generalized closed graph theorem}%
\label{sec:frechet}
%-----------------------------------------------------------------------

The open mapping and closed graph theorems generalize beyond Banach spaces to the
setting of \emph{Fr\'echet spaces}\index{Fr\'echet space}.

\begin{definition}\label{def:frechet-space}
\index{Fr\'echet space!definition}
A \textbf{Fr\'echet space} is a topological vector space $X$ whose topology is
induced by a \emph{complete, translation-invariant metric} $d$ and by a countable
family of seminorms\index{seminorm} $\{p_n\}_{n \geq 1}$ such that:
\[
  d(x,y) = \sum_{n=1}^{\infty} \frac{1}{2^n}
  \frac{p_n(x-y)}{1 + p_n(x-y)}.
\]
Equivalently, $X$ is a complete metrizable locally convex space.
\end{definition}

\begin{example}\label{ex:frechet-examples}
\index{Fr\'echet space!examples}
\begin{enumerate}[label=(\roman*)]
  \item Every Banach space is a Fr\'echet space (with a single seminorm equal to
    the norm).
  \item $C^\infty([0,1])$\index{$C^\infty([0,1])$}, the space of smooth functions
    on $[0,1]$, with seminorms $p_n(f) = \max_{0 \leq k \leq n} \norm{f^{(k)}}_\infty$.
  \item The Schwartz space\index{Schwartz space} $\mathcal{S}(\R^n)$ of rapidly
    decreasing functions, with seminorms
    $p_{\alpha,\beta}(f) = \sup_{x} \abs{x^\alpha D^\beta f(x)}$.
  \item $\mathcal{H}(\Omega)$\index{holomorphic functions}, the space of holomorphic
    functions on an open set $\Omega \subset \C$, with seminorms
    $p_n(f) = \sup_{K_n} \abs{f}$ for a compact exhaustion $\{K_n\}$ of $\Omega$.
  \item $\omega = \K^{\N}$\index{$\omega$ (sequence space)}, the space of all
    sequences, with seminorms $p_n(x) = \abs{x_n}$.
\end{enumerate}
\end{example}

\begin{theorem}[Open mapping theorem for Fr\'echet spaces]%
\label{thm:omt-frechet}
\index{open mapping theorem!Fr\'echet spaces}
Let $X$ and $Y$ be Fr\'echet spaces and $T \colon X \to Y$ a continuous linear
surjection. Then $T$ is open.
\end{theorem}

\begin{proof}[Proof sketch]
The proof follows the same pattern as for Banach spaces. The Baire category theorem
applies since $Y$ is a complete metric space. The approximation argument in Step~2
of Theorem~\ref{thm:omt} is adapted to the metric $d$ (or equivalently to the
countable family of seminorms), using the completeness of $X$ to sum the convergent
series.
\end{proof}

\begin{theorem}[Closed graph theorem for Fr\'echet spaces]%
\label{thm:cgt-frechet}
\index{closed graph theorem!Fr\'echet spaces}
Let $X$ and $Y$ be Fr\'echet spaces and $T \colon X \to Y$ a linear operator with
closed graph. Then $T$ is continuous.
\end{theorem}

\begin{proof}
The product $X \times Y$ is a Fr\'echet space. The graph $\Gamma(T)$ is a closed
subspace, hence a Fr\'echet space. The projection $\pi_1 \colon \Gamma(T) \to X$
is a continuous bijection between Fr\'echet spaces. By
Theorem~\ref{thm:omt-frechet}, $\pi_1$ is open, hence $\pi_1^{-1}$ is continuous.
Then $T = \pi_2 \circ \pi_1^{-1}$ is continuous.
\end{proof}

\begin{remark}\label{rem:webbed-spaces}
\index{webbed space}\index{De Wilde's theorem}
There are further generalizations. De Wilde (1971) proved the closed graph theorem
for operators from an \emph{ultrabornological}\index{ultrabornological space} space
to a \emph{webbed space}. The class of webbed spaces includes all Fr\'echet spaces,
countable inductive limits of Fr\'echet spaces (LF-spaces\index{LF-space}), and
their closed subspaces and quotients. This provides the most general framework in
which a ``closed graph implies continuous'' theorem holds without additional
hypotheses on the operator.
\end{remark}

\begin{example}[Failure of the closed graph theorem]\label{ex:cgt-failure}
\index{closed graph theorem!failure}
The closed graph theorem fails if $X$ is not complete. Let $X = (C^1([0,1]),
\norm{\cdot}_\infty)$ and $Y = (C([0,1]), \norm{\cdot}_\infty)$. The differentiation
operator $T = d/dx \colon X \to Y$ is closed (as shown in
Example~\ref{ex:differentiation-closed}) but unbounded. Here $X$ is not complete
in the $\norm{\cdot}_\infty$ norm.
\end{example}

%-----------------------------------------------------------------------
\section{Exercises for Chapter~\ref{ch:omt-cgt}}\label{sec:omt-cgt-exercises}
%-----------------------------------------------------------------------

\begin{exercise}[Open mapping theorem application; $\star$]\label{ex:omt-surjective-quotient}
\index{open mapping theorem!exercises}
Let $X$ be a Banach space and $M$ a closed subspace. Show that the quotient map
$\pi \colon X \to X/M$ is an open mapping. Conclude that the quotient topology on
$X/M$ coincides with the topology induced by the quotient norm.
\end{exercise}

\begin{exercise}[Two Banach space norms; $\star$]\label{ex:two-norms}
\index{equivalent norms!exercises}
Let $(X, \norm{\cdot}_1)$ and $(X, \norm{\cdot}_2)$ be Banach spaces with the same
underlying vector space. Suppose that $x_n \to 0$ in $\norm{\cdot}_1$ and
$x_n \to x$ in $\norm{\cdot}_2$ implies $x = 0$. Show that the two norms are
equivalent.

\noindent\textit{Hint:} Consider the identity map and use the closed graph theorem.
\end{exercise}

\begin{exercise}[Projection operators; $\star$]\label{ex:projection}
\index{projection operator!exercises}
Let $X$ be a Banach space and $P \in \mathcal{L}(X)$ a projection ($P^2 = P$).
Show that $X = \ran(P) \oplus \Ker(P)$ and both $\ran(P)$ and $\Ker(P)$ are closed.
Conversely, if $X = M \oplus N$ with $M, N$ closed subspaces, show that the
projection onto $M$ along $N$ is bounded.

\noindent\textit{Hint for the converse:} Use the closed graph theorem or the
Banach isomorphism theorem on the map $M \times N \to X$, $(m,n) \mapsto m + n$.
\end{exercise}

\begin{exercise}[Closed graph and sequential characterization; $\star$]%
\label{ex:cgt-sequential}
\index{closed graph theorem!exercises}
Let $X$ and $Y$ be Banach spaces and $T \colon X \to Y$ a linear operator.
Show that $T$ is bounded if and only if $x_n \weakto x$ in $X$ implies
$Tx_n \weakto Tx$ in $Y$ (i.e., $T$ is weakly sequentially continuous).
\end{exercise}

\begin{exercise}[Hellinger--Toeplitz; $\star\star$]\label{ex:hellinger-toeplitz}
\index{Hellinger--Toeplitz theorem!exercise}
Let $H$ be a Hilbert space. Suppose $T \colon H \to H$ is a linear map satisfying
$\inner{Tx}{y} = \inner{x}{Ty}$ for all $x, y \in H$. Prove that $T$ is bounded
\emph{without} using the closed graph theorem, by directly applying the uniform
boundedness principle.

\noindent\textit{Hint:} Fix $y \in H$ and consider the functionals
$f_y(x) = \inner{Tx}{y}$.
\end{exercise}

\begin{exercise}[Inverse of a perturbation; $\star\star$]\label{ex:neumann-series}
\index{Neumann series}\index{perturbation of invertible operator}
Let $X$ be a Banach space and $T \in \mathcal{L}(X)$ with $\norm{T} < 1$. Show that
$\Id - T$ is invertible with
\[
  (\Id - T)^{-1} = \sum_{n=0}^{\infty} T^n
\]
(the \textbf{Neumann series}), and $\norm{(\Id - T)^{-1}} \leq 1/(1-\norm{T})$.
Use this to show: if $S \in \mathcal{L}(X)$ is invertible and $\norm{U - S} <
1/\norm{S^{-1}}$, then $U$ is invertible.
\end{exercise}

\begin{exercise}[Closed range and duality; $\star\star$]\label{ex:closed-range-dual}
\index{closed range theorem!exercises}
Let $T \in \mathcal{L}(X,Y)$ with $X, Y$ Banach spaces.
\begin{enumerate}[label=(\alph*)]
  \item Show that if $T$ has closed range, then $\ran(T)$ is a Banach space isomorphic
    to $X/\Ker T$.
  \item Show that $T$ has closed range if and only if there exists $C > 0$ such that
    for every $y \in \ran(T)$, there exists $x \in X$ with $Tx = y$ and
    $\norm{x} \leq C\norm{y}$.
  \item Deduce that $T$ is surjective if and only if $T^*$ is injective and has
    closed range.
\end{enumerate}
\end{exercise}

\begin{exercise}[Fr\'echet space of smooth functions; $\star\star$]%
\label{ex:frechet-smooth}
\index{Fr\'echet space!exercises}
Let $X = C^\infty(\R)$ with the family of seminorms
$p_{n,K}(f) = \sup_{x \in [-n,n]} \abs{f^{(k)}(x)}$ for $k \leq n$.
\begin{enumerate}[label=(\alph*)]
  \item Show that $X$ is a Fr\'echet space.
  \item Show that the differentiation operator $D \colon X \to X$, $Df = f'$, is
    continuous.
  \item Show that $D$ is surjective (every smooth function has a smooth antiderivative)
    and apply the open mapping theorem.
\end{enumerate}
\end{exercise}

\begin{exercise}[Schauder's theorem; $\star\star\star$]\label{ex:schauder-theorem}
\index{Schauder's theorem}
Let $X, Y$ be Banach spaces and $T \in \mathcal{L}(X,Y)$. Prove that $T$ is
compact if and only if $T^*$ is compact.

\noindent\textit{Hint:} For one direction, use the Arzel\`a--Ascoli theorem. For the
other, use the closed range theorem and properties of compact operators.
\end{exercise}

\begin{exercise}[Operators on $\ell^p$; $\star\star$]\label{ex:operators-ell-p}
\index{operators on $\ell^p$}
Let $1 \leq p \leq \infty$ and let $A = (a_{ij})_{i,j \geq 1}$ be an infinite matrix
defining a linear map $T_A \colon \ell^p \to \ell^p$ by
$(T_A x)_i = \sum_j a_{ij} x_j$.
\begin{enumerate}[label=(\alph*)]
  \item Use the closed graph theorem to show: if $T_A$ maps $\ell^p$ into $\ell^p$
    (i.e., the series converges and the result is in $\ell^p$ for every
    $x \in \ell^p$), then $T_A$ is automatically bounded.
  \item For $p = 1$, show directly that $\norm{T_A} = \sup_j \sum_i \abs{a_{ij}}$.
  \item For $p = \infty$, show that $\norm{T_A} = \sup_i \sum_j \abs{a_{ij}}$.
\end{enumerate}
\end{exercise}

\begin{exercise}[Banach isomorphism and complementation; $\star\star\star$]%
\label{ex:complementation}
\index{complemented subspace}\index{Banach isomorphism theorem!exercises}
Let $X$ be a Banach space and $M, N$ closed subspaces with $M \cap N = \{0\}$
and $M + N = X$. Prove:
\begin{enumerate}[label=(\alph*)]
  \item The map $\Phi \colon M \times N \to X$, $\Phi(m,n) = m + n$ is a
    topological isomorphism.
  \item There exists $C > 0$ such that $\norm{m} + \norm{n} \leq C\norm{m+n}$
    for all $m \in M$, $n \in N$.
  \item Give an example of a Banach space $X$ and a closed subspace $M$ that is
    \emph{not} complemented. \textit{Hint:} $c_0$ in $\ell^\infty$.
\end{enumerate}
\end{exercise}

\begin{exercise}[The three-space property; $\star\star\star$]%
\label{ex:three-space}
\index{three-space property}
Let $X$ be a Banach space and $M$ a closed subspace such that both $M$ and $X/M$
are isomorphic (as Banach spaces) to Hilbert spaces. Must $X$ be isomorphic to a
Hilbert space? \textit{Hint:} Consider the Enflo--Lindenstrauss--Pisier example,
or prove the positive result under additional hypotheses.
\end{exercise}

\begin{exercise}[Uniform boundedness via open mapping; $\star\star$]%
\label{ex:ubp-via-omt}
\index{uniform boundedness principle!via open mapping}
Derive the uniform boundedness principle as a consequence of the open mapping
theorem.

\noindent\textit{Hint:} Given pointwise bounded $(T_n) \subset \mathcal{L}(X,Y)$,
consider the operator $\Phi \colon X \to \ell^\infty(Y)$ (or $Y^{\N}$) defined by
$\Phi(x) = (T_n x)_n$, show it has closed graph, and apply the closed graph theorem.
\end{exercise}

\begin{exercise}[Quotient maps are open; $\star$]\label{ex:quotient-open}
\index{quotient map!openness}
Let $X$ be a Banach space and $M \subset X$ a closed subspace. Prove directly
(without using the open mapping theorem) that the canonical projection
$\pi \colon X \to X/M$ is an open map. Then give a second proof using the open
mapping theorem.
\end{exercise}

\begin{exercise}[Isomorphic preduals; $\star\star\star$]\label{ex:isomorphic-preduals}
\index{predual!uniqueness}
Let $X$ and $Y$ be Banach spaces such that $X^* \cong Y^*$ (isometrically isomorphic).
Must $X \cong Y$? Investigate this question for the following cases:
\begin{enumerate}[label=(\alph*)]
  \item $X = \ell^1$ and $Y = \ell^1$. (Here $X^* = Y^* = \ell^\infty$.)
  \item $X = L^1([0,1])$. What is $X^*$? Find a non-isomorphic space with the
    same dual.
\end{enumerate}
\end{exercise}

\begin{exercise}[Closed graph for multilinear maps; $\star\star$]%
\label{ex:cgt-multilinear}
\index{closed graph theorem!multilinear maps}
Let $X_1, \ldots, X_n, Y$ be Banach spaces and
$T \colon X_1 \times \cdots \times X_n \to Y$ an $n$-linear map. Show that if
$T$ is \emph{separately} continuous in each variable, then $T$ is jointly continuous.

\noindent\textit{Hint:} Proceed by induction on $n$. For $n = 1$, this is trivial.
For the inductive step, fix all but one variable and apply the uniform boundedness
principle.
\end{exercise}

\begin{exercise}[Surjectivity and the open mapping theorem; $\star\star$]%
\label{ex:omt-not-surjective}
\index{open mapping theorem!non-surjective case}
Let $X, Y$ be Banach spaces and $T \in \mathcal{L}(X,Y)$ with $\ran(T)$ of
second category in $Y$. Show that $T$ is surjective.

\noindent\textit{Hint:} Revisit the proof of the open mapping lemma. If $\ran(T)$
is of second category, the same Baire argument shows that the closure of the image
of the unit ball has nonempty interior \emph{in $Y$}.
\end{exercise}

\begin{exercise}[Application to differential equations; $\star\star$]%
\label{ex:ode-closed-graph}
\index{differential equations!closed graph application}
Consider the Banach space $X = C([0,1])$ and the operator $T \colon D(T) \to X$
defined by $Tf = f'' + f$ with $D(T) = \{f \in C^2([0,1]) : f(0) = f(1) = 0\}$.
\begin{enumerate}[label=(\alph*)]
  \item Show that $T$ is a closed operator (with $D(T)$ equipped with the supremum
    norm from $X$).
  \item Show that $D(T)$ is not closed in $X$.
  \item Equip $D(T)$ with the graph norm. Show that $T$ becomes a bounded operator
    from $(D(T), \norm{\cdot}_T)$ to $X$.
  \item Determine $\Ker(T)$ explicitly.
\end{enumerate}
\end{exercise}

\begin{exercise}[The Mittag-Leffler theorem; $\star\star\star$]%
\label{ex:mittag-leffler}
\index{Mittag-Leffler theorem}
\emph{(A topological version.)} Let $(X_n, d_n)_{n \geq 0}$ be a sequence of
complete metric spaces and $\varphi_n \colon X_{n+1} \to X_n$ continuous maps
with dense range. Define the projective limit
$X = \varprojlim X_n = \{(x_n) \in \prod X_n : \varphi_n(x_{n+1}) = x_n\}$.
Prove that the projection $\pi_n \colon X \to X_n$ has dense range for each $n$.

\noindent\textit{Hint:} Use a Baire-category-style nested ball argument, constructing
compatible approximations in each $X_n$.
\end{exercise}

\begin{exercise}[Open mapping for non-complete spaces; $\star\star$]%
\label{ex:omt-counterexample}
\index{open mapping theorem!counterexample}
Let $X = (\ell^1, \norm{\cdot}_1)$ and define $Y = (\ell^1, \norm{\cdot}_2)$
where $\norm{x}_2 = \left(\sum_n \abs{x_n}^2\right)^{1/2}$. Note that $Y$ is
\emph{not} complete.
\begin{enumerate}[label=(\alph*)]
  \item Show that the identity $\Id \colon X \to Y$ is a well-defined, continuous,
    linear bijection.
  \item Show that $\Id^{-1}$ is \emph{not} continuous, hence $\Id$ is not open.
  \item Explain which hypothesis of the open mapping theorem fails.
\end{enumerate}
\end{exercise}

\begin{exercise}[Fredholm operators; $\star\star\star$]\label{ex:fredholm}
\index{Fredholm operator}\index{Fredholm index}
Let $X, Y$ be Banach spaces. An operator $T \in \mathcal{L}(X,Y)$ is called
\textbf{Fredholm} if $\Ker T$ is finite-dimensional, $\ran(T)$ is closed, and
$\codim(\ran T) < \infty$. The \textbf{Fredholm index} is
$\mathrm{ind}(T) = \dim \Ker T - \codim \ran T$.
\begin{enumerate}[label=(\alph*)]
  \item Show that $T$ is Fredholm if and only if $T^*$ is Fredholm, and
    $\mathrm{ind}(T^*) = -\mathrm{ind}(T)$.
  \item Show that if $T$ is Fredholm and $K \in \mathcal{L}(X,Y)$ is compact, then
    $T + K$ is Fredholm with $\mathrm{ind}(T+K) = \mathrm{ind}(T)$.
  \item Show that the set of Fredholm operators is open in $\mathcal{L}(X,Y)$
    and the index is locally constant.
\end{enumerate}
\end{exercise}
%%%%%%%%%%%%%%%%%%%%%%%%%%%%%%%%%%%%%%%%%%%%%%%%%%%%%%%%%%%%
%% CHAPTER 7 — Dual Spaces and Reflexivity
%%%%%%%%%%%%%%%%%%%%%%%%%%%%%%%%%%%%%%%%%%%%%%%%%%%%%%%%%%%%
\chapter{Dual Spaces and Reflexivity}\label{ch:dual-reflexivity}
\minitoc

\section{The Topological Dual}\label{sec:topological-dual}

Throughout this chapter, $E$ denotes a normed space over the field
$\K$ ($= \R$ or $\C$) unless otherwise stated.

\begin{definition}[Topological dual]\label{def:topological-dual}
\index{dual space!topological}
\index{topological dual}
The \textbf{topological dual}\index{dual!topological|see{topological dual}} (or
simply \textbf{dual}) of a normed space $E$ is the Banach space
\[
  E' \;=\; \mathcal{L}(E,\K)
  \;=\; \bigl\{\, f: E \to \K \;\big|\; f \text{ is continuous and linear}\bigr\},
\]
equipped with the operator norm
\[
  \norm{f}_{E'} \;=\; \sup_{\norm{x}\le 1} \abs{f(x)}.
\]
Elements of $E'$ are called \textbf{continuous linear functionals}\index{continuous linear functional}
(or simply \textbf{functionals}).
\end{definition}

\begin{remark}\label{rem:dual-banach}
\index{dual space!is Banach}
Since $\K$ is complete, $E' = \mathcal{L}(E,\K)$ is always a Banach space,
regardless of whether $E$ itself is complete.
\end{remark}

We recall the fundamental extension theorem whose consequences pervade the
entire theory.

\begin{theorem}[{Hahn--Banach, analytic form}]\label{thm:hahn-banach-recall}
\index{Hahn--Banach theorem}
Let $E$ be a real or complex vector space, $p:E\to\R$ a sublinear functional
(or seminorm in the complex case), $F\subset E$ a subspace, and $f:F\to\K$
a linear functional satisfying $\abs{f(x)}\le p(x)$ for all $x\in F$.  Then
there exists a linear functional $\tilde{f}:E\to\K$ extending $f$ such that
$\abs{\tilde{f}(x)}\le p(x)$ for all $x\in E$.
\end{theorem}

\begin{corollary}\label{cor:hb-norm-preserving}
\index{Hahn--Banach theorem!norm-preserving extension}
Let $F$ be a subspace of a normed space $E$ and let $f\in F'$.  Then there
exists $\tilde{f}\in E'$ such that $\tilde{f}\big|_F = f$ and
$\norm{\tilde{f}}_{E'} = \norm{f}_{F'}$.
\end{corollary}

\begin{corollary}\label{cor:hb-separation}
\index{Hahn--Banach theorem!separation of points}
For every $x\in E$ with $x\ne 0$, there exists $f\in E'$ with $\norm{f}=1$
and $f(x) = \norm{x}$.  In particular, $E'$ separates points of $E$:
\[
  \bigl(\forall f\in E',\; f(x)=0\bigr) \;\Longrightarrow\; x=0.
\]
\end{corollary}

\begin{proof}
Apply Corollary~\ref{cor:hb-norm-preserving} to the functional
$f_0: \K x \to \K$ defined by $f_0(\lambda x) = \lambda\norm{x}$, which
satisfies $\norm{f_0}=1$.
\end{proof}

% ------------------------------------------------------------------
\section{Duals of Classical Sequence Spaces}\label{sec:dual-classical}
% ------------------------------------------------------------------

\subsection{The dual of \texorpdfstring{$\ell^p$}{lp} for
\texorpdfstring{$1\le p<\infty$}{1<=p<infty}}

We denote by $p'$ the conjugate exponent\index{conjugate exponent} of $p$, defined by
$\frac{1}{p}+\frac{1}{p'}=1$ with the convention $1'=\infty$.

\begin{theorem}[{Dual of $\ell^p$, $1\le p<\infty$}]\label{thm:dual-lp}
\index{dual space!of $\ell^p$}
\index{$\ell^p$ space!dual of}
For $1\le p<\infty$, the map
\[
  \Phi: \ell^{p'} \longrightarrow (\ell^p)',
  \qquad \Phi(y)(x) = \sum_{n=1}^{\infty} x_n y_n,
\]
is an isometric isomorphism.
\end{theorem}

\begin{proof}
We divide the proof into three steps.

\medskip\noindent\textbf{Step 1: $\Phi$ is well-defined and $\norm{\Phi(y)} \le \norm{y}_{p'}$.}\index{H\"older's inequality}
By H\"older's inequality, for every $y\in\ell^{p'}$ and $x\in\ell^p$,
\[
  \Bigl|\sum_{n=1}^\infty x_n y_n\Bigr|
  \;\le\; \sum_{n=1}^\infty |x_n|\,|y_n|
  \;\le\; \norm{x}_p\,\norm{y}_{p'}.
\]
Hence $\Phi(y)\in(\ell^p)'$ and $\norm{\Phi(y)} \le \norm{y}_{p'}$.

\medskip\noindent\textbf{Step 2: $\norm{\Phi(y)} = \norm{y}_{p'}$ (isometry).}
We exhibit an element $x\in\ell^p$ with $\norm{x}_p = 1$ that achieves
equality.

\emph{Case $1 < p < \infty$} (so $1 < p' < \infty$).
Define
\[
  x_n = \begin{cases}
    |y_n|^{p'-2}\,\overline{y_n}\,\norm{y}_{p'}^{1-p'} & \text{if } y_n\ne 0,\\
    0 & \text{if } y_n=0.
  \end{cases}
\]
Then $|x_n|^p = |y_n|^{(p'-1)p}\,\norm{y}_{p'}^{(1-p')p}
= |y_n|^{p'}\,\norm{y}_{p'}^{-p'}$, so
$\norm{x}_p^p = \norm{y}_{p'}^{p'}\,\norm{y}_{p'}^{-p'} = 1$.
Moreover,
\[
  \Phi(y)(x) = \sum_n x_n y_n
  = \norm{y}_{p'}^{1-p'}\sum_n |y_n|^{p'}
  = \norm{y}_{p'}^{1-p'}\,\norm{y}_{p'}^{p'}
  = \norm{y}_{p'}.
\]

\emph{Case $p=1$} (so $p'=\infty$).
For any $\varepsilon>0$, pick $n_0$ with $|y_{n_0}|>\norm{y}_\infty - \varepsilon$.
Set $x = \overline{\sgn(y_{n_0})}\,e_{n_0}$ where $e_{n_0}$ is the $n_0$-th
standard basis vector.  Then $\norm{x}_1=1$ and
$\Phi(y)(x) = |y_{n_0}|>\norm{y}_\infty - \varepsilon$.  Letting
$\varepsilon\to 0$ gives $\norm{\Phi(y)}\ge\norm{y}_\infty$.

\medskip\noindent\textbf{Step 3: $\Phi$ is surjective.}
Let $f\in(\ell^p)'$.  Define $y_n = f(e_n)$ where $(e_n)_{n\ge 1}$ is the
standard basis\index{standard basis!of $\ell^p$}.

\emph{Case $1 < p < \infty$.}
For each $N$, define $x^{(N)}\in\ell^p$ by
\[
  x^{(N)}_n = \begin{cases}
    |y_n|^{p'-2}\,\overline{y_n} & \text{if } 1\le n\le N \text{ and } y_n\ne 0,\\
    0 & \text{otherwise}.
  \end{cases}
\]
Then $\norm{x^{(N)}}_p^p = \sum_{n=1}^N |y_n|^{(p'-1)p} = \sum_{n=1}^N |y_n|^{p'}$
and
\[
  f(x^{(N)}) = \sum_{n=1}^N |y_n|^{p'}.
\]
By the bound $|f(x^{(N)})|\le\norm{f}\,\norm{x^{(N)}}_p$,
\[
  \sum_{n=1}^N |y_n|^{p'}
  \;\le\; \norm{f}\,\Bigl(\sum_{n=1}^N |y_n|^{p'}\Bigr)^{1/p},
\]
hence $\bigl(\sum_{n=1}^N |y_n|^{p'}\bigr)^{1-1/p}\le\norm{f}$,
i.e.\ $\bigl(\sum_{n=1}^N |y_n|^{p'}\bigr)^{1/p'}\le\norm{f}$.
Letting $N\to\infty$, we get $y\in\ell^{p'}$ and $\norm{y}_{p'}\le\norm{f}$.

Now for any $x\in\ell^p$, the partial sums $s_N = \sum_{n=1}^N x_n e_n$
converge to $x$ in $\ell^p$, so
\[
  f(x) = \lim_{N\to\infty} f(s_N) = \lim_{N\to\infty} \sum_{n=1}^N x_n y_n
  = \sum_{n=1}^\infty x_n y_n = \Phi(y)(x).
\]
Hence $f = \Phi(y)$.

\emph{Case $p=1$.}
We have $|y_n| = |f(e_n)| \le \norm{f}\,\norm{e_n}_1 = \norm{f}$ for all $n$,
so $y\in\ell^\infty$ with $\norm{y}_\infty\le\norm{f}$.  The same density
argument as above shows $f = \Phi(y)$.
\end{proof}

\subsection{The dual of \texorpdfstring{$c_0$}{c0}}

\begin{theorem}[{Dual of $c_0$}]\label{thm:dual-c0}
\index{dual space!of $c_0$}
\index{$c_0$ space!dual of}
The map
\[
  \Phi: \ell^1 \longrightarrow c_0',
  \qquad \Phi(y)(x) = \sum_{n=1}^\infty x_n y_n,
\]
is an isometric isomorphism.
\end{theorem}

\begin{proof}
\textbf{Well-defined and contractive.}
For $y\in\ell^1$ and $x\in c_0$ (hence $x\in\ell^\infty$),
\[
  \bigl|\Phi(y)(x)\bigr| \le \sum_{n=1}^\infty |x_n|\,|y_n|
  \le \norm{x}_\infty\,\norm{y}_1,
\]
so $\Phi(y)\in c_0'$ with $\norm{\Phi(y)}\le\norm{y}_1$.

\medskip\noindent\textbf{Isometry.}
For $\varepsilon>0$, define $x^{(N)}\in c_0$ by
\[
  x^{(N)}_n = \begin{cases}
    \overline{\sgn(y_n)} & \text{if } 1\le n\le N, \\
    0 & \text{if } n > N.
  \end{cases}
\]
Then $\norm{x^{(N)}}_\infty \le 1$ and $\Phi(y)(x^{(N)}) = \sum_{n=1}^N |y_n|
\to \norm{y}_1$ as $N\to\infty$.  Hence $\norm{\Phi(y)} \ge \norm{y}_1$.

\medskip\noindent\textbf{Surjectivity.}
Let $f\in c_0'$ and set $y_n = f(e_n)$.  For each $N$ and each choice of
signs $\theta_n\in\{z\in\K : |z|=1\}$ with $\theta_n = \overline{\sgn(y_n)}$,
the vector $x^{(N)} = \sum_{n=1}^N \theta_n\,e_n \in c_0$ satisfies
$\norm{x^{(N)}}_\infty \le 1$ and
\[
  \sum_{n=1}^N |y_n| = f(x^{(N)}) \le \norm{f}.
\]
Letting $N\to\infty$ gives $y\in\ell^1$ with $\norm{y}_1\le\norm{f}$.
Since finite sequences are dense in $c_0$\index{density!finite sequences in $c_0$},
the same approximation argument as in Theorem~\ref{thm:dual-lp} shows $f=\Phi(y)$.
\end{proof}

\subsection{The dual of \texorpdfstring{$L^p$}{Lp}}

\begin{theorem}[{Dual of $L^p$, Riesz representation}]\label{thm:dual-Lp}
\index{dual space!of $L^p$}
\index{Riesz representation theorem!for $L^p$}
Let $(\Omega,\mathcal{A},\mu)$ be a $\sigma$-finite measure space and
$1\le p<\infty$.  Then the map
\[
  \Phi: L^{p'}(\mu) \longrightarrow (L^p(\mu))',
  \qquad \Phi(g)(f) = \int_\Omega f\,g\;\dd\mu,
\]
is an isometric isomorphism.
\end{theorem}

\begin{remark}
The complete proof uses the Radon--Nikod\'ym theorem\index{Radon--Nikod\'ym theorem}
and will be given in Chapter~11.  For $p=2$, the result follows immediately
from the Riesz representation theorem for Hilbert spaces.
\end{remark}

% ------------------------------------------------------------------
\section{The Canonical Injection and Bidual}\label{sec:bidual}
% ------------------------------------------------------------------

\begin{definition}[Bidual and canonical injection]\label{def:bidual}
\index{bidual}
\index{canonical injection}
\index{$J$, canonical injection}
The \textbf{bidual} of $E$ is $E'' = (E')'$.  The \textbf{canonical injection}
(or \textbf{evaluation map}) is
\[
  J: E \longrightarrow E'',
  \qquad J(x)(f) = f(x),\quad x\in E,\; f\in E'.
\]
\end{definition}

\begin{proposition}\label{prop:J-isometry}
\index{canonical injection!is isometry}
The canonical injection $J$ is a linear isometry from $E$ into $E''$.
\end{proposition}

\begin{proof}
Linearity is clear.  For each $x\in E$,
\[
  \norm{J(x)}_{E''} = \sup_{\norm{f}_{E'}\le 1} |J(x)(f)|
  = \sup_{\norm{f}\le 1} |f(x)|.
\]
On the one hand, $|f(x)|\le\norm{f}\,\norm{x}\le\norm{x}$ for $\norm{f}\le 1$,
so $\norm{J(x)}\le\norm{x}$.  On the other hand,
Corollary~\ref{cor:hb-separation} provides $f_0\in E'$ with $\norm{f_0}=1$ and
$f_0(x)=\norm{x}$, so $\norm{J(x)}\ge |f_0(x)| = \norm{x}$.
\end{proof}

\begin{remark}\label{rem:J-not-surjective}
The map $J$ is always injective (since it is an isometry), but it need not be
surjective.  When $J$ is surjective we say $E$ is \emph{reflexive}.
\end{remark}

% ------------------------------------------------------------------
\section{Reflexive Spaces}\label{sec:reflexive}
% ------------------------------------------------------------------

\begin{definition}[Reflexive space]\label{def:reflexive}
\index{reflexive space}
\index{space!reflexive}
A Banach space $E$ is called \textbf{reflexive} if the canonical injection
$J:E\to E''$ is surjective (hence an isometric isomorphism).
\end{definition}

\begin{remark}\label{rem:reflexive-isometric}
\index{James' space}
It is essential that the isomorphism onto $E''$ be the \emph{canonical} one.
James constructed a Banach space $J$ that is isometrically isomorphic to $J''$,
yet is not reflexive because the canonical injection is not surjective.
\end{remark}

\begin{example}\label{ex:reflexive-examples}
\index{reflexive space!examples}
\leavevmode
\begin{enumerate}[label=(\roman*)]
  \item Every finite-dimensional Banach space is reflexive (since $\dim E = \dim E' = \dim E''$).
  \item Every Hilbert space is reflexive (by the Riesz representation theorem).
  \item $L^p(\mu)$ is reflexive for $1 < p < \infty$ (as we shall prove below).
  \item $\ell^p$ is reflexive for $1 < p < \infty$.
  \item $\ell^1$, $\ell^\infty$, $c_0$, $L^1(\mu)$, and $L^\infty(\mu)$ are
        \textbf{not} reflexive (for non-trivial $\mu$).
\end{enumerate}
\end{example}

\begin{theorem}[$\ell^p$ is reflexive for $1<p<\infty$]\label{thm:lp-reflexive}
\index{$\ell^p$ space!reflexivity}
\index{reflexive space!$\ell^p$}
For $1<p<\infty$, the space $\ell^p$ is reflexive.
\end{theorem}

\begin{proof}
Let $p' = p/(p-1)$ be the conjugate exponent, so that $1<p'<\infty$ and
$(p')' = p$.  By Theorem~\ref{thm:dual-lp}, the isometric isomorphisms
\[
  (\ell^p)' \cong \ell^{p'}, \qquad
  (\ell^p)'' = ((\ell^p)')' \cong (\ell^{p'})' \cong \ell^{(p')'} = \ell^p
\]
hold.  We must verify that the composite isomorphism $(\ell^p)''\cong\ell^p$
is precisely the canonical injection $J$.

Let $x = (x_n)\in\ell^p$.  Under the identification
$\Phi: \ell^{p'}\xrightarrow{\sim}(\ell^p)'$, a functional $f\in(\ell^p)'$
corresponds to $y = (y_n)\in\ell^{p'}$ via $f(x)=\sum x_n y_n$.  Then
\[
  J(x)(f) = f(x) = \sum_{n=1}^\infty x_n y_n = \Psi(x)(y),
\]
where $\Psi:\ell^p\to(\ell^{p'})'$ is the canonical identification from
Theorem~\ref{thm:dual-lp} applied to $\ell^{p'}$.  Under the identification
$(\ell^p)'' \cong (\ell^{p'})' \cong \ell^p$, the element $J(x)$ corresponds
to $x$ itself.  Since $J$ is an isometry and the identification maps to all of
$\ell^p$, $J$ is surjective.
\end{proof}

\begin{theorem}[$\ell^1$ is not reflexive]\label{thm:l1-not-reflexive}
\index{$\ell^1$ space!not reflexive}
The space $\ell^1$ is not reflexive.
\end{theorem}

\begin{proof}
We have $(\ell^1)' \cong \ell^\infty$ by Theorem~\ref{thm:dual-lp}.
If $\ell^1$ were reflexive, then $(\ell^1)'' \cong \ell^1$ canonically,
whence $(\ell^\infty)' \cong \ell^1$.  But $\ell^\infty$ is non-separable
(it contains the uncountable set $\{(\varepsilon_n):\varepsilon_n\in\{0,1\}\}$
whose pairwise distances are all~$1$), while the dual of a separable space
need not be separable, but the predual of a separable space must be separable
(since a separable dual implies a separable predual is false in general).

More directly: we exhibit a functional in $(\ell^\infty)'$ not in $J(\ell^1)$.
Let $\mathcal{U}$ be a free ultrafilter\index{ultrafilter} on $\N$.  Define
$\Lambda:\ell^\infty\to\K$ by $\Lambda(x) = \lim_\mathcal{U} x_n$.  Then
$\Lambda$ is a bounded linear functional with $\norm{\Lambda}=1$, and
$\Lambda(e_n)=0$ for every standard basis vector $e_n$.  If $\Lambda = J(a)$
for some $a\in\ell^1$, then $a_n = \Lambda(e_n) = 0$ for all $n$, hence
$a=0$, contradicting $\norm{\Lambda}=1$.
\end{proof}

\begin{theorem}[$c_0$ is not reflexive]\label{thm:c0-not-reflexive}
\index{$c_0$ space!not reflexive}
The space $c_0$ is not reflexive.
\end{theorem}

\begin{proof}
We have $c_0' \cong \ell^1$ (Theorem~\ref{thm:dual-c0}) and
$c_0'' \cong (\ell^1)' \cong \ell^\infty$.  The canonical injection
$J:c_0\to\ell^\infty$ is the inclusion map.  Since $c_0\subsetneq\ell^\infty$
(for instance, the constant sequence $(1,1,1,\ldots)\in\ell^\infty\setminus c_0$),
$J$ is not surjective.
\end{proof}

% ------------------------------------------------------------------
\section{James' Characterization of Reflexivity}\label{sec:james-char}
% ------------------------------------------------------------------

The following deep result provides a purely geometric characterization of
reflexivity.

\begin{theorem}[{James, 1964}]\label{thm:james}
\index{James' theorem}
\index{reflexive space!James' characterization}
A Banach space $E$ is reflexive if and only if every continuous linear
functional $f\in E'$ attains its norm, i.e., there exists $x_0\in E$ with
$\norm{x_0}=1$ and $|f(x_0)|=\norm{f}$.
\end{theorem}

\begin{remark}
The ``only if'' direction is straightforward: if $E$ is reflexive, then $B_E$
is weakly compact (Kakutani's theorem, Theorem~\ref{thm:kakutani}), and any
$f\in E'$ is weakly continuous, hence attains its supremum on $B_E$.
The ``if'' direction is considerably harder; we refer to
\cite{James1964} for the proof.
\end{remark}

\begin{corollary}\label{cor:james-c0}
The space $c_0$ is not reflexive because the functional
$f(x) = \sum_{n=1}^\infty x_n/2^n$ satisfies $\norm{f}=1$ but does not
attain its norm on the closed unit ball of $c_0$.
\end{corollary}

\begin{proof}
We have $\norm{f} = \sum 2^{-n} = 1$.  If $x\in B_{c_0}$ with $\norm{x}_\infty\le 1$,
then $|f(x)|\le\sum |x_n|/2^n$.  Equality $|f(x)|=1$ requires $|x_n|=1$ for
all $n$, but then $x_n\not\to 0$, contradicting $x\in c_0$.
\end{proof}

% ------------------------------------------------------------------
\section{Subspaces and Quotients of Reflexive Spaces}\label{sec:sub-quot-reflexive}
% ------------------------------------------------------------------

\begin{theorem}\label{thm:subspace-reflexive}
\index{reflexive space!closed subspace}
Every closed subspace of a reflexive Banach space is reflexive.
\end{theorem}

\begin{proof}
Let $E$ be reflexive and $F\subset E$ a closed subspace.  We must show that
the canonical injection $J_F: F\to F''$ is surjective.

Let $\xi\in F''$.  Define $\eta\in E''$ by
\[
  \eta(f) = \xi(f\big|_F), \qquad f\in E'.
\]
This is well-defined: $f\big|_F\in F'$ for every $f\in E'$, and the restriction
map $R:E'\to F'$, $R(f)=f\big|_F$, is bounded with $\norm{R}\le 1$
(by Hahn--Banach it is surjective).  Moreover,
$|\eta(f)| = |\xi(f\big|_F)| \le \norm{\xi}\,\norm{f\big|_F}_{F'}
\le \norm{\xi}\,\norm{f}_{E'}$,
so $\eta\in E''$ with $\norm{\eta}\le\norm{\xi}$.

Since $E$ is reflexive, $\eta = J_E(x)$ for some $x\in E$, meaning
$f(x) = \eta(f) = \xi(f\big|_F)$ for all $f\in E'$.

We claim $x\in F$.  If $x\notin F$, by Hahn--Banach there exists $f\in E'$
with $f\big|_F = 0$ and $f(x) = 1$.  But then
$1 = f(x) = \xi(f\big|_F) = \xi(0) = 0$, a contradiction.

So $x\in F$, and for every $g\in F'$, extend $g$ to $\tilde{g}\in E'$ (Hahn--Banach):
\[
  J_F(x)(g) = g(x) = \tilde{g}(x) = \xi(\tilde{g}\big|_F) = \xi(g).
\]
Hence $J_F(x) = \xi$, proving surjectivity.
\end{proof}

\begin{theorem}\label{thm:quotient-reflexive}
\index{reflexive space!quotient}
\index{quotient space!of reflexive space}
If $E$ is a reflexive Banach space and $F\subset E$ is a closed subspace,
then the quotient space $E/F$ is reflexive.
\end{theorem}

\begin{proof}
Let $\pi:E\to E/F$ denote the quotient map\index{quotient map}.  The adjoint
$\pi':(E/F)'\to E'$ is an isometric embedding whose image is
$F^\perp = \{f\in E': f\big|_F = 0\}$\index{annihilator}.

Let $\xi\in(E/F)''$.  Define $\eta\in E''$ by
$\eta(f) = \xi(\psi)$ where $\psi\in(E/F)'$ is the unique functional with
$\pi'(\psi)=f$, whenever $f\in F^\perp$.  For general $f\in E'$, we note that
the composition $f\mapsto \xi\circ(\pi')^{-1}$ is defined on $F^\perp$ and extends
by Hahn--Banach; however, a cleaner approach is as follows.

Define $T:(E/F)'\to\K$ by $T = \xi$.  Consider the adjoint
$\pi'': (E/F)''\to E''$ of $\pi'$, defined by
$\pi''(\xi)(f) = \xi((\pi')^{-1}(f))$ for $f\in\ran(\pi')=F^\perp$---but
this requires care.  Instead, we use the transpose construction.

Let $\eta = \xi\circ\pi^{t}$ where $\pi^t: E'\to (E/F)'$ is defined by
$\pi^t(f)(\dot{x}) = f(x)$, valid when $f\in F^\perp$.  Since $\pi^t$ is not
defined on all of $E'$, we proceed differently.

Consider the map $\alpha: E' \to (E/F)'$ given by
$\alpha(f) = 0$ if we compose differently.  Let us use the direct approach.

Define $\zeta \in (F^\perp)' $ by $\zeta(f) = \xi(\psi_f)$ where
$\psi_f\in(E/F)'$ corresponds to $f\in F^\perp$ via the isometry $\pi'$.
Then $\zeta$ is bounded, and by Hahn--Banach we extend $\zeta$ to
$\eta\in E''$.  Since $E$ is reflexive, $\eta = J_E(x)$ for some $x\in E$,
meaning $f(x)=\zeta(f)$ for all $f\in F^\perp$.

Set $\dot{x} = \pi(x) = x + F \in E/F$.  For any $\psi\in(E/F)'$, let
$f = \pi'(\psi) = \psi\circ\pi \in F^\perp$.  Then
\[
  J_{E/F}(\dot{x})(\psi) = \psi(\dot{x}) = \psi(\pi(x)) = f(x)
  = \zeta(f) = \xi(\psi).
\]
Hence $J_{E/F}(\dot{x}) = \xi$ and $E/F$ is reflexive.
\end{proof}

\begin{corollary}\label{cor:reflexive-iff-dual}
\index{reflexive space!dual of}
A Banach space $E$ is reflexive if and only if $E'$ is reflexive.
\end{corollary}

\begin{proof}
($\Rightarrow$) Suppose $E$ is reflexive, so $J_E:E\xrightarrow{\sim}E''$.
Let $\Phi\in E'''$.  Define $f\in E'$ by $f(x) = \Phi(J_E(x))$ (well-defined
since $J_E$ is surjective).  Then for any $\xi\in E''$, writing
$\xi = J_E(x)$, we get $J_{E'}(f)(\xi) = \xi(f) = J_E(x)(f) = f(x)
= \Phi(J_E(x)) = \Phi(\xi)$.  So $J_{E'}(f) = \Phi$.

($\Leftarrow$) Suppose $E'$ is reflexive.  Since $J_E(E)$ is a closed subspace
of $E'' = (E')'$ (being the isometric image of a Banach space), and $E'$ is
reflexive, Theorem~\ref{thm:subspace-reflexive} shows every closed subspace
of $E''$ is reflexive.  But if $J_E(E)\ne E''$, then by Hahn--Banach there
exists $\Phi\in E'''\setminus\{0\}$ vanishing on $J_E(E)$.  Since $E'$ is
reflexive, $\Phi = J_{E'}(f)$ for some $f\in E'\setminus\{0\}$.
Then for all $x\in E$:
$0 = \Phi(J_E(x)) = J_{E'}(f)(J_E(x)) = J_E(x)(f) = f(x)$,
so $f=0$, contradicting $f\ne 0$.  Hence $J_E(E)=E''$.
\end{proof}

% ------------------------------------------------------------------
\section{The Eberlein--\v{S}mulian Theorem}\label{sec:eberlein-smulian}
% ------------------------------------------------------------------

\begin{theorem}[{Eberlein--\v{S}mulian}]\label{thm:eberlein-smulian}
\index{Eberlein--\v{S}mulian theorem}
\index{weak compactness!sequential characterization}
Let $E$ be a Banach space and $A\subset E$.  The following are equivalent:
\begin{enumerate}[label=(\roman*)]
  \item $A$ is relatively weakly compact (its weak closure is weakly compact).
  \item $A$ is relatively weakly sequentially compact (every sequence in $A$
        has a weakly convergent subsequence with limit in $E$).
  \item $A$ is weakly limit-point compact (every infinite subset of $A$ has
        a weak limit point in $E$).
\end{enumerate}
\end{theorem}

\begin{remark}\label{rem:eberlein-smulian-discussion}
This is remarkable because the weak topology on an infinite-dimensional space
is never metrizable on the whole space, so one cannot expect sequential and
topological compactness to coincide in general.  The Eberlein--\v{S}mulian
theorem\index{Eberlein--\v{S}mulian theorem!significance of} shows they do
coincide for the weak topology of a Banach space.  The proof is technical; see
\cite{DunfordSchwartz1958} or \cite{Megginson1998} for details.
\end{remark}

% ------------------------------------------------------------------
\section{Uniformly Convex Spaces and the Milman--Pettis Theorem}%
\label{sec:uniform-convexity}
% ------------------------------------------------------------------

\begin{definition}[Uniformly convex space]\label{def:uniformly-convex}
\index{uniformly convex space}
\index{space!uniformly convex}
A normed space $E$ is \textbf{uniformly convex} if for every $\varepsilon>0$
there exists $\delta>0$ such that
\[
  \norm{x}\le 1,\;\norm{y}\le 1,\;\norm{x-y}\ge\varepsilon
  \quad\Longrightarrow\quad
  \Bigl\|\frac{x+y}{2}\Bigr\| \le 1-\delta.
\]
The function $\delta_E(\varepsilon) = \inf\bigl\{1-\norm{\frac{x+y}{2}}:
\norm{x}\le 1,\norm{y}\le 1,\norm{x-y}\ge\varepsilon\bigr\}$
is called the \textbf{modulus of convexity}\index{modulus of convexity} of $E$.
\end{definition}

\begin{example}\label{ex:uniformly-convex}
\index{uniformly convex space!examples}
\leavevmode
\begin{enumerate}[label=(\roman*)]
  \item Every Hilbert space is uniformly convex.  Indeed, by the parallelogram
        law\index{parallelogram law},
        $\norm{x+y}^2 + \norm{x-y}^2 = 2\norm{x}^2+2\norm{y}^2 \le 4$,
        so $\norm{\frac{x+y}{2}}^2 \le 1-\frac{\varepsilon^2}{4}$, giving
        $\delta(\varepsilon) \ge 1-\sqrt{1-\varepsilon^2/4}$.
  \item $L^p(\mu)$ and $\ell^p$ are uniformly convex for $1<p<\infty$
        (Clarkson's inequalities\index{Clarkson's inequalities}).
  \item $L^1$, $L^\infty$, $\ell^1$, $\ell^\infty$ are \emph{not} uniformly convex.
\end{enumerate}
\end{example}

\begin{proposition}\label{prop:uc-strictly-convex}
\index{strictly convex space}
\index{uniformly convex space!is strictly convex}
Every uniformly convex space is strictly convex: if $\norm{x}=\norm{y}=1$
and $x\ne y$, then $\norm{\frac{x+y}{2}}<1$.
\end{proposition}

\begin{proof}
Set $\varepsilon = \norm{x-y}>0$.  By uniform convexity,
$\norm{\frac{x+y}{2}}\le 1-\delta(\varepsilon)<1$.
\end{proof}

\begin{theorem}[{Milman--Pettis, 1939}]\label{thm:milman-pettis}
\index{Milman--Pettis theorem}
\index{uniformly convex space!is reflexive}
Every uniformly convex Banach space is reflexive.
\end{theorem}

\begin{proof}
Let $E$ be a uniformly convex Banach space.  We must show that
$J:E\to E''$ is surjective.  Let $\xi\in E''$ with $\norm{\xi}=1$.
We will find $x\in E$ with $J(x)=\xi$.

\medskip\noindent\textbf{Step 1.}
By definition of the norm in $E''$ and the fact that $J(B_E)$ is dense in
$B_{E''}$ for the weak-$*$ topology of $E''$ (this is Goldstine's theorem,
Theorem~\ref{thm:goldstine}, proved in Chapter~8), for each $n\ge 1$ there
exists $x_n\in B_E$ such that
\begin{equation}\label{eq:mp-approx}
  |\xi(f) - f(x_n)| < \frac{1}{n} \qquad
  \text{for } f \text{ in a finite set } F_n\subset E'.
\end{equation}
However, we need a cleaner approach.  We use the following key fact.

\medskip\noindent\textbf{Step 1 (revised).}
Since $\norm{\xi}_{E''}=1$, for each $n$ there exists $f_n\in E'$ with
$\norm{f_n}=1$ and $\xi(f_n)>1-\frac{1}{n}$.  Choose $x_n\in E$ with
$\norm{x_n}\le 1$ and $f_n(x_n)>1-\frac{1}{n}$.

\medskip\noindent\textbf{Step 2: $(x_n)$ is Cauchy.}
Let $\varepsilon>0$ and let $\delta = \delta_E(\varepsilon)>0$ be the modulus
of uniform convexity.  We claim that for $m,n$ large enough,
$\norm{x_m-x_n}<\varepsilon$.

Suppose for contradiction that $\norm{x_m-x_n}\ge\varepsilon$ for infinitely
many pairs.  Then $\norm{\frac{x_m+x_n}{2}}\le 1-\delta$.  But for any
$f\in E'$ with $\norm{f}=1$,
\[
  f\Bigl(\frac{x_m+x_n}{2}\Bigr) = \frac{f(x_m)+f(x_n)}{2}.
\]
In particular, choosing appropriately: we take $f = f_n$ for large $n$ and
observe that we need $f$ to test \emph{both} $x_m$ and $x_n$.  Let us
refine the argument.

We use a single functional.  Since $\norm{\xi}=1$, there exists a net
$(x_\alpha)$ in $B_E$ with $J(x_\alpha)\to\xi$ weak-$*$ in $E''$.  Instead,
we use the following direct approach.

For $m,n$ large, pick $g\in E'$ with $\norm{g}=1$ and
$\xi(g)>1-\delta/2$.  Then for $n$ large enough (depending on $g$), since
$J(B_E)$ is weak-$*$ dense in $B_{E''}$, we can ensure $g(x_n)>1-\delta/2$
and similarly $g(x_m)>1-\delta/2$.  Hence
\[
  g\Bigl(\frac{x_m+x_n}{2}\Bigr) > 1-\frac{\delta}{2}.
\]
But $\norm{\frac{x_m+x_n}{2}}\le 1-\delta$ (by uniform convexity if
$\norm{x_m-x_n}\ge\varepsilon$), which gives
$g(\frac{x_m+x_n}{2})\le 1-\delta$, a contradiction.

More precisely: we construct the sequence as follows.  Pick $f\in E'$ with
$\norm{f}=1$ and $\xi(f)>1-\delta/4$.  By Goldstine's theorem, for each $n$
there exists $x_n\in B_E$ with $|f(x_n)-\xi(f)|<\frac{1}{n}$ (and also
with $|g(x_n)-\xi(g)|<\frac{1}{n}$ for finitely many $g$'s we pre-select).
In particular, $f(x_n)\to\xi(f)>1-\delta/4$, so for large $n$,
$\operatorname{Re} f(x_n)>1-\delta/2$.

If $\norm{x_m-x_n}\ge\varepsilon$, then $\norm{\frac{x_m+x_n}{2}}\le 1-\delta$, so
\[
  \operatorname{Re} f\Bigl(\frac{x_m+x_n}{2}\Bigr) \le \Bigl\|\frac{x_m+x_n}{2}\Bigr\| \le 1-\delta.
\]
But $\operatorname{Re} f(\frac{x_m+x_n}{2}) = \frac{1}{2}(\operatorname{Re} f(x_m)+\operatorname{Re} f(x_n))
> 1-\delta/2$ for $m,n$ large, a contradiction.

\medskip\noindent\textbf{Step 3: Conclusion.}
Since $E$ is complete, $x_n\to x$ for some $x\in E$ with $\norm{x}\le 1$.
For any $f\in E'$, $J(x)(f)=f(x)=\lim f(x_n)$.  By construction (choosing
the sequence so that $f(x_n)\to\xi(f)$ for all $f$ in a countable dense
subset of $E'$, and using the fact that $\norm{J(x_n)-\xi}\to 0$ in $E''$
follows from the Cauchy property), we obtain $J(x)=\xi$.

Let us make Step 3 precise.  We have shown that any sequence $(x_n)\subset B_E$
with $f_n(x_n)\to 1$ for some sequence $(f_n)$ in the unit sphere of $E'$
with $\xi(f_n)\to 1$ must be Cauchy.  Now use a diagonal argument: let
$(g_k)_{k\ge 1}$ be a sequence dense in $S_{E'}$ (or any countable family).
By Goldstine's theorem, for each $n$ we can find $x_n\in B_E$ with
$|g_k(x_n)-\xi(g_k)|<\frac{1}{n}$ for $k=1,\ldots,n$.  The Cauchy argument
shows $x_n\to x$.  Then $g_k(x) = \lim g_k(x_n) = \xi(g_k)$ for all $k$.
By density and continuity, $f(x)=\xi(f)$ for all $f\in E'$, i.e.\ $J(x)=\xi$.

(Note: if $E'$ is not separable, one modifies the argument to work with nets,
or uses a different approach.  The argument works for arbitrary $E$ because
uniform convexity forces Cauchy behavior without needing separability of $E'$;
one applies the argument to any single $\xi\in E''$.)
\end{proof}

\begin{corollary}\label{cor:Lp-reflexive}
\index{$L^p$ space!reflexivity}
For $1<p<\infty$, the spaces $L^p(\mu)$ and $\ell^p$ are reflexive.
\end{corollary}

\begin{proof}
These spaces are uniformly convex by Clarkson's inequalities, hence reflexive
by the Milman--Pettis theorem.
\end{proof}

% ------------------------------------------------------------------
\section{Exercises}\label{sec:exercises-dual-reflexivity}
% ------------------------------------------------------------------

\begin{exercise}[$\star$]\label{exo:dual-finite-dim}
\index{dual space!finite-dimensional}
Let $E$ be a finite-dimensional normed space.  Show that $E'\cong E$ (as
vector spaces) and that every linear functional on $E$ is continuous.
\end{exercise}

\begin{exercise}[$\star$]\label{exo:annihilator}
\index{annihilator}
Let $E$ be a normed space and $A\subset E$.  The \emph{annihilator} of $A$ is
$A^\perp = \{f\in E': f(a)=0\;\forall a\in A\}$.
\begin{enumerate}[label=(\alph*)]
  \item Show that $A^\perp$ is a closed subspace of $E'$.
  \item If $F$ is a closed subspace of $E$, show that $F^\perp\cong (E/F)'$
        isometrically.
  \item Show that $E/F^\perp \cong F'$ if $E$ is reflexive.
\end{enumerate}
\end{exercise}

\begin{exercise}[$\star\star$]\label{exo:separable-dual}
\index{separable space!dual of}
Prove that if $E'$ is separable, then $E$ is separable.
\textit{Hint:} For each $f_n$ in a dense subset of $S_{E'}$, choose $x_n$
with $|f_n(x_n)|>\frac{1}{2}\norm{f_n}$.  Show that $\spn\{x_n\}$ is dense.
\end{exercise}

\begin{exercise}[$\star\star$]\label{exo:reflexive-dual-separable}
\index{reflexive space!separability}
Show that a reflexive Banach space is separable if and only if its dual is
separable.
\end{exercise}

\begin{exercise}[$\star$]\label{exo:subspace-quotient-reflexive}
Let $E$ be a Banach space and $F$ a closed subspace.  Show that $E$ is
reflexive if and only if both $F$ and $E/F$ are reflexive.
\end{exercise}

\begin{exercise}[$\star\star$]\label{exo:linfty-dual}
\index{$\ell^\infty$ space!not reflexive}
\begin{enumerate}[label=(\alph*)]
  \item Show that $\ell^\infty$ is not separable.
  \item Deduce that $\ell^1$ is not reflexive (without using ultrafilters).
  \textit{Hint:} Use Exercise~\ref{exo:separable-dual}.
\end{enumerate}
\end{exercise}

\begin{exercise}[$\star\star$]\label{exo:norm-attaining}
\index{norm-attaining functional}
Let $E$ be a Banach space and $f\in E'$ with $\norm{f}=1$.  We say $f$
\emph{attains its norm} if there exists $x\in B_E$ with $|f(x)|=1$.
\begin{enumerate}[label=(\alph*)]
  \item Show that in a reflexive space, every $f\in E'$ attains its norm.
  \item Find an explicit $f\in c_0'$ that does not attain its norm.
\end{enumerate}
\end{exercise}

\begin{exercise}[$\star\star$]\label{exo:uc-unique-nearest}
\index{uniformly convex space!nearest point}
Let $E$ be a uniformly convex Banach space and $C\subset E$ a nonempty closed
convex subset.  Show that for every $x\in E$ there exists a unique $y\in C$
with $\norm{x-y}=\inf_{z\in C}\norm{x-z}$.
\end{exercise}

\begin{exercise}[$\star\star\star$]\label{exo:day}
\index{Day's theorem}
(Day) Show that $\ell^1$ admits no equivalent uniformly convex norm.
\textit{Hint:} Use the fact that $\ell^1$ is not reflexive and the
Milman--Pettis theorem.
\end{exercise}

\begin{exercise}[$\star$]\label{exo:dual-direct-sum}
\index{dual space!of direct sum}
Let $E_1,E_2$ be Banach spaces and $E = E_1\oplus_p E_2$ with the $\ell^p$-norm
($1\le p<\infty$).  Show that $E'\cong E_1'\oplus_{p'} E_2'$ isometrically.
\end{exercise}

\begin{exercise}[$\star\star$]\label{exo:natural-embedding-range}
\index{canonical injection!range}
Let $E$ be a normed space.  Show that $J(E)$ is dense in $E''$ if and only if
$E'$ is separating (which it always is, by Hahn--Banach).  Conclude that
$J(E)$ is always dense in $E''$ for the weak-$*$ topology $\sigma(E'',E')$.
\end{exercise}

\begin{exercise}[$\star\star\star$]\label{exo:james-space}
\index{James' space}
(James' space) Consider the space $\mathcal{J}$ of all sequences $x=(x_n)$ with
$x_n\to 0$ and finite \emph{James norm}:
\[
  \norm{x}_{\mathcal{J}} = \sup\Bigl\{
    \Bigl(\sum_{i=1}^{k-1}(x_{p_{i+1}}-x_{p_i})^2\Bigr)^{1/2}:
    k\ge 2,\; p_1<p_2<\cdots<p_k
  \Bigr\}.
\]
\begin{enumerate}[label=(\alph*)]
  \item Verify that $\mathcal{J}$ is a Banach space.
  \item Show that $\mathcal{J}$ is isometrically isomorphic to $\mathcal{J}''$
        but is not reflexive ($\codim J(\mathcal{J})=1$ in $\mathcal{J}''$).
\end{enumerate}
\end{exercise}


%%%%%%%%%%%%%%%%%%%%%%%%%%%%%%%%%%%%%%%%%%%%%%%%%%%%%%%%%%%%
%% CHAPTER 8 — Weak Topologies
%%%%%%%%%%%%%%%%%%%%%%%%%%%%%%%%%%%%%%%%%%%%%%%%%%%%%%%%%%%%
\chapter{Weak Topologies}\label{ch:weak-topologies}
\minitoc

\section{Initial Topologies: Review}\label{sec:initial-topology}

We begin with a review of the general construction that underlies both the
weak topology and the weak-$*$ topology.

\begin{definition}[Initial topology]\label{def:initial-topology}
\index{initial topology}
\index{topology!initial}
Let $X$ be a set and $\{(Y_i,\tau_i)\}_{i\in I}$ a family of topological
spaces.  For each $i\in I$, let $\varphi_i: X\to Y_i$ be a map.  The
\textbf{initial topology} on $X$ with respect to the family
$(\varphi_i)_{i\in I}$ is the coarsest topology on $X$ making each
$\varphi_i$ continuous.  A subbasis for this topology is given by
\[
  \bigl\{\,\varphi_i^{-1}(U_i) : i\in I,\; U_i\in\tau_i\,\bigr\}.
\]
\end{definition}

\begin{proposition}\label{prop:initial-convergence}
\index{initial topology!convergence in}
In the initial topology, a net $(x_\alpha)_{\alpha\in A}$ converges to $x$
in $X$ if and only if $\varphi_i(x_\alpha)\to\varphi_i(x)$ in $Y_i$ for
every $i\in I$.
\end{proposition}

\begin{proposition}\label{prop:initial-hausdorff}
\index{initial topology!Hausdorff}
The initial topology is Hausdorff if and only if the family $(\varphi_i)_{i\in I}$
separates points of $X$: for every $x\ne y$ in $X$, there exists $i\in I$
with $\varphi_i(x)\ne\varphi_i(y)$.
\end{proposition}

% ------------------------------------------------------------------
\section{The Weak Topology \texorpdfstring{$\sigma(E,E')$}{sigma(E,E')}}%
\label{sec:weak-topology}
% ------------------------------------------------------------------

\begin{definition}[Weak topology]\label{def:weak-topology}
\index{weak topology}
\index{topology!weak}
\index{$\sigma(E,E')$}
Let $E$ be a normed space.  The \textbf{weak topology} on $E$, denoted
$\sigma(E,E')$, is the initial topology on $E$ with respect to the family
of all continuous linear functionals:
\[
  \sigma(E,E') = \text{initial topology w.r.t.\ } \{f:E\to\K\}_{f\in E'}.
\]
\end{definition}

\begin{proposition}[Subbasis and neighborhood base]\label{prop:weak-nbhd}
\index{weak topology!neighborhood base}
A subbasis of open sets for $\sigma(E,E')$ is given by
\[
  \bigl\{f^{-1}(U) : f\in E',\; U\subset\K \text{ open}\bigr\}.
\]
A base of neighborhoods of a point $x_0\in E$ is given by the sets
\[
  V(x_0; f_1,\ldots,f_n;\varepsilon)
  = \bigl\{x\in E : |f_k(x)-f_k(x_0)| < \varepsilon,\; k=1,\ldots,n\bigr\}
\]
where $n\ge 1$, $f_1,\ldots,f_n\in E'$, and $\varepsilon>0$.
\end{proposition}

\begin{theorem}\label{thm:weak-hausdorff}
\index{weak topology!is Hausdorff}
The weak topology $\sigma(E,E')$ is Hausdorff.
\end{theorem}

\begin{proof}
By Proposition~\ref{prop:initial-hausdorff}, it suffices to show that $E'$
separates points of $E$.  This is exactly the content of
Corollary~\ref{cor:hb-separation} (a consequence of the Hahn--Banach theorem).
\end{proof}

\begin{proposition}[Weak topology is coarser]\label{prop:weak-coarser}
\index{weak topology!coarser than norm topology}
The weak topology $\sigma(E,E')$ is coarser than (or equal to) the norm topology.
If $\dim E = \infty$, the inclusion is strict: $\sigma(E,E') \subsetneq \tau_{\norm{\cdot}}$.
\end{proposition}

\begin{proof}
Every $f\in E'$ is norm-continuous, so every weakly open set is norm-open.
For the strict inclusion, note that the open unit ball $B(0,1)$ is norm-open.
If it were weakly open, it would contain a basic weak neighborhood
$V(0;f_1,\ldots,f_n;\varepsilon)$.  But
$\bigcap_{k=1}^n \Ker f_k$ is a subspace of codimension at most $n$, hence
of infinite dimension when $\dim E=\infty$.  Any nonzero $v$ in this
intersection satisfies $tv\in V(0;f_1,\ldots,f_n;\varepsilon)$ for all $t$,
so $V$ contains the entire line through $v$, which is not contained in
$B(0,1)$.
\end{proof}

\begin{proposition}\label{prop:weak-non-metrizable}
\index{weak topology!non-metrizable}
If $E$ is an infinite-dimensional normed space, the weak topology on $E$ is
not metrizable (on the whole space $E$).
\end{proposition}

\begin{proof}
In a metrizable space, a point in the closure of a set $A$ is the limit of a
sequence from $A$.  Consider the weak topology on an infinite-dimensional space.
Every weak neighborhood of $0$ contains a subspace of finite codimension
(as shown above), hence is unbounded.  In particular, $B_E$ is not weakly
open.

More formally: if the weak topology were metrizable, then since $E$ is a
topological vector space under the weak topology, it would be first countable.
But a locally convex space is first countable (at the origin) if and only if
its topology is generated by countably many seminorms, which for $\sigma(E,E')$
would require $E'$ to be the span of countably many functionals.  When
$\dim E = \infty$, $E'$ is infinite-dimensional (by Hahn--Banach), and in fact
$E'$ cannot be the countable union of finite-dimensional subspaces (by Baire's
theorem if $E'$ is a Banach space---which it is---then it cannot be countable-dimensional).
\end{proof}

% ------------------------------------------------------------------
\section{Weak Convergence}\label{sec:weak-convergence}
% ------------------------------------------------------------------

\begin{definition}[Weak convergence of sequences]\label{def:weak-convergence}
\index{weak convergence}
\index{convergence!weak}
\index{$\weakto$}
A sequence $(x_n)$ in a normed space $E$ \textbf{converges weakly}\index{weak convergence!of sequences}
to $x\in E$ if
\[
  f(x_n) \to f(x) \qquad \text{for every } f\in E'.
\]
We write $x_n\weakto x$.
\end{definition}

\begin{proposition}\label{prop:weak-vs-norm}
\index{weak convergence!vs.\ norm convergence}
\leavevmode
\begin{enumerate}[label=(\roman*)]
  \item Norm convergence implies weak convergence: if $\norm{x_n-x}\to 0$,
        then $x_n\weakto x$.
  \item The converse is false in general (in infinite-dimensional spaces).
  \item In finite-dimensional spaces, weak and norm convergence coincide.
\end{enumerate}
\end{proposition}

\begin{proof}
(i) If $\norm{x_n-x}\to 0$ and $f\in E'$, then
$|f(x_n)-f(x)|\le\norm{f}\norm{x_n-x}\to 0$.

(ii) In $\ell^2$, the standard basis vectors $e_n\weakto 0$ (since for any
$f\in(\ell^2)'\cong\ell^2$ given by $y=(y_k)$, $f(e_n)=y_n\to 0$), but
$\norm{e_n}=1\not\to 0$.

(iii) In $\R^d$, the functionals $f_k(x)=x_k$ ($k=1,\ldots,d$) span $E'$,
so weak convergence is equivalent to coordinate-wise convergence, which is
equivalent to norm convergence.
\end{proof}

\begin{example}\label{ex:weak-convergence-examples}
\index{weak convergence!examples}
\leavevmode
\begin{enumerate}[label=(\roman*)]
  \item In $\ell^p$ ($1<p<\infty$): $e_n\weakto 0$ but $\norm{e_n}_p=1$.
  \item In $L^2([0,1])$: $\sqrt{2}\sin(2\pi n\,\cdot\,)\weakto 0$
        (Riemann--Lebesgue lemma\index{Riemann--Lebesgue lemma}).
  \item In $\ell^1$: $e_n\weakto 0$ as well (for any $y\in\ell^\infty=(\ell^1)'$,
        $y_n = y(e_n)$; but since we only need $y_n\to 0$ for $y\in c_0$...
        actually $(\ell^1)'=\ell^\infty$ so we need $y_n\to 0$ for all
        $y\in\ell^\infty$, which fails for $y=(1,1,1,\ldots)$).
        Hence $e_n\not\weakto 0$ in $\ell^1$.
  \item In $\ell^1$: the sequence $x_n = e_1+e_2+\cdots+e_n$ does not converge
        weakly (take $f=(1,1,1,\ldots)\in\ell^\infty$).
\end{enumerate}
\end{example}

The next result is fundamental and illustrates the power of the
Banach--Steinhaus theorem in the study of weak convergence.

\begin{theorem}[Weakly convergent sequences are bounded]\label{thm:weak-bounded}
\index{weak convergence!boundedness}
\index{Banach--Steinhaus theorem!application to weak convergence}
Let $E$ be a Banach space and $(x_n)$ a sequence in $E$ with
$x_n\weakto x$.  Then:
\begin{enumerate}[label=(\roman*)]
  \item $\sup_n\norm{x_n}<\infty$.
  \item $\norm{x}\le\liminf_{n\to\infty}\norm{x_n}$.
\end{enumerate}
\end{theorem}

\begin{proof}
(i) Consider the family $(J(x_n))_{n\ge 1}$ in $E'' = \mathcal{L}(E',\K)$.
For each $f\in E'$, the sequence $(J(x_n)(f))_{n\ge 1} = (f(x_n))_{n\ge 1}$
converges (to $f(x)$), hence is bounded:
$\sup_n |J(x_n)(f)| < \infty$.
By the Banach--Steinhaus theorem\index{Banach--Steinhaus theorem} (uniform
boundedness principle), applied to the family of bounded linear operators
$J(x_n):E'\to\K$,
\[
  \sup_n \norm{J(x_n)}_{E''} < \infty.
\]
Since $J$ is an isometry, $\norm{J(x_n)}=\norm{x_n}$, giving
$\sup_n\norm{x_n}<\infty$.

(ii) For every $f\in E'$ with $\norm{f}\le 1$,
\[
  |f(x)| = \lim_n |f(x_n)| \le \liminf_n \norm{f}\,\norm{x_n}
  \le \liminf_n \norm{x_n}.
\]
Taking the supremum over $\norm{f}\le 1$ gives
$\norm{x}\le\liminf_n\norm{x_n}$.
\end{proof}

\begin{remark}\label{rem:weak-lsc}
\index{norm!weak lower semicontinuity}
Part (ii) says that the norm is \emph{weakly lower semicontinuous}\index{lower semicontinuity!of norm}.
This is extremely useful in the calculus of variations.
\end{remark}

\begin{proposition}\label{prop:weak-plus-norm-convergence}
\index{weak convergence!and norm convergence}
Let $E$ be a Banach space.  If $x_n\weakto x$ and $\norm{x_n}\to\norm{x}$,
and if $E$ is uniformly convex, then $x_n\to x$ in norm.
\end{proposition}

\begin{proof}
If $\norm{x}=0$, then $\norm{x_n}\to 0$ and we are done.
Assume $\norm{x}>0$.  Let $u_n = x_n/\norm{x_n}$ and $u=x/\norm{x}$ (for
$n$ large enough).  Then $\norm{u_n}=\norm{u}=1$.  Since $E$ is uniformly
convex, if $u_n\not\to u$ in norm, there exist $\varepsilon>0$ and a
subsequence with $\norm{u_{n_k}-u}\ge\varepsilon$, giving
$\norm{\frac{u_{n_k}+u}{2}}\le 1-\delta$.

By Hahn--Banach, pick $f\in E'$ with $\norm{f}=1$ and $f(u)=1$.  Then
$f(u_{n_k})\to f(u) = 1$ (weak convergence and $\norm{x_n}\to\norm{x}$), so
$f(\frac{u_{n_k}+u}{2})\to 1$, contradicting
$\norm{\frac{u_{n_k}+u}{2}}\le 1-\delta$.
\end{proof}

% ------------------------------------------------------------------
\section{The Weak-\texorpdfstring{$*$}{*} Topology
\texorpdfstring{$\sigma(E',E)$}{sigma(E',E)}}%
\label{sec:weak-star-topology}
% ------------------------------------------------------------------

\begin{definition}[Weak-$*$ topology]\label{def:weak-star}
\index{weak-$*$ topology}
\index{topology!weak-$*$}
\index{$\sigma(E',E)$}
Let $E$ be a normed space.  The \textbf{weak-$*$ topology} on $E'$, denoted
$\sigma(E',E)$, is the initial topology on $E'$ with respect to the family
of evaluation maps $\{\hat{x}: E'\to\K\}_{x\in E}$, where
$\hat{x}(f)=f(x)$.
\end{definition}

\begin{remark}\label{rem:weak-star-vs-weak}
\index{weak-$*$ topology!vs.\ weak topology}
The weak-$*$ topology on $E'$ is coarser than or equal to the weak topology
$\sigma(E',(E')')=\sigma(E',E'')$ on $E'$.  In $\sigma(E',E)$, one only
requires convergence on evaluation functionals $\hat{x}$ for $x\in E$,
whereas $\sigma(E',E'')$ requires convergence against all of $E''$.
They coincide if and only if $E$ is reflexive.
\end{remark}

\begin{definition}[Weak-$*$ convergence]\label{def:weak-star-convergence}
\index{weak-$*$ convergence}
\index{convergence!weak-$*$}
\index{$\weakstarto$}
A net $(f_\alpha)$ in $E'$ converges to $f\in E'$ in the weak-$*$ topology
if and only if $f_\alpha(x)\to f(x)$ for every $x\in E$.  For sequences, we
write $f_n\weakstarto f$.
\end{definition}

\begin{proposition}\label{prop:weak-star-hausdorff}
\index{weak-$*$ topology!is Hausdorff}
The weak-$*$ topology is Hausdorff.
\end{proposition}

\begin{proof}
The evaluation maps $\hat{x}:E'\to\K$ separate points of $E'$: if $f\ne g$,
there exists $x\in E$ with $f(x)\ne g(x)$, i.e.\ $\hat{x}(f)\ne\hat{x}(g)$.
\end{proof}

\begin{proposition}\label{prop:weak-star-nbhd}
\index{weak-$*$ topology!neighborhood base}
A base of neighborhoods of $f_0\in E'$ for the weak-$*$ topology is
\[
  W(f_0;x_1,\ldots,x_n;\varepsilon) = \bigl\{f\in E': |f(x_k)-f_0(x_k)|<\varepsilon,
  \;k=1,\ldots,n\bigr\},
\]
where $n\ge 1$, $x_1,\ldots,x_n\in E$, $\varepsilon>0$.
\end{proposition}

% ------------------------------------------------------------------
\section{The Banach--Alaoglu Theorem}\label{sec:banach-alaoglu}
% ------------------------------------------------------------------

This is one of the most important compactness results in functional analysis.

\begin{theorem}[{Banach--Alaoglu, 1940}]\label{thm:banach-alaoglu}
\index{Banach--Alaoglu theorem}
\index{compactness!weak-$*$}
Let $E$ be a normed space.  The closed unit ball
\[
  B_{E'} = \{f\in E': \norm{f}\le 1\}
\]
is compact in the weak-$*$ topology $\sigma(E',E)$.
\end{theorem}

\begin{proof}
For each $x\in E$, define
\[
  D_x = \{\lambda\in\K : |\lambda|\le\norm{x}\},
\]
which is a compact subset of $\K$.  Consider the product space
\[
  P = \prod_{x\in E} D_x,
\]
which is compact by Tychonoff's theorem\index{Tychonoff's theorem}.  An element
of $P$ is a function $\omega: E\to\K$ with $|\omega(x)|\le\norm{x}$ for all $x$.

Define the map
\[
  \iota: B_{E'} \longrightarrow P, \qquad \iota(f) = (f(x))_{x\in E}.
\]
This is well-defined since $|f(x)|\le\norm{f}\,\norm{x}\le\norm{x}$ for
$f\in B_{E'}$.

\medskip\noindent\textbf{Step 1: $\iota$ is injective.}
If $\iota(f)=\iota(g)$, then $f(x)=g(x)$ for all $x\in E$, so $f=g$.

\medskip\noindent\textbf{Step 2: $\iota$ is a homeomorphism onto its image.}
The product topology on $P$ is the topology of pointwise convergence.
The weak-$*$ topology on $B_{E'}$ is also the topology of pointwise convergence
(convergence at each $x\in E$).  Hence $\iota$ and $\iota^{-1}$ (on $\iota(B_{E'})$)
are both continuous.

\medskip\noindent\textbf{Step 3: $\iota(B_{E'})$ is closed in $P$.}
Let $(\omega_\alpha)$ be a net in $\iota(B_{E'})$ converging to $\omega\in P$.
Each $\omega_\alpha = \iota(f_\alpha)$ for some $f_\alpha\in B_{E'}$.
Convergence in $P$ means $f_\alpha(x)\to\omega(x)$ for every $x\in E$.

We must show $\omega$ is linear.  For $x,y\in E$ and $\lambda\in\K$:
\begin{align*}
  \omega(\lambda x+y) &= \lim_\alpha f_\alpha(\lambda x+y)
  = \lim_\alpha \bigl[\lambda f_\alpha(x)+f_\alpha(y)\bigr] \\
  &= \lambda\lim_\alpha f_\alpha(x) + \lim_\alpha f_\alpha(y)
  = \lambda\,\omega(x)+\omega(y).
\end{align*}
So $\omega$ is linear and $|\omega(x)|\le\norm{x}$ for all $x$ (since
$\omega\in P$), hence $\omega\in B_{E'}$ and $\iota(\omega)=\omega$.

\medskip\noindent\textbf{Conclusion.}
$\iota(B_{E'})$ is a closed subset of the compact space $P$, hence compact.
Since $\iota$ is a homeomorphism onto its image, $B_{E'}$ is weak-$*$ compact.
\end{proof}

\begin{corollary}\label{cor:alaoglu-bounded}
\index{Banach--Alaoglu theorem!for bounded sets}
For any $r>0$, the ball $\{f\in E': \norm{f}\le r\}$ is weak-$*$ compact.
\end{corollary}

\begin{proof}
Apply the theorem to the rescaled space, or note that the map $f\mapsto rf$
is a weak-$*$ homeomorphism.
\end{proof}

\begin{corollary}\label{cor:weak-star-subsequence}
\index{weak-$*$ convergence!subsequences in separable spaces}
If $E$ is a separable normed space, then every bounded sequence in $E'$ has
a weak-$*$ convergent subsequence.
\end{corollary}

\begin{proof}
Let $(f_n)\subset E'$ with $\sup_n\norm{f_n}\le M$.  Then $(f_n)\subset MB_{E'}$,
which is weak-$*$ compact by Corollary~\ref{cor:alaoglu-bounded}.

It remains to show that the weak-$*$ topology on $MB_{E'}$ is metrizable
when $E$ is separable.  Let $(x_k)_{k\ge 1}$ be a dense sequence in $E$.
Define
\[
  d(f,g) = \sum_{k=1}^\infty \frac{1}{2^k}\,
  \frac{|f(x_k)-g(x_k)|}{1+|f(x_k)-g(x_k)|}.
\]
This is a metric on $E'$, and on bounded subsets of $E'$, it induces the
weak-$*$ topology\index{weak-$*$ topology!metrizability on bounded sets}.
Indeed, if $f_n\weakstarto f$ in $E'$ with $\norm{f_n}\le M$, then
$f_n(x_k)\to f(x_k)$ for each $k$, hence $d(f_n,f)\to 0$ by dominated
convergence.  Conversely, if $d(f_n,f)\to 0$, then $f_n(x_k)\to f(x_k)$
for each $k$; by density and uniform boundedness, $f_n(x)\to f(x)$ for all
$x\in E$.

Since $MB_{E'}$ is weak-$*$ compact and metrizable, it is sequentially compact.
\end{proof}

\begin{remark}\label{rem:alaoglu-vs-heine-borel}
\index{Banach--Alaoglu theorem!vs.\ Heine--Borel}
The Banach--Alaoglu theorem can be viewed as an infinite-dimensional analogue of
the Heine--Borel theorem: closed bounded sets in $\R^n$ are compact, and the
unit ball in $E'$ is compact---but only in the weak-$*$ topology, not the
norm topology (unless $E$ is finite-dimensional).
\end{remark}

% ------------------------------------------------------------------
\section{Goldstine's Theorem}\label{sec:goldstine}
% ------------------------------------------------------------------

\begin{theorem}[{Goldstine, 1938}]\label{thm:goldstine}
\index{Goldstine's theorem}
Let $E$ be a normed space and $J:E\to E''$ the canonical injection.  Then
$J(B_E)$ is weak-$*$ dense in $B_{E''}$\,:
\[
  \overline{J(B_E)}^{\sigma(E'',E')} = B_{E''}.
\]
\end{theorem}

\begin{proof}
Let $\xi\in B_{E''}$ and let $W = W(\xi;f_1,\ldots,f_n;\varepsilon)$ be a
basic weak-$*$ neighborhood of $\xi$ in $E''$:
\[
  W = \bigl\{\eta\in E'': |\eta(f_k)-\xi(f_k)|<\varepsilon,\;k=1,\ldots,n\bigr\}.
\]
We need to find $x\in B_E$ with $J(x)\in W$, i.e.,
$|f_k(x)-\xi(f_k)|<\varepsilon$ for $k=1,\ldots,n$.

Consider the linear map $T:E\to\K^n$ defined by $T(x)=(f_1(x),\ldots,f_n(x))$.
Set $a = (\xi(f_1),\ldots,\xi(f_n))\in\K^n$.  We need to show $a\in T(B_E) + B(0,\varepsilon)$
in $\K^n$, or equivalently that $a$ is in the closure of $T(B_E)$ in $\K^n$.

\medskip\noindent\textbf{Claim:} $a\in\overline{T(B_E)}$ (closure in $\K^n$).

Since $T$ is linear and continuous, $T(B_E)$ is a convex subset of $\K^n$.
In finite-dimensional spaces, the closure of a convex set equals its weak
closure, so it suffices to show $a\in\overline{T(B_E)}$.

Suppose not.  Then by the Hahn--Banach separation theorem in $\K^n$ (a
finite-dimensional space), there exists a linear functional $\ell:\K^n\to\R$
and $\gamma\in\R$ with $\ell(a)>\gamma\ge\ell(T(x))$ for all $x\in B_E$.

Writing $\ell(\lambda_1,\ldots,\lambda_n) = \operatorname{Re}\sum_{k=1}^n \alpha_k\lambda_k$
for some $\alpha_k\in\K$, and setting $g = \sum_{k=1}^n \alpha_k f_k\in E'$,
we get
\[
  \operatorname{Re}\,g(x) = \ell(T(x)) \le \gamma < \ell(a)
  = \operatorname{Re}\sum_k \alpha_k\xi(f_k) = \operatorname{Re}\,\xi(g)
\]
for all $x\in B_E$.  Taking the supremum over $x\in B_E$:
\[
  \norm{g} = \sup_{\norm{x}\le 1}\operatorname{Re}\,g(x) \le \gamma
  < \operatorname{Re}\,\xi(g) \le |\xi(g)| \le \norm{\xi}\,\norm{g}\le\norm{g},
\]
a contradiction.  Hence $a\in\overline{T(B_E)}$, and we can find $x\in B_E$
with $T(x)$ arbitrarily close to $a$, which means $J(x)\in W$.
\end{proof}

\begin{corollary}\label{cor:goldstine-density}
\index{Goldstine's theorem!consequences}
$J(E)$ is weak-$*$ dense in $E''$.
\end{corollary}

\begin{proof}
For any $\xi\in E''$ with $\norm{\xi}=r>0$, applying Goldstine's theorem to
$\xi/r\in B_{E''}$ and scaling back gives $\xi\in\overline{J(rB_E)}^{w^*}
\subset\overline{J(E)}^{w^*}$.
\end{proof}

% ------------------------------------------------------------------
\section{Kakutani's Theorem}\label{sec:kakutani}
% ------------------------------------------------------------------

\begin{theorem}[{Kakutani, 1938}]\label{thm:kakutani}
\index{Kakutani's theorem}
\index{reflexive space!Kakutani characterization}
A Banach space $E$ is reflexive if and only if $B_E$ is weakly compact.
\end{theorem}

\begin{proof}
($\Rightarrow$) Assume $E$ is reflexive, so $J:E\to E''$ is a surjective
isometry.  Under $J$, $B_E$ is mapped onto $B_{E''}$.  By the Banach--Alaoglu
theorem (Theorem~\ref{thm:banach-alaoglu}), $B_{E''}$ is weak-$*$ compact in
$\sigma(E'',E')$.  Since $J$ is surjective, the weak topology $\sigma(E,E')$
on $E$ corresponds (via $J$) to the weak-$*$ topology $\sigma(E'',E')$ on $E''$.
More precisely, for a net $(x_\alpha)$ in $E$:
\begin{align*}
  x_\alpha\xrightarrow{\sigma(E,E')} x
  &\;\Longleftrightarrow\; f(x_\alpha)\to f(x)\;\forall f\in E' \\
  &\;\Longleftrightarrow\; J(x_\alpha)(f)\to J(x)(f)\;\forall f\in E' \\
  &\;\Longleftrightarrow\; J(x_\alpha)\xrightarrow{\sigma(E'',E')} J(x).
\end{align*}
So $J:(E,\sigma(E,E'))\to(E'',\sigma(E'',E'))$ is a homeomorphism, and
$J(B_E) = B_{E''}$ is $\sigma(E'',E')$-compact.  Hence $B_E$ is
$\sigma(E,E')$-compact.

($\Leftarrow$) Assume $B_E$ is weakly compact.  Then $J(B_E)$ is compact
in $\sigma(E'',E')$ (since $J$ is a weak-to-weak-$*$ homeomorphism onto its
image, as shown above).  In particular, $J(B_E)$ is $\sigma(E'',E')$-closed.

By Goldstine's theorem, $J(B_E)$ is $\sigma(E'',E')$-dense in $B_{E''}$.
Since $J(B_E)$ is also $\sigma(E'',E')$-closed, we conclude
$J(B_E) = B_{E''}$.  This implies $J(E) = E''$ (since every $\xi\in E''$ with
$\norm{\xi}=r$ satisfies $\xi/r\in B_{E''} = J(B_E)$), so $E$ is reflexive.
\end{proof}

\begin{remark}\label{rem:kakutani-significance}
\index{Kakutani's theorem!significance}
Kakutani's theorem converts the algebraic condition of reflexivity ($J$ surjective)
into a topological condition (weak compactness of $B_E$).  This is extremely
useful because compactness arguments (extracting convergent subsequences,
attaining suprema of continuous functions, etc.)\ become available.
\end{remark}

% ------------------------------------------------------------------
\section{Mazur's Theorem}\label{sec:mazur}
% ------------------------------------------------------------------

\begin{theorem}[{Mazur, 1933}]\label{thm:mazur}
\index{Mazur's theorem}
\index{convex set!weak closure}
Let $E$ be a normed space and $C\subset E$ a convex set.  Then $C$ is
norm-closed if and only if $C$ is weakly closed.  Equivalently,
\[
  \overline{C}^{\,\norm{\cdot}} = \overline{C}^{\,\sigma(E,E')}.
\]
\end{theorem}

\begin{proof}
Since the weak topology is coarser than the norm topology, every weakly closed
set is norm-closed.  So we only need to show: if $C$ is convex and norm-closed,
then $C$ is weakly closed.

Let $x_0\notin C$.  We show there is a weakly open set containing $x_0$ that
is disjoint from $C$.  Since $C$ is norm-closed and convex, by the geometric
Hahn--Banach theorem\index{Hahn--Banach theorem!geometric form} (strict separation
of a point from a closed convex set), there exist $f\in E'$ and $\gamma\in\R$
such that
\[
  \operatorname{Re}\,f(x_0) > \gamma > \sup_{c\in C}\operatorname{Re}\,f(c).
\]
The set $U = \{x\in E: \operatorname{Re}\,f(x)>\gamma\}$ is weakly open (since $f$ is
weakly continuous), contains $x_0$, and satisfies $U\cap C = \emptyset$.
Hence $C^c$ is weakly open, so $C$ is weakly closed.

For the equivalent formulation: the norm closure $\overline{C}$ is convex
(closure of a convex set is convex) and norm-closed, hence weakly closed.
So $\overline{C}^{w}\subset\overline{C}$.  The reverse inclusion
$\overline{C}\subset\overline{C}^{w}$ holds because the weak topology is
coarser.
\end{proof}

\begin{corollary}[Weakly convergent sequences and convex combinations]\label{cor:mazur-convex-comb}
\index{Mazur's theorem!convex combinations}
\index{weak convergence!and convex combinations}
Let $E$ be a Banach space and $x_n\weakto x$.  Then $x$ is in the
norm-closed convex hull of $\{x_n:n\ge 1\}$.  In other words, there exist
convex combinations $y_k = \sum_{n\in F_k}\lambda_n^{(k)}x_n$ (with
$F_k$ finite, $\lambda_n^{(k)}\ge 0$, $\sum\lambda_n^{(k)}=1$)
such that $y_k\to x$ in norm.
\end{corollary}

\begin{proof}
Let $C = \conv\{x_n:n\ge 1\}$.  Since $x_n\weakto x$, the point $x$ belongs
to the weak closure of $\{x_n\}$, hence to the weak closure of $C$.
By Mazur's theorem, $\overline{C}^{w} = \overline{C}^{\norm{\cdot}}$, so
$x\in\overline{C}^{\norm{\cdot}} = \overline{\conv\{x_n\}}^{\norm{\cdot}}$.
\end{proof}

\begin{remark}\label{rem:mazur-significance}
Mazur's theorem is the key bridge between weak and norm topologies for convex
sets.  It says that for convex sets, the distinction between weak and norm
topologies disappears.  This is not true for non-convex sets: for instance,
the unit sphere $S_E = \{x:\norm{x}=1\}$ is norm-closed but weakly dense
in $B_E$ when $\dim E=\infty$ (a consequence of the fact that no weakly open
set is norm-bounded).
\end{remark}

% ------------------------------------------------------------------
\section{Application: Existence of Minimizers}\label{sec:minimizers}
% ------------------------------------------------------------------

\begin{theorem}[Direct method of the calculus of variations]\label{thm:direct-method}
\index{calculus of variations!direct method}
\index{direct method}
\index{minimizer!existence of}
Let $E$ be a reflexive Banach space, $C\subset E$ a nonempty, bounded, closed,
and convex subset, and $\Phi: C\to\R$ a function that is sequentially weakly
lower semicontinuous (i.e., $x_n\weakto x$ implies
$\Phi(x)\le\liminf_n\Phi(x_n)$).  Then $\Phi$ attains its minimum on $C$:
there exists $x_0\in C$ with $\Phi(x_0)=\inf_{x\in C}\Phi(x)$.
\end{theorem}

\begin{proof}
Let $m = \inf_{x\in C}\Phi(x)\in[-\infty,+\infty)$.  Choose a minimizing
sequence $(x_n)\subset C$ with $\Phi(x_n)\to m$.

Since $E$ is reflexive and $C$ is bounded, closed, and convex, $C$ is weakly
compact by Kakutani's theorem (Theorem~\ref{thm:kakutani}) and the fact that
a closed convex subset of a weakly compact set is weakly compact (closed in
the weak topology by Mazur's theorem, hence weakly compact as a weakly closed
subset of the weakly compact set $\overline{B(0,R)}^{w}$ for some $R$ large
enough).

By the Eberlein--\v{S}mulian theorem (Theorem~\ref{thm:eberlein-smulian}),
$(x_n)$ has a subsequence $(x_{n_k})$ with $x_{n_k}\weakto x_0$ for some
$x_0\in C$ (since $C$ is weakly closed, hence weakly sequentially closed).

By weak lower semicontinuity,
\[
  \Phi(x_0) \le \liminf_{k\to\infty}\Phi(x_{n_k}) = m.
\]
Since $x_0\in C$, $\Phi(x_0)\ge m$, hence $\Phi(x_0)=m$.
\end{proof}

\begin{example}\label{ex:minimizer-norm}
\index{minimizer!of norm}
Let $E$ be a reflexive Banach space and $C\subset E$ a nonempty, closed, convex
subset.  For any $y\in E$, the function $\Phi(x) = \norm{x-y}$ is weakly lower
semicontinuous (Theorem~\ref{thm:weak-bounded}(ii)).  If $C$ is bounded, the
direct method gives a minimizer.  Even if $C$ is unbounded, one can restrict
to $C\cap \overline{B}(y,r)$ for $r=\inf_{c\in C}\norm{c-y}+1$, and the
direct method still applies.  Hence every nonempty closed convex subset of a
reflexive Banach space has a \emph{nearest point} to any given $y$.
\end{example}

\begin{example}[Minimization of an energy functional]\label{ex:energy-minimizer}
\index{energy functional}
\index{calculus of variations!energy minimization}
Let $\Omega\subset\R^d$ be a bounded open set and consider the Sobolev space
$H^1_0(\Omega)$\index{Sobolev space}, which is a reflexive Banach space (in fact
a Hilbert space).  For $f\in L^2(\Omega)$, define the energy functional
\[
  \mathcal{E}(u) = \frac{1}{2}\int_\Omega |\nabla u|^2\,\dd x
  - \int_\Omega fu\,\dd x.
\]
Then $\mathcal{E}$ is weakly lower semicontinuous and coercive on $H^1_0(\Omega)$
(by the Poincar\'e inequality\index{Poincar\'e inequality},
$\mathcal{E}(u)\to+\infty$ as $\norm{u}_{H^1_0}\to\infty$), so $\mathcal{E}$
attains its minimum.  The minimizer is the weak solution of $-\Delta u = f$.
\end{example}

% ------------------------------------------------------------------
\section{Additional Results}\label{sec:weak-additional}
% ------------------------------------------------------------------

\begin{proposition}[Characterization of weak-$*$ continuous functionals]%
\label{prop:weak-star-continuous}
\index{weak-$*$ topology!continuous functionals}
A linear functional $\Phi:E'\to\K$ is $\sigma(E',E)$-continuous if and only if
there exists $x\in E$ such that $\Phi(f)=f(x)$ for all $f\in E'$.
\end{proposition}

\begin{proof}
The ``if'' direction is clear: $\Phi = \hat{x}$ is weak-$*$ continuous
by definition.

For the ``only if'' direction: if $\Phi$ is $\sigma(E',E)$-continuous, there
exist $x_1,\ldots,x_n\in E$ and $\varepsilon>0$ such that
$|\Phi(f)|<1$ whenever $|f(x_k)|<\varepsilon$ for $k=1,\ldots,n$.
This implies $\Ker\Phi\supset\bigcap_{k=1}^n\Ker\hat{x}_k$.  By a standard
lemma on linear functionals\index{linear functional!kernel containment},
$\Phi\in\spn\{\hat{x}_1,\ldots,\hat{x}_n\}$, i.e.,
$\Phi = \sum_{k=1}^n\alpha_k\hat{x}_k = \widehat{\sum\alpha_k x_k}$.
Setting $x=\sum\alpha_k x_k$ gives $\Phi(f)=f(x)$.
\end{proof}

\begin{proposition}[Weak-$*$ compactness and reflexivity]\label{prop:wstar-compact-reflexive}
\index{weak-$*$ topology!and reflexivity}
Let $E$ be a Banach space.  Then $B_E$ is weak-$*$ compact in $E''$
(i.e., $J(B_E)$ is $\sigma(E'',E')$-compact) if and only if $E$ is reflexive.
\end{proposition}

\begin{proof}
If $E$ is reflexive, $J(B_E) = B_{E''}$ is weak-$*$ compact by Banach--Alaoglu.
Conversely, if $J(B_E)$ is $\sigma(E'',E')$-compact, then it is
$\sigma(E'',E')$-closed.  By Goldstine's theorem, $J(B_E)$ is dense in
$B_{E''}$, so $J(B_E)=B_{E''}$, giving reflexivity.
\end{proof}

\begin{theorem}[Weak-$*$ sequential compactness in separable duals]%
\label{thm:wstar-sep-sequential}
\index{weak-$*$ convergence!in separable spaces}
If $E$ is a separable Banach space, then every bounded sequence in $E'$ has
a weak-$*$ convergent subsequence.
\end{theorem}

\begin{proof}
This was established in Corollary~\ref{cor:weak-star-subsequence}: when $E$ is
separable, the weak-$*$ topology on bounded subsets of $E'$ is metrizable,
and bounded subsets are relatively weak-$*$ compact by Banach--Alaoglu, hence
sequentially compact.
\end{proof}

\begin{proposition}[Weak-$*$ closed subspaces]\label{prop:wstar-closed-subspace}
\index{weak-$*$ topology!closed subspaces}
A subspace $F\subset E'$ is $\sigma(E',E)$-closed if and only if
$F = G^\perp$ for some subset $G\subset E$, where
$G^\perp = \{f\in E': f(x)=0\;\forall x\in G\}$.
\end{proposition}

\begin{proof}
($\Leftarrow$) $G^\perp = \bigcap_{x\in G}\Ker\hat{x}$, which is an
intersection of weak-$*$ closed sets, hence weak-$*$ closed.

($\Rightarrow$) Let $F$ be a weak-$*$ closed subspace.  Set
$G = {}^\perp F = \{x\in E: f(x)=0\;\forall f\in F\}$.  Then $F\subset G^\perp$
and we claim $F = G^\perp$.  If $g\in G^\perp\setminus F$, then by the Hahn--Banach
theorem for the locally convex space $(E',\sigma(E',E))$, there exists a
$\sigma(E',E)$-continuous functional $\Phi$ with $\Phi\big|_F=0$ and $\Phi(g)\ne 0$.
By Proposition~\ref{prop:weak-star-continuous}, $\Phi = \hat{x}$ for some
$x\in E$.  Then $f(x)=0$ for all $f\in F$, so $x\in{}^\perp F = G$, giving
$g(x)=0$ (since $g\in G^\perp$), i.e.\ $\Phi(g)=0$, a contradiction.
\end{proof}

% ------------------------------------------------------------------
\section{Exercises}\label{sec:exercises-weak-topologies}
% ------------------------------------------------------------------

\begin{exercise}[$\star$]\label{exo:weak-convergence-lp}
\index{weak convergence!in $\ell^p$}
Show that $e_n\weakto 0$ in $\ell^p$ for $1<p<\infty$ but $e_n\not\weakto 0$
in $\ell^1$.
\end{exercise}

\begin{exercise}[$\star$]\label{exo:weak-topology-subbase}
\index{weak topology!subbasis}
Let $E$ be a normed space.
\begin{enumerate}[label=(\alph*)]
  \item Show that the weak topology is the coarsest topology making every
        $f\in E'$ continuous.
  \item Show that a linear map $T:E\to F$ between normed spaces is
        weakly continuous (continuous from $(E,\sigma(E,E'))$ to $(F,\sigma(F,F'))$)
        if and only if it is norm-continuous.
\end{enumerate}
\end{exercise}

\begin{exercise}[$\star\star$]\label{exo:weak-closed-not-seq-closed}
\index{weak topology!closure vs.\ sequential closure}
(This exercise shows that in non-separable spaces, weak sequential closure
may differ from weak closure.)
Let $E = \ell^\infty$ and consider the set $A = \{e_n : n\ge 1\}$.
\begin{enumerate}[label=(\alph*)]
  \item Show that $0$ is in the weak closure of $A$.
  \item Show that no subsequence of $(e_n)$ converges weakly to $0$ in
        $\ell^\infty$.
\end{enumerate}
\textit{Hint for (b):} Use the fact that $(\ell^\infty)'$ is much larger
than $\ell^1$.
\end{exercise}

\begin{exercise}[$\star$]\label{exo:weak-star-basic}
\index{weak-$*$ convergence!examples}
\begin{enumerate}[label=(\alph*)]
  \item In $\ell^\infty = (\ell^1)'$, show that $e_n\weakstarto 0$.
  \item In $L^\infty([0,1]) = (L^1([0,1]))'$, show that $\sin(2\pi n\,\cdot\,)
        \weakstarto 0$.
\end{enumerate}
\end{exercise}

\begin{exercise}[$\star\star$]\label{exo:alaoglu-necessity}
\index{Banach--Alaoglu theorem!necessity of weak-$*$}
Show that the closed unit ball of an infinite-dimensional Banach space is
\textbf{not} compact in the norm topology.
\textit{Hint:} Construct an infinite sequence with no convergent subsequence
(e.g., using Riesz's lemma\index{Riesz's lemma}).
\end{exercise}

\begin{exercise}[$\star\star$]\label{exo:mazur-application}
\index{Mazur's theorem!application}
Let $(x_n)$ be a sequence in a Banach space $E$ with $x_n\weakto x$.
\begin{enumerate}[label=(\alph*)]
  \item Use Mazur's theorem to show that there exist convex combinations
        $y_N = \sum_{n=N}^{M_N}\lambda_n x_n$ (with $\lambda_n\ge 0$,
        $\sum\lambda_n=1$) such that $\norm{y_N-x}\to 0$.
  \item Give an explicit construction when $E = L^2([0,1])$ and
        $x_n(t) = \sqrt{2}\sin(2\pi nt)$.
\end{enumerate}
\end{exercise}

\begin{exercise}[$\star\star$]\label{exo:reflexive-wstar-weak}
\index{weak-$*$ topology!equals weak on reflexive}
Let $E$ be a Banach space.  Show that the following are equivalent:
\begin{enumerate}[label=(\roman*)]
  \item $E$ is reflexive.
  \item The weak and weak-$*$ topologies on $E'$ coincide:
        $\sigma(E',E) = \sigma(E',E'')$.
  \item $B_{E'}$ is weakly compact (in the weak topology $\sigma(E',(E')')$).
\end{enumerate}
\end{exercise}

\begin{exercise}[$\star\star\star$]\label{exo:schur}
\index{Schur's theorem}
\index{$\ell^1$ space!Schur property}
(Schur's theorem) Show that in $\ell^1$, weak convergence implies norm
convergence: if $x_n\weakto x$ in $\ell^1$, then $\norm{x_n-x}_1\to 0$.
\textit{Hint:} Argue by contradiction.  If $\norm{x_n-x}_1\not\to 0$, pass
to a subsequence and use a gliding hump argument to construct
$f\in\ell^\infty = (\ell^1)'$ with $f(x_n)\not\to f(x)$.
\end{exercise}

\begin{exercise}[$\star\star$]\label{exo:kakutani-application}
\index{Kakutani's theorem!application}
Use Kakutani's theorem to give another proof that $c_0$ is not reflexive.
\textit{Hint:} Show that the sequence $(e_1+e_2+\cdots+e_n)/n$ is in $B_{c_0}$
but has no weakly convergent subsequence, using Schur's theorem and the
identification $c_0''\cong\ell^\infty$.
\end{exercise}

\begin{exercise}[$\star\star$]\label{exo:goldstine-application}
\index{Goldstine's theorem!application}
Let $E$ be a Banach space and $f\in E'$.  Use Goldstine's theorem to show
that $\norm{f}_{E'} = \sup\{|f(x)|: x\in B_E\}$ can also be written as
$\norm{f}_{E'} = \sup\{|\xi(f)|: \xi\in B_{E''}\}$.
\end{exercise}

\begin{exercise}[$\star\star$]\label{exo:weak-compact-bounded}
\index{weak compactness!implies boundedness}
Show that every weakly compact subset of a normed space is bounded.
\textit{Hint:} Use the uniform boundedness principle.
\end{exercise}

\begin{exercise}[$\star\star\star$]\label{exo:dunford-pettis}
\index{Dunford--Pettis theorem}
\index{weak compactness!in $L^1$}
(Dunford--Pettis criterion) Let $(\Omega,\mu)$ be a finite measure space and
$K\subset L^1(\mu)$.  State (without proof) the Dunford--Pettis criterion for
$K$ to be relatively weakly compact.  Verify it for $K = \{f_n\}$ where
$f_n = n\,\mathbf{1}_{[0,1/n^2]}$ in $L^1([0,1])$.
\end{exercise}

\begin{exercise}[$\star$]\label{exo:weak-vs-norm-closed}
\index{weak topology!closed sets}
Give an example of a subset of a Banach space that is norm-closed but not
weakly closed.
\textit{Hint:} Consider the unit sphere in an infinite-dimensional space.
\end{exercise}

\begin{exercise}[$\star\star$]\label{exo:weak-lower-semicontinuous-convex}
\index{convex function!weak lower semicontinuity}
Let $E$ be a Banach space and $\Phi:E\to\R\cup\{+\infty\}$ a convex,
lower semicontinuous (for the norm topology) function.  Show that $\Phi$ is
weakly lower semicontinuous.
\textit{Hint:} The sublevel sets $\{x:\Phi(x)\le\alpha\}$ are convex and
norm-closed; apply Mazur's theorem.
\end{exercise}

\begin{exercise}[$\star\star\star$]\label{exo:minimizer-Lp}
\index{minimizer!in $L^p$}
\index{calculus of variations!in $L^p$}
Let $1<p<\infty$, $\Omega\subset\R^d$ bounded and open, $g\in L^{p'}(\Omega)$.
Consider the problem of minimizing
\[
  \Phi(u) = \frac{1}{p}\int_\Omega |u|^p\,\dd x - \int_\Omega gu\,\dd x
\]
over $u\in L^p(\Omega)$.
\begin{enumerate}[label=(\alph*)]
  \item Show that $\Phi$ is weakly lower semicontinuous and coercive.
  \item Prove that a minimizer exists and is unique.
  \item Find the minimizer explicitly.
\end{enumerate}
\end{exercise}

\begin{exercise}[$\star\star$]\label{exo:weak-compact-operators}
\index{compact operator!and weak convergence}
Let $E,F$ be Banach spaces and $T\in\mathcal{L}(E,F)$.  Show that $T$ is
compact if and only if $x_n\weakto 0$ in $E$ implies $Tx_n\to 0$ in norm
in $F$.
\textit{Hint:} Use the Eberlein--\v{S}mulian theorem and the fact that
$T$ maps weakly convergent sequences to weakly convergent sequences.
\end{exercise}

\begin{exercise}[$\star$]\label{exo:bidual-chain}
\index{bidual!iterated}
Let $E$ be a Banach space.
\begin{enumerate}[label=(\alph*)]
  \item Show that $E$ is isometrically embedded in $E''$, which is isometrically
        embedded in $E^{(4)} = (E'')''$, etc.
  \item If $E$ is reflexive, show that all iterated duals $E^{(n)}$ are
        reflexive and canonically isomorphic to $E$ (for $n$ even) or $E'$
        (for $n$ odd).
  \item If $E$ is not reflexive, show that $\dim E^{(2n)} < \dim E^{(2n+2)}$
        for all $n\ge 0$ (in the sense that the canonical injection is not
        surjective).
\end{enumerate}
\end{exercise}
% === START OF CHUNK 5: Chapters 9--12 + End Matter ===
% ============================================================
%  Functional Analysis — A Graduate Course
%  Chapters 9--12, Appendix, Index, \end{document}
% ============================================================

%%%%%%%%%%%%%%%%%%%%%%%%%%%%%%%%%%%%%%%%%%%%%%%%%%%%%%%%%%%%
%%  CHAPTER 9: Spectral Theory of Compact Operators
%%%%%%%%%%%%%%%%%%%%%%%%%%%%%%%%%%%%%%%%%%%%%%%%%%%%%%%%%%%%
\chapter{Spectral Theory of Compact Operators}\label{ch:compact-spectral}
\minitoc

\section{Compact Operators: Review and Examples}\label{sec:compact-review}

\index{compact operator|(}
\index{operator!compact|(}

We begin with a systematic study of the spectral properties of compact operators,
building on the foundational Banach and Hilbert space theory developed in earlier
chapters. The class of compact operators is arguably the most tractable
infinite-dimensional analogue of matrices, and their spectral theory mirrors the
finite-dimensional situation remarkably closely.

\begin{definition}[Compact operator]\label{def:compact-op}
\index{compact operator!definition}
Let $E$ and $F$ be Banach spaces. A bounded linear operator $T \colon E \to F$ is
called \emph{compact} (or \emph{completely continuous}) if for every bounded
sequence $(x_n)_{n \ge 1}$ in $E$, the sequence $(Tx_n)_{n \ge 1}$ has a
convergent subsequence in $F$.

Equivalently, $T$ is compact if and only if the image $T(B_E)$ of the closed
unit ball $B_E = \{x \in E : \norm{x} \le 1\}$ is relatively compact in $F$
(i.e., $\overline{T(B_E)}$ is compact).
\end{definition}

We denote by $\mathcal{K}(E, F)$ the set of all compact operators from $E$ to $F$,
\index{K(E,F)@$\mathcal{K}(E,F)$}
and write $\mathcal{K}(E) = \mathcal{K}(E, E)$.

\begin{example}[Finite-rank operators]\label{ex:finite-rank-compact}
\index{finite-rank operator}
\index{operator!finite-rank}
An operator $T \colon E \to F$ is of \emph{finite rank} if $\dim(\im T) < \infty$.
Every finite-rank bounded operator is compact: if $(x_n)$ is bounded in $E$,
then $(Tx_n)$ lies in the finite-dimensional subspace $\im T$, and by the
Bolzano--Weierstrass theorem\index{Bolzano--Weierstrass theorem} it admits a
convergent subsequence.
\end{example}

\begin{example}[Integral operators]\label{ex:integral-compact}
\index{integral operator!compact}
\index{operator!integral}
Let $\Omega \subset \R^d$ be a bounded measurable set and let
$K \in L^2(\Omega \times \Omega)$. The operator $T_K \colon L^2(\Omega) \to L^2(\Omega)$
defined by
\[
  (T_K f)(x) = \int_\Omega K(x, y)\, f(y)\, \dd y
\]
is compact. Indeed, $T_K$ is a Hilbert--Schmidt operator with
$\norm{T_K}_{\mathrm{HS}} = \norm{K}_{L^2(\Omega \times \Omega)}$, and every
Hilbert--Schmidt operator is compact (see \S\ref{sec:hilbert-schmidt}).
\end{example}

\begin{example}[Diagonal operators]\label{ex:diagonal-compact}
\index{operator!diagonal}
Let $H = \ell^2(\N)$ and let $(\lambda_n)_{n \ge 1}$ be a bounded sequence of
scalars. Define $T \colon H \to H$ by $T(e_n) = \lambda_n e_n$, where $(e_n)$
is the standard orthonormal basis. Then $T$ is compact if and only if
$\lambda_n \to 0$ as $n \to \infty$.
\end{example}

\begin{example}[The inclusion $\ell^1 \hookrightarrow \ell^2$ is not compact]
\label{ex:non-compact}
The natural inclusion $\iota \colon \ell^1(\N) \hookrightarrow \ell^2(\N)$
is bounded (with $\norm{\iota} = 1$) but not compact. The standard basis
vectors $e_n$ satisfy $\norm{e_n}_{\ell^1} = 1$ and
$\norm{e_m - e_n}_{\ell^2} = \sqrt{2}$ for $m \ne n$, so $(e_n)$ has no
convergent subsequence in $\ell^2$.
\end{example}

\section{The Space of Compact Operators}\label{sec:compact-space}

\index{compact operator!space of}

\begin{proposition}\label{prop:compact-ideal}
Let $E$, $F$, and $G$ be Banach spaces.
\begin{enumerate}[label=(\roman*)]
  \item $\mathcal{K}(E, F)$ is a closed linear subspace of $\mathcal{L}(E, F)$.
  \item If $T \in \mathcal{K}(E, F)$ and $S \in \mathcal{L}(F, G)$, then
        $ST \in \mathcal{K}(E, G)$.
  \item If $T \in \mathcal{L}(E, F)$ and $S \in \mathcal{K}(F, G)$, then
        $ST \in \mathcal{K}(E, G)$.
\end{enumerate}
In particular, $\mathcal{K}(E)$ is a closed two-sided ideal\index{ideal!of compact operators}
in the Banach algebra $\mathcal{L}(E)$.
\end{proposition}

\begin{proof}
Linearity of $\mathcal{K}(E, F)$ and properties (ii) and (iii) are straightforward
exercises. We prove that $\mathcal{K}(E, F)$ is closed.

Let $(T_n)_{n \ge 1} \subset \mathcal{K}(E, F)$ with $T_n \to T$ in
$\mathcal{L}(E, F)$. We show $T$ is compact. Let $(x_k)_{k \ge 1}$
be a bounded sequence in $E$ with $\norm{x_k} \le M$ for all $k$.

We use a diagonal argument.\index{diagonal argument} Since $T_1$ is compact,
extract a subsequence $(x_k^{(1)})$ of $(x_k)$ such that $(T_1 x_k^{(1)})$
converges. From $(x_k^{(1)})$, extract a further subsequence $(x_k^{(2)})$ such
that $(T_2 x_k^{(2)})$ converges, and so on. The diagonal sequence
$y_k := x_k^{(k)}$ satisfies: $(T_n y_k)_{k \ge 1}$ converges for every
$n \ge 1$.

We claim $(Ty_k)$ is Cauchy. Given $\varepsilon > 0$, choose $n$ so that
$\norm{T - T_n} < \varepsilon / (3M)$. Since $(T_n y_k)$ converges, there
exists $K$ such that for all $j, k \ge K$,
$\norm{T_n y_j - T_n y_k} < \varepsilon / 3$. Then for $j, k \ge K$:
\begin{align*}
  \norm{Ty_j - Ty_k}
  &\le \norm{Ty_j - T_n y_j} + \norm{T_n y_j - T_n y_k} + \norm{T_n y_k - Ty_k} \\
  &\le \norm{T - T_n}\norm{y_j} + \frac{\varepsilon}{3} + \norm{T_n - T}\norm{y_k} \\
  &< \frac{\varepsilon}{3M} \cdot M + \frac{\varepsilon}{3} + \frac{\varepsilon}{3M} \cdot M
  = \varepsilon.
\end{align*}
Since $F$ is complete, $(Ty_k)$ converges, so $T$ is compact.
\end{proof}

\begin{remark}\label{rem:approx-property}
\index{approximation property}
A Banach space $E$ has the \emph{approximation property} if every compact operator
$T \in \mathcal{K}(E)$ is the operator-norm limit of finite-rank operators. All
Hilbert spaces and all $L^p$ spaces ($1 \le p \le \infty$) have the approximation
property. Enflo (1973)\index{Enflo counterexample} constructed a separable Banach
space failing this property.
\end{remark}

\section{The Fredholm Alternative}\label{sec:fredholm-alt}

\index{Fredholm alternative|(}

The Fredholm alternative is the cornerstone result connecting the solvability of
equations involving compact operators to a finite-dimensional eigenvalue problem.

\begin{theorem}[Fredholm alternative]\label{thm:fredholm-alt}
\index{Fredholm alternative!theorem}
Let $E$ be a Banach space and $T \in \mathcal{K}(E)$. Exactly one of the
following holds:
\begin{enumerate}[label=(\alph*)]
  \item The operator $\Id - T$ is bijective, $(\Id - T)^{-1} \in \mathcal{L}(E)$,
        and for every $y \in E$ the equation $x - Tx = y$ has a unique solution.
  \item $\Ker(\Id - T) \ne \{0\}$, i.e., $\lambda = 1$ is an eigenvalue of $T$.
        In this case, $\dim \Ker(\Id - T) < \infty$,
        $\im(\Id - T)$ is closed with
        $\codim \im(\Id - T) = \dim \Ker(\Id - T^*) = \dim \Ker(\Id - T) \ge 1$.
\end{enumerate}
\end{theorem}

We prepare the proof with several lemmas.

\begin{lemma}\label{lem:kernel-finite-dim}
\index{kernel!of Id minus compact}
If $T \in \mathcal{K}(E)$, then $\Ker(\Id - T)$ is finite-dimensional.
\end{lemma}

\begin{proof}
Let $N = \Ker(\Id - T)$. On $N$ we have $T|_N = \Id_N$, so $T$ maps
the unit ball $B_N = B_E \cap N$ onto itself. Since $T(B_N)$ is relatively compact
and $T(B_N) = B_N$, the closed unit ball of $N$ is compact. By Riesz's
theorem\index{Riesz theorem!finite dimension}, $\dim N < \infty$.
\end{proof}

\begin{lemma}\label{lem:range-closed}
\index{range!closed}
If $T \in \mathcal{K}(E)$, then $\im(\Id - T)$ is closed in $E$.
\end{lemma}

\begin{proof}
Let $y_n = (\Id - T)x_n \to y$. We must find $x$ with $(\Id - T)x = y$.
Let $N = \Ker(\Id - T)$ and let $d_n = d(x_n, N)$. Since $N$ is finite-dimensional
(hence closed), choose $z_n \in N$ with $\norm{x_n - z_n} \le d_n + 1/n$ and set
$\tilde{x}_n = x_n - z_n$. Then $(\Id - T)\tilde{x}_n = y_n$.

\emph{Claim:} $(\tilde{x}_n)$ is bounded. Suppose not; then (passing to a
subsequence) $\norm{\tilde{x}_n} \to \infty$. Set
$u_n = \tilde{x}_n / \norm{\tilde{x}_n}$. Then
$(\Id - T)u_n = y_n / \norm{\tilde{x}_n} \to 0$. Since $T$ is compact and
$(u_n)$ is bounded, a subsequence satisfies $Tu_{n_k} \to w$. Hence
$u_{n_k} = (\Id - T)u_{n_k} + Tu_{n_k} \to 0 + w = w$, so $\norm{w} = 1$ and
$(\Id - T)w = 0$, i.e., $w \in N$. But
\[
  d(u_n, N) = \frac{d(\tilde{x}_n, N)}{\norm{\tilde{x}_n}}
  = \frac{d_n}{\norm{\tilde{x}_n}}
  \ge \frac{d_n}{d_n + 1/n} \to 1
\]
(since $d_n \ge \norm{\tilde{x}_n} - \norm{0} \to \infty$ when
$\norm{\tilde{x}_n} \to \infty$ and $\tilde{x}_n$ has distance $d_n$ to $N$
with $\norm{\tilde{x}_n} \le d_n + 1/n$). Yet $u_{n_k} \to w \in N$ gives
$d(u_{n_k}, N) \to 0$, a contradiction. So $(\tilde{x}_n)$ is bounded.

Since $(\tilde{x}_n)$ is bounded and $T$ is compact, passing to a subsequence,
$T\tilde{x}_{n_k} \to v$. Then
$\tilde{x}_{n_k} = y_{n_k} + T\tilde{x}_{n_k} \to y + v =: x$, and
$(\Id - T)x = y$.
\end{proof}

\begin{lemma}[Riesz's lemma on ascending chains]\label{lem:ascending-chains}
\index{Riesz lemma!ascending chains}
Let $T \in \mathcal{K}(E)$ and define $N_k = \Ker(\Id - T)^k$ for $k \ge 0$.
Then $N_0 \subset N_1 \subset N_2 \subset \cdots$ and there exists $p \ge 0$
such that $N_p = N_{p+1} = N_{p+2} = \cdots$
\end{lemma}

\begin{proof}
The inclusions are clear. Suppose for contradiction that $N_k \subsetneq N_{k+1}$
for all $k$. By Riesz's lemma\index{Riesz lemma}, for each $k \ge 1$ there
exists $x_k \in N_k$ with $\norm{x_k} = 1$ and $d(x_k, N_{k-1}) \ge 1/2$.
For $j < k$, one checks that
\[
  Tx_k - Tx_j = x_k - \bigl[(\Id - T)x_k + x_j - (\Id - T)x_j\bigr],
\]
and the bracketed term lies in $N_{k-1}$. Hence
$\norm{Tx_k - Tx_j} \ge d(x_k, N_{k-1}) \ge 1/2$, so $(Tx_k)$ has no convergent
subsequence, contradicting compactness.
\end{proof}

\begin{proof}[Proof of Theorem~\ref{thm:fredholm-alt}]
\index{Fredholm alternative!proof}
Finite-dimensionality of $\Ker(\Id - T)$ is Lemma~\ref{lem:kernel-finite-dim},
and closedness of $\im(\Id - T)$ is Lemma~\ref{lem:range-closed}.

\textbf{Claim:} If $\Id - T$ is injective, then it is surjective.

Set $R_n = \im(\Id - T)^n$. Then $E = R_0 \supset R_1 \supset R_2 \supset \cdots$
and each $R_n$ is closed. If $\Id - T$ is not surjective, then
$R_0 \supsetneq R_1$, and injectivity forces $R_n \supsetneq R_{n+1}$ for all $n$.
By Riesz's lemma, pick $x_n \in R_n$ with $\norm{x_n} = 1$ and
$d(x_n, R_{n+1}) \ge 1/2$. Then for $m > n$,
$Tx_n - Tx_m = x_n - [(\Id - T)x_n + x_m - (\Id - T)x_m]$ where the bracket
lies in $R_{n+1}$, so $\norm{Tx_n - Tx_m} \ge 1/2$. This contradicts compactness.

\textbf{Index zero:} When $E$ is a Hilbert space $H$, note that $T^*$ is also
compact. Since $\im(\Id - T)$ is closed,
$\im(\Id - T)^\perp = \Ker(\Id - T^*)$. Hence
$\codim \im(\Id - T) = \dim \Ker(\Id - T^*)$. The Fredholm index
$\operatorname{ind}(\Id - T) = \dim \Ker(\Id - T) - \dim \Ker(\Id - T^*)$
equals zero; this follows by a homotopy argument: the map
$t \mapsto \operatorname{ind}(\Id - tT)$ is continuous and integer-valued (hence
constant) for $t \in [0,1]$, and at $t = 0$ the index is $0$.

For general Banach spaces, one uses the dual $T^* \in \mathcal{K}(E^*)$ and the
Hahn--Banach theorem\index{Hahn--Banach theorem} to establish
$\im(\Id - T) = {}^\perp\!\Ker(\Id - T^*)$ and the analogous index computation.
\end{proof}

\index{Fredholm alternative|)}

\begin{corollary}\label{cor:fredholm-equations}
\index{Fredholm equation}
Let $T \in \mathcal{K}(E)$ and $\lambda \ne 0$. Then either $\lambda\Id - T$ is
bijective (and $(\lambda\Id - T)^{-1} \in \mathcal{L}(E)$), or $\lambda$ is an
eigenvalue of $T$ with finite-dimensional eigenspace. The equation
$\lambda x - Tx = y$ is solvable for all $y$ if and only if $\lambda x = Tx$
implies $x = 0$.
\end{corollary}

\begin{proof}
Apply Theorem~\ref{thm:fredholm-alt} to the compact operator $T/\lambda$.
\end{proof}

\section{Riesz--Schauder Theory}\label{sec:riesz-schauder}

\index{Riesz--Schauder theory|(}
\index{spectrum!of compact operator|(}

\begin{theorem}[Riesz--Schauder]\label{thm:riesz-schauder}
\index{Riesz--Schauder theorem}
Let $E$ be an infinite-dimensional Banach space and $T \in \mathcal{K}(E)$. Then:
\begin{enumerate}[label=(\roman*)]
  \item $0 \in \sigma(T)$.
  \item $\sigma(T) \setminus \{0\} = \sigma_p(T) \setminus \{0\}$: every nonzero
        spectral value is an eigenvalue.\index{eigenvalue!of compact operator}
  \item Each nonzero eigenvalue has finite multiplicity.\index{eigenvalue!finite multiplicity}
  \item $\sigma(T)$ is at most countable, with $0$ as the only possible
        accumulation point.\index{spectrum!countable}
  \item For every $\varepsilon > 0$, only finitely many eigenvalues satisfy
        $\abs{\lambda} \ge \varepsilon$.
\end{enumerate}
\end{theorem}

\begin{proof}
\textbf{(i)} If $T$ were invertible, $\Id = T^{-1}T$ would be compact (as a
composition of a bounded with a compact operator), contradicting
infinite-dimensionality of $E$.

\textbf{(ii)} If $\lambda \ne 0$ and $\lambda \in \sigma(T)$, then $\lambda\Id - T$
is not bijective. By Corollary~\ref{cor:fredholm-equations}, it is not injective,
so $\lambda \in \sigma_p(T)$.

\textbf{(iii)} Follows from Lemma~\ref{lem:kernel-finite-dim} applied to $T/\lambda$.

\textbf{(iv)--(v)} Suppose there are infinitely many distinct eigenvalues
$(\lambda_n)_{n \ge 1}$ with $\abs{\lambda_n} \ge \varepsilon > 0$.
Let $x_n$ be a corresponding unit eigenvector and
$E_n = \spn\{x_1, \ldots, x_n\}$. Since eigenvalues are distinct, the
eigenvectors are linearly independent and $E_{n-1} \subsetneq E_n$.
By Riesz's lemma, there exist $y_n \in E_n$ with $\norm{y_n} = 1$ and
$d(y_n, E_{n-1}) \ge 1/2$.

Now $Ty_n = \lambda_n y_n + w_n$ where $w_n \in E_{n-1}$ (since $T$ maps $E_n$
into $E_n$ and the ``leading term'' is $\lambda_n y_n$). For $m < n$:
\[
  Ty_n - Ty_m = \lambda_n y_n + (\text{element of } E_{n-1}),
\]
hence $\norm{Ty_n - Ty_m} \ge \abs{\lambda_n}\, d(y_n, E_{n-1})
\ge \varepsilon / 2 > 0$. Thus $(Ty_n)$ has no convergent subsequence,
contradicting compactness.
\end{proof}

\index{Riesz--Schauder theory|)}
\index{spectrum!of compact operator|)}

\section{Spectral Theorem for Compact Self-Adjoint Operators}
\label{sec:compact-sa-spectral}

\index{spectral theorem!compact self-adjoint|(}

We now specialize to Hilbert spaces and self-adjoint compact operators.

\begin{theorem}[Spectral theorem for compact self-adjoint operators]
\label{thm:compact-sa-spectral}
\index{spectral theorem!compact self-adjoint!statement}
Let $H$ be a separable Hilbert space and $T \in \mathcal{K}(H)$ self-adjoint.
Then there exists an orthonormal basis $(e_n)_{n \ge 1}$ of $H$ consisting of
eigenvectors of $T$:
\[
  Te_n = \lambda_n e_n, \qquad \lambda_n \in \R, \qquad \lambda_n \to 0.
\]
For every $x \in H$:
\[
  Tx = \sum_{n=1}^\infty \lambda_n \inner{x}{e_n}\, e_n.
\]
The eigenvalues are real, and eigenspaces for distinct eigenvalues are orthogonal.
\end{theorem}

\begin{lemma}\label{lem:sa-norm-eigenvalue}
\index{norm!attained on spectrum}
If $T \in \mathcal{L}(H)$ is self-adjoint, then
$\norm{T} = \sup_{\norm{x}=1} \abs{\inner{Tx}{x}}$.
If additionally $T$ is compact and $T \ne 0$, then either $\norm{T}$ or $-\norm{T}$
is an eigenvalue of $T$.
\end{lemma}

\begin{proof}
Set $m = \sup_{\norm{x}=1} \abs{\inner{Tx}{x}}$. Clearly $m \le \norm{T}$.
For the reverse, the polarization identity\index{polarization identity} gives
\[
  4\operatorname{Re}\inner{Tx}{y} = \inner{T(x+y)}{x+y} - \inner{T(x-y)}{x-y}.
\]
Taking $\norm{x} = \norm{y} = 1$ and using $\abs{\inner{Tu}{u}} \le m\norm{u}^2$
yields $\abs{\operatorname{Re}\inner{Tx}{y}} \le m$ after optimizing over the
phase of $y$, hence $\norm{T} \le m$.

Now let $T$ be compact with $T \ne 0$. Choose $(x_n)$ with $\norm{x_n} = 1$ and
$\inner{Tx_n}{x_n} \to \mu$ where $\abs{\mu} = \norm{T}$. By compactness, pass
to a subsequence with $Tx_n \to y$. Compute:
\[
  \norm{Tx_n - \mu x_n}^2 = \norm{Tx_n}^2 - 2\mu\inner{Tx_n}{x_n} + \mu^2
  \to \norm{y}^2 - 2\mu^2 + \mu^2.
\]
Since $\norm{y} = \lim\norm{Tx_n} \le \norm{T} = \abs{\mu}$ and
$\norm{Tx_n}^2 = \inner{Tx_n}{Tx_n} = \inner{T^2 x_n}{x_n}$ with
$\abs{\inner{T^2 x_n}{x_n}} \le \norm{T}^2$ (by applying $m \le \norm{T}$ to
$T^2$, but more carefully:
$\norm{Tx_n}^2 \le \norm{T}\norm{Tx_n}\norm{x_n}$ so $\norm{Tx_n} \le \norm{T}$),
we get $\norm{Tx_n}^2 \to \norm{y}^2$ and we need $\norm{y} = \abs{\mu}$.
In fact, $\norm{Tx_n - \mu x_n}^2 = \norm{Tx_n}^2 - 2\mu\inner{Tx_n}{x_n} + \mu^2
\le \mu^2 - 2\mu\inner{Tx_n}{x_n} + \mu^2 \to 2\mu^2 - 2\mu^2 = 0$.
Hence $Tx_n - \mu x_n \to 0$, so $\mu x_n \to y$, giving
$x_n \to e := y/\mu$ with $\norm{e} = 1$ and $Te = \mu e$.
\end{proof}

\begin{proof}[Proof of Theorem~\ref{thm:compact-sa-spectral}]
\index{spectral theorem!compact self-adjoint!proof}
We construct the eigenvectors iteratively.

\textbf{Step 1.} By Lemma~\ref{lem:sa-norm-eigenvalue}, if $T \ne 0$ then either
$\norm{T}$ or $-\norm{T}$ is an eigenvalue $\lambda_1$ with eigenvector $e_1$
($\norm{e_1} = 1$, $\abs{\lambda_1} = \norm{T}$).

\textbf{Step 2.} Set $H_1 = \{e_1\}^\perp$. Since $T = T^*$, $H_1$ is
$T$-invariant: for $x \perp e_1$,
$\inner{Tx}{e_1} = \inner{x}{Te_1} = \lambda_1\inner{x}{e_1} = 0$.
The restriction $T_1 = T|_{H_1}$ is compact and self-adjoint on $H_1$ with
$\norm{T_1} \le \norm{T}$. Apply the lemma to get $\lambda_2$, $e_2$ with
$\abs{\lambda_2} \le \abs{\lambda_1}$.

\textbf{Step 3.} Iterate: $H_n = \{e_1, \ldots, e_n\}^\perp$, and
$\abs{\lambda_{n+1}} = \norm{T|_{H_n}}$.

\textbf{Step 4.} $\lambda_n \to 0$: if $\abs{\lambda_n} \ge \delta > 0$ for all $n$,
then $(e_n)$ is orthonormal with
$\norm{Te_m - Te_n}^2 = \lambda_m^2 + \lambda_n^2 \ge 2\delta^2$ for $m \ne n$,
contradicting compactness.

\textbf{Step 5 (Completeness).} Let $H_\infty = \bigcap_n H_n$. Then $T|_{H_\infty}$
is compact self-adjoint with $\norm{T|_{H_\infty}} = \lim_n \abs{\lambda_{n+1}} = 0$,
so $T|_{H_\infty} = 0$. Every vector in $H_\infty$ is an eigenvector for
eigenvalue $0$. Choose an orthonormal basis of $H_\infty$ and adjoin it to
$(e_n)$ to obtain a complete orthonormal basis of $H$, each element being an
eigenvector of $T$.
\end{proof}

\index{spectral theorem!compact self-adjoint|)}

\section{The Hilbert--Schmidt Theorem}\label{sec:hilbert-schmidt}

\index{Hilbert--Schmidt operator|(}
\index{Hilbert--Schmidt theorem|(}

\begin{definition}[Hilbert--Schmidt operator]\label{def:hilbert-schmidt}
\index{Hilbert--Schmidt operator!definition}
Let $H$ be a separable Hilbert space with orthonormal basis $(e_n)$. An operator
$T \in \mathcal{L}(H)$ is \emph{Hilbert--Schmidt} if
\[
  \norm{T}_{\mathrm{HS}}^2 := \sum_{n=1}^\infty \norm{Te_n}^2 < \infty.
\]
This quantity is independent of the choice of basis.
\end{definition}

\begin{theorem}[Hilbert--Schmidt]\label{thm:hilbert-schmidt}
\index{Hilbert--Schmidt theorem!statement}
Every Hilbert--Schmidt operator is compact. The space $\mathcal{S}_2(H)$ of
Hilbert--Schmidt operators is a two-sided $*$-ideal in $\mathcal{L}(H)$, and
$(\mathcal{S}_2(H), \norm{\cdot}_{\mathrm{HS}})$ is a Hilbert space with
inner product
$\inner{S}{T}_{\mathrm{HS}} = \sum_n \inner{Se_n}{Te_n}$.
\end{theorem}

\begin{proof}
\index{Hilbert--Schmidt theorem!proof}
Define $T_N x = \sum_{n=1}^N \inner{x}{e_n}\, Te_n$, a finite-rank operator. Then
\[
  \norm{(T - T_N)x}^2
  = \norm[\bigg]{\sum_{n > N} \inner{x}{e_n}\, Te_n}^2
  \le \Bigl(\sum_{n > N} \norm{Te_n}^2\Bigr) \norm{x}^2
\]
by Cauchy--Schwarz. Hence $\norm{T - T_N} \le (\sum_{n > N} \norm{Te_n}^2)^{1/2} \to 0$.
Since $\mathcal{K}(H)$ is closed, $T$ is compact.

The ideal and Hilbert space properties are verified by direct computation using
Parseval's identity\index{Parseval's identity}.
\end{proof}

\begin{example}[Hilbert--Schmidt integral operator]\label{ex:hs-integral}
\index{Hilbert--Schmidt operator!integral}
If $K \in L^2(\Omega \times \Omega)$, the integral operator $T_K$ of
Example~\ref{ex:integral-compact} is Hilbert--Schmidt with
$\norm{T_K}_{\mathrm{HS}} = \norm{K}_{L^2}$. Indeed, for any ONB $(e_n)$ of
$L^2(\Omega)$:
\[
  \sum_n \norm{T_K e_n}^2
  = \int_\Omega \sum_n \abs{\inner{K(x,\cdot)}{e_n}}^2\, \dd x
  = \int_\Omega \norm{K(x,\cdot)}_{L^2}^2\, \dd x
  = \norm{K}_{L^2(\Omega \times \Omega)}^2.
\]
\end{example}

\index{Hilbert--Schmidt operator|)}
\index{Hilbert--Schmidt theorem|)}

\section{Application: Fredholm Integral Equations}\label{sec:fredholm-integral}

\index{Fredholm integral equation|(}

\begin{theorem}[Solvability of Fredholm integral equations]
\label{thm:fredholm-integral-solve}
\index{Fredholm integral equation!solvability}
Let $K \in L^2(\Omega \times \Omega)$ be Hermitian ($K(x,y) = \overline{K(y,x)}$
a.e.) and $(\lambda_n, e_n)$ the eigenpairs from Theorem~\ref{thm:compact-sa-spectral}.
Given $g \in L^2(\Omega)$ and $\mu \ne 0$, consider the equation
$f - \mu T_K f = g$.
\begin{enumerate}[label=(\roman*)]
  \item If $1/\mu \notin \{\lambda_n\}$, the unique solution is
  \[
    f = g + \sum_{n \ge 1} \frac{\mu\lambda_n}{1 - \mu\lambda_n}\, \inner{g}{e_n}\, e_n.
  \]
  \item If $1/\mu = \lambda_k$ for some $k$, a solution exists if and only if
        $g \perp \Ker(\Id - \mu T_K)$.
\end{enumerate}
\end{theorem}

\begin{proof}
Expand $f = \sum_n c_n e_n$ and $g = \sum_n g_n e_n$. The equation yields
$c_n(1 - \mu\lambda_n) = g_n$. If $1 - \mu\lambda_n \ne 0$ for all $n$,
then $c_n = g_n/(1-\mu\lambda_n)$ and
\[
  f = \sum_n \frac{g_n}{1 - \mu\lambda_n}\, e_n
  = g + \sum_n \frac{\mu\lambda_n}{1 - \mu\lambda_n}\, g_n\, e_n.
\]
Convergence holds since $\lambda_n \to 0$. When $1 - \mu\lambda_k = 0$,
solvability requires $g_k = 0$, i.e., $g \perp e_k$.
\end{proof}

\index{Fredholm integral equation|)}

\section{Exercises}

\begin{exercise}[$\star$]\label{exo:compact-weak-strong}
\index{compact operator!weakly convergent sequences}
Let $H$ be a Hilbert space and $T \in \mathcal{K}(H)$. Show that if
$x_n \weakto x$, then $Tx_n \to Tx$ in norm.
\end{exercise}

\begin{exercise}[$\star$]\label{exo:compact-product}
Let $T \in \mathcal{L}(E,F)$ and $S \in \mathcal{K}(F,G)$. Show $ST$ is compact.
Give an example where neither $S$ nor $T$ is compact but $ST$ is.
\end{exercise}

\begin{exercise}[$\star\star$]\label{exo:volterra-compact}
\index{Volterra operator}
Let $(Vf)(x) = \int_0^x f(t)\, \dd t$ on $L^2([0,1])$. Show $V$ is compact
with $\sigma(V) = \{0\}$ and $\sigma_p(V) = \emptyset$.
\end{exercise}

\begin{exercise}[$\star\star$]\label{exo:trace-class}
\index{trace class operator}
Show that every trace class operator is Hilbert--Schmidt. If $T$ is trace class
and self-adjoint, show $\Tr(T) = \sum_n \inner{Te_n}{e_n}$ is basis-independent.
\end{exercise}

\begin{exercise}[$\star\star$]\label{exo:mercer}
\index{Mercer's theorem}
(Mercer's theorem.) Let $K \colon [0,1]^2 \to \R$ be continuous, symmetric, with
$T_K \ge 0$. Show $K(x,y) = \sum_n \lambda_n e_n(x) e_n(y)$ uniformly.
\end{exercise}

\begin{exercise}[$\star\star\star$]\label{exo:schauder-fixed-pt}
\index{Schauder fixed-point theorem}
(Schauder.) Let $C$ be closed, bounded, convex in a Banach space $E$ and
$T \colon C \to C$ compact. Prove $T$ has a fixed point.
\end{exercise}

\begin{exercise}[$\star\star$]\label{exo:eigenvalue-variational}
\index{min-max principle}
Let $T \in \mathcal{K}(H)$ be positive self-adjoint. Prove the min-max principle:
\[
  \lambda_n = \min_{\codim V = n-1} \max_{\substack{x \in V \\ \norm{x}=1}}
  \inner{Tx}{x}
  = \max_{\dim W = n} \min_{\substack{x \in W \\ \norm{x}=1}} \inner{Tx}{x}.
\]
\end{exercise}

\begin{exercise}[$\star$]\label{exo:compact-spectrum-example}
Give an example of a bounded operator on $\ell^2$ whose spectrum is $[0,1]$.
\end{exercise}

\begin{exercise}[$\star\star\star$]\label{exo:fredholm-index}
\index{Fredholm operator!index}
Show: (a) $\Id - K$ is Fredholm of index zero for all $K \in \mathcal{K}(E)$;
(b) the set of Fredholm operators is open and the index is locally constant.
\end{exercise}

\index{compact operator|)}
\index{operator!compact|)}

%%%%%%%%%%%%%%%%%%%%%%%%%%%%%%%%%%%%%%%%%%%%%%%%%%%%%%%%%%%%
%%  CHAPTER 10: The Spectral Theorem for Self-Adjoint Operators
%%%%%%%%%%%%%%%%%%%%%%%%%%%%%%%%%%%%%%%%%%%%%%%%%%%%%%%%%%%%
\chapter{The Spectral Theorem for Self-Adjoint Operators}
\label{ch:spectral-thm}
\minitoc

\section{Continuous Functional Calculus}\label{sec:cont-func-calc}

\index{functional calculus!continuous|(}
\index{continuous functional calculus|(}

The spectral theorem for compact self-adjoint operators expresses $T$ as a
(norm-convergent) sum involving its eigenvalues and eigenprojections. For a
general bounded self-adjoint operator, eigenvalues need not exist (consider
multiplication by $x$ on $L^2([0,1])$), and we must develop a more sophisticated
framework.

\begin{notation}\label{not:spectrum-sa}
For a bounded self-adjoint operator $T$ on a Hilbert space $H$, we write
$\sigma(T) \subset \R$ for its spectrum and set
$m = \inf_{\norm{x}=1}\inner{Tx}{x}$, $M = \sup_{\norm{x}=1}\inner{Tx}{x}$,
so that $\sigma(T) \subset [m, M]$.
\end{notation}

\begin{theorem}[Continuous functional calculus]\label{thm:cont-func-calc}
\index{continuous functional calculus!theorem}
\index{functional calculus!continuous!theorem}
Let $T \in \mathcal{L}(H)$ be self-adjoint. There exists a unique isometric
$*$-algebra homomorphism
\[
  \Phi_T \colon C(\sigma(T)) \longrightarrow \mathcal{L}(H),
  \qquad f \longmapsto f(T),
\]
satisfying:
\begin{enumerate}[label=(\roman*)]
  \item $\Phi_T(1) = \Id$ and $\Phi_T(\mathrm{id}) = T$,
        where $\mathrm{id}(\lambda) = \lambda$.
  \item $\Phi_T(\bar{f}) = \Phi_T(f)^*$.
  \item $\norm{f(T)} = \norm{f}_{C(\sigma(T))} := \sup_{\lambda \in \sigma(T)} \abs{f(\lambda)}$.
  \item $\sigma(f(T)) = f(\sigma(T))$ (spectral mapping theorem).
        \index{spectral mapping theorem}
\end{enumerate}
\end{theorem}

\begin{proof}
\textbf{Step 1 (Polynomials).} For a polynomial $p(\lambda) = \sum_{k=0}^n a_k \lambda^k$,
define $p(T) = \sum_{k=0}^n a_k T^k$. We must show that
$\norm{p(T)} = \sup_{\lambda \in \sigma(T)} \abs{p(\lambda)}$.

Since $T$ is self-adjoint, $p(T)$ is normal (it commutes with its adjoint).
For a self-adjoint operator $S$, $\norm{S} = r(S)$ (the spectral radius).
For a normal operator $N$, $\norm{N^2} = \norm{N^*N} = \norm{N}^2$, hence
by induction $\norm{N^{2^k}} = \norm{N}^{2^k}$, giving
$\norm{N} = r(N) = \sup_{\mu \in \sigma(N)} \abs{\mu}$.
Since $\sigma(p(T)) = p(\sigma(T))$ by the polynomial spectral mapping theorem,
\[
  \norm{p(T)} = r(p(T)) = \sup_{\mu \in \sigma(p(T))} \abs{\mu}
  = \sup_{\lambda \in \sigma(T)} \abs{p(\lambda)}.
\]

\textbf{Step 2 (Extension).} The map $p \mapsto p(T)$ is an isometry from
$(C_{\mathrm{poly}}(\sigma(T)), \norm{\cdot}_\infty)$ to $\mathcal{L}(H)$.
By the Weierstrass approximation theorem\index{Weierstrass approximation theorem},
polynomials are dense in $C(\sigma(T))$ (since $\sigma(T)$ is a compact subset
of $\R$). The isometry extends uniquely to a map
$\Phi_T \colon C(\sigma(T)) \to \mathcal{L}(H)$.

\textbf{Step 3 (Properties).} The $*$-homomorphism properties follow by
continuity from the polynomial case. The spectral mapping theorem:
$\mu \in \sigma(f(T))$ iff $f(T) - \mu\Id$ is not invertible. By the
isometry, this happens iff $f - \mu$ has no inverse in $C(\sigma(T))$,
i.e., iff $f(\lambda) = \mu$ for some $\lambda \in \sigma(T)$.
\end{proof}

\index{functional calculus!continuous|)}
\index{continuous functional calculus|)}

\section{Spectral Measures and Resolution of the Identity}
\label{sec:spectral-measures}

\index{spectral measure|(}
\index{resolution of the identity|(}

\begin{definition}[Spectral measure]\label{def:spectral-measure}
\index{spectral measure!definition}
Let $H$ be a Hilbert space and $(\Omega, \mathcal{B})$ a measurable space. A
\emph{spectral measure} (or \emph{projection-valued measure}) on $(\Omega, \mathcal{B})$
with values in $\mathcal{L}(H)$ is a map
$E \colon \mathcal{B} \to \mathcal{L}(H)$ satisfying:
\begin{enumerate}[label=(\roman*)]
  \item $E(\Delta)$ is an orthogonal projection for each $\Delta \in \mathcal{B}$.
  \item $E(\emptyset) = 0$ and $E(\Omega) = \Id$.
  \item $E(\Delta_1 \cap \Delta_2) = E(\Delta_1) E(\Delta_2)$ for all
        $\Delta_1, \Delta_2 \in \mathcal{B}$.
  \item If $(\Delta_n)$ are pairwise disjoint, then
        $E(\bigcup_n \Delta_n) = \sum_n E(\Delta_n)$ (strong operator convergence).
\end{enumerate}
\end{definition}

\begin{remark}\label{rem:scalar-measure}
\index{spectral measure!scalar}
For any $x, y \in H$, the map $\Delta \mapsto \inner{E(\Delta)x}{y}$ defines a
complex Borel measure $\mu_{x,y}$ on $\Omega$. When $x = y$, the measure
$\mu_{x,x}(\Delta) = \norm{E(\Delta)x}^2$ is positive. These are called the
\emph{scalar spectral measures}.
\end{remark}

\begin{definition}[Resolution of the identity]\label{def:resolution-identity}
\index{resolution of the identity!definition}
For a bounded self-adjoint operator $T$ with $\sigma(T) \subset [m, M]$, a
\emph{resolution of the identity} is a family of projections $(E_\lambda)_{\lambda \in \R}$
satisfying:
\begin{enumerate}[label=(\roman*)]
  \item $E_\lambda \le E_\mu$ (i.e., $\im E_\lambda \subset \im E_\mu$)
        for $\lambda \le \mu$.
  \item $E_\lambda = 0$ for $\lambda < m$ and $E_\lambda = \Id$ for $\lambda \ge M$.
  \item $E_\lambda$ is strongly right-continuous:
        $E_{\lambda^+} = E_\lambda$.
\end{enumerate}
\end{definition}

\section{The Spectral Theorem (Spectral Measure Version)}
\label{sec:spectral-thm-measure}

\index{spectral theorem!bounded self-adjoint|(}

\begin{theorem}[Spectral theorem for bounded self-adjoint operators]
\label{thm:spectral-thm-bounded-sa}
\index{spectral theorem!bounded self-adjoint!statement}
Let $T \in \mathcal{L}(H)$ be self-adjoint. There exists a unique spectral measure
$E$ on $(\sigma(T), \mathcal{B}(\sigma(T)))$ such that
\[
  T = \int_{\sigma(T)} \lambda\, \dd E(\lambda),
\]
meaning $\inner{Tx}{y} = \int_{\sigma(T)} \lambda\, \dd\mu_{x,y}(\lambda)$
for all $x, y \in H$. More generally, for every bounded Borel function
$f \colon \sigma(T) \to \C$:
\[
  f(T) = \int_{\sigma(T)} f(\lambda)\, \dd E(\lambda) \in \mathcal{L}(H).
\]
\end{theorem}

\begin{proof}
\index{spectral theorem!bounded self-adjoint!proof}
\textbf{Step 1 (From continuous to Borel functions).}
By Theorem~\ref{thm:cont-func-calc}, we have a $*$-homomorphism
$\Phi_T \colon C(\sigma(T)) \to \mathcal{L}(H)$. For each $x \in H$, the map
$f \mapsto \inner{f(T)x}{x}$ is a positive linear functional on $C(\sigma(T))$
(positive because $f \ge 0$ implies $f(T) \ge 0$).

By the Riesz representation theorem\index{Riesz representation theorem!for measures},
there exists a unique positive Borel measure $\mu_x$ on $\sigma(T)$ such that
\[
  \inner{f(T)x}{x} = \int_{\sigma(T)} f(\lambda)\, \dd\mu_x(\lambda)
  \quad \text{for all } f \in C(\sigma(T)).
\]
Note $\mu_x(\sigma(T)) = \inner{\Id\, x}{x} = \norm{x}^2$.

\textbf{Step 2 (Complex measures).}
By polarization\index{polarization identity}, for $x, y \in H$ define the complex
measure $\mu_{x,y}$ via
\[
  \mu_{x,y} = \tfrac{1}{4}\bigl(\mu_{x+y} - \mu_{x-y}
  + i\mu_{x+iy} - i\mu_{x-iy}\bigr).
\]
Then $\inner{f(T)x}{y} = \int f\, \dd\mu_{x,y}$ for all $f \in C(\sigma(T))$.

\textbf{Step 3 (Extension to bounded Borel functions).}
For a bounded Borel function $f$, the map $(x, y) \mapsto \int f\, \dd\mu_{x,y}$
is a bounded sesquilinear form. By the Riesz representation theorem for bounded
sesquilinear forms, there exists a unique $f(T) \in \mathcal{L}(H)$ with
$\inner{f(T)x}{y} = \int f\, \dd\mu_{x,y}$.

The map $f \mapsto f(T)$ is a $*$-homomorphism from $B(\sigma(T))$ (bounded Borel
functions) to $\mathcal{L}(H)$. The multiplicativity $f(T)g(T) = (fg)(T)$ follows
from the continuous case by a monotone class argument\index{monotone class theorem}.

\textbf{Step 4 (Spectral measure).}
For $\Delta \in \mathcal{B}(\sigma(T))$, set $E(\Delta) = \mathbf{1}_\Delta(T)$.
Since $\mathbf{1}_\Delta^2 = \mathbf{1}_\Delta = \overline{\mathbf{1}_\Delta}$,
we have $E(\Delta)^2 = E(\Delta) = E(\Delta)^*$, so $E(\Delta)$ is an orthogonal
projection. The spectral measure axioms follow from the properties of the
$*$-homomorphism:
\begin{itemize}
  \item $E(\sigma(T)) = \mathbf{1}_{\sigma(T)}(T) = \Id$.
  \item $E(\Delta_1 \cap \Delta_2) = \mathbf{1}_{\Delta_1 \cap \Delta_2}(T)
        = \mathbf{1}_{\Delta_1}(T)\mathbf{1}_{\Delta_2}(T) = E(\Delta_1)E(\Delta_2)$.
  \item Countable additivity in the strong operator topology follows from the
        dominated convergence theorem applied to $\mu_{x,x}$.
\end{itemize}

\textbf{Uniqueness.} If $E'$ is another spectral measure with
$T = \int \lambda\, \dd E'(\lambda)$, then $p(T) = \int p\, \dd E'$ for every
polynomial $p$. By density of polynomials and uniqueness in the Riesz representation
theorem, $\mu'_{x,x} = \mu_{x,x}$ for all $x$, hence $E' = E$.
\end{proof}

\index{spectral theorem!bounded self-adjoint|)}

\section{Borel Functional Calculus}\label{sec:borel-func-calc}

\index{functional calculus!Borel|(}
\index{Borel functional calculus|(}

The spectral theorem gives us a powerful extension of the continuous functional
calculus.

\begin{definition}[Borel functional calculus]\label{def:borel-func-calc}
\index{Borel functional calculus!definition}
For $T$ self-adjoint and $f \colon \sigma(T) \to \C$ a bounded Borel function,
we define
\[
  f(T) = \int_{\sigma(T)} f(\lambda)\, \dd E(\lambda),
\]
where $E$ is the spectral measure of $T$.
\end{definition}

\begin{proposition}\label{prop:borel-calc-properties}
\index{Borel functional calculus!properties}
Let $T$ be bounded self-adjoint with spectral measure $E$, and let
$f, g \in B(\sigma(T))$.
\begin{enumerate}[label=(\roman*)]
  \item $(\alpha f + \beta g)(T) = \alpha f(T) + \beta g(T)$.
  \item $(fg)(T) = f(T)\, g(T)$.
  \item $\bar{f}(T) = f(T)^*$.
  \item $\norm{f(T)} \le \norm{f}_\infty$.
  \item If $f_n \to f$ pointwise with $\sup_n \norm{f_n}_\infty < \infty$, then
        $f_n(T) \to f(T)$ in the strong operator topology.
  \item $\sigma(f(T)) \subset \overline{f(\sigma(T))}$.
\end{enumerate}
\end{proposition}

\begin{proof}
Properties (i)--(iv) follow from the construction in Theorem~\ref{thm:spectral-thm-bounded-sa}.
For~(v), let $x \in H$; then
$\norm{f_n(T)x - f(T)x}^2 = \int \abs{f_n - f}^2\, \dd\mu_{x,x}$.
By dominated convergence (the integrand is bounded by $(2\sup_n\norm{f_n}_\infty)^2$
and $\mu_{x,x}$ is finite), this tends to $0$.
Property~(vi) follows from (ii): if $\mu \notin \overline{f(\sigma(T))}$, then
$h = 1/(f - \mu)$ is bounded and Borel on $\sigma(T)$, so
$h(T)(f(T) - \mu\Id) = \Id$.
\end{proof}

\index{functional calculus!Borel|)}
\index{Borel functional calculus|)}

\section{Positive Operators and Square Root}\label{sec:positive-sqrt}

\index{positive operator|(}
\index{square root!of positive operator|(}

\begin{definition}\label{def:positive-op}
\index{positive operator!definition}
A self-adjoint operator $T \in \mathcal{L}(H)$ is \emph{positive} (written
$T \ge 0$) if $\inner{Tx}{x} \ge 0$ for all $x \in H$. Equivalently,
$\sigma(T) \subset [0, \infty)$.
\end{definition}

\begin{theorem}[Square root]\label{thm:square-root}
\index{square root!existence and uniqueness}
If $T \ge 0$, there exists a unique $S \ge 0$ with $S^2 = T$. We write
$S = T^{1/2}$ or $S = \sqrt{T}$. Moreover, $S$ commutes with every operator
that commutes with $T$.
\end{theorem}

\begin{proof}
\textbf{Existence.} Apply the continuous functional calculus with
$f(\lambda) = \sqrt{\lambda}$ on $\sigma(T) \subset [0, \norm{T}]$. Set
$S = f(T) = \sqrt{\cdot}(T)$. Then $S^2 = (\sqrt{\cdot})^2(T) = \mathrm{id}(T) = T$,
$S^* = \overline{\sqrt{\cdot}}(T) = \sqrt{\cdot}(T) = S$, and
$\sigma(S) = f(\sigma(T)) \subset [0, \infty)$, so $S \ge 0$.

\textbf{Uniqueness.} Let $R \ge 0$ with $R^2 = T$. Then $R$ commutes with $T = R^2$,
hence with any polynomial in $T$, and by continuity with $S = \sqrt{T}$.
Let $N = \Ker(R - S)^\perp$ and decompose $H = \Ker(R - S) \oplus N$. On $N$,
$R - S$ is injective. Now $(R - S)(R + S) = R^2 - S^2 = 0$ (using $RS = SR$),
so for $x \in N$, $(R + S)x \in \Ker(R - S)$, hence
$\inner{(R+S)x}{x} = 0$ for $x \in N$ (since $N \perp \Ker(R - S)$). But
$R + S \ge 0$, so $\inner{Rx}{x} + \inner{Sx}{x} = 0$ with both terms
non-negative forces $Rx = Sx = 0$ for $x \in N$, hence $(R - S)x = 0$. Thus
$N \subset \Ker(R - S)$, so $N = \{0\}$ and $R = S$.

\textbf{Commuting property.} If $A$ commutes with $T$, then $A$ commutes with
$p(T)$ for every polynomial $p$, hence with $\sqrt{T}$ by continuity of
$\Phi_T$.
\end{proof}

\index{positive operator|)}
\index{square root!of positive operator|)}

\section{Polar Decomposition}\label{sec:polar-decomp}

\index{polar decomposition|(}

\begin{definition}\label{def:abs-value-op}
\index{absolute value!of operator}
For $T \in \mathcal{L}(H)$, define $\abs{T} = (T^*T)^{1/2}$. Note $T^*T \ge 0$,
so $\abs{T}$ exists by Theorem~\ref{thm:square-root}.
\end{definition}

\begin{theorem}[Polar decomposition]\label{thm:polar-decomp}
\index{polar decomposition!theorem}
For every $T \in \mathcal{L}(H)$, there exists a unique partial isometry
$U \in \mathcal{L}(H)$\index{partial isometry} such that $T = U\abs{T}$ and
$\Ker U = \Ker T$. Moreover:
\begin{enumerate}[label=(\roman*)]
  \item $U$ is an isometry from $(\Ker T)^\perp = \overline{\im \abs{T}}$ onto
        $\overline{\im T}$.
  \item $\abs{T} = U^* T$.
  \item $T^* = \abs{T}\, U^*$ and $\abs{T^*} = U\abs{T}U^*$.
\end{enumerate}
\end{theorem}

\begin{proof}
\textbf{Construction of $U$.} For $x \in H$, observe
$\norm{Tx}^2 = \inner{T^*Tx}{x} = \inner{\abs{T}^2 x}{x} = \norm{\abs{T}x}^2$.
Hence the map $\abs{T}x \mapsto Tx$ is a well-defined isometry from
$\im\abs{T}$ to $\im T$. Extend by continuity to an isometry
$U_0 \colon \overline{\im\abs{T}} \to \overline{\im T}$, and set $U = 0$
on $(\overline{\im\abs{T}})^\perp = \Ker\abs{T} = \Ker T$ (the last equality
because $\norm{\abs{T}x} = \norm{Tx}$). Then $U$ is a partial isometry with
$T = U\abs{T}$ and $\Ker U = \Ker T$.

\textbf{Uniqueness.} If $T = V\abs{T}$ with $V$ a partial isometry and
$\Ker V = \Ker T$, then $V$ agrees with $U$ on $\im\abs{T}$ (since
$V\abs{T}x = Tx = U\abs{T}x$) and on $\Ker\abs{T} = \Ker T$ (since both
vanish). By density, $V = U$.

\textbf{Property (ii):} $U^*T = U^*U\abs{T} = P_{\overline{\im\abs{T}}}\abs{T} = \abs{T}$
since $\im\abs{T} \subset \overline{\im\abs{T}}$.
\end{proof}

\index{polar decomposition|)}

\section{Unbounded Operators}\label{sec:unbounded-ops}

\index{unbounded operator|(}

Many important operators in analysis and physics (differential operators, for
instance) are not bounded. We develop the basic framework.

\begin{definition}\label{def:unbounded-op}
\index{unbounded operator!definition}
An \emph{unbounded operator} on a Hilbert space $H$ is a linear map
$T \colon \dom(T) \to H$ where $\dom(T)$\index{domain} is a linear subspace of $H$
(the \emph{domain} of $T$), not necessarily all of $H$.
\end{definition}

\begin{definition}[Graph, closed operator]\label{def:graph-closed}
\index{graph!of operator}
\index{closed operator}
\index{operator!closed}
The \emph{graph} of $T$ is $\Gamma(T) = \{(x, Tx) : x \in \dom(T)\} \subset H \times H$.
The operator $T$ is \emph{closed} if $\Gamma(T)$ is closed in $H \times H$, i.e.,
whenever $x_n \in \dom(T)$, $x_n \to x$, and $Tx_n \to y$, we have
$x \in \dom(T)$ and $Tx = y$.
\end{definition}

\begin{definition}[Closable operator]\label{def:closable}
\index{closable operator}
$T$ is \emph{closable} if the closure $\overline{\Gamma(T)}$ is the graph of an
operator (equivalently, if $x_n \in \dom(T)$, $x_n \to 0$, $Tx_n \to y$ implies
$y = 0$). The resulting operator $\bar{T}$ is the \emph{closure} of $T$.
\end{definition}

\begin{definition}[Symmetric, self-adjoint]\label{def:sa-unbounded}
\index{self-adjoint operator!unbounded}
\index{symmetric operator}
Let $T$ be densely defined ($\overline{\dom(T)} = H$).
\begin{itemize}
  \item $T$ is \emph{symmetric} if $\inner{Tx}{y} = \inner{x}{Ty}$ for all
        $x, y \in \dom(T)$.
  \item The \emph{adjoint} $T^*$ has domain
        $\dom(T^*) = \{y \in H : x \mapsto \inner{Tx}{y} \text{ is bounded on } \dom(T)\}$,
        and $T^*y$ is the unique element such that $\inner{Tx}{y} = \inner{x}{T^*y}$.
  \item $T$ is \emph{self-adjoint} if $T = T^*$ (meaning $\dom(T) = \dom(T^*)$
        and $Tx = T^*x$ for all $x \in \dom(T)$).
\end{itemize}
A symmetric operator satisfies $T \subset T^*$ (i.e., $\dom(T) \subset \dom(T^*)$
and $T^*|_{\dom(T)} = T$). Self-adjointness is the stronger condition
$\dom(T) = \dom(T^*)$.
\end{definition}

\begin{remark}\label{rem:sa-vs-symmetric}
\index{symmetric vs.\ self-adjoint}
The distinction between symmetric and self-adjoint is crucial. A symmetric operator
may have many self-adjoint extensions, exactly one, or none. The spectral theorem
requires genuine self-adjointness.
\end{remark}

\begin{example}[The Laplacian]\label{ex:laplacian}
\index{Laplacian}
\index{operator!Laplacian}
On $H = L^2(\R^d)$, define $T = -\Delta$ with $\dom(T) = H^2(\R^d)$
(the Sobolev space\index{Sobolev space}). This is self-adjoint. Indeed,
for $f, g \in H^2(\R^d)$:
\[
  \inner{-\Delta f}{g}_{L^2} = \int_{\R^d} \nabla f \cdot \overline{\nabla g}\, \dd x
  = \inner{f}{-\Delta g}_{L^2}
\]
by integration by parts. Self-adjointness (not just symmetry) can be proved via
the Fourier transform:\index{Fourier transform} under $\mathcal{F}$, $-\Delta$
becomes multiplication by $\abs{\xi}^2$, which is self-adjoint on
$\dom = \{f \in L^2 : \abs{\xi}^2 \hat{f} \in L^2\} = H^2(\R^d)$.
\end{example}

\begin{example}[Momentum operator]\label{ex:momentum}
\index{momentum operator}
On $L^2(\R)$, define $P = -i\frac{\dd}{\dd x}$ with $\dom(P) = H^1(\R)$. Then
$P$ is self-adjoint. Via Fourier transform, $P$ becomes multiplication by $\xi$.
\end{example}

\begin{proposition}[Criteria for self-adjointness]\label{prop:sa-criteria}
\index{self-adjoint operator!criteria}
Let $T$ be symmetric and densely defined. The following are equivalent:
\begin{enumerate}[label=(\roman*)]
  \item $T$ is self-adjoint.
  \item $T$ is closed and $\Ker(T^* \pm i) = \{0\}$.
  \item $\ran(T \pm i) = H$.
\end{enumerate}
\end{proposition}

\begin{proof}
(i)$\Rightarrow$(ii): If $T = T^*$, then $T$ is closed (since $T^*$ is always
closed). If $T^* y = iy$, then $\inner{Ty}{y} = \inner{y}{T^*y} = \inner{y}{iy} = -i\norm{y}^2$
but also $\inner{Ty}{y} = \inner{T^*y}{y} = \inner{iy}{y} = i\norm{y}^2$, giving
$y = 0$.

(ii)$\Rightarrow$(iii): We show $\ran(T - i)$ is closed and dense. Closedness:
if $x \in \dom(T)$, $\norm{(T - i)x}^2 = \norm{Tx}^2 + \norm{x}^2$ (since
$\operatorname{Re}\inner{Tx}{ix} = 0$ for symmetric $T$), so $T - i$ is bounded
below, forcing $\ran(T - i)$ closed. Density: if $y \perp \ran(T - i)$, then
$y \in \dom(T^*)$ and $T^*y = iy$, hence $y = 0$ by hypothesis. Similarly for
$\ran(T + i)$.

(iii)$\Rightarrow$(i): Let $y \in \dom(T^*)$. We show $y \in \dom(T)$.
Since $\ran(T + i) = H$, there exists $x \in \dom(T)$ with
$(T + i)x = (T^* + i)y$. Since $T \subset T^*$, we have
$T^*(x - y) = -i(x - y)$, i.e., $(x - y) \in \Ker(T^* + i)$.
But $\ran(T - i) = H$ implies $\Ker(T^* + i) = \ran(T - i)^\perp = \{0\}$.
Hence $x = y$ and $y \in \dom(T)$.
\end{proof}

\section{Spectral Theorem for Unbounded Self-Adjoint Operators}
\label{sec:spectral-unbounded}

\index{spectral theorem!unbounded self-adjoint|(}

\begin{theorem}[Spectral theorem, unbounded version]\label{thm:spectral-unbounded}
\index{spectral theorem!unbounded self-adjoint!statement}
Let $T$ be a self-adjoint operator on $H$ (possibly unbounded). There exists a
unique spectral measure $E$ on $(\R, \mathcal{B}(\R))$ with
$E(\R \setminus \sigma(T)) = 0$ and
\[
  T = \int_{-\infty}^{\infty} \lambda\, \dd E(\lambda),
\]
in the sense that
\[
  \dom(T) = \Bigl\{x \in H : \int_{-\infty}^{\infty} \lambda^2\, \dd\mu_{x,x}(\lambda)
  < \infty\Bigr\}
\]
and $\inner{Tx}{y} = \int \lambda\, \dd\mu_{x,y}(\lambda)$ for $x \in \dom(T)$,
$y \in H$.

For every Borel measurable function $f \colon \R \to \C$, the operator
\[
  f(T) = \int f(\lambda)\, \dd E(\lambda)
\]
is defined on $\dom(f(T)) = \{x \in H : \int \abs{f}^2\, \dd\mu_{x,x} < \infty\}$.
\end{theorem}

The proof uses the Cayley transform\index{Cayley transform}
$U = (T - i)(T + i)^{-1}$, which is a unitary operator when $T$ is self-adjoint
(by Proposition~\ref{prop:sa-criteria}). One applies the spectral theorem for
unitary operators (a variant of the bounded case) to $U$, then transfers back
via $T = i(\Id + U)(\Id - U)^{-1}$.

\index{spectral theorem!unbounded self-adjoint|)}

\section{Application: Quantum Mechanics}\label{sec:quantum-mech}

\index{quantum mechanics|(}

The spectral theorem provides the mathematical foundation for quantum mechanics.

\begin{remark}[Observables as self-adjoint operators]
\index{observable}
In the quantum-mechanical formalism:\index{quantum mechanics!formalism}
\begin{itemize}
  \item The \emph{state space} is a separable Hilbert space $H$ (typically $L^2(\R^3)$).
  \item \emph{Observables} (position, momentum, energy, etc.) are represented by
        self-adjoint operators on $H$.
  \item If $T$ is an observable with spectral measure $E$ and the system is in
        state $\psi$ ($\norm{\psi} = 1$), the probability of measuring $T$ in a
        Borel set $\Delta \subset \R$ is
        $\Pr(T \in \Delta) = \inner{E(\Delta)\psi}{\psi} = \mu_{\psi,\psi}(\Delta)$.
  \item The expected value is
        $\langle T \rangle_\psi = \inner{T\psi}{\psi} = \int \lambda\, \dd\mu_{\psi,\psi}(\lambda)$.
  \item The Heisenberg uncertainty principle\index{Heisenberg uncertainty principle}
        $\Delta T \cdot \Delta S \ge \frac{1}{2}\abs{\inner{[T,S]\psi}{\psi}}$
        (where $[T,S] = TS - ST$) follows from the Cauchy--Schwarz inequality.
\end{itemize}
\end{remark}

\index{quantum mechanics|)}

\section{Exercises}

\begin{exercise}[$\star$]\label{exo:spectral-projection}
\index{spectral projection}
Let $T$ be bounded self-adjoint with spectral measure $E$. Show that
$T$ has an eigenvalue $\lambda$ if and only if $E(\{\lambda\}) \ne 0$, and in
that case $E(\{\lambda\})$ is the orthogonal projection onto $\Ker(T - \lambda)$.
\end{exercise}

\begin{exercise}[$\star\star$]\label{exo:func-calc-composition}
Let $T$ be self-adjoint, $f$ continuous, and $g$ bounded Borel. Show $g(f(T)) = (g \circ f)(T)$.
\end{exercise}

\begin{exercise}[$\star$]\label{exo:positive-sqrt-compute}
\index{square root!computation}
Let $A = \begin{pmatrix} 2 & 1 \\ 1 & 2 \end{pmatrix}$ on $\R^2$. Compute $A^{1/2}$.
\end{exercise}

\begin{exercise}[$\star\star$]\label{exo:spectral-theorem-multiplication}
\index{multiplication operator!spectral measure}
Let $g \in L^\infty(\R)$ be real-valued and $M_g$ the multiplication operator on
$L^2(\R)$. Describe the spectral measure of $M_g$ explicitly.
\end{exercise}

\begin{exercise}[$\star\star$]\label{exo:unbounded-sa-extension}
\index{self-adjoint extension}
Let $T = -\frac{\dd^2}{\dd x^2}$ on $\dom(T) = C_c^\infty((0,1))$ in $L^2((0,1))$.
Show $T$ is symmetric but not self-adjoint. Find its self-adjoint extensions
(parametrized by boundary conditions).
\end{exercise}

\begin{exercise}[$\star\star\star$]\label{exo:stone-theorem}
\index{Stone's theorem}
(Stone's theorem.) Let $(U(t))_{t \in \R}$ be a strongly continuous one-parameter
unitary group on $H$. Show there exists a self-adjoint operator $A$ such that
$U(t) = e^{itA}$ for all $t$.
\end{exercise}

\begin{exercise}[$\star\star$]\label{exo:polar-normal}
Show that $T$ is normal if and only if $\abs{T^*} = \abs{T}$ (equivalently,
$U$ in the polar decomposition is unitary on $(\Ker T)^\perp$).
\end{exercise}

\begin{exercise}[$\star\star\star$]\label{exo:kato-rellich}
\index{Kato--Rellich theorem}
(Kato--Rellich.) Let $A$ be self-adjoint and $B$ symmetric with $\dom(A) \subset \dom(B)$.
Suppose there exist $a < 1$ and $b \ge 0$ with $\norm{Bx} \le a\norm{Ax} + b\norm{x}$
for all $x \in \dom(A)$. Prove that $A + B$ is self-adjoint on $\dom(A)$.
\end{exercise}

\begin{exercise}[$\star$]\label{exo:hydrogen-atom}
\index{hydrogen atom}
Describe the Hamiltonian of the hydrogen atom $H = -\Delta - 1/\abs{x}$ on
$L^2(\R^3)$ and explain why the Kato--Rellich theorem guarantees it is
self-adjoint on $H^2(\R^3)$.
\end{exercise}


%%%%%%%%%%%%%%%%%%%%%%%%%%%%%%%%%%%%%%%%%%%%%%%%%%%%%%%%%%%%
%%  CHAPTER 11: L^p Spaces and Duality
%%%%%%%%%%%%%%%%%%%%%%%%%%%%%%%%%%%%%%%%%%%%%%%%%%%%%%%%%%%%
\chapter{$L^p$ Spaces and Duality}\label{ch:Lp-duality}
\minitoc

\section{H\"older and Minkowski Inequalities}\label{sec:holder-minkowski}

\index{Holder inequality@H\"older inequality|(}
\index{Minkowski inequality|(}

Throughout this chapter, $(\Omega, \mathcal{A}, \mu)$ denotes a $\sigma$-finite
measure space and $1 \le p \le \infty$. We write $p'$ (or $q$) for the conjugate
exponent:\index{conjugate exponent} $\frac{1}{p} + \frac{1}{p'} = 1$ (with
$1' = \infty$ and $\infty' = 1$).

\begin{lemma}[Young's inequality for products]\label{lem:young-ineq}
\index{Young's inequality!for products}
For $a, b \ge 0$ and $1 < p < \infty$ with conjugate $q$:
\[
  ab \le \frac{a^p}{p} + \frac{b^q}{q}.
\]
Equality holds if and only if $a^p = b^q$.
\end{lemma}

\begin{proof}
The function $t \mapsto e^t$ is convex, so
$e^{\alpha s + \beta t} \le \alpha e^s + \beta e^t$ for $\alpha, \beta > 0$,
$\alpha + \beta = 1$. Take $\alpha = 1/p$, $\beta = 1/q$, $s = p\ln a$,
$t = q\ln b$ (for $a, b > 0$; the case $a = 0$ or $b = 0$ is trivial).
\end{proof}

\begin{theorem}[H\"older's inequality]\label{thm:holder}
\index{Holder inequality@H\"older inequality!statement}
Let $1 \le p \le \infty$ and $q = p'$. If $f \in L^p(\mu)$ and $g \in L^q(\mu)$,
then $fg \in L^1(\mu)$ and
\[
  \norm{fg}_{L^1} \le \norm{f}_{L^p}\, \norm{g}_{L^q}.
\]
\end{theorem}

\begin{proof}
The cases $p = 1$ (or $p = \infty$) are immediate. For $1 < p < \infty$, we may
assume $\norm{f}_p, \norm{g}_q > 0$. Set $\tilde{f} = f/\norm{f}_p$ and
$\tilde{g} = g/\norm{g}_q$. By Young's inequality:
\[
  \abs{\tilde{f}(x)\tilde{g}(x)} \le \frac{\abs{\tilde{f}(x)}^p}{p}
  + \frac{\abs{\tilde{g}(x)}^q}{q}.
\]
Integrating: $\int \abs{\tilde{f}\tilde{g}}\, \dd\mu \le \frac{1}{p} + \frac{1}{q} = 1$.
Hence $\norm{fg}_1 \le \norm{f}_p\, \norm{g}_q$.
\end{proof}

\begin{theorem}[Minkowski's inequality]\label{thm:minkowski}
\index{Minkowski inequality!statement}
For $1 \le p \le \infty$ and $f, g \in L^p(\mu)$:
\[
  \norm{f + g}_{L^p} \le \norm{f}_{L^p} + \norm{g}_{L^p}.
\]
\end{theorem}

\begin{proof}
For $p = 1$ and $p = \infty$ this is the triangle inequality. For $1 < p < \infty$:
\begin{align*}
  \norm{f+g}_p^p
  &= \int \abs{f+g}^p\, \dd\mu
  \le \int \abs{f+g}^{p-1}\abs{f}\, \dd\mu + \int \abs{f+g}^{p-1}\abs{g}\, \dd\mu.
\end{align*}
Since $(p-1)q = p$, the function $\abs{f+g}^{p-1} \in L^q$ with
$\norm{\abs{f+g}^{p-1}}_q = \norm{f+g}_p^{p/q}$. Apply H\"older to each term:
\[
  \norm{f+g}_p^p \le \norm{f+g}_p^{p/q}(\norm{f}_p + \norm{g}_p).
\]
Divide by $\norm{f+g}_p^{p/q}$ (assuming it is nonzero) and use $p - p/q = 1$.
\end{proof}

\index{Holder inequality@H\"older inequality|)}
\index{Minkowski inequality|)}

\section{Completeness of $L^p$: The Riesz--Fischer Theorem}
\label{sec:Lp-completeness}

\index{Riesz--Fischer theorem|(}
\index{Lp space@$L^p$ space!completeness|(}

\begin{theorem}[Riesz--Fischer]\label{thm:riesz-fischer}
\index{Riesz--Fischer theorem!statement}
For $1 \le p \le \infty$, the space $L^p(\Omega, \mu)$ is a Banach space.
\end{theorem}

\begin{proof}[Proof for $1 \le p < \infty$]
\index{Riesz--Fischer theorem!proof}
Let $(f_n)$ be Cauchy in $L^p$. We show it converges by extracting a subsequence
that converges both in $L^p$ and almost everywhere.

\textbf{Step 1 (Subsequence).} Choose $n_1 < n_2 < \cdots$ such that
$\norm{f_{n_{k+1}} - f_{n_k}}_p < 2^{-k}$. Set $g_k = f_{n_k}$ and
\[
  S_N(x) = \sum_{k=1}^{N} \abs{g_{k+1}(x) - g_k(x)}, \qquad
  S(x) = \sum_{k=1}^{\infty} \abs{g_{k+1}(x) - g_k(x)}.
\]
By Minkowski's inequality, $\norm{S_N}_p \le \sum_{k=1}^N 2^{-k} < 1$.
By the monotone convergence theorem\index{monotone convergence theorem},
$\norm{S}_p \le 1 < \infty$, so $S(x) < \infty$ a.e.

\textbf{Step 2 (Pointwise limit).} Where $S(x) < \infty$, the telescoping series
$g_1(x) + \sum_{k=1}^\infty (g_{k+1}(x) - g_k(x))$ converges absolutely.
Define $f(x) = \lim_{k \to \infty} g_k(x)$ where this limit exists, and
$f(x) = 0$ elsewhere.

\textbf{Step 3 ($L^p$ convergence).} We have $\abs{g_k(x) - f(x)}^p \to 0$ a.e.
and $\abs{g_k - f}^p \le (2S)^p \in L^1$. By dominated convergence\index{dominated convergence theorem},
$\norm{g_k - f}_p \to 0$.

\textbf{Step 4 (Full sequence).} Since $(f_n)$ is Cauchy and the subsequence
$(g_k) = (f_{n_k})$ converges to $f$ in $L^p$, the full sequence $f_n \to f$
in $L^p$.
\end{proof}

\begin{proof}[Proof for $p = \infty$]
If $(f_n)$ is Cauchy in $L^\infty$, then for each pair $m, n$,
$\abs{f_m(x) - f_n(x)} \le \norm{f_m - f_n}_\infty$ outside a $\mu$-null set
$N_{m,n}$. The union $N = \bigcup_{m,n} N_{m,n}$ is still $\mu$-null, and on
$\Omega \setminus N$ the sequence $(f_n(x))$ is uniformly Cauchy. Define
$f(x) = \lim f_n(x)$ on $\Omega \setminus N$, $f = 0$ on $N$. Then
$\norm{f_n - f}_\infty \to 0$.
\end{proof}

\index{Riesz--Fischer theorem|)}
\index{Lp space@$L^p$ space!completeness|)}

\section{Density of $C_c^\infty$ in $L^p$}\label{sec:density-Cc}

\index{density!of smooth functions|(}
\index{Cc infty@$C_c^\infty$!density in $L^p$|(}

\begin{theorem}\label{thm:Cc-dense}
\index{Cc infty@$C_c^\infty$!density in $L^p$!statement}
Let $\Omega \subset \R^d$ be open and $1 \le p < \infty$. Then $C_c^\infty(\Omega)$
is dense in $L^p(\Omega)$.
\end{theorem}

The proof uses convolution with an approximation to the identity.

\begin{definition}[Mollifier]\label{def:mollifier}
\index{mollifier}
\index{approximation to the identity}
Let $\rho \in C_c^\infty(\R^d)$ with $\rho \ge 0$, $\supp\rho \subset B(0,1)$,
and $\int \rho = 1$. For $\varepsilon > 0$, set
$\rho_\varepsilon(x) = \varepsilon^{-d}\rho(x/\varepsilon)$, so that
$\supp\rho_\varepsilon \subset B(0,\varepsilon)$ and $\int \rho_\varepsilon = 1$.
The family $(\rho_\varepsilon)$ is called an \emph{approximation to the identity}
(or \emph{mollifier}).
\end{definition}

\section{Convolution and Approximation}\label{sec:convolution-approx}

\index{convolution|(}

\begin{definition}\label{def:convolution}
\index{convolution!definition}
For $f \in L^p(\R^d)$ and $g \in L^1(\R^d)$, the \emph{convolution} is
\[
  (f * g)(x) = \int_{\R^d} f(x - y)\, g(y)\, \dd y.
\]
\end{definition}

\begin{proposition}[Young's convolution inequality]\label{prop:young-conv}
\index{Young's inequality!for convolutions}
If $f \in L^p(\R^d)$ and $g \in L^1(\R^d)$, then $f * g \in L^p(\R^d)$ with
$\norm{f * g}_p \le \norm{f}_p\, \norm{g}_1$.
\end{proposition}

\begin{proof}
For $p = 1$ this is Fubini--Tonelli\index{Fubini theorem}. For $1 < p < \infty$:
write $\abs{g} = \abs{g}^{1/q}\abs{g}^{1/p}$ (with $1/p + 1/q = 1$) and apply
H\"older:
\begin{align*}
  \abs{(f*g)(x)}
  &\le \int \abs{f(x-y)}\abs{g(y)}\, \dd y \\
  &\le \Bigl(\int \abs{g(y)}\, \dd y\Bigr)^{1/q}
       \Bigl(\int \abs{f(x-y)}^p\abs{g(y)}\, \dd y\Bigr)^{1/p}.
\end{align*}
Raising to the $p$-th power and integrating in $x$:
$\norm{f*g}_p^p \le \norm{g}_1^{p/q} \int\!\int \abs{f(x-y)}^p\abs{g(y)}\, \dd y\, \dd x
= \norm{g}_1^{p/q}\norm{f}_p^p\norm{g}_1 = \norm{f}_p^p\norm{g}_1^p$.
\end{proof}

\begin{proposition}[Approximation by mollifiers]\label{prop:mollifier-approx}
\index{mollifier!approximation}
If $f \in L^p(\R^d)$ with $1 \le p < \infty$, then
$f * \rho_\varepsilon \in C^\infty(\R^d)$ and
$f * \rho_\varepsilon \to f$ in $L^p$ as $\varepsilon \to 0$.
\end{proposition}

\begin{proof}
Smoothness: differentiating under the integral sign (justified by dominated
convergence and the compactness of $\supp\rho_\varepsilon$) gives
$D^\alpha(f * \rho_\varepsilon) = f * D^\alpha\rho_\varepsilon \in C^\infty$.

$L^p$ convergence: by Minkowski's integral inequality\index{Minkowski integral inequality},
\[
  \norm{f * \rho_\varepsilon - f}_p
  = \norm[\bigg]{\int [\tau_y f - f]\, \rho_\varepsilon(y)\, \dd y}_p
  \le \int \norm{\tau_y f - f}_p\, \rho_\varepsilon(y)\, \dd y,
\]
where $(\tau_y f)(x) = f(x - y)$. Since translations are continuous in $L^p$
for $p < \infty$, given $\delta > 0$ there exists $\eta > 0$ with
$\norm{\tau_y f - f}_p < \delta$ for $\abs{y} < \eta$. For
$\varepsilon < \eta$, $\supp\rho_\varepsilon \subset B(0,\eta)$, so the
integral above is $< \delta$.
\end{proof}

\begin{proof}[Proof of Theorem~\ref{thm:Cc-dense}]
\textbf{Step 1.} Approximate $f \in L^p(\Omega)$ by $f_n = f\mathbf{1}_{K_n}$
where $K_n \nearrow \Omega$ are compact. Then $f_n \to f$ in $L^p$ by dominated
convergence.

\textbf{Step 2.} Each $f_n$ (extended by zero to $\R^d$) lies in $L^p(\R^d)$.
By Proposition~\ref{prop:mollifier-approx},
$g_\varepsilon = f_n * \rho_\varepsilon \in C^\infty$ and
$g_\varepsilon \to f_n$ in $L^p$. For $\varepsilon < d(K_n, \Omega^c)$,
$\supp g_\varepsilon \subset \Omega$ and is compact, so
$g_\varepsilon \in C_c^\infty(\Omega)$.
\end{proof}

\index{convolution|)}
\index{density!of smooth functions|)}
\index{Cc infty@$C_c^\infty$!density in $L^p$|)}

\section{The Dual of $L^p$}\label{sec:Lp-dual}

\index{dual space!of $L^p$|(}
\index{Lp space@$L^p$ space!dual|(}

The following is one of the central results in the theory of $L^p$ spaces.

\begin{theorem}[Dual of $L^p$]\label{thm:Lp-dual}
\index{dual space!of $L^p$!theorem}
Let $(\Omega, \mathcal{A}, \mu)$ be $\sigma$-finite and $1 \le p < \infty$. The
map
\[
  \Phi \colon L^q(\mu) \longrightarrow (L^p(\mu))', \qquad
  \Phi(g)(f) = \int_\Omega f\, \bar{g}\, \dd\mu,
\]
where $q = p'$, is an isometric isomorphism. That is,
$(L^p(\mu))' \cong L^q(\mu)$.
\end{theorem}

\begin{proof}
\index{dual space!of $L^p$!proof}
\textbf{Step 1 ($\Phi$ is an isometry).}
By H\"older, $\abs{\Phi(g)(f)} \le \norm{f}_p\norm{g}_q$, so
$\norm{\Phi(g)} \le \norm{g}_q$. For the reverse, take
$f = \abs{g}^{q-2}\bar{g}$ (when $g \ne 0$; here $\abs{g}^{q-2}\bar{g}$ has
$\abs{f}^p = \abs{g}^{(q-1)p} = \abs{g}^q$, so $f \in L^p$). Then
$\Phi(g)(f) = \int \abs{g}^q\, \dd\mu = \norm{g}_q^q$ and
$\norm{f}_p = \norm{g}_q^{q/p}$, giving
$\norm{\Phi(g)} \ge \norm{g}_q^q / \norm{g}_q^{q/p} = \norm{g}_q^{q - q/p} = \norm{g}_q$.

\textbf{Step 2 ($\Phi$ is surjective).} Let $\varphi \in (L^p(\mu))'$.

\emph{Case $p = 1$ ($q = \infty$):} For measurable $A$ with $\mu(A) < \infty$,
define $\nu(A) = \varphi(\mathbf{1}_A)$. Then $\nu$ is a signed (or complex) measure
absolutely continuous with respect to $\mu$ (since $\mu(A) = 0$ implies
$\mathbf{1}_A = 0$ in $L^1$). By the Radon--Nikodym theorem\index{Radon--Nikodym theorem},
$\dd\nu = g\, \dd\mu$ for some measurable $g$. One checks $g \in L^\infty$ with
$\norm{g}_\infty \le \norm{\varphi}$.

\emph{Case $1 < p < \infty$:} Define $\nu$ as above. The key estimate is:
for any measurable set $A$ with $\mu(A) > 0$,
\[
  \abs{\nu(A)} = \abs{\varphi(\mathbf{1}_A)} \le \norm{\varphi}\, \norm{\mathbf{1}_A}_p
  = \norm{\varphi}\, \mu(A)^{1/p},
\]
so $\nu \ll \mu$. By Radon--Nikodym, $\dd\nu = g\, \dd\mu$. We need $g \in L^q$.

For simple functions $s = \sum_k c_k \mathbf{1}_{A_k}$:
$\varphi(s) = \int s\, g\, \dd\mu$. By density of simple functions in $L^p$,
$\varphi(f) = \int f\, g\, \dd\mu$ for all $f \in L^p$ (once we verify integrability).

To show $g \in L^q$: let $g_n = g\, \mathbf{1}_{\{\abs{g} \le n\}}$. Set
$f_n = \abs{g_n}^{q-2}\overline{g_n}$, so $\abs{f_n}^p = \abs{g_n}^{(q-1)p} = \abs{g_n}^q$.
Then
\[
  \int \abs{g_n}^q\, \dd\mu = \int f_n\, g\, \dd\mu = \varphi(f_n)
  \le \norm{\varphi}\, \norm{f_n}_p = \norm{\varphi}\, \Bigl(\int \abs{g_n}^q\, \dd\mu\Bigr)^{1/p}.
\]
Hence $\norm{g_n}_q^{q/q} = \norm{g_n}_q^q / \norm{g_n}_q^{q/p} \le \norm{\varphi}$,
i.e., $\norm{g_n}_q \le \norm{\varphi}$. By monotone convergence,
$\norm{g}_q \le \norm{\varphi} < \infty$.
\end{proof}

\begin{remark}\label{rem:L-infty-dual}
\index{dual space!of $L^\infty$}
\index{Linfty@$L^\infty$!dual}
The dual of $L^\infty$ is \emph{not} $L^1$ in general; $(L^\infty(\mu))'$ is
strictly larger, consisting of all finitely additive signed measures absolutely
continuous with respect to $\mu$. One has $L^1 \hookrightarrow (L^\infty)'$
isometrically, but the inclusion is proper when $\mu$ is not purely atomic.
\end{remark}

\begin{remark}\label{rem:reflexive-Lp}
\index{Lp space@$L^p$ space!reflexive}
\index{reflexive!$L^p$ spaces}
Since $(L^p)' \cong L^q$ and $(L^q)' \cong L^p$ for $1 < p < \infty$,
the spaces $L^p$ are \emph{reflexive} for $1 < p < \infty$. The spaces
$L^1$ and $L^\infty$ are not reflexive (except in trivial cases).
\end{remark}

\section{$L^2$ as a Hilbert Space}\label{sec:L2-hilbert}

\index{L2 space@$L^2$ space!Hilbert space}

\begin{proposition}\label{prop:L2-hilbert}
$L^2(\Omega, \mu)$ is a Hilbert space with inner product
$\inner{f}{g} = \int_\Omega f\, \bar{g}\, \dd\mu$.
The Riesz representation theorem\index{Riesz representation theorem!Hilbert space}
for Hilbert spaces gives $(L^2)' \cong L^2$ directly, consistent with
Theorem~\ref{thm:Lp-dual} at $p = q = 2$.
\end{proposition}

\section{The Radon--Nikodym Theorem}\label{sec:radon-nikodym}

\index{Radon--Nikodym theorem|(}

We state the theorem used in the proof of Theorem~\ref{thm:Lp-dual} and give a
proof using Hilbert space methods (following von Neumann).

\begin{theorem}[Radon--Nikodym]\label{thm:radon-nikodym}
\index{Radon--Nikodym theorem!statement}
Let $\mu$ and $\nu$ be $\sigma$-finite measures on $(\Omega, \mathcal{A})$ with
$\nu \ll \mu$ (i.e., $\mu(A) = 0 \Rightarrow \nu(A) = 0$). Then there exists
a measurable function $g \ge 0$ (unique $\mu$-a.e.) such that
$\dd\nu = g\, \dd\mu$, i.e.,
$\nu(A) = \int_A g\, \dd\mu$ for all $A \in \mathcal{A}$.
The function $g = \frac{\dd\nu}{\dd\mu}$ is the \emph{Radon--Nikodym derivative}.
\index{Radon--Nikodym derivative}
\end{theorem}

\begin{proof}[Proof (von Neumann)]
\index{Radon--Nikodym theorem!proof}
Set $\lambda = \mu + \nu$. Then $L^2(\lambda) \to \R$,
$f \mapsto \int f\, \dd\nu$, is a bounded linear functional (since
$\abs{\int f\, \dd\nu} \le \int \abs{f}\, \dd\nu \le \int \abs{f}\, \dd\lambda$
and by Cauchy--Schwarz $\le \lambda(\Omega)^{1/2}\norm{f}_{L^2(\lambda)}$ when
$\lambda(\Omega) < \infty$; the $\sigma$-finite case requires a standard exhaustion
argument).

By the Riesz representation theorem for $L^2(\lambda)$, there exists
$h \in L^2(\lambda)$ with $\int f\, \dd\nu = \int fh\, \dd\lambda$ for all
$f \in L^2(\lambda)$. Taking $f = \mathbf{1}_A$:
$\nu(A) = \int_A h\, \dd\lambda$. Since $0 \le \nu(A) \le \lambda(A)$, we
get $0 \le h \le 1$ $\lambda$-a.e.

Now $\int f\, \dd\nu = \int fh\, \dd\mu + \int fh\, \dd\nu$, so
$\int f(1 - h)\, \dd\nu = \int fh\, \dd\mu$. On $\{h = 1\}$, this gives
$\mu(\{h = 1\}) = 0$ (take $f = \mathbf{1}_{\{h=1\}}$; the left side is $0$). Set
$g = h/(1-h)$ on $\{h < 1\}$ and $g = 0$ on $\{h = 1\}$. Then for any measurable
$A$, taking $f = \mathbf{1}_A/(1-h)$ (which is integrable) yields
$\nu(A) = \int_A g\, \dd\mu$.
\end{proof}

\index{Radon--Nikodym theorem|)}

\section{Interpolation: The Riesz--Thorin Theorem}\label{sec:riesz-thorin}

\index{Riesz--Thorin theorem|(}
\index{interpolation|(}

\begin{theorem}[Riesz--Thorin]\label{thm:riesz-thorin}
\index{Riesz--Thorin theorem!statement}
Let $(\Omega_1, \mu_1)$ and $(\Omega_2, \mu_2)$ be $\sigma$-finite measure spaces.
Let $1 \le p_0, p_1, q_0, q_1 \le \infty$ and let $T$ be a linear operator
defined on $L^{p_0}(\mu_1) + L^{p_1}(\mu_1)$ satisfying
\[
  \norm{Tf}_{q_0} \le M_0\norm{f}_{p_0}, \qquad
  \norm{Tf}_{q_1} \le M_1\norm{f}_{p_1}.
\]
Then for every $0 < \theta < 1$, with
\[
  \frac{1}{p} = \frac{1-\theta}{p_0} + \frac{\theta}{p_1}, \qquad
  \frac{1}{q} = \frac{1-\theta}{q_0} + \frac{\theta}{q_1},
\]
we have $\norm{Tf}_q \le M_0^{1-\theta} M_1^\theta\, \norm{f}_p$ for all
$f \in L^p(\mu_1)$.
\end{theorem}

\begin{proof}
\index{Riesz--Thorin theorem!proof}
We use the Hadamard three-lines lemma\index{Hadamard three-lines lemma}. Consider
the strip $S = \{z \in \C : 0 \le \operatorname{Re} z \le 1\}$.

\textbf{Step 1 (Reduction).} By duality, it suffices to show that for simple
functions $f = \sum_j a_j \mathbf{1}_{A_j}$ and $g = \sum_k b_k \mathbf{1}_{B_k}$
with $\norm{f}_p = \norm{g}_{q'} = 1$:
\[
  \abs[\bigg]{\int (Tf)\, \bar{g}\, \dd\mu_2} \le M_0^{1-\theta} M_1^\theta.
\]

\textbf{Step 2 (Analytic family).} Define, for $z \in S$:
\[
  f_z(x) = \abs{f(x)}^{p(\frac{1-z}{p_0} + \frac{z}{p_1})}
  \frac{f(x)}{\abs{f(x)}}, \qquad
  g_z(y) = \abs{g(y)}^{q'(\frac{1-z}{q_0'} + \frac{z}{q_1'})}
  \frac{g(y)}{\abs{g(y)}}.
\]
(Set $f_z = 0$ where $f = 0$, similarly for $g_z$.) At $z = \theta$:
$f_\theta = f$, $g_\theta = g$.

Define $F(z) = \int (Tf_z)\, \overline{g_z}\, \dd\mu_2$. Then $F$ is continuous
on $S$, analytic on the interior (since $f$ and $g$ are simple functions), and
bounded on $S$.

\textbf{Step 3 (Boundary estimates).} On $\operatorname{Re} z = 0$:
$\norm{f_{it}}_{p_0} = 1$ and $\norm{g_{it}}_{q_0'} = 1$, so
$\abs{F(it)} \le \norm{Tf_{it}}_{q_0}\norm{g_{it}}_{q_0'} \le M_0$.
Similarly, $\abs{F(1 + it)} \le M_1$ on $\operatorname{Re} z = 1$.

\textbf{Step 4 (Three-lines lemma).} By the Hadamard three-lines lemma (if
$F$ is bounded and continuous on $S$, analytic inside, with $\abs{F} \le M_j$
on the line $\operatorname{Re} z = j$ for $j = 0, 1$, then
$\abs{F(z)} \le M_0^{1 - \operatorname{Re} z} M_1^{\operatorname{Re} z}$):
\[
  \abs[\bigg]{\int (Tf)\bar{g}\, \dd\mu_2}
  = \abs{F(\theta)} \le M_0^{1-\theta} M_1^\theta. \qedhere
\]
\end{proof}

\index{Riesz--Thorin theorem|)}
\index{interpolation|)}

\section{Application: Fourier Transform on $L^1$ and $L^2$}
\label{sec:fourier-L1-L2}

\index{Fourier transform|(}

\begin{definition}\label{def:fourier-L1}
\index{Fourier transform!on $L^1$}
For $f \in L^1(\R^d)$, the \emph{Fourier transform} is
\[
  \hat{f}(\xi) = \int_{\R^d} f(x)\, e^{-2\pi i x \cdot \xi}\, \dd x, \qquad \xi \in \R^d.
\]
\end{definition}

\begin{proposition}\label{prop:fourier-L1}
The Fourier transform $\mathcal{F} \colon L^1(\R^d) \to C_0(\R^d)$ is a bounded
linear map with $\norm{\hat{f}}_\infty \le \norm{f}_1$.
\end{proposition}

\begin{theorem}[Plancherel]\label{thm:plancherel}
\index{Plancherel theorem}
\index{Fourier transform!on $L^2$}
The Fourier transform extends uniquely to a unitary operator
$\mathcal{F} \colon L^2(\R^d) \to L^2(\R^d)$:
\[
  \norm{\hat{f}}_{L^2} = \norm{f}_{L^2} \qquad \text{for all } f \in L^2(\R^d).
\]
\end{theorem}

\begin{proof}
For $f \in L^1 \cap L^2$ (which is dense in $L^2$), one verifies
$\norm{\hat{f}}_2 = \norm{f}_2$ by computing
$\inner{\hat{f}}{\hat{g}}_2 = \inner{f}{g}_2$ for Schwartz functions and
extending by density.

Alternatively, by the Riesz--Thorin theorem: $\mathcal{F}$ maps $L^1 \to L^\infty$
with norm $1$ and (once established on a dense subspace) $L^2 \to L^2$ with norm $1$.
By interpolation at $\theta = 1/2$ (giving $p = 4/3$, $q = 4$), one obtains the
Hausdorff--Young inequality\index{Hausdorff--Young inequality}
$\norm{\hat{f}}_q \le \norm{f}_p$ for $1 \le p \le 2$, $q = p'$.
\end{proof}

\index{Fourier transform|)}

\section{Exercises}

\begin{exercise}[$\star$]\label{exo:holder-equality}
Characterize equality in H\"older's inequality.
\end{exercise}

\begin{exercise}[$\star$]\label{exo:Lp-inclusions}
\index{Lp space@$L^p$ space!inclusions}
Let $\mu(\Omega) < \infty$. Show that $L^q(\Omega) \subset L^p(\Omega)$ for
$1 \le p \le q \le \infty$ with $\norm{f}_p \le \mu(\Omega)^{1/p - 1/q}\norm{f}_q$.
Give a counterexample when $\mu(\Omega) = \infty$.
\end{exercise}

\begin{exercise}[$\star\star$]\label{exo:Lp-separable}
\index{Lp space@$L^p$ space!separability}
Show that $L^p(\R^d)$ is separable for $1 \le p < \infty$ but $L^\infty(\R^d)$
is not.
\end{exercise}

\begin{exercise}[$\star\star$]\label{exo:weak-Lp}
\index{weak convergence!in $L^p$}
Let $1 < p < \infty$ and $f_n \weakto f$ in $L^p$. Show $\norm{f}_p \le \liminf \norm{f_n}_p$.
If additionally $\norm{f_n}_p \to \norm{f}_p$, show $f_n \to f$ in $L^p$.
\end{exercise}

\begin{exercise}[$\star\star$]\label{exo:dual-ell-1}
\index{dual space!of $\ell^1$}
Prove directly that $(\ell^1)' \cong \ell^\infty$.
\end{exercise}

\begin{exercise}[$\star\star\star$]\label{exo:dunford-pettis}
\index{Dunford--Pettis theorem}
(Dunford--Pettis.) A subset $K \subset L^1(\mu)$ is weakly sequentially compact
if and only if it is bounded and uniformly integrable. Prove this.
\end{exercise}

\begin{exercise}[$\star$]\label{exo:convolution-Lp}
Show $f * g * h = f * (g * h)$ for appropriate $f$, $g$, $h$ in $L^p$ spaces.
\end{exercise}

\begin{exercise}[$\star\star$]\label{exo:fourier-schwartz}
\index{Fourier transform!on Schwartz space}
Show the Fourier transform is a bijection on the Schwartz space $\mathcal{S}(\R^d)$
and satisfies $\widehat{D^\alpha f}(\xi) = (2\pi i \xi)^\alpha \hat{f}(\xi)$.
\end{exercise}

\begin{exercise}[$\star\star\star$]\label{exo:marcinkiewicz}
\index{Marcinkiewicz interpolation}
(Marcinkiewicz interpolation.) Prove the weak-type version of Riesz--Thorin:
if $T$ is weak-$(p_0, q_0)$ and weak-$(p_1, q_1)$, then $T$ is strong-$(p, q)$
for $p_0 < p < p_1$.
\end{exercise}


%%%%%%%%%%%%%%%%%%%%%%%%%%%%%%%%%%%%%%%%%%%%%%%%%%%%%%%%%%%%
%%  CHAPTER 12: Banach Algebras and C*-Algebras
%%%%%%%%%%%%%%%%%%%%%%%%%%%%%%%%%%%%%%%%%%%%%%%%%%%%%%%%%%%%
\chapter{Banach Algebras and $C^*$-Algebras --- Introduction}
\label{ch:banach-cstar}
\minitoc

\section{Banach Algebras}\label{sec:banach-alg-def}

\index{Banach algebra|(}

\begin{definition}[Banach algebra]\label{def:banach-algebra}
\index{Banach algebra!definition}
A \emph{Banach algebra} is a Banach space $(A, \norm{\cdot})$ equipped with a
bilinear multiplication $A \times A \to A$ satisfying:
\begin{enumerate}[label=(\roman*)]
  \item $\norm{ab} \le \norm{a}\norm{b}$ for all $a, b \in A$ (submultiplicativity).
  \item Associativity: $(ab)c = a(bc)$.
\end{enumerate}
If there exists a unit $e \in A$ with $ea = ae = a$ for all $a$ and
$\norm{e} = 1$, we say $A$ is \emph{unital}\index{Banach algebra!unital}.
\end{definition}

\begin{example}\label{ex:banach-algebras}
\index{Banach algebra!examples}
\begin{enumerate}[label=(\alph*)]
  \item $\mathcal{L}(E)$ with operator composition, for any Banach space $E$.
  \item $C(K)$ with pointwise multiplication, for a compact Hausdorff space $K$.
        This is a commutative unital Banach algebra.\index{C(K)@$C(K)$!Banach algebra}
  \item $\ell^1(\Z)$ with convolution: $(a * b)_n = \sum_{k \in \Z} a_k b_{n-k}$.
        \index{ell1(Z)@$\ell^1(\Z)$!Banach algebra}
        The unit is $\delta_0$. This is a commutative unital Banach algebra.
  \item $L^1(\R)$ with convolution. This is a commutative Banach algebra
        \emph{without} unit.\index{L1(R)@$L^1(\R)$!Banach algebra}
\end{enumerate}
\end{example}

\section{Spectrum and Spectral Radius}\label{sec:spectrum-banach-alg}

\index{spectrum!in Banach algebra|(}
\index{spectral radius|(}

\begin{definition}\label{def:spectrum-algebra}
\index{spectrum!in Banach algebra!definition}
Let $A$ be a unital Banach algebra and $a \in A$. The \emph{spectrum} of $a$ is
\[
  \Spec(a) = \{\lambda \in \C : \lambda e - a \text{ is not invertible in } A\}.
\]
The \emph{spectral radius} is $r(a) = \sup\{\abs{\lambda} : \lambda \in \Spec(a)\}$.
\end{definition}

\begin{proposition}\label{prop:neumann-series}
\index{Neumann series}
If $\norm{a} < 1$ in a unital Banach algebra, then $e - a$ is invertible with
$(e - a)^{-1} = \sum_{n=0}^\infty a^n$ and
$\norm{(e-a)^{-1}} \le (1 - \norm{a})^{-1}$.
\end{proposition}

\begin{proof}
The series converges absolutely since $\sum \norm{a^n} \le \sum \norm{a}^n < \infty$.
Then $(e - a)\sum_{n=0}^N a^n = e - a^{N+1} \to e$.
\end{proof}

\begin{theorem}[Spectral radius formula]\label{thm:spectral-radius}
\index{spectral radius!formula}
For $a$ in a unital Banach algebra:
\[
  r(a) = \lim_{n \to \infty} \norm{a^n}^{1/n} = \inf_{n \ge 1} \norm{a^n}^{1/n}.
\]
Moreover, $\Spec(a)$ is a nonempty compact subset of the disk $\{\abs{\lambda} \le \norm{a}\}$.
\end{theorem}

\begin{proof}
\index{spectral radius!formula!proof}
\textbf{Step 1 ($\Spec(a)$ is compact and nonempty).}
If $\abs{\lambda} > \norm{a}$, then $\lambda e - a = \lambda(e - a/\lambda)$ is
invertible by Proposition~\ref{prop:neumann-series}, so $\lambda \notin \Spec(a)$.
Hence $\Spec(a) \subset \overline{B}(0, \norm{a})$.

The map $\lambda \mapsto \lambda e - a$ is continuous $\C \to A$, and the set
of invertible elements $A^\times$ is open (if $b$ is invertible and
$\norm{c - b} < \norm{b^{-1}}^{-1}$, then $c = b(e - b^{-1}(b - c))$ is
invertible). Hence $\C \setminus \Spec(a)$ is open, so $\Spec(a)$ is closed,
thus compact.

\emph{Nonemptiness:} for any $\varphi \in A'$, the function
$f(\lambda) = \varphi((\lambda e - a)^{-1})$ is analytic on $\C \setminus \Spec(a)$.
If $\Spec(a) = \emptyset$, $f$ is entire. For $\abs{\lambda} > \norm{a}$:
$\abs{f(\lambda)} \le \norm{\varphi}\, \norm{(\lambda e - a)^{-1}}
\le \norm{\varphi}/(\abs{\lambda} - \norm{a}) \to 0$ as $\abs{\lambda} \to \infty$.
By Liouville's theorem\index{Liouville's theorem}, $f \equiv 0$. Since this holds
for all $\varphi \in A'$, Hahn--Banach gives $(\lambda e - a)^{-1} = 0$, a
contradiction. Hence $\Spec(a) \ne \emptyset$.

\textbf{Step 2 ($r(a) = \lim \norm{a^n}^{1/n}$).}
If $\lambda \in \Spec(a)$, then $\lambda^n \in \Spec(a^n)$ (since $\lambda^n e - a^n$
factors), so $\abs{\lambda}^n \le \norm{a^n}$, giving
$r(a) \le \inf_n \norm{a^n}^{1/n}$.

For the reverse, let $R = \limsup_n \norm{a^n}^{1/n}$. The resolvent
$R(\lambda) = (\lambda e - a)^{-1} = \sum_{n=0}^\infty a^n/\lambda^{n+1}$
converges for $\abs{\lambda} > R$ (by the root test). Since this Laurent series
converges on $\C \setminus \Spec(a)$ and diverges at points of $\Spec(a)$,
we have $r(a) \ge R$.

To see $R = \lim$ (not just $\limsup$): let $\alpha = \inf_n \norm{a^n}^{1/n}$.
For any $n$, writing $m = nq + r$ with $0 \le r < n$:
$\norm{a^m}^{1/m} \le \norm{a^n}^{q/m}\norm{a}^{r/m}$. Taking $\limsup$ as
$m \to \infty$ (so $q/m \to 1/n$, $r/m \to 0$):
$R \le \norm{a^n}^{1/n}$ for all $n$, hence $R \le \alpha$. Combined with
$\alpha \le R$, we get $\alpha = R = \lim_n \norm{a^n}^{1/n}$.
\end{proof}

\index{spectrum!in Banach algebra|)}
\index{spectral radius|)}

\section{Ideals, Quotients, and Morphisms}\label{sec:ideals-quotients}

\index{Banach algebra!ideal|(}
\index{Banach algebra!quotient|(}

\begin{definition}\label{def:ideal-algebra}
\index{ideal!in Banach algebra}
A \emph{(two-sided) ideal} in a Banach algebra $A$ is a closed subspace $I$ such
that $aI \subset I$ and $Ia \subset I$ for all $a \in A$. An ideal is
\emph{maximal}\index{maximal ideal} if it is proper ($I \ne A$) and not contained
in any larger proper ideal.
\end{definition}

\begin{proposition}\label{prop:quotient-algebra}
\index{Banach algebra!quotient!construction}
If $I$ is a closed ideal in $A$, then $A/I$ is a Banach algebra with
$\norm{a + I} = \inf_{x \in I} \norm{a + x}$, the quotient norm.
If $A$ is unital, so is $A/I$ (with unit $e + I$).
\end{proposition}

\begin{definition}\label{def:algebra-morphism}
\index{Banach algebra!morphism}
A \emph{morphism} of Banach algebras is a bounded linear map
$\phi \colon A \to B$ satisfying $\phi(ab) = \phi(a)\phi(b)$. If $A$ and $B$ are
unital and $\phi(e_A) = e_B$, we say $\phi$ is \emph{unital}.
\end{definition}

\index{Banach algebra!ideal|)}
\index{Banach algebra!quotient|)}

\section{Commutative Banach Algebras: Characters and Gelfand Space}
\label{sec:gelfand-space}

\index{Gelfand theory|(}
\index{character|(}

\begin{definition}\label{def:character}
\index{character!definition}
Let $A$ be a commutative unital Banach algebra. A \emph{character} (or
\emph{multiplicative linear functional}\index{multiplicative linear functional})
is a nonzero algebra homomorphism $\chi \colon A \to \C$, i.e., $\chi(ab) = \chi(a)\chi(b)$
and $\chi \ne 0$.
\end{definition}

\begin{proposition}\label{prop:character-properties}
\index{character!properties}
Let $\chi$ be a character of a commutative unital Banach algebra $A$.
\begin{enumerate}[label=(\roman*)]
  \item $\chi$ is continuous with $\norm{\chi} = 1$.
  \item $\chi(e) = 1$.
  \item $\Ker\chi$ is a maximal ideal.
  \item $\chi(a) \in \Spec(a)$ for all $a \in A$.
\end{enumerate}
\end{proposition}

\begin{proof}
(ii) $\chi(e) = \chi(e \cdot e) = \chi(e)^2$ and $\chi \ne 0$, so $\chi(e) = 1$.

(iv) $\chi(a) e - a \in \Ker\chi$ which is a proper ideal, hence
$\chi(a)e - a$ is not invertible.

(i) From (iv), $\abs{\chi(a)} \le r(a) \le \norm{a}$, so $\norm{\chi} \le 1$.
Since $\chi(e) = 1$, $\norm{\chi} = 1$.

(iii) $\Ker\chi$ has codimension $1$ (since $\chi$ is surjective onto $\C$) and
is an ideal, hence maximal.
\end{proof}

\begin{definition}[Gelfand space]\label{def:gelfand-space}
\index{Gelfand space}
The \emph{Gelfand space} (or \emph{maximal ideal space}, or \emph{character space})
of $A$ is the set $\hat{A}$ of all characters of $A$, equipped with the weak-$*$
topology\index{weak-* topology!on Gelfand space} inherited from $A'$.
\end{definition}

\begin{proposition}\label{prop:gelfand-space-compact}
\index{Gelfand space!compactness}
$\hat{A}$ is a compact Hausdorff space (in the weak-$*$ topology).
\end{proposition}

\begin{proof}
By Proposition~\ref{prop:character-properties}, $\hat{A} \subset \{\varphi \in A' :
\norm{\varphi} \le 1\}$, which is weak-$*$ compact by Banach--Alaoglu\index{Banach--Alaoglu theorem}.
It suffices to show $\hat{A} \cup \{0\}$ is weak-$*$ closed. If
$\chi_\alpha \to \varphi$ weak-$*$ with $\chi_\alpha \in \hat{A}$, then
$\varphi(ab) = \lim \chi_\alpha(ab) = \lim \chi_\alpha(a)\chi_\alpha(b) = \varphi(a)\varphi(b)$.
So $\varphi$ is multiplicative, hence $\varphi \in \hat{A} \cup \{0\}$. Since
$\chi(e) = 1$ for all $\chi \in \hat{A}$ and $\varphi(e) = \lim \chi_\alpha(e) = 1$,
$\varphi \ne 0$, so actually $\hat{A}$ is itself closed, hence compact.
\end{proof}

\section{The Gelfand Transform}\label{sec:gelfand-transform}

\index{Gelfand transform|(}

\begin{definition}\label{def:gelfand-transform}
\index{Gelfand transform!definition}
The \emph{Gelfand transform} is the map
$\Gamma \colon A \to C(\hat{A})$ defined by
\[
  \hat{a}(\chi) = \chi(a), \qquad a \in A,\ \chi \in \hat{A}.
\]
\end{definition}

\begin{theorem}\label{thm:gelfand-transform}
\index{Gelfand transform!properties}
The Gelfand transform $\Gamma \colon a \mapsto \hat{a}$ is a unital algebra
homomorphism from $A$ to $C(\hat{A})$ with:
\begin{enumerate}[label=(\roman*)]
  \item $\widehat{ab} = \hat{a}\hat{b}$ and $\widehat{\alpha a + \beta b} = \alpha\hat{a} + \beta\hat{b}$.
  \item $\norm{\hat{a}}_\infty = r(a)$ for all $a \in A$.
  \item $\hat{a}(\hat{A}) = \Spec(a)$ for all $a$.
  \item $\Ker\Gamma = \{a \in A : r(a) = 0\}$ (the \emph{radical}\index{radical} of $A$).
\end{enumerate}
\end{theorem}

\begin{proof}
(i) is immediate: $\widehat{ab}(\chi) = \chi(ab) = \chi(a)\chi(b) = \hat{a}(\chi)\hat{b}(\chi)$.

(iii) $\lambda \in \Spec(a)$ iff $\lambda e - a$ is not invertible. In a commutative
unital Banach algebra, a noninvertible element is contained in a maximal ideal,
and every maximal ideal is $\Ker\chi$ for some $\chi \in \hat{A}$ (by Zorn's lemma
and the correspondence between maximal ideals and characters). Hence
$\lambda \in \Spec(a)$ iff $\chi(\lambda e - a) = 0$ for some $\chi$, iff
$\lambda = \chi(a)$ for some $\chi$.

(ii) $\norm{\hat{a}}_\infty = \sup_\chi \abs{\chi(a)} = \sup_{\lambda \in \Spec(a)} \abs{\lambda} = r(a)$.

(iv) $\hat{a} = 0$ iff $\norm{\hat{a}}_\infty = 0$ iff $r(a) = 0$.
\end{proof}

\index{Gelfand transform|)}
\index{character|)}
\index{Gelfand theory|)}

\section{$C^*$-Algebras}\label{sec:cstar-def}

\index{C*-algebra@$C^*$-algebra|(}

\begin{definition}[$C^*$-algebra]\label{def:cstar-algebra}
\index{C*-algebra@$C^*$-algebra!definition}
A \emph{$C^*$-algebra} is a Banach algebra $A$ equipped with an
\emph{involution}\index{involution} $* \colon A \to A$ (an antilinear map with
$(a^*)^* = a$ and $(ab)^* = b^*a^*$) satisfying the \emph{$C^*$-identity}:
\[
  \norm{a^*a} = \norm{a}^2 \qquad \text{for all } a \in A.
\]
\end{definition}

\begin{example}\label{ex:cstar-examples}
\index{C*-algebra@$C^*$-algebra!examples}
\begin{enumerate}[label=(\alph*)]
  \item $\mathcal{L}(H)$ with $T^* = $ adjoint, for a Hilbert space $H$.
  \item $C(K)$ with $f^* = \bar{f}$, for compact Hausdorff $K$.
  \item Any norm-closed $*$-subalgebra of $\mathcal{L}(H)$.
\end{enumerate}
\end{example}

\begin{proposition}\label{prop:cstar-spectral-radius}
\index{C*-algebra@$C^*$-algebra!spectral radius}
In a $C^*$-algebra: $\norm{a} = r(a)$ for every normal element ($a^*a = aa^*$).
In particular, $\norm{a} = r(a)$ for self-adjoint $a$.
\end{proposition}

\begin{proof}
If $a$ is normal, $\norm{a^2}^2 = \norm{(a^2)^*(a^2)} = \norm{(a^*a)^2}
= \norm{a^*a}^2 = \norm{a}^4$, so $\norm{a^2} = \norm{a}^2$.
By induction, $\norm{a^{2^n}} = \norm{a}^{2^n}$, hence
$r(a) = \lim \norm{a^n}^{1/n} = \norm{a}$.
\end{proof}

\section{The Commutative Gelfand--Naimark Theorem}\label{sec:gelfand-naimark}

\index{Gelfand--Naimark theorem!commutative|(}

\begin{theorem}[Commutative Gelfand--Naimark]\label{thm:gelfand-naimark}
\index{Gelfand--Naimark theorem!commutative!statement}
Let $A$ be a commutative unital $C^*$-algebra. The Gelfand transform
$\Gamma \colon A \to C(\hat{A})$ is an isometric $*$-isomorphism.
\end{theorem}

\begin{proof}
\index{Gelfand--Naimark theorem!commutative!proof}
\textbf{Step 1 ($\Gamma$ is a $*$-homomorphism).}
We know $\Gamma$ is an algebra homomorphism. We must show
$\widehat{a^*}(\chi) = \overline{\hat{a}(\chi)}$, i.e., $\chi(a^*) = \overline{\chi(a)}$.

Every element of a $C^*$-algebra can be written $a = h + ik$ with
$h = (a + a^*)/2$ and $k = (a - a^*)/(2i)$ self-adjoint. It suffices to show
$\chi(h) \in \R$ for self-adjoint $h$.

For self-adjoint $h$ and $t \in \R$, set $u_t = e^{ith}$ (defined by the power
series, which converges). One checks $u_t^* = e^{-ith} = u_t^{-1}$, so $u_t$
is unitary: $\norm{u_t}^2 = \norm{u_t^* u_t} = \norm{e} = 1$.
Hence $\abs{\chi(u_t)} \le \norm{u_t} = 1$ and
$\abs{\chi(u_t)^{-1}} = \abs{\chi(u_{-t})} \le 1$, so $\abs{\chi(u_t)} = 1$.

Now $\chi(u_t) = e^{it\chi(h)}$ (since $\chi$ is an algebra homomorphism
preserving the exponential series). Hence $\abs{e^{it\chi(h)}} = 1$ for all
$t \in \R$, which forces $\chi(h) \in \R$.

\textbf{Step 2 ($\Gamma$ is isometric).}
Since $A$ is commutative, every element is normal (as $a^*a = aa^*$ when
multiplication commutes). By Proposition~\ref{prop:cstar-spectral-radius},
$\norm{a} = r(a) = \norm{\hat{a}}_\infty$.

\textbf{Step 3 ($\Gamma$ is surjective).}
By the isometry, $\Gamma(A)$ is a closed $*$-subalgebra of $C(\hat{A})$
containing the constants (since $\hat{e} = 1$) and separating points of
$\hat{A}$ (if $\chi_1 \ne \chi_2$, there exists $a$ with
$\chi_1(a) \ne \chi_2(a)$, i.e., $\hat{a}(\chi_1) \ne \hat{a}(\chi_2)$).
Since $\Gamma$ is a $*$-map, $\Gamma(A)$ is closed under conjugation.
By the Stone--Weierstrass theorem\index{Stone--Weierstrass theorem},
$\Gamma(A) = C(\hat{A})$.
\end{proof}

\index{Gelfand--Naimark theorem!commutative|)}

\section{Continuous Functional Calculus in $C^*$-Algebras}
\label{sec:func-calc-cstar}

\index{functional calculus!in $C^*$-algebras|(}

\begin{theorem}\label{thm:func-calc-cstar}
\index{functional calculus!in $C^*$-algebras!theorem}
Let $A$ be a unital $C^*$-algebra and $a \in A$ normal. There exists a unique
isometric $*$-isomorphism $\Phi_a \colon C(\Spec(a)) \to C^*(a, e)$ (the
$C^*$-subalgebra generated by $a$ and $e$) with $\Phi_a(\mathrm{id}) = a$ and
$\Phi_a(1) = e$.
\end{theorem}

\begin{proof}
$C^*(a, e)$ is a commutative $C^*$-algebra (since $a$ is normal). By the
Gelfand--Naimark theorem, $C^*(a, e) \cong C(\widehat{C^*(a,e)})$.
The evaluation map $\chi \mapsto \chi(a)$ is a homeomorphism
$\widehat{C^*(a,e)} \to \Spec(a)$ (it is continuous and injective---since
$\chi$ is determined by $\chi(a)$---and surjective by
Theorem~\ref{thm:gelfand-transform}(iii)). The composition gives the desired
isomorphism $C(\Spec(a)) \cong C^*(a, e)$.
\end{proof}

\index{functional calculus!in $C^*$-algebras|)}

\section{States and Positive Forms}\label{sec:states}

\index{state!on $C^*$-algebra|(}
\index{positive linear functional|(}

\begin{definition}\label{def:positive-form}
\index{positive linear functional!definition}
A linear functional $\omega \colon A \to \C$ on a $C^*$-algebra is
\emph{positive} if $\omega(a^*a) \ge 0$ for all $a \in A$.
\end{definition}

\begin{definition}\label{def:state}
\index{state!definition}
A \emph{state} on a unital $C^*$-algebra $A$ is a positive linear functional
$\omega$ with $\omega(e) = 1$. The set of all states is the \emph{state
space}\index{state space} $S(A)$.
\end{definition}

\begin{proposition}\label{prop:state-properties}
\index{state!properties}
Let $\omega$ be a state on a unital $C^*$-algebra $A$.
\begin{enumerate}[label=(\roman*)]
  \item $\omega(a^*) = \overline{\omega(a)}$.
  \item $\abs{\omega(a)}^2 \le \omega(a^*a)$ (Cauchy--Schwarz).
  \item $\norm{\omega} = \omega(e) = 1$.
  \item $S(A)$ is convex and weak-$*$ compact.
\end{enumerate}
\end{proposition}

\begin{proof}
(i)--(ii): The map $(a,b) \mapsto \omega(b^*a)$ is a positive semidefinite
sesquilinear form, so (i) and (ii) follow from standard properties.

(iii): $\abs{\omega(a)}^2 \le \omega(a^*a) \le \norm{\omega}\norm{a^*a}
= \norm{\omega}\norm{a}^2$, hence $\norm{\omega} \le$ itself (circular unless
we note $\omega(e) = 1$ gives $\norm{\omega} \ge 1$, and the Cauchy--Schwarz
bound gives $\abs{\omega(a)} \le \omega(e)^{1/2}\omega(a^*a)^{1/2}
\le \omega(e)^{1/2}\norm{a^*a}^{1/2}\norm{\omega}^{1/2}$... more cleanly:
$\abs{\omega(a)} \le \norm{a}$ follows from $\abs{\omega(a)}^2 \le \omega(a^*a)
\le \norm{a^*a} = \norm{a}^2$ once we know $\omega$ is bounded, which follows
from positivity on a $C^*$-algebra). Thus $\norm{\omega} = 1$.

(iv): Convexity: if $\omega_1, \omega_2$ are states and $0 \le t \le 1$, then
$t\omega_1 + (1-t)\omega_2$ is positive with value $1$ at $e$. Weak-$*$
compactness: $S(A) \subset \{\varphi \in A' : \norm{\varphi} \le 1\}$ is weak-$*$
closed (being defined by the conditions $\omega(e) = 1$ and $\omega(a^*a) \ge 0$),
hence compact by Banach--Alaoglu.
\end{proof}

\index{state!on $C^*$-algebra|)}
\index{positive linear functional|)}

\section{The GNS Construction}\label{sec:gns}

\index{GNS construction|(}
\index{Gelfand--Naimark--Segal construction|see{GNS construction}}

The Gelfand--Naimark--Segal (GNS) construction produces a Hilbert space
representation from a state.

\begin{theorem}[GNS construction]\label{thm:gns}
\index{GNS construction!theorem}
Let $A$ be a unital $C^*$-algebra and $\omega$ a state on $A$. There exist a
Hilbert space $H_\omega$, a unit vector $\xi_\omega \in H_\omega$, and a
$*$-representation $\pi_\omega \colon A \to \mathcal{L}(H_\omega)$ such that:
\begin{enumerate}[label=(\roman*)]
  \item $\omega(a) = \inner{\pi_\omega(a)\xi_\omega}{\xi_\omega}$ for all $a \in A$.
  \item $\xi_\omega$ is cyclic: $\overline{\pi_\omega(A)\xi_\omega} = H_\omega$.
        \index{cyclic vector}
\end{enumerate}
The triple $(H_\omega, \pi_\omega, \xi_\omega)$ is unique up to unitary equivalence.
\end{theorem}

\begin{proof}
\index{GNS construction!proof}
\textbf{Step 1 (Inner product on $A$).} Define the sesquilinear form on $A$:
$\inner{a}{b}_\omega = \omega(b^*a)$. This is positive semidefinite:
$\inner{a}{a}_\omega = \omega(a^*a) \ge 0$.

\textbf{Step 2 (Quotient by null space).} Let
$N_\omega = \{a \in A : \omega(a^*a) = 0\}$. By Cauchy--Schwarz
($\abs{\omega(b^*a)}^2 \le \omega(a^*a)\omega(b^*b)$), $N_\omega$ is a left
ideal of $A$ (if $a \in N_\omega$, then for any $b \in A$:
$\omega((ba)^*ba) = \omega(a^*b^*ba) \le \norm{b^*b}\omega(a^*a) = 0$
by positivity and $C^*$-identity reasoning). The quotient $A / N_\omega$
inherits an inner product:
$\inner{a + N_\omega}{b + N_\omega} = \omega(b^*a)$.

\textbf{Step 3 (Completion).} Let $H_\omega$ be the completion of $A / N_\omega$.

\textbf{Step 4 (Representation).} For $a \in A$, define
$\pi_\omega(a)(b + N_\omega) = ab + N_\omega$. This is well-defined since
$N_\omega$ is a left ideal. We check boundedness:
\begin{align*}
  \norm{\pi_\omega(a)(b + N_\omega)}^2
  &= \omega(b^*a^*ab) \le \norm{a^*a}\, \omega(b^*b)
  = \norm{a}^2\, \norm{b + N_\omega}^2,
\end{align*}
where the inequality uses $a^*a \le \norm{a}^2 e$ (in the order of self-adjoint
elements) and positivity of $\omega$. Hence $\norm{\pi_\omega(a)} \le \norm{a}$
and $\pi_\omega(a)$ extends to $H_\omega$.

One verifies $\pi_\omega(ab) = \pi_\omega(a)\pi_\omega(b)$ and
$\pi_\omega(a^*) = \pi_\omega(a)^*$ (from the definition of the inner product).

\textbf{Step 5 (Cyclic vector).} Set $\xi_\omega = e + N_\omega \in H_\omega$.
Then $\pi_\omega(a)\xi_\omega = a + N_\omega$, so
$\pi_\omega(A)\xi_\omega = A/N_\omega$, which is dense in $H_\omega$ by
construction.

Also, $\inner{\pi_\omega(a)\xi_\omega}{\xi_\omega} = \omega(e^*\cdot a \cdot e) = \omega(a)$.

\textbf{Uniqueness.} If $(H', \pi', \xi')$ is another GNS triple, define
$U \colon \pi_\omega(a)\xi_\omega \mapsto \pi'(a)\xi'$. This is isometric
(since both inner products equal $\omega(a^*a)$) and extends to a unitary
$U \colon H_\omega \to H'$ with $U\pi_\omega(a) = \pi'(a)U$ and $U\xi_\omega = \xi'$.
\end{proof}

\begin{corollary}[Gelfand--Naimark embedding]\label{cor:gn-embedding}
\index{Gelfand--Naimark theorem!general}
Every $C^*$-algebra admits an isometric $*$-representation on a Hilbert space.
\end{corollary}

\begin{proof}
For each $a \in A$ with $a \ne 0$, there exists a state $\omega$ with
$\omega(a^*a) = \norm{a}^2$ (by the Hahn--Banach theorem applied to the
$C^*$-subalgebra generated by $a^*a$). The direct sum
$\pi = \bigoplus_{\omega \in S(A)} \pi_\omega$ on
$H = \bigoplus_\omega H_\omega$ gives a faithful $*$-representation.
\end{proof}

\index{GNS construction|)}

\section{Application: Representation Theory and Quantum Mechanics}
\label{sec:cstar-qm}

\index{quantum mechanics!algebraic formulation|(}

\begin{remark}[Algebraic quantum mechanics]
\index{observable!algebra of}
In the algebraic approach to quantum mechanics:
\begin{itemize}
  \item Observables form a $C^*$-algebra $A$ (self-adjoint elements represent
        physical quantities).
  \item States are positive normalized linear functionals $\omega \in S(A)$.
  \item The GNS construction recovers the Hilbert space formulation: given a
        state $\omega$, the observables act on $H_\omega$ via $\pi_\omega$.
  \item Different states may give rise to unitarily inequivalent representations,
        which is related to the phenomenon of \emph{spontaneous symmetry breaking}
        \index{symmetry breaking} and \emph{superselection sectors}
        \index{superselection sectors} in quantum field theory.
\end{itemize}
\end{remark}

\index{quantum mechanics!algebraic formulation|)}

\section{Exercises}

\begin{exercise}[$\star$]\label{exo:spectrum-examples}
\index{spectrum!computation}
Compute $\Spec(a)$ for: (a) the identity $e$; (b) an idempotent $p = p^2$;
(c) a nilpotent element $a^n = 0$.
\end{exercise}

\begin{exercise}[$\star\star$]\label{exo:gelfand-ell1}
\index{Gelfand transform!of $\ell^1(\Z)$}
Identify the Gelfand space of $\ell^1(\Z)$ (with convolution product) with the
unit circle $\mathbb{T}$, and the Gelfand transform with the Fourier series.
\end{exercise}

\begin{exercise}[$\star$]\label{exo:cstar-norm-unique}
\index{C*-algebra@$C^*$-algebra!uniqueness of norm}
Show that a $C^*$-algebra has a unique $C^*$-norm: if $\norm{\cdot}_1$ and
$\norm{\cdot}_2$ are two norms making $A$ a $C^*$-algebra, then
$\norm{a}_1 = \norm{a}_2$ for all $a$.
\end{exercise}

\begin{exercise}[$\star\star$]\label{exo:cstar-positive}
\index{positive element!in $C^*$-algebra}
Let $A$ be a $C^*$-algebra and $a \in A$ self-adjoint. Show $a \ge 0$ if and
only if $\Spec(a) \subset [0, \infty)$.
\end{exercise}

\begin{exercise}[$\star\star$]\label{exo:pure-state}
\index{pure state}
A state $\omega$ is \emph{pure} if it is an extreme point of $S(A)$. Show that
$\omega$ is pure if and only if $\pi_\omega$ is irreducible.
\end{exercise}

\begin{exercise}[$\star\star\star$]\label{exo:gelfand-naimark-noncommutative}
\index{Gelfand--Naimark theorem!noncommutative}
Using the GNS construction and Corollary~\ref{cor:gn-embedding}, prove that
every $C^*$-algebra $A$ is isometrically $*$-isomorphic to a norm-closed
$*$-subalgebra of $\mathcal{L}(H)$ for some Hilbert space $H$.
\end{exercise}

\begin{exercise}[$\star\star$]\label{exo:wiener-lemma}
\index{Wiener's lemma}
(Wiener's lemma.) Let $f \in \ell^1(\Z)$ have absolutely convergent Fourier
series $\hat{f}(\theta) = \sum_n a_n e^{in\theta}$. Show that if $\hat{f}$
never vanishes on $\mathbb{T}$, then $1/\hat{f}$ also has absolutely convergent
Fourier series. \emph{Hint:} use the Gelfand theory of $\ell^1(\Z)$.
\end{exercise}

\begin{exercise}[$\star\star$]\label{exo:kadison}
\index{Kadison--Schwarz inequality}
(Kadison's inequality.) Let $\pi \colon A \to \mathcal{L}(H)$ be a unital
$*$-homomorphism. Show that $\pi(a)^*\pi(a) \le \pi(a^*a)$ for all $a$.
\end{exercise}

\begin{exercise}[$\star\star\star$]\label{exo:spectral-permanence}
\index{spectral permanence}
(Spectral permanence.) Let $B$ be a $C^*$-subalgebra of a $C^*$-algebra $A$
with the same unit. Show that $\Spec_B(b) = \Spec_A(b)$ for all $b \in B$.
\end{exercise}

\index{C*-algebra@$C^*$-algebra|)}
\index{Banach algebra|)}


%%%%%%%%%%%%%%%%%%%%%%%%%%%%%%%%%%%%%%%%%%%%%%%%%%%%%%%%%%%%
%%  APPENDIX
%%%%%%%%%%%%%%%%%%%%%%%%%%%%%%%%%%%%%%%%%%%%%%%%%%%%%%%%%%%%
\appendix
\chapter{Review of Analysis and Topology}\label{app:review}
\minitoc

\index{appendix|(}

This appendix collects fundamental results from analysis and topology that are
used throughout the text. Proofs are included for the reader's convenience.

\section{Baire's Category Theorem}\label{sec:baire}

\index{Baire category theorem|(}

\begin{theorem}[Baire category theorem]\label{thm:baire}
\index{Baire category theorem!statement}
Let $X$ be a complete metric space (or, more generally, a locally compact
Hausdorff space). If $(U_n)_{n \ge 1}$ is a sequence of dense open subsets of
$X$, then $\bigcap_{n=1}^\infty U_n$ is dense in $X$.

Equivalently, $X$ cannot be written as a countable union of closed sets with
empty interior: if $X = \bigcup_n F_n$ with $F_n$ closed, then at least one
$F_n$ has nonempty interior.
\end{theorem}

\begin{proof}
Let $B_0 \subset X$ be a nonempty open ball. Since $U_1$ is dense,
$B_0 \cap U_1 \ne \emptyset$; choose a closed ball
$\overline{B}_1 \subset B_0 \cap U_1$ with radius $r_1 < 1$. Since $U_2$
is dense, $B_1 \cap U_2 \ne \emptyset$; choose $\overline{B}_2 \subset B_1 \cap U_2$
with $r_2 < 1/2$. Continue inductively: $\overline{B}_{n+1} \subset B_n \cap U_{n+1}$,
$r_n < 1/n$.

The centers $(x_n)$ form a Cauchy sequence (since $x_m, x_n \in B_N$ for
$m, n \ge N$, with $\diam B_N < 2/N$). By completeness, $x_n \to x$.
Since $x \in \overline{B}_n$ for all $n$ (as the nested intersection of closed
sets), $x \in \overline{B}_n \subset U_n$ for all $n$. Hence
$x \in B_0 \cap \bigcap_n U_n$, proving density.
\end{proof}

\begin{remark}\label{rem:baire-applications}
The three pillars of functional analysis---the uniform boundedness principle,
the open mapping theorem, and the closed graph theorem---all rely on Baire's
theorem. In this text, it is invoked in Chapters~2, 3, and~4.
\end{remark}

\index{Baire category theorem|)}

\section{Zorn's Lemma}\label{sec:zorn}

\index{Zorn's lemma|(}

\begin{theorem}[Zorn's lemma]\label{thm:zorn}
\index{Zorn's lemma!statement}
Let $(P, \le)$ be a nonempty partially ordered set in which every totally
ordered subset (chain) has an upper bound. Then $P$ has at least one maximal
element.
\end{theorem}

Zorn's lemma is equivalent to the axiom of choice\index{axiom of choice} and
is used throughout functional analysis, notably in proving the Hahn--Banach
theorem, the existence of bases (Hamel bases), and the existence of maximal
ideals.

\index{Zorn's lemma|)}

\section{Tychonoff's Theorem}\label{sec:tychonoff}

\index{Tychonoff's theorem|(}

\begin{theorem}[Tychonoff]\label{thm:tychonoff}
\index{Tychonoff's theorem!statement}
An arbitrary product of compact topological spaces is compact (in the product
topology).
\end{theorem}

\begin{proof}[Proof sketch using Alexander's subbasis theorem]
\index{Alexander's subbasis theorem}
By Alexander's subbasis theorem, it suffices to show that every cover by
subbasis elements has a finite subcover. The subbasis for the product topology
$\prod_\alpha X_\alpha$ consists of sets $\pi_\alpha^{-1}(U_\alpha)$ where
$U_\alpha \subset X_\alpha$ is open. Suppose $\prod X_\alpha \subset
\bigcup_{i \in I} \pi_{\alpha_i}^{-1}(U_{\alpha_i})$ has no finite subcover.
For each index $\alpha$, consider the open sets $U_{\alpha_i}$ with
$\alpha_i = \alpha$. If these cover $X_\alpha$, a finite subcover of $X_\alpha$
(by compactness) gives a finite subcover of the product restricted to that
coordinate---one checks this yields a contradiction. Hence for each $\alpha$,
$X_\alpha$ is not covered, and one can choose $x_\alpha \notin \bigcup_i U_{\alpha_i}$
(for those $i$ with $\alpha_i = \alpha$). The point $(x_\alpha)_\alpha$ then
lies in no set of the cover, contradicting the assumption that it was a cover.
\end{proof}

\begin{remark}
Tychonoff's theorem is equivalent to the axiom of choice. It is used in the
proof of the Banach--Alaoglu theorem\index{Banach--Alaoglu theorem} (weak-$*$
compactness of the unit ball in a dual space).
\end{remark}

\index{Tychonoff's theorem|)}

\section{Lebesgue Measure and Integration}\label{sec:lebesgue-review}

\index{Lebesgue measure|(}
\index{Lebesgue integral|(}

We recall the essential facts. Let $(\Omega, \mathcal{A}, \mu)$ be a measure space.

\begin{definition}\label{def:measurable-function}
\index{measurable function}
A function $f \colon \Omega \to \overline{\R}$ (or $\C$) is \emph{measurable} if
$f^{-1}(B) \in \mathcal{A}$ for every Borel set $B$.
\end{definition}

\begin{theorem}[Monotone convergence]\label{thm:monotone-conv}
\index{monotone convergence theorem}
If $0 \le f_1 \le f_2 \le \cdots$ are measurable with $f_n \nearrow f$ pointwise,
then $\int f_n\, \dd\mu \nearrow \int f\, \dd\mu$.
\end{theorem}

\begin{theorem}[Dominated convergence]\label{thm:dominated-conv}
\index{dominated convergence theorem}
If $f_n \to f$ a.e., $\abs{f_n} \le g$ a.e.\ for some $g \in L^1(\mu)$, then
$f \in L^1(\mu)$ and $\int f_n\, \dd\mu \to \int f\, \dd\mu$.
\end{theorem}

\begin{proof}
By Fatou's lemma\index{Fatou's lemma} applied to $g + f_n \ge 0$ and $g - f_n \ge 0$:
\[
  \int g + \int f \le \liminf \int (g + f_n) = \int g + \liminf \int f_n,
\]
\[
  \int g - \int f \le \liminf \int (g - f_n) = \int g - \limsup \int f_n.
\]
Hence $\int f \le \liminf \int f_n \le \limsup \int f_n \le \int f$.
\end{proof}

\begin{theorem}[Fubini--Tonelli]\label{thm:fubini}
\index{Fubini theorem}
\index{Tonelli theorem}
Let $(X, \mu)$ and $(Y, \nu)$ be $\sigma$-finite measure spaces.
\begin{enumerate}[label=(\roman*)]
  \item (Tonelli) If $f \ge 0$ is measurable on $X \times Y$, then
  \[
    \int_{X \times Y} f\, \dd(\mu \otimes \nu)
    = \int_X \left(\int_Y f(x,y)\, \dd\nu(y)\right) \dd\mu(x)
    = \int_Y \left(\int_X f(x,y)\, \dd\mu(x)\right) \dd\nu(y).
  \]
  \item (Fubini) If $f \in L^1(\mu \otimes \nu)$, the same iterated integral
        equalities hold and the inner integrals define integrable functions a.e.
\end{enumerate}
\end{theorem}

\begin{remark}
Fubini's theorem is used extensively in Chapter~\ref{ch:Lp-duality} (for
convolution and duality proofs) and in Chapter~\ref{ch:compact-spectral} (for
Hilbert--Schmidt operators).
\end{remark}

\index{Lebesgue measure|)}
\index{Lebesgue integral|)}

\section{Additional Useful Results}\label{sec:additional-results}

\begin{theorem}[Urysohn's lemma]\label{thm:urysohn}
\index{Urysohn's lemma}
If $X$ is a normal topological space and $A, B \subset X$ are disjoint closed
sets, there exists a continuous function $f \colon X \to [0,1]$ with $f|_A = 0$
and $f|_B = 1$.
\end{theorem}

\begin{theorem}[Partition of unity]\label{thm:partition-unity}
\index{partition of unity}
Let $X$ be a paracompact Hausdorff space and $(U_\alpha)$ an open cover. There
exists a partition of unity $(\varphi_\alpha)$ subordinate to $(U_\alpha)$:
$\varphi_\alpha \ge 0$, $\supp\varphi_\alpha \subset U_\alpha$,
$\sum_\alpha \varphi_\alpha = 1$ (locally finite sum).
\end{theorem}

\begin{theorem}[Arzel\`a--Ascoli]\label{thm:arzela-ascoli}
\index{Arzela--Ascoli theorem@Arzel\`a--Ascoli theorem}
Let $K$ be a compact metric space. A subset $\mathcal{F} \subset C(K)$ is
relatively compact if and only if it is bounded and equicontinuous.
\end{theorem}

\begin{theorem}[Stone--Weierstrass]\label{thm:stone-weierstrass-app}
\index{Stone--Weierstrass theorem}
Let $K$ be a compact Hausdorff space and $A \subset C(K, \R)$ a subalgebra that
separates points and contains the constants. Then $A$ is dense in $C(K, \R)$.
The complex version requires $A$ to be closed under conjugation.
\end{theorem}

\index{appendix|)}


%%%%%%%%%%%%%%%%%%%%%%%%%%%%%%%%%%%%%%%%%%%%%%%%%%%%%%%%%%%%
%%  BIBLIOGRAPHY
%%%%%%%%%%%%%%%%%%%%%%%%%%%%%%%%%%%%%%%%%%%%%%%%%%%%%%%%%%%%
\begin{thebibliography}{99}

\bibitem{dales2000}
H.\,G. Dales,
\emph{Banach Algebras and Automatic Continuity},
London Mathematical Society Monographs~24, Clarendon Press, 2000.

\bibitem{James1964}
R.\,C. James,
Weakly compact sets,
\emph{Transactions of the AMS} \textbf{113} (1964), 129--140.

\bibitem{DunfordSchwartz1958}
N. Dunford and J.\,T. Schwartz,
\emph{Linear Operators, Part I: General Theory},
Interscience Publishers, 1958.

\bibitem{Megginson1998}
R.\,E. Megginson,
\emph{An Introduction to Banach Space Theory},
Graduate Texts in Mathematics~183, Springer-Verlag, 1998.

\end{thebibliography}

%%%%%%%%%%%%%%%%%%%%%%%%%%%%%%%%%%%%%%%%%%%%%%%%%%%%%%%%%%%%
%%  INDEX AND END
%%%%%%%%%%%%%%%%%%%%%%%%%%%%%%%%%%%%%%%%%%%%%%%%%%%%%%%%%%%%
\printindex
\end{document}
